%% LyX 2.4.4 created this file.  For more info, see https://www.lyx.org/.
%% Do not edit unless you really know what you are doing.
\documentclass[english,hebrew]{article}
\usepackage[HE8,T1]{fontenc}
\usepackage[utf8]{inputenc}
\setcounter{secnumdepth}{-2}
\setlength{\parindent}{0bp}
\usepackage{xcolor}
\usepackage{verbatim}
\usepackage{calc}
\usepackage{amsmath}
\usepackage{amssymb}
\PassOptionsToPackage{normalem}{ulem}
\usepackage{ulem}

\makeatletter

%%%%%%%%%%%%%%%%%%%%%%%%%%%%%% LyX specific LaTeX commands.
\DeclareTextSymbolDefault{\textquotedbl}{T1}
%% Because html converters don't know tabularnewline
\providecommand{\tabularnewline}{\\}

\@ifundefined{date}{}{\date{}}
%%%%%%%%%%%%%%%%%%%%%%%%%%%%%% User specified LaTeX commands.
\usepackage{graphicx}
\usepackage{geometry}
\usepackage{amsmath}
\usepackage{amssymb}
\newcommand{\limi}{\underset{n\rightarrow\infty}{\lim}}
\newcommand{\limxz}{\underset{x\rightarrow x_{0}}{\lim}}
\newcommand{\limfxz}{\underset{x\rightarrow x_{0}}{\lim}f(x)}
\newcommand{\limfxm}{\underset{x\rightarrow x_{0}^{-}}{\lim}f(x)}
\newcommand{\limfxp}{\underset{x\rightarrow x_{0}^{+}}{\lim}f(x)}
\newcommand{\limxm}{\underset{x\rightarrow x_{0}^{-}}{\lim}}
\newcommand{\limxp}{\underset{x\rightarrow x_{0}^{+}}{\lim}}
\newcommand{\sqrm}{M_{m\times m}(\mathbb{F})}
\newcommand{\seqa}{(a_{n})_{n=1}^{\infty}}
\newcommand{\seqx}{(x_{n})_{n=1}^{\infty}}
\newcommand{\funcdr}{f:D\rightarrow\mathbb{R}}
\newcommand{\matrmn}{M_{m\times n}(\mathbb{F})}
\newcommand{\sqrmn}{M_{n\times n}(\mathbb{F})}
\newcommand{\spn}{\text{Span }}

\makeatother

\usepackage{babel}
\begin{document}
\newgeometry{top=1cm, bottom=2cm, left=1cm, right=1cm}
\title{אלגברה ליניארית {\beginL 1\endL} - הסיכום המוכלל}
\author{נעם לב}

\maketitle
\noindent\fbox{\begin{minipage}[t]{1\columnwidth - 2\fboxsep - 2\fboxrule}%
\begin{center}
האוניברסיטה העברית, אלגברה ליניארית )1( - \L{$80134$}, סמסטר{\small{}
}{\beginL 2026\endL}{\small{} }א' )תשפ\textquotedblleft ו א'(.
\par\end{center}
\begin{center}
מרצה: ד\textquotedblright ר אלכסנדר )אלכס( גורביץ', מתרגל: מר יואב
כהן.
\par\end{center}
\begin{center}
הסיכום הוא לא רשמי ולא עבר אישור של הסגל, על כן אין להשתמש בו )ובין
השאר במספור שלו או בשמות הטענות שבו( כמקור רשמי. אם ישנן טעויות, הן
שלי ולא של הסגל.
\centerline{\textcolor{magenta}{גרסה \L{$1.5$}}}
\par\end{center}%
\end{minipage}}

\medskip{}

\begin{itemize}
\item בשחור - החומר עצמו. {]}שמות בסוגריים כאלו הם שמות שאני נתתי לטענות{[}.
\item \textcolor{blue}{בכחול - הערות והדגשות.}
\item \textcolor{violet}{בסגול - הערות והסברים מ-\L{$AI$}.}
\item \textcolor{teal}{בצבע הזה - הערות והסברים נוספים.}
\end{itemize}
\vfill{}

\par\noindent\centerline{%
  \L{\includegraphics[width=0.45\textheight]{\detokenize{./gen_sum_vol_2_figs/gauss.jpeg}}}%
}\par

\begin{comment}
שלבים בהכנה:
\begin{enumerate}
\item העתקה של תוכן הרשומות מהרצאות {\beginL 1-14\endL} בעזרת ג'מיני
\item הוספה של שאר ההרצאות שכבר כתבתי בליך
\item צפיה וסיכום של ההרצאות החסרות
\item וריפיקציה של ההרצאות שהעתקתי עם ליך
\item הוספה של תרגולים, מטלות בית, וכולי
\item הוספת הערות לחומרים שלא מבין
\end{enumerate}
מספר הרצאות כולל: {\beginL 29\endL}\\

\textcolor{red}{לתת שמות לכל הטענות והמשפטים וכו' שלא הספקתי לתת להם!}
\end{comment}

\newpage{}

\tableofcontents{}

\newpage{}

\noindent\fbox{\begin{minipage}[t]{1\columnwidth - 2\fboxsep - 2\fboxrule}%
\begin{center}
אלגברה ליניארית )1( - \L{$80134$} \textbar{} האונ' העברית \textbar{}
תשפ\textquotedblright ו א' \textbar{} מרצה: ד\textquotedblright ר
אלכס גורביץ' \textbar{} מתרגל: מר יואב כהן \textbar{} ליקוט ועריכה:
נעם לב
\par\end{center}%
\end{minipage}}

\part{הרצאה {\beginL 1\endL} - {\beginL 20.10.2025\endL}}

\section{כללי}
\begin{itemize}
\item בכל שבוע יש תרגיל בית להגשה. 
\item מומלץ מאוד להגיע להרצאות ולתרגולים. 
\item יש הרצאות, תרגולים ותרגולים פעילים.
\item נוכחות בתרגולים הפעילים מהווה חלק מהציון. 
\item חשוב לקבל משוב על הפתרונות שלנו וזה בעיקר בתרגולים הפעילים. 
\item למידע נוסף כדאי לקרוא את קובץ \textquotedblleft נהלי הקורס\textquotedblright{}
במודל. 
\end{itemize}
\noindent\rule{1\linewidth}{1pt}

אלגברה לינארית נועדה להכליל את השיטות והמושגים של ההנדסה האנליטית.

\section{נושאי לימוד}
\begin{itemize}
\item פרק 1 - אובייקטים קונקרטיים בהנדסה אנליטית \\
 )2 מימדים, 3 מימדים, תיאור ע"י משוואות( 
\item פרק 2 - מרחבים וקטוריים )הכללה של מישור( \\
 תלות לינארית \\
 בסיס 
\item פרק 3 - העתקות לינאריות 
\item פרק 4 - דטרמיננטות 
\end{itemize}
\noindent\rule{1\linewidth}{1pt}

\section{סקירת מושגים בסיסיים מתורת הקבוצות: )סימונים בסיסיים(}

\textbf{קבוצה} - אוסף של איברים.

כל איבר או שייך או לא שייך לקבוצה ולא שניהם ביחד.

אם איבר \L{$a$} \textbf{שייך} לקבוצה \L{$S$} רושמים \L{$a\in S$}.
אחרת רושמים \L{$a\notin S$}.\\


\subsection{דרכים להגדיר קבוצה}

1. ע\textquotedblright י רשימת האיברים \L{%
\[
S=\{1,5,7\}
\]
}%

2. ע\textquotedblright י קבוצה נתונה ותנאי )כך ש- ( \L{%
\[
\{x\in S\mid x>2\}=\{5,7\}
\]
}%
דהיינו - כל \L{$x$} שב-\L{$S$} כך ש-\L{$x$} גדול מ-2.\\

3. ע\textquotedblright י קבוצה נתונה וביטוי מתמטי \L{%
\[
A=\{(x-3)^{2}\mid x\in S\}=\{4,16\}
\]
}%
דהיינו - כל איבר מהצורה \L{$(x-3)^{2}$} כך ש- \L{$x\in S$}.

נשים לב שסדר הכתיבה משתנה בדוגמאות 2 ו-3 - זה קריטי.

\subsection{הכלה}

תהיינה \L{$S_{1},S_{2}$} קבוצות.

אומרים כי \L{$S_{1}$} מוכלת ב- \L{$S_{2}$} )או ש- \L{$S_{1}$} היא
תת-קבוצה של \L{$S_{2}$}(

ורושמים \L{$S_{1}\subseteq S_{2}$} כאשר כל אחד מהאיברים של \L{$S_{1}$}
שייך ל- \L{$S_{2}$}.

\subsection{פעולות על קבוצות}

\textbf{איחוד} של \L{$S_{1}$} ו- \L{$S_{2}$} מוגדר ע\textquotedblright י
\L{%
\[
S_{1}\cup S_{2}=\{x\mid x\in S_{1}\lor x\in S_{2}\}
\]
}%
נזכור ש- \textquotedblright או\textquotedblright{} במתמטיקה הוא \L{$\text{Inclusive Or}$},
דהיינו, או א' או ב' או שניהם.

\textbf{חיתוך} של \L{$S_{1}$} ו- \L{$S_{2}$} מוגדר ע\textquotedblright י
\L{%
\[
S_{1}\cap S_{2}=\{x\mid x\in S_{1}\land x\in S_{2}\}
\]
}%

\textbf{הפרש} של \L{$S_{1}$} ו- \L{$S_{2}$} מוגדר ע\textquotedblright י
\L{%
\[
S_{1}\setminus S_{2}=\{x\mid x\in S_{1}\land x\notin S_{2}\}
\]
}%

הקבוצה שלה לא שייך אף איבר נקראת \textbf{הקבוצה הריקה} ומסומנת ב-
\L{$\emptyset$}.

\L{%
\[
\emptyset=\{\}
\]
}%


\subsection{קבוצות מספרים}

קבוצת המספרים השלמים: \L{%
\[
\mathbb{Z}=\{\dots,-2,-1,0,1,2,3,\dots\}
\]
}%

קבוצת המספרים הטבעיים: \L{%
\[
\mathbb{N}=\{x\in\mathbb{Z}\mid x>0\}
\]
}%
)יש מקומות בהם מגדירים \L{$x\ge0$} אך לא בקורס זה(

קבוצת המספרים הרציונליים: \L{%
\[
\mathbb{Q}=\left\{ \frac{m}{n}\mid m,n\in\mathbb{Z}\ ,\ n\neq0\right\} 
\]
}%

קבוצת המספרים הממשיים: \L{$\mathbb{R}$}

לא נגדיר במדויק בקורס זה.

אפשר להבין אותה אינטואיטיבית כציר המספרים

\L{%
\[
\begin{array}{ccccc}
\rule{2cm}{0.5pt} & | & \rule{1cm}{0.5pt} & | & \rule{4cm}{0.5pt}\\
 & -1 &  & 0
\end{array}
\]
}%

לכל מספר ממשי מתאימה נק' על ציר המספרים. \\

\textcolor{green}{לגבי ההרצאה בימי חמישי - בפועל לרוב לא תהיה הרצאה.
רק במקרים מיוחדים כשיהיה צורך לשיעור השלמה. יודיעו על כך במודל. כדאי
להסתכל גם איזה מרצה שלח את ההודעה כי כל המרצים יכולים לקבוע זאת.}

\section{\L{$n$}-יה סדורה}

בקבוצה \L{$\mathbb{R}$} מוגדרות פעולות חיבור וכפל שמקיימות חוקים
מוכרים.\\

יהי \L{$n\in\mathbb{N}$}.

\textbf{\L{$n$}-יה סדורה} של מספרים ממשיים היא ביטוי מהצורה \L{$\begin{pmatrix}a_{1}\\
a_{2}\\
\vdots\\
a_{n}
\end{pmatrix}$} כאשר \L{$a_{1},\dots,a_{n}\in\mathbb{R}$}.

בניגוד לקבוצות, ב-\L{$n$}-יה סדר האיברים חשוב.

לדוגמה: \L{$\{1,2\}=\{2,1\}$} אבל \L{$\begin{pmatrix}1\\
2
\end{pmatrix}\neq\begin{pmatrix}2\\
1
\end{pmatrix}$}.

בקבוצות לחזרות אין משמעות, ב-\L{$n$}-יה לחזרות יש משמעות.

לדוגמה: \L{$\{1,1\}=\{1\}$} אבל \L{$\begin{pmatrix}1\\
1
\end{pmatrix}\neq(1)$}.

הקבוצה של כל ה-\L{$n$}-יות מסומנת ב- \L{$\mathbb{R}^{n}$}

כלומר \L{%
\[
\mathbb{R}^{n}=\left\{ \begin{pmatrix}a_{1}\\
\vdots\\
a_{n}
\end{pmatrix}\,\middle|\,a_{1},\dots,a_{n}\in\mathbb{R}\right\} 
\]
}%

כאשר \L{$n=2$} מקבלים \L{%
\[
\mathbb{R}^{2}=\left\{ \begin{pmatrix}a_{1}\\
a_{2}
\end{pmatrix}\,\middle|\,a_{1},a_{2}\in\mathbb{R}\right\} 
\]
}%

שזה למעשה המישור הקרטזי.

\medskip{}

\par\noindent\centerline{%
  \L{\includegraphics[width=0.3\textheight]{\detokenize{./gen_sum_vol_2_figs/fig1.png}}}%
}\par

\medskip{}

ב- \L{$\mathbb{R}^{n}$} מוגדרות הפעולות הבאות: )מקבעים את \L{$n$}
היות ואי אפשר לחבר \L{$n$}-יות מסדר שונה.(

1( \textbf{חיבור} \L{%
\[
\begin{pmatrix}a_{1}\\
\vdots\\
a_{n}
\end{pmatrix}+\begin{pmatrix}b_{1}\\
\vdots\\
b_{n}
\end{pmatrix}=\begin{pmatrix}a_{1}+b_{1}\\
\vdots\\
a_{n}+b_{n}
\end{pmatrix}
\]
}%

2( \textbf{כפל של \L{$n$}-יה במספר )כפל בסקלר(}

הסדר חשוב, כותבים את המספר בצד שמאל.

עבור \L{$c\in\mathbb{R}$} \L{%
\[
c\begin{pmatrix}a_{1}\\
\vdots\\
a_{n}
\end{pmatrix}=\begin{pmatrix}ca_{1}\\
\vdots\\
ca_{n}
\end{pmatrix}
\]
}%


\section{משוואות לינאריות}

יהי \L{$n\in\mathbb{N}$}.

משוואה לינארית ב-\L{$n$} נעלמים היא משוואה מהצורה \L{%
\[
a_{1}x_{1}+\dots+a_{n}x_{n}=b\quad\quad\quad\qquad\qquad(I)
\]
}%
כאשר \L{$x_{1},\dots,x_{n}$} הם הנעלמים ו- \L{$a_{1},\dots,a_{n},b\in\mathbb{R}$}
הם המקדמים של המשוואה. \L{$b$} נקרא מקדם חופשי.

\L{$n$}-יה \L{$\begin{pmatrix}d_{1}\\
\vdots\\
d_{n}
\end{pmatrix}\in\mathbb{R}^{n}$} נקראת פתרון של המשוואה )I( כאשר \L{$a_{1}d_{1}+...+a_{n}d_{n}=b$}.

\textbf{דוגמה:}

עם \L{$n=2$}.

\L{$x_{1}-x_{2}=1\quad\quad\quad\qquad\qquad(II)$}

כאן \L{$a_{1}=1,a_{2}=-1,b=1$}.

\L{$\begin{pmatrix}1\\
0
\end{pmatrix}$} הוא פתרון של )II( כי \L{$1-0=1$}.

גם \L{$\begin{pmatrix}2\\
1
\end{pmatrix}$} הוא פתרון של )II( כי \L{$2-1=1$}.

\L{$\begin{pmatrix}0\\
1
\end{pmatrix}$} אינו פתרון של )II( כי \L{$0-1\neq1$}

עכשיו נסתכל על הקבוצה של כל הפתרונות של משוואה מסוימת.

קבוצת הפתרונות של )I( היא

\L{%
\[
\left\{ \begin{pmatrix}x_{1}\\
\vdots\\
x_{n}
\end{pmatrix}\in\mathbb{R}^{n}\,\middle|\,a_{1}x_{1}+\dots+a_{n}x_{n}=b\right\} 
\]
}%

דוגמה: \L{$x_{1}-x_{2}=1\quad(II)$}

קבוצת הפתרונות של )II( היא

\L{%
\[
\left\{ \begin{pmatrix}x_{1}\\
x_{2}
\end{pmatrix}\in\mathbb{R}^{2}\,\middle|\,x_{1}-x_{2}=1\right\} 
\]
}%

\L{%
\[
=\left\{ \begin{pmatrix}x_{1}\\
x_{2}
\end{pmatrix}\in\mathbb{R}^{2}\,\middle|\,x_{1}=x_{2}+1\right\} 
\]
}%

\L{%
\[
=\left\{ \begin{pmatrix}x_{2}+1\\
x_{2}
\end{pmatrix}\,\middle|\,x_{2}\in\mathbb{R}\right\} 
\]
}%

\L{%
\[
=\left\{ \begin{pmatrix}1\\
0
\end{pmatrix}+\begin{pmatrix}x_{2}\\
x_{2}
\end{pmatrix}\,\middle|\,x_{2}\in\mathbb{R}\right\} 
\]
}%

\L{%
\[
=\left\{ \begin{pmatrix}1\\
0
\end{pmatrix}+x_{2}\begin{pmatrix}1\\
1
\end{pmatrix}\,\middle|\,x_{2}\in\mathbb{R}\right\} 
\]
}%

הנוסחה האחרונה נקראת הצגה פרמטרית לקבוצת הפתרונות של המשוואה )II(.

\newpage{}

\noindent\fbox{\begin{minipage}[t]{1\columnwidth - 2\fboxsep - 2\fboxrule}%
\begin{center}
אלגברה ליניארית )1( - \L{$80134$} \textbar{} האונ' העברית \textbar{}
תשפ\textquotedblright ו א' \textbar{} מרצה: ד\textquotedblright ר
אלכס גורביץ' \textbar{} מתרגל: מר יואב כהן \textbar{} ליקוט ועריכה:
נעם לב
\par\end{center}%
\end{minipage}}

\part{הרצאה {\beginL 2\endL} - {\beginL 21.10.2025\endL}}

\section{תזכורת}

התחלנו לדבר על משוואות לינאריות. בסוף השיעור הקודם ראינו דוגמה למשוואה
לינארית וכתבנו את קבוצת הפתרונות שלה בהצגה פרמטרית.

\L{%
\[
(II)\quad x_{1}-x_{2}=1
\]
}%

קבוצת הפתרונות של )II( היא

\L{%
\[
\left\{ \begin{pmatrix}x_{1}\\
x_{2}
\end{pmatrix}\in\mathbb{R}^{2}\,\middle|\,x_{1}-x_{2}=1\right\} 
\]
}%

\L{%
\[
=\left\{ \begin{pmatrix}1\\
0
\end{pmatrix}+x_{2}\begin{pmatrix}1\\
1
\end{pmatrix}\,\middle|\,x_{2}\in\mathbb{R}\right\} 
\]
}%

להצגה הפרמטרית יש גם משמעות גיאומטרית,

\medskip{}

\par\noindent\centerline{%
  \L{\includegraphics[width=0.3\textheight]{\detokenize{./gen_sum_vol_2_figs/fig2.png}}}%
}\par

\medskip{}

לוקחים את הנקודה \L{$\begin{pmatrix}1\\
0
\end{pmatrix}$}. מעבירים דרכה ישר שמקביל לישר שעובר בראשית הצירים ובנקודה שמתאימה
ל- \L{$x_{2}\begin{pmatrix}1\\
1
\end{pmatrix}$}. כאן זה הישר הירוק. זה מייצג את קבוצת הפתרונות של המשוואה. זה למעשה
\L{$y=x-1$}. )זו רק דוגמה, זה לא קריטי להבין(.

\noindent\textcolor{black}{\rule{1\linewidth}{1.5pt}}

\section{צורה כללית של הצגה פרמטרית}

הצגה פרמטרית של קבוצת הפתרונות של משוואה לינארית ב- \L{$n$} נעלמים
היא הצגה מהצורה:

\L{%
\[
\left\{ u_{0}+t_{1}u_{1}+\dots+t_{k}u_{k}\mid t_{1},\dots,t_{k}\in\mathbb{R}\right\} 
\]
}%

כאשר \L{$k\in\mathbb{N}\cup\{0\}$} ו- \L{$u_{0},u_{1},\dots,u_{k}\in\mathbb{R}^{n}$}.

דוגמה: בהצגה פרמטרית של קבוצת הפתרונות של )II(: \L{$k=1$} , \L{$t_{1}=x_{2}$}
, \L{$u_{1}=\begin{pmatrix}1\\
1
\end{pmatrix}$} , \L{$u_{0}=\begin{pmatrix}1\\
0
\end{pmatrix}$}.

\section{מערכות משוואות לינאריות}

\subsection{הגדרה - מערכת משוואות לינאריות )ממ\textquotedblright ל(}

מערכת של \L{$m$} משוואות לינאריות ב- \L{$n$} נעלמים היא )נסמן את
המבנה הבא ב- \L{$(*)$}(

\L{%
\[
\begin{cases}
a_{11}x_{1}+...+a_{1n}x_{n}=b_{1}\\
a_{21}x_{1}+...+a_{2n}x_{n}=b_{2}\\
\vdots\\
a_{m1}x_{1}+...+a_{mn}x_{n}=b_{m}
\end{cases}
\]
}%

כאשר \L{$x_{1},\dots,x_{n}$} הם הנעלמים ו- \L{$a_{i,j},b_{i}\in\mathbb{R}$}
המקדמים )ה-\L{$b$}-ים הם המקדמים החופשיים(.\\

דוגמה: \L{$m=2,n=2$}:

\L{%
\[
\begin{cases}
x_{1}-x_{2}=1\\
x_{2}=2
\end{cases}\quad(III)
\]
}%

כאן מתקיים:

\L{$a_{11}=1\ ,\ a_{12}=-1\ ,\ a_{21}=0$},

\L{$a_{22}=1\ ,\ b_{1}=1\ ,\ b_{2}=2$}

\subsection{פתרון של מערכת משוואות לינאריות )הגדרה(}

\L{$n$}-יה \L{$\begin{pmatrix}d_{1}\\
\vdots\\
d_{n}
\end{pmatrix}\in\mathbb{R}^{n}$} נקראת פתרון של ממ\textquotedblright ל כאשר היא מהווה פתרון של כל
אחת ממשוואות המערכת. לפיכך, הקבוצה של כל הפתרונות של המערכת היא החיתוך
של קבוצת הפתרונות של משוואות המערכת.

למשל, עבור )III( , קבוצת הפתרונות היא \L{$\left\{ \begin{pmatrix}3\\
2
\end{pmatrix}\right\} $}.

\medskip{}

\par\noindent\centerline{%
  \L{\includegraphics[width=0.3\textheight]{\detokenize{./gen_sum_vol_2_figs/fig3.png}}}%
}\par

\medskip{}


\section{מציאת הצגה פרמטרית עבור קבוצת הפתרונות של ממ\textquotedblright ל}

המטרה: בהינתן ממ\textquotedblright ל, למצוא הצגה פרמטרית עבור קבוצת
הפתרונות שלה.

\subsection{הגדרה - מטריצה של ממ\textquotedblright ל}

המטריצה של מערכת מהצורה \L{$(*)$} מוגדרת להיות

\L{%
\[
\begin{pmatrix}a_{11} & ... & a_{1n}\\
a_{21} & ... & a_{2n}\\
\vdots &  & \vdots\\
a_{m1} & ... & a_{mn}
\end{pmatrix}
\]
}%


\subsection{הגדרה - מטריצה \underline{מורחבת} של ממ\textquotedblright ל}

המטריצה המורחבת של מערכת מהצורה \L{$(*)$} מוגדרת להיות

\L{%
\[
\left(\begin{array}{ccc|c}
a_{11} & \dots & a_{1n} & b_{1}\\
a_{21} & \dots & a_{2n} & b_{2}\\
 & \vdots &  & \vdots\\
a_{m1} & \dots & a_{mn} & b_{m}
\end{array}\right)
\]
}%

\textbf{דוגמה:} המטריצה המורחבת של הממ\textquotedblright ל )III( היא:

\L{%
\[
\left(\begin{array}{cc|c}
1 & -1 & 1\\
0 & 1 & 2
\end{array}\right)
\]
}%

כמות העמודות במטריצה מורחבת היא \L{$n+1$}.

\section{מטריצה מדורגת מצומצמת}

\subsection{הגדרה - האיבר המוביל}

האיבר המוביל בשורה כלשהי של מטריצה, האיבר המוביל הוא האיבר השמאלי
ביותר בשורה אשר שונה מ-0.

\L{%
\[
\left(\begin{array}{cc|c}
\underline{1} & -1 & 1\\
0 & \underline{1} & 2
\end{array}\right)
\]
}%

רק בשורה שכולה אפסים אין איבר מוביל.

יתכן שהאיבר המוביל יהיה בעמודה האחרונה )כלומר, מקדם חופשי(.

\subsection{מטריצה מדורגת מצומצמת}

\textbf{מטריצה נקראת מדורגת מצומצמת כאשר:} 
\begin{enumerate}
\item האיבר המוביל של כל שורה )במידה וקיים - שורה אפסים מתקבלת( שווה ל-1. 
\item כל המספרים בעמודה של איבר מוביל של שורה כלשהי )חוץ ממנו( שווים ל-0. 
\item אם בשורה כלשהי )חוץ מהשורה הראשונה( קיים איבר מוביל, אז גם בשורה שמעליה
קיים איבר מוביל והוא נמצא בצד שמאל מהאיבר המוביל של השורה המקורית. 
\end{enumerate}
\textbf{דוגמה:} \L{%
\[
\left(\begin{array}{ccccc|c}
0 & 1 & -2 & 0 & -1 & 3\\
0 & 0 & 0 & 1 & 2 & -1\\
0 & 0 & 0 & 0 & 0 & 0
\end{array}\right)
\]
}%
)3 משוואות, 6 מקדמים(. זו מטריצה מדורגת מצומצמת.

\bigskip{}


\section{שורת סתירה}

אם המטריצה המורחבת של ממ\textquotedblright ל כלשהי היא מדורגת מצומצמת
וקיימת בה שורה שבה איבר מוביל בעמודה האחרונה, אז אחת המשוואות של המערכת
היא: \L{%
\[
0x_{1}+\dots+0x_{n}=1
\]
}%
ולמערכת \textbf{אין פתרונות}.

שורה כזאת )שורה עם איבר מוביל בעמודה האחרונה של מטריצה מורחבת מדורגת
מצומצמת( נקראת \textbf{שורת סתירה}.

\bigskip{}


\section{נעלמים קשורים וחופשיים}

\textbf{אחרת )אם למטריצה מורחבת מדורגת מצומצמת של ממ\textquotedblright ל
אין שורת סתירה( למערכת יש פתרון:}

אם המטריצה המורחבת של ממ\textquotedblright ל היא מדורגת מצומצמת ואין
בה שורה עם איבר מוביל בעמודה האחרונה, נקרא לנעלמים המתאימים לעמודות
עם איבר מוביל \textbf{קשורים} ולנעלמים המתאימים לעמודות בלי איבר מוביל
\textbf{חופשיים}. נבטא את הנעלמים הקשורים דרך החופשיים ובצורה זו נקבל
הצגה פרמטרית לקבוצת הפתרונות של הממ\textquotedblright ל.

\subsection{דוגמה }

\L{%
\[
\begin{cases}
x_{2}-2x_{3}-x_{5}=3\\
x_{4}+2x_{5}=-1\\
0=0
\end{cases}
\]
}%

המטריצה המתאימה: \L{%
\[
\left(\begin{array}{ccccc|c}
0 & 1 & -2 & 0 & -1 & 3\\
0 & 0 & 0 & 1 & 2 & -1\\
0 & 0 & 0 & 0 & 0 & 0
\end{array}\right)
\]
}%

נעלמים קשורים: \L{$x_{2},x_{4}$}

נעלמים חופשיים: \L{$x_{1},x_{3},x_{5}$}

נבטא את הנעלמים הקשורים דרך החופשיים:

\L{%
\[
\begin{cases}
x_{2}=3+2x_{3}+x_{5}\\
x_{4}=-1-2x_{5}
\end{cases}
\]
}%

אז קבוצת הפתרונות היא:

\L{%
\[
\left\{ \begin{pmatrix}x_{1}\\
x_{2}\\
x_{3}\\
x_{4}\\
x_{5}
\end{pmatrix}\in\mathbb{R}^{5}\,\middle|\,\begin{array}{l}
x_{2}=3+2x_{3}+x_{5}\\
x_{4}=-1-2x_{5}
\end{array}\right\} 
\]
}%

\L{%
\[
=\left\{ \begin{pmatrix}x_{1}\\
3+2x_{3}+x_{5}\\
x_{3}\\
-1-2x_{5}\\
x_{5}
\end{pmatrix}\,\middle|\,x_{1},x_{3},x_{5}\in\mathbb{R}\right\} 
\]
}%

\L{%
\[
=\left\{ \begin{pmatrix}0\\
3\\
0\\
-1\\
0
\end{pmatrix}+\begin{pmatrix}x_{1}\\
0\\
0\\
0\\
0
\end{pmatrix}+\begin{pmatrix}0\\
2x_{3}\\
x_{3}\\
0\\
0
\end{pmatrix}+\begin{pmatrix}0\\
x_{5}\\
0\\
-2x_{5}\\
x_{5}
\end{pmatrix}\,\middle|\,x_{1},x_{3},x_{5}\in\mathbb{R}\right\} 
\]
}%

\L{%
\[
=\left\{ \underbrace{\begin{pmatrix}0\\
3\\
0\\
-1\\
0
\end{pmatrix}}_{u_{0}}+x_{1}\underbrace{\begin{pmatrix}1\\
0\\
0\\
0\\
0
\end{pmatrix}}_{u_{1}}+x_{3}\underbrace{\begin{pmatrix}0\\
2\\
1\\
0\\
0
\end{pmatrix}}_{u_{2}}+x_{5}\underbrace{\begin{pmatrix}0\\
1\\
0\\
-2\\
1
\end{pmatrix}}_{u_{3}}\,\middle|\,x_{1},x_{3},x_{5}\in\mathbb{R}\right\} 
\]
}%

\L{%
\[
t_{1}=x_{1},\ t_{2}=x_{3},\ t_{3}=x_{5},\ k=3
\]
}%

\newpage{}

\noindent\fbox{\begin{minipage}[t]{1\columnwidth - 2\fboxsep - 2\fboxrule}%
\begin{center}
אלגברה ליניארית )1( - \L{$80134$} \textbar{} האונ' העברית \textbar{}
תשפ\textquotedblright ו א' \textbar{} מרצה: ד\textquotedblright ר
אלכס גורביץ' \textbar{} מתרגל: מר יואב כהן \textbar{} ליקוט ועריכה:
נעם לב
\par\end{center}%
\end{minipage}}

\part{הרצאה {\beginL 3\endL} - {\beginL 27.10.2025\endL}}

נלמד פעולות על מטריצות שלא משנות את הקבוצת הפתרונות של המערכת, אף
שהן משנות את מערכת המשוואות.

נקראות פעולות שורה אלמנטריות.

\section{פעולות שורה אלמנטריות}

כל אחת מהפעולות שומרת על מימדי המטריצה.

יהיו \L{$m,n\in\mathbb{N}$}. פעולת שורה אלמנטריות )פש\textquotedblright א(
הן פעולות שמתאימות לכל מטריצה בגודל \L{$m\times n$} מטריצה אחרת באותו
גודל בהתאם לכללים מסוימים. ישנם 3 סוגים:

\subsection{{\beginL 1\endL}. כפל בסקלר}

כפל של השורה ה-\L{$i$} ב-\L{$0\neq c\in\mathbb{R}$}.

\L{%
\[
R_{i}\to cR_{i}
\]
}%

כאשר \L{$R_{i}$} זו השורה ה-\L{$i$} של המטריצה, מלשון \L{$\text{Row}$}.

\subsection{{\beginL 2\endL}. הוספה של שורה אחרת כפול סקלר}

הוספה של השורה ה-\L{$j$} המוכפלת ב-\L{$c\in\mathbb{R}$} לשורה ה-\L{$i$}
)\L{$i\neq j$}(

\L{%
\[
R_{i}\to R_{i}+cR_{j}
\]
}%


\subsection{{\beginL 3\endL}. החלפת שורות}

החלפת השורה ה-\L{$i$} והשורה ה-\L{$j$} )\L{$i\neq j$}(

\L{%
\[
R_{i}\leftrightarrow R_{j}
\]
}%


\subsubsection{פעולה הפוכה}

תהי \L{$A$} מטריצה בגודל \L{$m\times n$}.

יהי \L{$\mathcal{E}$} פש\textquotedblright א, המטריצה המתקבלת מ-\L{$A$}
ע\textquotedblright י ביצוע של \L{$\mathcal{E}$} מסומנת ב- \L{$\mathcal{E}(A)$}.

פש\textquotedblright א \L{$\delta,\mathcal{E}$} נקראות הפוכות זו
לזו כאשר לכל מטריצה \L{$A$} מתקיים \L{$\delta(\mathcal{E}(A))=A$}
וגם \L{$\mathcal{E}(\delta(A))=A$}.\\

מתקיים: לכל פש\textquotedblright א קיימת פש\textquotedblright א הפוכה
)ויחידה, אבל לא הוכחנו זאת(.

אם \L{$\mathcal{E}=(R_{i}\to cR_{i})$} אז פש\textquotedblright א
ההפוכה היא \L{$\delta=(R_{i}\to\frac{1}{c}R_{i})$}.

אם \L{$\mathcal{E}=(R_{i}\to R_{i}+cR_{j})$} אז פש\textquotedblright א
ההפוכה היא \L{$\delta=(R_{i}\to R_{i}-cR_{j})$}.

אם \L{$\mathcal{E}=(R_{i}\leftrightarrow R_{j})$} אז פש\textquotedblright א
ההפוכה היא \L{$\delta=\mathcal{E}$}.\\

\textcolor{teal}{)העשרה שהיא הצצה לשלב יותר מתקדם בקורס: פש\textquotedblright א
שקולה לכפל משמאל במטריצה אלמנטרית ופש\textquotedblright א הפוכה שקולה
לכפל משמאל במטריצה האלמנטרית ההפוכה. היות ולכל מטריצה הפיכה קיימת
מטריצה הפוכה אחת בדיוק, גם לכל פש\textquotedblright א קיימת פש\textquotedblright א
הפוכה אחת בדיוק(.}

\section{טענה {]}אינווריאנטיות קבוצת הפתרונות של ממ\textquotedblright ל תחת
פש\textquotedblright א{[}}

נתונות שתי ממ\textquotedblright ל בעלות \L{$m$} משוואות ו-\L{$n$}
נעלמים כך שהמטריצה המורחבת של הממ\textquotedblright ל השניה מתקבלת
מהמטריצה המורחבת של הממ\textquotedblright ל הראשונה ע\textquotedblright י
ביצוע פש\textquotedblright א כלשהי. אז קבוצת הפתרונות של שתי המערכות
האלה זהות.

\subsection{הוכחה}

{]}נוכיח עבור פש\textquotedblright א אחת, השאר כתרגיל לבית{[}

נוכיח את הטענה במקרה שבו \L{$\mathcal{E}=R_{i}\to R_{i}+cR_{j}$}
)שאר המקרים - תרגיל(.

כדי להראות שקבוצות הפתרונות של שתי ממ\textquotedblright ל האלה זהות
מספיק להראות כי כל פתרון של ממ\textquotedblright ל הראשונה הוא פתרון
של ממ\textquotedblright ל השניה, כי המטריצה של הממ\textquotedblright ל
הראשונה מתקבלת מהמטריצה של הממ\textquotedblright ל השניה ע\textquotedblright י
ביצוע פש\textquotedblright א הפוכה, לפש\textquotedblright א הנתונה.

נסמן את המטריצה המורחבת של הממ\textquotedblright ל הראשונה ב-

\L{%
\[
A=\left(\begin{array}{ccc|c}
a_{11} & \dots & a_{1n} & b_{1}\\
\vdots &  & \vdots & \vdots\\
\vdots &  & \vdots & \vdots\\
a_{m1} & \dots & a_{mn} & b_{m}
\end{array}\right)
\]
}%

יהי \L{$\begin{pmatrix}d_{1}\\
\vdots\\
d_{n}
\end{pmatrix}\in\mathbb{R}^{n}$} פתרון של הממ\textquotedblright ל הראשונה. אז בפרט הוא פתרון של המשוואות
ה-\L{$i$} וה-\L{$j$} של הממ\textquotedblright ל הראשונה, כלומר

\L{%
\[
a_{i1}d_{1}+\dots+a_{in}d_{n}=b_{i}
\]
}%
\L{%
\[
a_{j1}d_{1}+\dots+a_{jn}d_{n}=b_{j}
\]
}%

כל משוואה פרט למשוואה ה-\L{$i$} בממ\textquotedblright ל השניה היא
כמו בממ\textquotedblright ל הראשונה. המשוואה ה-\L{$i$} בממ\textquotedblright ל
השניה היא

\L{%
\[
(a_{i1}+ca_{j1})x_{1}+...+(a_{in}+ca_{jn})x_{n}=b_{i}+cb_{j}
\]
}%

נציב את \L{$\begin{pmatrix}x_{1}\\
\vdots\\
x_{n}
\end{pmatrix}=\begin{pmatrix}d_{1}\\
\vdots\\
d_{n}
\end{pmatrix}$} ונקבל

\L{%
\[
(a_{i1}+ca_{j1})d_{1}+\dots+(a_{in}+ca_{jn})d_{n}=
\]
}%
\L{%
\[
(a_{i1}d_{1}+\dots+a_{in}d_{n})+c(a_{j1}d_{1}+\dots+a_{jn}d_{n})=b_{i}+cb_{j}
\]
}%

זאת אומרת ש- \L{$\begin{pmatrix}d_{1}\\
\vdots\\
d_{n}
\end{pmatrix}$} הוא פתרון של הממ\textquotedblright ל השניה וסיימנו.
\begin{flushright}
\L{$\square$}
\par\end{flushright}

\section{הגדרה - שקילות שורה}

תהיינה \L{$A,B$} מטריצות בגודל \L{$m\times n$}. אומרים כי \L{$B$}
שקולת שורה ל-\L{$A$} כאשר ניתן לקבל את \L{$B$} מ-\L{$A$} ע\textquotedblright י
ביצוע של סדרת פש\textquotedblright א.

\subsection{מסקנה}

אם נתונות שתי ממ\textquotedblright ל כך שהמטריצות המורחבות שלהן שקולות
שורה, אז קבוצות הפתרונות של הממ\textquotedblright ל האלה זהות.

\medskip{}


\section{אלגוריתם דירוג מטריצות}

\textcolor{blue}{החלק הבא יעסוק באלגוריתם למציאת מטריצה מדורגת מצומצמת
השקולת שורה למטריצה נתונה.}

\subsection{עובדה )שלא נוכיח(}

לכל מטריצה קיימת מטריצה מדורגת מצומצמת שקולת שורה לה והיא יחידה.

\subsection{האלגוריתם}

1( נשתמש בפש\textquotedblright א מהסוג \L{$R_{i}\to cR_{i}$} כדי
להשיג שכל איבר מוביל יהיה שווה ל-1.

2( נשתמש בפש\textquotedblright א מהסוג \L{$R_{i}\to R_{i}+cR_{j}$}
כדי להשיג שכל המספרים בעמודה של איבר מוביל )חוץ ממנו( שווים ל-0.

3( נשתמש בפש\textquotedblright א מהסוג \L{$R_{i}\leftrightarrow R_{j}$}
כדי להבטיח קיום של התנאי השלישי בהגדרה של מטריצה מדורגת מצומצמת )ה-\textquotedblright דירוג\textquotedblright (.

\subsection{דוגמה}

\L{%
\[
\left(\begin{array}{cccc}
0 & 2 & 4 & -2\\
3 & 2 & -2 & -5\\
1 & -1 & -4 & 1
\end{array}\right)\xrightarrow{R_{1}\to\frac{1}{2}R_{1}}
\]
}%

\L{%
\[
\left(\begin{array}{cccc}
0 & 1 & 2 & -1\\
3 & 2 & -2 & -5\\
1 & -1 & -4 & 1
\end{array}\right)\xrightarrow{R_{2}\to R_{2}-2R_{1}}
\]
}%

\L{%
\[
\left(\begin{array}{cccc}
0 & 1 & 2 & -1\\
3 & 0 & -6 & -3\\
1 & -1 & -4 & 1
\end{array}\right)\xrightarrow{R_{3}\to R_{3}+R_{1}}
\]
}%

\L{%
\[
\left(\begin{array}{cccc}
0 & 1 & 2 & -1\\
3 & 0 & -6 & -3\\
1 & 0 & -2 & 0
\end{array}\right)\xrightarrow{R_{2}\to\frac{1}{3}R_{2}}
\]
}%

\L{%
\[
\left(\begin{array}{cccc}
0 & 1 & 2 & -1\\
1 & 0 & -2 & -1\\
1 & 0 & -2 & 0
\end{array}\right)\xrightarrow{R_{3}\to R_{3}-R_{2}}
\]
}%

\L{%
\[
\left(\begin{array}{cccc}
0 & 1 & 2 & -1\\
1 & 0 & -2 & -1\\
0 & 0 & 0 & 1
\end{array}\right)\xrightarrow{R_{1}\to R_{1}+R_{3}}
\]
}%

\L{%
\[
\left(\begin{array}{cccc}
0 & 1 & 2 & 0\\
1 & 0 & -2 & 0\\
0 & 0 & 0 & 1
\end{array}\right)\xrightarrow{R_{1}\leftrightarrow R_{2}}
\]
}%

\L{%
\[
\left(\begin{array}{cccc}
1 & 0 & -2 & 0\\
0 & 1 & 2 & 0\\
0 & 0 & 0 & 1
\end{array}\right)
\]
}%

\textcolor{blue}{שזו מטריצה מדורגת מצומצמת השקולת שורה למטריצה המקורית.}

\textcolor{blue}{נציין שהדרך שעברנו היא לא הדרך היחידה לעשות זאת,
אבל בסוף תמיד מתקבלת אותה מטריצה מדורגת מצומצמת )לא נוכיח או נשתמש
בזה(.}

\rule{1\linewidth}{1.5pt}

\section*{\textcolor{violet}{סיכום: דירוג מטריצות ואלגוריתם גאוס-ז'ורדן}}

\subsection*{\textcolor{violet}{הגדרה ומטרה}}

\textcolor{violet}{אלגוריתם גאוס-ז'ורדן )\L{$\text{Gauss-Jordan Elimination}$}(
הוא שיטה אלגוריתמית במערכות ליניאריות לביצוע פעולות שורה אלמנטריות
על מטריצה נתונה, במטרה להביאה ל}\textbf{\textcolor{violet}{צורה מדורגת
קנונית}}\textcolor{violet}{{} )\L{$\text{Reduced Row Echelon Form - RREF}$}(.
שיטה זו משמשת לפתרון מערכות משוואות ליניאריות, למציאת המטריצה ההופכית,
ולחישוב הדרגה )\L{$Rank$}( של המטריצה {]}מושגים שנלמד בהמשך{[}.}

\subsection*{\textcolor{violet}{פעולות שורה אלמנטריות}}

\textcolor{violet}{במהלך הדירוג, מותר לבצע שלוש פעולות הפיכות על שורות
המטריצה, אשר אינן משנות את מרחב הפתרונות של המערכת: }
\begin{enumerate}
\item \textbf{\textcolor{violet}{החלפת שורות}}\textcolor{violet}{: החלפה
בין שורה \L{$R_{i}$} לשורה \L{$R_{j}$}. מסומן: \L{$R_{i}\leftrightarrow R_{j}$}. }
\item \textbf{\textcolor{violet}{כפל בסקלר}}\textcolor{violet}{: כפל כל
איברי השורה \L{$R_{i}$} בסקלר \L{$\lambda\neq0$}. מסומן: \L{$R_{i}\to\lambda R_{i}$}. }
\item \textbf{\textcolor{violet}{הוספת שורה}}\textcolor{violet}{: הוספת
כפולה של שורה \L{$R_{j}$} לשורה אחרת \L{$R_{i}$}. מסומן: \L{$R_{i}\to R_{i}+\lambda R_{j}$}. }
\end{enumerate}

\subsection*{\textcolor{violet}{צורה מדורגת קנונית )\L{$\text{RREF}$}(}}

\textcolor{violet}{מטריצה תיקרא }\textbf{\textcolor{violet}{מדורגת
קנונית}}\textcolor{violet}{{} אם מתקיימים התנאים הבאים: }
\begin{enumerate}
\item \textcolor{violet}{כל שורות האפסים )אם ישנן( מופיעות בתחתית המטריצה. }
\item \textcolor{violet}{האיבר הראשון שאינו אפס בכל שורה )הנקרא }\textbf{\textcolor{violet}{איבר
מוביל}}\textcolor{violet}{{} או \L{$\text{Pivot}$}( הוא \L{$1$}. }
\item \textcolor{violet}{האיבר המוביל של כל שורה נמצא מימין לאיבר המוביל
של השורה שמעליו. }
\item \textcolor{violet}{בעמודה שבה יש איבר מוביל )\L{$1$}(, כל שאר האיברים
הם אפסים )גם מעל וגם מתחת(. }
\end{enumerate}

\subsection*{\textcolor{violet}{תיאור האלגוריתם}}

\textcolor{violet}{תהי \L{$A$} מטריצה מסדר \L{$m\times n$}. האלגוריתם
עובד באופן איטרטיבי על עמודות ושורות המטריצה משמאל לימין: }
\begin{enumerate}
\item \textcolor{violet}{מצא את העמודה הראשונה שאינה מכילה רק אפסים )זוהי
עמודת הציר הנוכחית(. }
\item \textcolor{violet}{אם האיבר העליון בעמודה זו הוא \L{$0$}, בצע החלפת
שורות כדי להביא איבר שונה מאפס לראש העמודה. }
\item \textcolor{violet}{בצע פעולת כפל בסקלר כדי להפוך את האיבר המוביל ל-\L{$1$}. }
\item \textcolor{violet}{השתמש באיבר המוביל )\L{$1$}( כדי לאפס את כל האיברים
האחרים באותה עמודה )מעליו ומתחתיו( על ידי פעולת הוספת שורות. }
\item \textcolor{violet}{התעלם מהשורה הנוכחית ומהעמודה הנוכחית, וחזור על
התהליך עבור תת-המטריצה הנותרת. }
\end{enumerate}

\subsection*{\textcolor{violet}{דוגמה מספרית}}

\textcolor{violet}{נתונה המטריצה \L{$A$}: \L{%
\[
A=\begin{pmatrix}0 & 2 & 4\\
1 & 3 & 3\\
3 & 7 & 12
\end{pmatrix}
\]
}%
}\textbf{\textcolor{violet}{שלב 1: החלפת שורות ליצירת איבר מוביל ראשון.}}\textcolor{violet}{{}
נחליף את \L{$R_{1}$} עם \L{$R_{2}$} כדי לקבל \L{$1$} בפינה השמאלית
העליונה: \L{%
\[
\xrightarrow{R_{1}\leftrightarrow R_{2}}\begin{pmatrix}1 & 3 & 3\\
0 & 2 & 4\\
3 & 7 & 12
\end{pmatrix}
\]
}%
}\textbf{\textcolor{violet}{שלב 2: איפוס מתחת לאיבר המוביל הראשון.}}\textcolor{violet}{{}
נאפס את ה-\L{$3$} בשורה השלישית על ידי הפעולה \L{$R_{3}\to R_{3}-3R_{1}$}:
\L{%
\[
\xrightarrow{R_{3}\to R_{3}-3R_{1}}\begin{pmatrix}1 & 3 & 3\\
0 & 2 & 4\\
0 & -2 & 3
\end{pmatrix}
\]
}%
}\textbf{\textcolor{violet}{שלב 3: טיפול באיבר המוביל השני.}}\textcolor{violet}{{}
במטריצה שנותרה )שורות 2-3(, האיבר המוביל הוא \L{$2$}.}

\textcolor{violet}{ננרמל אותו ל-\L{$1$} על ידי חלוקה ב-2 )\L{$R_{2}\to\frac{1}{2}R_{2}$}(:
\L{%
\[
\xrightarrow{R_{2}\to\frac{1}{2}R_{2}}\begin{pmatrix}1 & 3 & 3\\
0 & 1 & 2\\
0 & -2 & 3
\end{pmatrix}
\]
}%
}\textbf{\textcolor{violet}{שלב 4: איפוס העמודה השנייה )מעל ומתחת(.}}

\textcolor{violet}{נאפס את ה-\L{$-2$} בשורה השלישית )\L{$R_{3}\to R_{3}+2R_{2}$}(
ואת ה-\L{$3$} בשורה הראשונה )\L{$R_{1}\to R_{1}-3R_{2}$}(: \L{%
\[
\xrightarrow{\begin{subarray}{l}
R_{3}\to R_{3}+2R_{2}\\
R_{1}\to R_{1}-3R_{2}
\end{subarray}}\begin{pmatrix}1 & 0 & -3\\
0 & 1 & 2\\
0 & 0 & 7
\end{pmatrix}
\]
}%
}\textbf{\textcolor{violet}{שלב 5: טיפול באיבר המוביל השלישי וסיום.}}\textcolor{violet}{{}
ננרמל את האיבר המוביל בשורה השלישית ל-\L{$1$} )\L{$R_{3}\to\frac{1}{7}R_{3}$}(:
\L{%
\[
\xrightarrow{R_{3}\to\frac{1}{7}R_{3}}\begin{pmatrix}1 & 0 & -3\\
0 & 1 & 2\\
0 & 0 & 1
\end{pmatrix}
\]
}%
כעת נאפס את האיברים מעל האיבר המוביל השלישי: \L{%
\[
\xrightarrow{\begin{subarray}{l}
R_{2}\to R_{2}-2R_{3}\\
R_{1}\to R_{1}+3R_{3}
\end{subarray}}\begin{pmatrix}1 & 0 & 0\\
0 & 1 & 0\\
0 & 0 & 1
\end{pmatrix}
\]
}%
התקבלה מטריצת היחידה, שהיא הצורה המדורגת קנונית של \L{$A$}.}

\rule{1\linewidth}{1.5pt}

\section{ממ\textquotedblright ל הומוגנית}

\textbf{{]}סוג חשוב של ממ\textquotedblright ל:{[}}

ממ\textquotedblright ל עם מטריצה מורחבת

\L{%
\[
\left(\begin{array}{ccc|c}
a_{11} & \dots & a_{1n} & b_{1}\\
\vdots &  & \vdots & \vdots\\
\vdots &  & \vdots & \vdots\\
a_{m1} & \dots & a_{mn} & b_{m}
\end{array}\right)
\]
}%

נקראת \underline{הומוגנית} כאשר \L{%
\[
b_{1}=\dots=b_{m}=0
\]
}%


\section{תכונה {]}הפתרון הטריוויאלי של ממ\textquotedblright ל הומוגנית{[}}

נתונה ממ\textquotedblright ל הומוגנית ב-\L{$n$} נעלמים. אז \L{$\begin{pmatrix}0\\
\vdots\\
0
\end{pmatrix}\in\mathbb{R}^{n}$} היא תמיד פתרון של הממ\textquotedblright ל.

הפתרון הזה נקרא \textbf{הפתרון הטריוויאלי}.

\textcolor{blue}{הערה: הפתרון הטריוויאלי יכול להיות יחיד, ויכול להיות
לא יחיד )כלומר יכולים להיות עוד פתרונות לממ\textquotedblright ל הומוגנית(.
אבל בכל מקרה תמיד יש לפחות פתרון אחד שהוא הפתרון הטריוויאלי.}

\newpage{}

\noindent\fbox{\begin{minipage}[t]{1\columnwidth - 2\fboxsep - 2\fboxrule}%
\begin{center}
אלגברה ליניארית )1( - \L{$80134$} \textbar{} האונ' העברית \textbar{}
תשפ\textquotedblright ו א' \textbar{} מרצה: ד\textquotedblright ר
אלכס גורביץ' \textbar{} מתרגל: מר יואב כהן \textbar{} ליקוט ועריכה:
נעם לב
\par\end{center}%
\end{minipage}}

\part{הרצאה {\beginL 4\endL} - {\beginL 28.10.2025\endL}}

\section{תזכורת}

ממ\textquotedblright ל עם מטריצה מורחבת \L{$\left(\begin{array}{ccc|c}
a_{11} & \dots & a_{1n} & b_{1}\\
\vdots &  & \vdots & \vdots\\
a_{m1} & \dots & a_{mn} & b_{m}
\end{array}\right)$} נקראת הומוגנית כאשר \L{$b_{1}=\dots=b_{m}=0$}.

לממ\textquotedblright ל הומוגנית ב-\L{$n$} נעלמים, \L{$\begin{pmatrix}0\\
\vdots\\
0
\end{pmatrix}\in\mathbb{R}^{n}$} הוא תמיד פתרון )הפתרון הטריוויאלי(.

\rule{1\linewidth}{1.5pt}

\textcolor{blue}{ננסה לסווג מתי הפתרון הטריוויאלי זה הפתרון היחיד
של ממ\textquotedblright ל הומוגנית ומתי לא.}

\section{תכונה )חשובה( {]}שמירת ההומוגניות תחת שקילות שורה{[}}

אם מטריצות מורחבות של שתי ממ\textquotedblright ל שקולות שורה, ואחת
מהממ\textquotedblright ל האלה היא הומוגנית, אז גם הממ\textquotedblright ל
השניה הומוגנית.

\section{טענה {]}אפיון פתרון יחיד )טריוויאלי( לממ\textquotedblright ל הומוגנית
- איבר מוביל בכל עמודה שאינה ההרחבה{[}}

נתונה ממ\textquotedblright ל הומוגנית עם מטריצה מורחבת מדורגת מצומצמת.

אם בכל אחת מהעמודות של המטריצה המורחבת )חוץ מהעמודה הימנית( קיים איבר
מוביל, אז כל אחד מהמשתנים הוא קשור, והפתרון הטריוויאלי הוא הפתרון
היחיד של הממ\textquotedblright ל.

דוגמה: \L{%
\[
\left(\begin{array}{cc|c}
1 & 0 & 0\\
0 & 1 & 0\\
0 & 0 & 0
\end{array}\right)
\]
}%

\L{%
\[
\begin{cases}
x_{1}=0\\
x_{2}=0\\
0=0
\end{cases}
\]
}%

אחרת, כלומר אם קיימת עמודה של המטריצה המורחבת )שאינה העמודה הימנית(
ובה אין איבר מוביל, אז הנעלם המתאים לעמודה הוא נעלם חופשי, ולכן למערכת
יהיה גם פתרון לא טריוויאלי.

\subsection{מסקנה )חשובה( {]}בממ\textquotedblright ל הומוגנית עם יותר נעלמים
ממשוואות, יש פתרון לא-טריוויאלי{[}}

נתונה ממ\textquotedblright ל הומוגנית שבה מספר המשוואות קטן ממספר
הנעלמים, אז לממ\textquotedblright ל הזאת יהיה פתרון לא טריוויאלי )ולכן
אינסוף פתרונות(.

דוגמה: \L{%
\[
\left(\begin{array}{ccc|c}
1 & 2 & 3 & 0\\
4 & 5 & 6 & 0
\end{array}\right)
\]
}%

\textcolor{blue}{עד כאן הדיון על ממ\textquotedblright ל הומוגניות.
נלמד כעת על דרך נוספת לכתוב ממ\textquotedblright ל.}

\section{\textcolor{violet}{תכונה {]}שלוש האפשרויות לכמות הפתרונות{[}}}

\textcolor{violet}{באלגברה ליניארית )מעל שדות אינסופיים כמו המספרים
הממשיים \L{$\mathbb{R}$}(, למערכת משוואות ליניאריות יש בדיוק שלוש
אפשרויות:}
\begin{enumerate}
\item \textcolor{violet}{אין פתרון.}
\item \textcolor{violet}{פתרון יחיד.}
\item \textcolor{violet}{אינסוף פתרונות.}
\end{enumerate}


\section{צירוף לינארי}

תהיינה \L{$v_{1},\dots,v_{n}\in\mathbb{R}^{m}$}, \L{$a_{1},\dots,a_{n}\in\mathbb{R}$}.

\textbf{צירוף לינארי} של \L{$v_{1},\dots,v_{n}$} עם מקדמים \L{$a_{1},\dots,a_{n}$}
הוא \L{$a_{1}v_{1}+\dots+a_{n}v_{n}\in\mathbb{R}^{m}$}.

\section{כפל של מטריצה ב-\L{$n$}-יה}

תהיינה \L{$v_{1},\dots,v_{n}\in\mathbb{R}^{m}$}, \L{$a_{1},\dots,a_{n}\in\mathbb{R}$}.

תהי \L{$A$} מטריצה בגודל \L{$m\times n$} )כלומר \L{$m$} שורות ו-
\L{$n$} עמודות(. יהי \L{$d=\begin{pmatrix}d_{1}\\
\vdots\\
d_{n}
\end{pmatrix}\in\mathbb{R}^{n}$} )מוגדר רק כאשר זה אותו \L{$n$}(.

נסמן את העמודות של \L{$A$} ב- \L{$C_{1}^{A},\dots,C_{n}^{A}$} )מלשון
\L{$\text{Column}$}( )נשים לב שכל אחת מהן היא \L{$m$}-יה סדורה(.\\

נגדיר \L{$Ad=d_{1}C_{1}^{A}+\dots+d_{n}C_{n}^{A}\in\mathbb{R}^{m}$}.\\

אם \L{$A=\begin{pmatrix}a_{11} & \dots & a_{1n}\\
\vdots &  & \vdots\\
a_{m1} & \dots & a_{mn}
\end{pmatrix}$} אז

\L{%
\[
Ad=\begin{pmatrix}d_{1}a_{11}+\dots+d_{n}a_{1n}\\
\vdots\\
d_{1}a_{m1}+\dots+d_{n}a_{mn}
\end{pmatrix}\in\mathbb{R}^{m}
\]
}%

\medskip{}

תהי \L{$b=\begin{pmatrix}b_{1}\\
\vdots\\
b_{m}
\end{pmatrix}\in\mathbb{R}^{m}$}. כדי ש- \L{$d=\begin{pmatrix}d_{1}\\
\vdots\\
d_{n}
\end{pmatrix}\in\mathbb{R}^{n}$} תהיה פתרון של ממ\textquotedblright ל עם מטריצה מורחבת \L{$(A|b)$}
צריך להתקיים

\L{%
\[
\begin{cases}
a_{11}d_{1}+\dots+a_{1n}d_{n}=b_{1}\\
\vdots\\
a_{m1}d_{1}+\dots+a_{mn}d_{n}=b_{m}
\end{cases}
\]
}%

\textbf{כלומר \L{$Ad=b$}.}

לכן ניתן לכתוב את הממ\textquotedblright ל עם המטריצה המורחבת \L{$(A|b)$}
באופן הבא: \L{$Ax=b$} כאשר \L{$x=\begin{pmatrix}x_{1}\\
\vdots\\
x_{n}
\end{pmatrix}$} הוא ה-\L{$n$}-ית הנעלמים.

\textcolor{teal}{כפל של מטריצה ב-\L{$n$}-יה זה כמו אקורדיון.}

\subsection{תכונות של כפל מטריצה ב-\L{$n$}-יה}

1( לכל מטריצה בגודל \L{$m\times n$}: \L{%
\[
A\begin{pmatrix}0\\
\vdots\\
0
\end{pmatrix}=\begin{pmatrix}0\\
\vdots\\
0
\end{pmatrix}
\]
}%
)נשים לב: ה-\L{$\vec{0}$} השמאלי ב-\L{$\mathbb{R}^{n}$}, הימני ב-\L{$\mathbb{R}^{m}$}(.\\

2( לכל מטריצה בגודל \L{$m\times n$} ולכל \L{$d,h\in\mathbb{R}^{n}$}:
\L{%
\[
A(d+h)=Ad+Ah
\]
}%

3( לכל מטריצה בגודל \L{$m\times n$} ולכל \L{$c\in\mathbb{R},d\in\mathbb{R}^{n}$}:
\L{%
\[
A(cd)=c(Ad)
\]
}%

4( לכל \L{$d\in\mathbb{R}^{n}$}, אם \L{$A$} היא מטריצה בגודל \L{$m\times n$}
כך שהשורה ה-\L{$i$} של \L{$A$} היא שורת אפסים אז הקואורדינאטה ה-\L{$i$}
של \L{$Ad$} היא 0.

5( נסמן \L{$e_{i}=\begin{pmatrix}0\\
\vdots\\
1\\
\vdots\\
0
\end{pmatrix}\in\mathbb{R}^{n}$} )ה-1 במקום ה-\L{$i$}(. לדוגמה: \L{$e_{2}=\begin{pmatrix}0\\
1\\
0
\end{pmatrix}\in\mathbb{R}^{3}$}. \textcolor{violet}{זה וקטור היחידה הסטנדרטי.}

\L{$Ae_{i}=C_{i}^{A}$}. כלומר מטריצה \L{$A$} כפול \L{$e_{i}$} שווה
לעמודה ה-\L{$i$} של \L{$A$}.

6( נסמן ב-\L{$I_{n}$} את המטריצה בגודל \L{$n\times n$} כך שהעמודה
ה-\L{$i$} שלה היא \L{$e_{i}$}. לדוגמה \L{$n=3$}: \L{%
\[
I_{3}=\begin{pmatrix}1 & 0 & 0\\
0 & 1 & 0\\
0 & 0 & 1
\end{pmatrix}
\]
}%
\L{$I_{n}$} נקראת \textbf{מטריצת היחידה}.

\textcolor{violet}{\L{%
\[
I_{n}=\Big[e_{1}\ \big|\ e_{2}\ \big|\ \dots\ \big|\ e_{n}\Big]
\]
}%
}

לכל \L{$d=\begin{pmatrix}d_{1}\\
\vdots\\
d_{n}
\end{pmatrix}\in\mathbb{R}^{n}$}:

\L{%
\[
I_{n}d=d_{1}e_{1}+\dots+d_{n}e_{n}=\begin{pmatrix}d_{1}\\
0\\
\vdots\\
0
\end{pmatrix}+\begin{pmatrix}0\\
d_{2}\\
\vdots\\
0
\end{pmatrix}+\dots+\begin{pmatrix}0\\
\vdots\\
0\\
d_{n}
\end{pmatrix}=\begin{pmatrix}d_{1}\\
d_{2}\\
\vdots\\
d_{n}
\end{pmatrix}=d
\]
}%

\bigskip{}

\textcolor{blue}{וסיימנו את הפרק הראשון של הקורס.}

\textcolor{blue}{הקדמה לפרק השני: נגדיר מושגים יותר מופשטים בצורה
פחות מפורשת בעזרת התכונות שלהם )בעזרת אקסיומות(. נלמד על שדות.}

\newpage{}

\noindent\fbox{\begin{minipage}[t]{1\columnwidth - 2\fboxsep - 2\fboxrule}%
\begin{center}
אלגברה ליניארית )1( - \L{$80134$} \textbar{} האונ' העברית \textbar{}
תשפ\textquotedblright ו א' \textbar{} מרצה: ד\textquotedblright ר
אלכס גורביץ' \textbar{} מתרגל: מר יואב כהן \textbar{} ליקוט ועריכה:
נעם לב
\par\end{center}%
\end{minipage}}

\part{הרצאה {\beginL 5\endL} - {\beginL 03.11.2025\endL}}

\section{שדות}

קבוצה \L{$\mathbb{F}$} נקראת שדה כאשר ב-\L{$\mathbb{F}$} מוגדרות
פעולות חיבור )\L{$+$}( וכפל )\L{$\cdot$}( אשר לכל זוג איברים מ-\L{$\mathbb{F}$}
מתאימות איבר מ-\L{$\mathbb{F}$} כך שמתקיימות האקסיומות הבאות:\\

1( \L{%
\[
\forall a,b\in\mathbb{F},\ a+b=b+a
\]
}%
 

2( \L{%
\[
\forall a,b,c\in\mathbb{F},\ (a+b)+c=a+(b+c)
\]
}%

3( \L{%
\[
\exists0_{\mathbb{F}}\in\mathbb{F},\forall a\in\mathbb{F},\ a+0_{\mathbb{F}}=a
\]
}%
 

4( \L{%
\[
\forall a\in\mathbb{F},\exists-a\in\mathbb{F},\ a+(-a)=0_{\mathbb{F}}
\]
}%
 

5( \L{%
\[
\forall a,b\in\mathbb{F},\ ab=ba
\]
}%
 

6( \L{%
\[
\forall a,b,c\in\mathbb{F},\ (ab)c=a(bc)
\]
}%
 

7( \L{%
\[
\exists1_{\mathbb{F}}\in\mathbb{F},\ 1_{\mathbb{F}}\neq0_{\mathbb{F}},\forall a\in\mathbb{F},\ a\cdot1_{\mathbb{F}}=a
\]
}%
 

8( \L{%
\[
\forall_{0\neq a\in\mathbb{F}},\exists_{\frac{1}{a}\in\mathbb{F}},\ a\cdot(\frac{1}{a})=1_{\mathbb{F}}
\]
}%
 

9( \L{%
\[
\forall_{a,b,c\in\mathbb{F}},\ a\cdot(b+c)=a\cdot b+a\cdot c
\]
}%


\subsection{דוגמאות}
\begin{enumerate}
\item \L{$\mathbb{R}$} שדה של מספרים ממשיים עם חיבור וכפל.
\item \L{$\mathbb{C}$} שדה של מספרים מרוכבים עם חיבור וכפל.
\item \L{$\mathbb{Q}$} שדה של מספרים רציונליים עם חיבור וכפל. \L{$\mathbb{Q}=\left\{ \frac{a}{b}\,\middle|\,a\in\mathbb{Z},b\in\mathbb{N}\right\} $}
\item \L{$\mathbb{N},\mathbb{Z}$} \emph{לא} מהווים שדה עם חיבור וכפל.
\item \L{$\mathbb{F}_{2}=\{0,1\}$} שדה עם חיבור וכפל ב-\L{$\mathbb{F}_{2}$}
שנגדיר באופן הבא:
\end{enumerate}
\begin{center}
\begin{tabular}{c|c|c}
\R{\L{$+$} } & \R{0 } & \R{1 }\tabularnewline
\hline 
\R{0 } & \R{0 } & \R{1 }\tabularnewline
\hline 
\R{1 } & \R{1 } & \R{0 }\tabularnewline
\end{tabular}\quad{}\quad{}%
\begin{tabular}{c|c|c}
\R{\L{$\cdot$} } & \R{0 } & \R{1 }\tabularnewline
\hline 
\R{0 } & \R{0 } & \R{0 }\tabularnewline
\hline 
\R{1 } & \R{0 } & \R{1 }\tabularnewline
\end{tabular}
\par\end{center}

\subsection{תכונות של שדות}

)תכונות שנובעות מהאקסיומות(\\

1( \underline{יחידות האיבר הנייטרלי ביחס לחיבור}

אם \L{$z\in\mathbb{F}$} כך ש- \L{$\forall a\in\mathbb{F},a+z=a$}
אז \L{$z=0_{\mathbb{F}}$} )יש נייטרלי יחיד ביחס לחיבור(.

\underline{הוכחה:} \L{%
\[
z\underset{\text{axiom 3}}{=}z+0_{\mathbb{F}}\underset{\text{axiom 1}}{=}0_{\mathbb{F}}+z\underset{\text{given}}{=}0_{\mathbb{F}}
\]
}%

2( \underline{יחידות האיבר הנגדי}

אם \L{$a,b\in\mathbb{F}$} כך ש- \L{$a+b=0_{\mathbb{F}}$}, אז \L{$b=-a$}.

\underline{הוכחה:}

\L{%
\[
-a\underset{\text{axiom 3}}{=}(-a)+0_{\mathbb{F}}\underset{\text{axiom 1}}{=}0_{\mathbb{F}}+(-a)\underset{\text{given}}{=}(a+b)+(-a)
\]
}%

\L{%
\[
\underset{\text{axiom 1}}{=}(b+a)+(-a)\underset{\text{axiom 2}}{=}b+(a+(-a))\underset{\text{axiom 4}}{=}b+0_{\mathbb{F}}\underset{\text{axiom 3}}{=}b
\]
}%

ובסה\textquotedblright כ קיבלנו \L{$b=-a$}.\\

3( אם \L{$a\in\mathbb{F}$} כך ש- \L{$a+a=a$}, אז \L{$a=0_{\mathbb{F}}$}.

\underline{הוכחה:}

\L{%
\[
0_{\mathbb{F}}\underset{\text{axiom 4}}{=}a+(-a)\underset{\text{given}}{=}(a+a)+(-a)\underset{\text{axiom 2}}{=}a+(a+(-a))\underset{\text{axiom 4}}{=}a+0_{\mathbb{F}}\underset{\text{axiom 3}}{=}a
\]
}%

ובסה\textquotedblright כ קיבלנו \L{$0_{\mathbb{F}}=a$} ולכן לפי אקסיומה
{\beginL 1\endL}, \L{$a=0_{\mathbb{F}}$}.\\

4( \L{$\forall a\in\mathbb{F},\ a\cdot0_{\mathbb{F}}=0_{\mathbb{F}}$}

\underline{הוכחה:}

\L{%
\[
a\cdot0_{\mathbb{F}}\underset{\text{axiom 3}}{=}a\cdot(0_{\mathbb{F}}+0_{\mathbb{F}})\underset{\text{axiom 9}}{=}a\cdot0_{\mathbb{F}}+a\cdot0_{\mathbb{F}}
\]
}%

קיבלנו שלאיבר \L{$x=a\cdot0_{\mathbb{F}}$} מתקיים \L{$x=x+x$}. לפי
תכונה 3 שהוכחנו לעיל, מתקיים:

\L{%
\[
a\cdot0_{\mathbb{F}}=0_{\mathbb{F}}
\]
}%


\section{מרחבים וקטוריים}

\underline{הבחנה:} כל מה שלמדנו על ממ\textquotedblright ל תקף
כאשר במקום מספרים ממשיים משתמשים באיברים משדה כלשהו.\\

בפרט, בהינתן שדה \L{$\mathbb{F}$} ניתן להגדיר

\L{%
\[
\mathbb{F}^{n}=\left\{ \begin{pmatrix}x_{1}\\
\vdots\\
x_{n}
\end{pmatrix}\,\middle|\,x_{1},\dots,x_{n}\in\mathbb{F}\right\} 
\]
}%

בדומה ל-\L{$\mathbb{R}^{n}$}, ב-\L{$\mathbb{F}^{n}$} מוגדרות פעולות
של חיבור וכפל בסקלר מ-\L{$\mathbb{F}$}.

\subsection{הגדרה - מרחב וקטורי )מ\textquotedblright ו(}

יהי \L{$\mathbb{F}$} שדה.

קבוצה \L{$V$} נקראת מרחב וקטורי מעל \L{$\mathbb{F}$} כאשר ב-\L{$V$}
מוגדרות:
\begin{enumerate}
\item פעולת חיבור אשר לכל זוג איברים מ-\L{$V$} מתאימה איבר כלשהו מ-\L{$V$}.
\item פעולת כפל בסקלר אשר לכל איבר מ-\L{$V$} ולכל סקלר מ-\L{$\mathbb{F}$}
מתאימה איבר כלשהו מ-\L{$V$}.
\end{enumerate}
\textbf{כך שמתקיימות האקסיומות הבאות:}\\

\begin{enumerate}
\item \L{%
\[
\forall_{u,v\in V}\quad u+v=v+u
\]
}%
\item \L{%
\[
\forall_{u,v,w\in V}\quad(u+v)+w=u+(v+w)
\]
}%
\item \L{%
\[
\exists_{0_{V}\in V}\quad\forall_{v\in V}\quad v+0_{V}=v
\]
}%
 
\item \L{%
\[
\forall_{v\in V}\quad\exists_{(-v)\in V}\quad v+(-v)=0_{V}
\]
}%
\item \L{%
\[
\exists_{1_{\mathbb{F}}\in\mathbb{F}}:\forall_{v\in V},\quad1_{\mathbb{F}}\cdot v=v
\]
}%
\item \L{%
\[
\forall_{v\in V},\forall_{a,b\in\mathbb{F}},\ (ab)v=a(bv)
\]
}%
\item \L{%
\[
\forall_{v\in V}\quad\forall_{a,b\in\mathbb{F}}\quad(a+b)v=av+bv
\]
}%
\item \L{%
\[
\forall_{u,v\in V}\quad\forall_{a\in\mathbb{F}}\quad a(u+v)=au+av
\]
}%
\end{enumerate}
לאיברים של \L{$\mathbb{F}$} נקרא \emph{סקלרים} ולאיברים של \L{$V$}
נקרא \emph{וקטורים}.

\newpage{}

\noindent\fbox{\begin{minipage}[t]{1\columnwidth - 2\fboxsep - 2\fboxrule}%
\begin{center}
אלגברה ליניארית )1( - \L{$80134$} \textbar{} האונ' העברית \textbar{}
תשפ\textquotedblright ו א' \textbar{} מרצה: ד\textquotedblright ר
אלכס גורביץ' \textbar{} מתרגל: מר יואב כהן \textbar{} ליקוט ועריכה:
נעם לב
\par\end{center}%
\end{minipage}}

\part{הרצאה {\beginL 6\endL} - {\beginL 04.11.2025\endL}}

\section{דוגמאות למרחבים וקטוריים}

יהי \L{$\mathbb{F}$} שדה.\\

1( המרחב הוקטורי של קבוצת ה-\L{$n$}-יות משדה \L{$\mathbb{F}$}. זהו
המרחב הוקטורי שמעליו פותרים ממ\textquotedblright ל.

\L{%
\[
\mathbb{F}^{n}=\left\{ \begin{pmatrix}x_{1}\\
\vdots\\
x_{n}
\end{pmatrix}\,\middle|\,x_{1},\dots,x_{n}\in\mathbb{F}\right\} 
\]
}%

2( המרחב הוקטורי של המטריצות בעלות \L{$m$} שורות ו-\L{$n$} עמודות,
שאיבריהן נמצאים ב-\L{$\mathbb{F}$}. זהו המרחב הוקטורי המגדיר פעולות
על מטריצות.

עבור \L{$m,n\in\mathbb{N}$}

\L{%
\[
M_{m\times n}(\mathbb{F})=\left\{ \begin{pmatrix}a_{11} & \dots & a_{1n}\\
\vdots &  & \vdots\\
a_{m1} & \dots & a_{mn}
\end{pmatrix}\,\middle|\,\forall_{1\le i\le m}\quad\forall_{1\le j\le n}\quad a_{ij}\in\mathbb{F}\right\} 
\]
}%

פעולות ב-\L{$M_{m\times n}(\mathbb{F})$} מוגדרות ע\textquotedblright י:

\L{%
\[
\begin{pmatrix}a_{11} & \dots & a_{1n}\\
\vdots &  & \vdots\\
a_{m1} & \dots & a_{mn}
\end{pmatrix}+\begin{pmatrix}b_{11} & \dots & b_{1n}\\
\vdots &  & \vdots\\
b_{m1} & \dots & b_{mn}
\end{pmatrix}=
\]
}%
\L{%
\[
\begin{pmatrix}a_{11}+b_{11} & \dots & a_{1n}+b_{1n}\\
\vdots &  & \vdots\\
a_{m1}+b_{m1} & \dots & a_{mn}+b_{mn}
\end{pmatrix}
\]
}%

לכל \L{$c\in\mathbb{F}$} \L{%
\[
c\begin{pmatrix}a_{11} & \dots & a_{1n}\\
\vdots &  & \vdots\\
a_{m1} & \dots & a_{mn}
\end{pmatrix}=\begin{pmatrix}ca_{11} & \dots & ca_{1n}\\
\vdots &  & \vdots\\
ca_{m1} & \dots & ca_{mn}
\end{pmatrix}
\]
}%
\\

\textcolor{blue}{בפועל הדוגמה הראשונה היא מקרה פרטי של הדוגמה השניה
כי \L{$n$}-יה זו מטריצה עם עמודה 1 ו-\L{$n$} שורות.}\\

3. \textcolor{blue}{המרחב הוקטורי של קבוצת הסדרות האינסופיות משדה
\L{$\mathbb{F}$}. זהו המ\textquotedblright ו של ה-\L{$n$}-יות עבור
\L{$n=\infty$}.}

\L{%
\[
\mathbb{F}^{\mathbb{N}}=\left\{ (x_{1},x_{2},\dots)\mid\forall_{i\in\mathbb{N}}\ (x_{i}\in\mathbb{F})\right\} 
\]
}%
פעולות מוגדרות ע\textquotedblright י:

\L{%
\[
(x_{1},x_{2},\dots)+(y_{1},y_{2},\dots)=(x_{1}+y_{1},x_{2}+y_{2},\dots)
\]
}%

\L{%
\[
c(x_{1},x_{2},\dots)=(cx_{1},cx_{2},\dots)
\]
}%

4( \textcolor{blue}{המרחב הוקטורי של כל הפולינומים מעל שדה \L{$\mathbb{F}$}
והמ\textquotedblright ו של כל הפולינומים ממעלה קטנה או שווה ל-\L{$n$}
מעל שדה \L{$\mathbb{F}$}.}

בהינתן פולינום \L{$P=a_{n}x^{n}+...+a_{1}x+a_{0}$} עם \L{$a_{n}\neq0_{\mathbb{F}}$},

\L{$a_{n},\dots,a_{0}$} נקראים המקדמים של הפולינום ו-\L{$n$} נקרא
המעלה של הפולינום. רושמים \L{$\deg P=n$}. בנוסף מגדירים \L{$\deg0_{\mathbb{F}}=-\infty$}.

\begin{comment}
להכניס כאן אינטואיציה מדוע זה מינוס אינסוף שקיבלתי אי שם בקנדה עליון
בתחילת השנה
\end{comment}

\textbf{המ\textquotedblright ו של כל הפולינומים:}

\L{%
\[
\mathbb{F}[x]=\left\{ a_{n}x^{n}+...+a_{1}x+a_{0}\,\middle|\,n\in\mathbb{N}\cup\{0\},\ a_{0},\dots,a_{n}\in\mathbb{F},a_{n}\ne0_{\mathbb{F}}\right\} \cup\{0_{\mathbb{F}}\}
\]
}%

\textbf{\textcolor{black}{המ\textquotedblright ו של כל הפולינומים
ממעלה קטנה או שווה ל-\L{$n$}:}}

\L{%
\[
\mathbb{F}[x]_{\le n}=\left\{ a_{n}x^{n}+\dots+a_{1}x+a_{0}\,\middle|\,a_{0},a_{1},\dots,a_{n}\in\mathbb{F}\right\} 
\]
}%

\L{$\mathbb{F}[x]_{\le n}$} הוא תמ\textquotedblright ו של \L{$\mathbb{F}[x]$},
נלמד בהמשך מה זה אומר. \\

\uline{פעולות ב-\L{\mbox{$\mathbb{F}[x]$}} מוגדרות באופן הבא:}

\textbf{חיבור של פולינומים:}

\L{%
\[
(a_{n}x^{n}+\dots+a_{1}x+a_{0})+(b_{m}x^{m}+\dots+b_{1}x+b_{0})
\]
}%
\L{%
\[
=(a_{k}+b_{k})x^{k}+\dots+(a_{1}+b_{1})x+(a_{0}+b_{0})
\]
}%

כאשר \L{$k=\max(m,n)$}. עבור \L{$i>n$} מניחים \L{$a_{i}=0$}. עבור
\L{$j>m$} מניחים \L{$b_{j}=0$}.

הערה: במידה ויצא \L{$a_{k}+b_{k}=0$} מוחקים את המחובר \L{$(a_{k}+b_{k})x^{k}$}
וכן הלאה, כלומר מעלת הפולינום עשויה לרדת.

דוגמאות לחיבור פולינומים: \L{%
\[
(2x^{2}+3x+4)+(5x+6)=2x^{2}+8x+10
\]
}%
\L{%
\[
(3x+4)+(-3x+5)=9
\]
}%

\textbf{כפל בסקלר של פולינומים:}

\L{%
\[
c(a_{n}x^{n}+\dots+a_{1}x+a_{0})=
\]
}%
\L{%
\[
(ca_{n}x^{n}+\dots+ca_{1}x+ca_{0})
\]
}%

\L{%
\[
0(a_{n}x^{n}+\dots+a_{1}x+a_{0})=0
\]
}%

\noindent\fbox{\begin{minipage}[t]{1\columnwidth - 2\fboxsep - 2\fboxrule}%
\textcolor{teal}{נשים לב: קבוצת הפולינומים שהם רק מדרגה \L{$n$} היא
בד\textquotedblright כ אינה מ\textquotedblright ו, שכן היא לא סגורה
לחיבור. }

\textcolor{teal}{למשל, האם קבוצת הפולינומים הממשיים מדרגה 2,}\\
\textcolor{teal}{\L{$\mathbb{R}[2]=\left\{ a_{2}x^{2}+a_{1}x+a_{0}|a_{2},a_{1},a_{0}\in\mathbb{R},a_{2}\neq0_{\mathbb{\mathbb{R}}}\right\} \cup\{0_{\mathbb{\mathbb{R}}}\}$},
היא מ\textquotedblleft ו?}\\
\textcolor{teal}{לא, כי היא לא סגורה לחיבור. למשל, \L{$x^{2}\in\mathbb{R}[2]\wedge-x^{2}+5x\in\mathbb{R}[2]$}
אבל}\\
\textcolor{teal}{\L{$x^{2}+(-x^{2}+5x)=5x\notin\mathbb{R}[2]$}. אבל
אפשר לתקן את זה:}

\textcolor{teal}{קבוצת הפולינומים הממשיים עד דרגה 2 היא כן מ\textquotedblright ו:}\\
\textcolor{teal}{\L{$P_{2}(\mathbb{R})=\left\{ a_{2}x^{2}+a_{1}x+a_{0}|a_{0},a_{1},a_{2}\in\mathbb{R}\right\} $}.}\\
\textcolor{teal}{}\\
\textcolor{teal}{ההגדרה של קבוצת הפולינומים שמעלתם היא בדיוק \L{$n$}:}

\textcolor{teal}{\L{%
\[
\mathbb{F}[x]_{=n}=\left\{ a_{n}x^{n}+\dots+a_{1}x+a_{0}\,\middle|\,a_{n},\dots,a_{0}\in\mathbb{F},a_{n}\neq0_{\mathbb{F}}\right\} 
\]
}%
}%
\end{minipage}}

\medskip{}

5. הדוגמה הכי חשובה - המרחב הוקטורי של הפונקציות מקבוצה לשדה.

תהי \L{$S$} קבוצה כלשהי.פונקציה מ-\L{$S$} לשדה \L{$\mathbb{F}$}
היא התאמה שלכל איבר של \L{$S$} מתאימה איבר )יחיד( מ-\L{$\mathbb{F}$}.

כאשר \L{$f$} היא פונקציה מ-\L{$S$} ל-\L{$\mathbb{F}$} רושמים \L{$f:S\to\mathbb{F}$}.

נסתכל על המרחב הוקטורי: \L{%
\[
\mathbb{F}^{S}=\{f:S\to\mathbb{F}\mid f\text{ is a function}\}
\]
}%

פעולות ב-\L{$\mathbb{F}^{S}$} מוגדרות באופן הבא:

אם \L{$f,g:S\to\mathbb{F}$} נגדיר \L{$f+g:S\to\mathbb{F}$} ע\textquotedblright י
\L{%
\[
\forall s\in S,\ (f+g)(s)=f(s)+g(s)
\]
}%

אם \L{$f:S\to\mathbb{F},\ c\in\mathbb{F}$} נגדיר \L{$c\cdot f:S\to\mathbb{F}$}
ע\textquotedblright י \L{%
\[
\forall s\in S,\ (cf)(s)=c\cdot f(s)
\]
}%


\section{תכונות של מרחבים וקטוריים}

יהי \L{$\mathbb{F}$} שדה. יהי \L{$V$} מרחב וקטורי )מ\textquotedblright ו(
מעל \L{$\mathbb{F}$}.
\begin{enumerate}
\item אם \L{$w\in V$} כך ש- \L{$\forall u\in V,u+w=u$} אז \L{$w=0_{V}$}.
\item אם \L{$v,w\in V$} כך ש- \L{$v+w=0_{V}$}, אז \L{$w=-v$}.
\item אם \L{$v\in V$} כך ש- \L{$v+v=v$}, אז \L{$v=0_{V}$}.
\item \L{$\forall v\in V,\ 0_{\mathbb{F}}\cdot v=0_{V}$}
\item \L{$\forall a\in\mathbb{F},\ a\cdot0_{V}=0_{V}$}
\item \L{$\forall v\in V,\ (-1_{\mathbb{F}})\cdot v=-v$}
\end{enumerate}
כאשר \L{$-1_{\mathbb{F}}\in\mathbb{F}$} הוא הנגדי של \L{$1_{\mathbb{F}}\in\mathbb{F}$}.
כלומר \L{$1_{\mathbb{F}}+(-1_{\mathbb{F}})=0_{\mathbb{F}}$}. )למשל,
אם \L{$\mathbb{F}=\{0,1\}$} אז \L{$-1_{\mathbb{F}}=1_{\mathbb{F}}$}(.
\\

\underline{הוכחה:} )לטענה 6(

יהי \L{$v\in V$}. \L{%
\[
v+(-1_{\mathbb{F}})v\underset{\text{axiom 5 }}{=}1_{\mathbb{F}}\cdot v+(-1_{\mathbb{F}})v\underset{\text{axiom 7}}{=}(1_{\mathbb{F}}+(-1_{\mathbb{F}}))v\underset{\text{axiom 4}}{=}0_{\mathbb{F}}\cdot v\underset{\text{Property 4}}{=}0_{V}
\]
}%

לפי תכונה 2 מתקיים \L{$(-1_{\mathbb{F}})v=-v$}.

\section{תת-מרחב וקטורי )תמ\textquotedblright ו(}

יהי \L{$\mathbb{F}$} שדה. יהי \L{$V$} מרחב וקטורי מעל \L{$\mathbb{F}$}.

תהי \L{$U$} תת-קבוצה של \L{$V$}, כלומר \L{$U\subseteq V$}.

\L{$U$} נקרא תת-מרחב של \L{$V$} כאשר מתקיימים התנאים הבאים )מכך
שזו תת-קבוצה של מ\textquotedblright ו ברגע שמוודאים את שלושת התנאים
הללו אפשר לדעת שזה גם מ\textquotedblright ו ובהתאם מקיים את כל האקסיומות
של מ\textquotedblright ו(:
\begin{enumerate}
\item \L{$0_{V}\in U$}
\item אם \L{$u,w\in U$} אז \L{$u+w\in U$} )סגירות ביחס לחיבור(
\item אם \L{$u\in U$} ו- \L{$c\in\mathbb{F}$} אז \L{$cu\in U$} )סגירות
ביחס לכפל בסקלר(
\end{enumerate}
הבחנה: אם \L{$U$} הוא תת-מרחב של \L{$V$}, אז \L{$U$} הוא גם מרחב
וקטורי מעל \L{$\mathbb{F}$}.

\newpage{}

\noindent\fbox{\begin{minipage}[t]{1\columnwidth - 2\fboxsep - 2\fboxrule}%
\begin{center}
אלגברה ליניארית )1( - \L{$80134$} \textbar{} האונ' העברית \textbar{}
תשפ\textquotedblright ו א' \textbar{} מרצה: ד\textquotedblright ר
אלכס גורביץ' \textbar{} מתרגל: מר יואב כהן \textbar{} ליקוט ועריכה:
נעם לב
\par\end{center}%
\end{minipage}}

\part{הרצאה {\beginL 7\endL} - {\beginL 10.11.2025\endL}}

\textcolor{blue}{בשיעור הקודם התחלנו לדבר על מרחבים וקטוריים. ראינו
דוגמאות של מרחבים וקטוריים שונים וגם הגדרנו תמ\textquotedblright ו.}

\textcolor{blue}{ראינו שכל תמ\textquotedblright ו הוא בעצמו מ\textquotedblright ו. }

\section{דוגמאות לתת-מרחבים וקטוריים}

\textbf{דוגמה 1:} יהי \L{$V=\mathbb{F}^{n}$} )\L{$n$}-יות עם קואורדינאטות
מ-\L{$\mathbb{F}$}( מ\textquotedblright ו מעל שדה \L{$\mathbb{F}$}.
נגדיר: \L{%
\[
U=\left\{ \begin{pmatrix}x_{1}\\
\vdots\\
x_{n}
\end{pmatrix}\in\mathbb{F}^{n}\,\middle|\,x_{n}=0_{\mathbb{F}}\right\} 
\]
}%

)למשל, אם \L{$V=\mathbb{R}^{2}$} אז \L{$U$} זה ציר ה-\L{$x$}(

נוודא שהדרישות של תת-מרחב מתקיימות:

1( כאשר \L{$V=\mathbb{F}^{n}$}, וקטור האפס הוא \L{$0_{V}=\begin{pmatrix}0_{\mathbb{F}}\\
\vdots\\
0_{\mathbb{F}}
\end{pmatrix}$}. \L{$\begin{pmatrix}0_{\mathbb{F}}\\
\vdots\\
0_{\mathbb{F}}
\end{pmatrix}$} בפרט מקיים \L{$x_{n}=0_{\mathbb{F}}$} לכן \L{$0_{V}\in U$}.

2( יהיו \L{$\begin{pmatrix}x_{1}\\
\vdots\\
x_{n}
\end{pmatrix},\begin{pmatrix}y_{1}\\
\vdots\\
y_{n}
\end{pmatrix}\in U$}, אז \L{$x_{n}=y_{n}=0_{\mathbb{F}}$}. מכאן \L{$x_{n}+y_{n}=0_{\mathbb{F}}+0_{\mathbb{F}}=0_{\mathbb{F}}$}
לכן: \L{%
\[
\begin{pmatrix}x_{1}\\
\vdots\\
x_{n}
\end{pmatrix}+\begin{pmatrix}y_{1}\\
\vdots\\
y_{n}
\end{pmatrix}=\begin{pmatrix}x_{1}+y_{1}\\
\vdots\\
x_{n}+y_{n}
\end{pmatrix}=\begin{pmatrix}x_{1}+y_{1}\\
\vdots\\
0_{\mathbb{F}}
\end{pmatrix}\in U
\]
}%

3( יהיו \L{$\begin{pmatrix}x_{1}\\
\vdots\\
x_{n}
\end{pmatrix}\in U$} ו- \L{$c\in\mathbb{F}$}. אז \L{$x_{n}=0_{\mathbb{F}}$}. מכאן \L{$cx_{n}=0_{\mathbb{F}}$}.
\L{%
\[
c\begin{pmatrix}x_{1}\\
\vdots\\
x_{n}
\end{pmatrix}=\begin{pmatrix}cx_{1}\\
\vdots\\
cx_{n}
\end{pmatrix}=\begin{pmatrix}cx_{1}\\
\vdots\\
0_{\mathbb{F}}
\end{pmatrix}\in U
\]
}%

\textbf{דוגמה 2:} \L{$V=\mathbb{F}[x]$} )מרחב הפולינומים(, \L{$n\in\mathbb{N}\cup\{0\}$},
\L{%
\[
U=\{P\in\mathbb{F}[x]\mid\deg P\le n\}=\mathbb{F}[x]_{\le n}
\]
}%
)כל הפולינומים שהמעלה שלהם קטנה או שווה ל-\L{$n$}(.

נוודא שהדרישות של תת-מרחב מתקיימות:

{\beginL 1\endL}( \L{$0_{\mathbb{F}}\in U$} כי \L{$\deg0_{\mathbb{F}}=-\infty\le n$}.

{\beginL 2\endL}( יהיו \L{$P,Q\in U$}. אז \L{$\deg P,\deg Q\le n$}.
\L{%
\[
\deg(P+Q)\le\max(\deg(P),\deg(Q))\le n
\]
}%
לכן \L{$P+Q\in U$}.

{\beginL 3\endL}( יהי \L{$P\in U,c\in\mathbb{F}$}. אז \L{$\deg(P)\le n$}.
מכאן: \L{%
\[
\deg(cP)\le\deg(P)\le n
\]
}%
)שווה כש- \L{$c\ne0_{\mathbb{F}}$} וקטנה כש- \L{$c=0_{\mathbb{F}}$}(.
לכן \L{$cP\in U$}.

\section{אנטי דוגמאות}
\begin{itemize}
\item דוגמה לתת-קבוצה שהיא לא תת-מרחב: \L{$\emptyset$}. כי הרי \L{$0_{V}\notin\emptyset$}.
\item \L{$V=\mathbb{R}^{2}$} מעל שדה \L{$\mathbb{R}$}, \L{$U=\left\{ \begin{pmatrix}x_{1}\\
x_{2}
\end{pmatrix}\in\mathbb{R}^{2}\;\middle|\;x_{2}\ge0\right\} $}. )המחשה: כל ציר ה-\L{$x$} ו"מעליו" בציר ה-\L{$y$}(. הדרישה הראשונה
והשנייה מתקיימות, אבל השלישית לא: \L{$U$} אינה סגורה ביחס לכפל בסקלר.
כי למשל \L{$\begin{pmatrix}0\\
1
\end{pmatrix}\in U,c=-1\in\mathbb{R}$}, אבל: \L{$-1\begin{pmatrix}0\\
1
\end{pmatrix}=\begin{pmatrix}0\\
-1
\end{pmatrix}\notin U$}.
\item \L{$V=\mathbb{R}^{2}$} מ"ו מעל \L{$\mathbb{R}$}, \L{$U=\left\{ \begin{pmatrix}x_{1}\\
x_{2}
\end{pmatrix}\in\mathbb{R}^{2}\;\middle|\;x_{1}x_{2}=0\right\} $}. )המחשה: כל ציר ה-\L{$x$} וכל ציר ה-\L{$y$}(. הדרישות הראשונה והשלישית
מתקיימות, אבל השנייה לא: \L{$U$} אינה סגורה ביחס לחיבור כי למשל \L{$\begin{pmatrix}0\\
1
\end{pmatrix},\begin{pmatrix}1\\
0
\end{pmatrix}\in U$} אבל \L{$\begin{pmatrix}0\\
1
\end{pmatrix}+\begin{pmatrix}1\\
0
\end{pmatrix}=\begin{pmatrix}1\\
1
\end{pmatrix}\notin U$} )שכן \L{$1\cdot1=1\ne0$}(.
\item \L{$n\in\mathbb{N},V=\mathbb{F}[x]$}. \L{%
\[
U=\mathbb{F}[x]_{=n}\cup\{0_{\mathbb{F}}\}=\{P\in\mathbb{F}[x]\mid\deg P=n\}\cup\{0_{\mathbb{F}}\}
\]
}%
\L{$U$} אינה סגורה ביחס לחיבור כי למשל: \L{%
\[
x^{n}+1\in U,-x^{n}\in U,(x^{n}+1)+(-x^{n})=1\notin U
\]
}%
כי \L{$\deg1=0\ne n$}. 
\end{itemize}

\section{\textcolor{blue}{טרמינולוגיה}}

\textcolor{blue}{איברים של שדה - סקלרים. איברים של מ"ו - וקטורים.
שזה )מ"ו( רק קבוצה עם 2 פעולות )חיבור וכפל בסקלר( שמקיימות את ה-8
אקסיומות שראינו. אנחנו מדברים באופן מופשט לרוב ולאו דווקא יודעים איך
האיברים נראים.}

\section{הגדרה - צירוף ליניארי )צ\textquotedblright ל(}

יהי \L{$V$} מ"ו מעל \L{$\mathbb{F}$} ויהיו \L{$v_{1},\dots,v_{n}\in V$},
\L{$c_{1},\dots,c_{n}\in\mathbb{F}$}. צירוף ליניארי של \L{$v_{1},\dots,v_{n}$}
עם מקדמים \L{$c_{1},\dots,c_{n}$} מוגדר ע"י \L{%
\[
\sum_{i=1}^{n}c_{i}v_{i}=c_{1}v_{1}+\dots+c_{n}v_{n}\in V
\]
}%


\section{תכונה )של תמ\textquotedblright ו( {]}סגירות ביחס לצ\textquotedblright ל{[}}

יהי \L{$V$} מ"ו מעל \L{$\mathbb{F}$}, \L{$U$} תמ\textquotedblright ו
של \L{$V$}. אז כל צירוף לינארי של וקטורים מ-\L{$U$} שייך ל-\L{$U$}.

כלומר, \L{$U$} סגור לכל צירוף ליניארי של וקטורים מ-\L{$U$}.

\subsection{הוכחה}

יהיו \L{$u_{1},\dots,u_{n}\in U$} ויהיו \L{$c_{1},\dots,c_{n}\in\mathbb{F}$}.

מאחר ש-\L{$U$} סגור לכפל בסקלר לכל \L{$1\le i\le n$} מתקיים \L{$c_{i}u_{i}\in U$}.

מאחר ש-\L{$U$} סגור ביחס לחיבור מתקיים \L{$c_{1}u_{1}+\dots+c_{n}u_{n}\in U$}
וסיימנו.

\section{טענה {]}חיתוך תת-מרחבים הוא תמ\textquotedblright ו{[}}

יהי \L{$V$} מ"ו מעל שדה \L{$\mathbb{F}$}. אם \L{$U,W$} תת-מרחבים
של \L{$V$}, אז \L{$U\cap W$} גם הוא תת-מרחב של \L{$V$}.

\subsection{הוכחה}
\begin{enumerate}
\item \L{$0_{V}\in U\cap W$} כי \L{$0_{V}\in U$} וגם \L{$0_{V}\in W$}
)מהיותם תת-מרחבים(. 
\item יהיו \L{$v_{1},v_{2}\in U\cap W$} אז \L{$v_{1},v_{2}\in U$} וגם
\L{$v_{1},v_{2}\in W$}. מאחר ש-\L{$U$} סגור ביחס לחיבור, \L{$v_{1}+v_{2}\in U$}.
מאחר ש-\L{$W$} סגור ביחס לחיבור, \L{$v_{1}+v_{2}\in W$} ולכן \L{$v_{1}+v_{2}\in U\cap W$}. 
\item יהיו \L{$v_{1}\in U\cap W,c\in\mathbb{F}$}. אז \L{$v_{1}\in U$}
וגם \L{$v_{1}\in W$}. מאחר ש-\L{$U$} סגור ביחס לכפל בסקלר, \L{$cv_{1}\in U$}.
באותו אופן, \L{$cv_{1}\in W$}. לכן \L{$cv_{1}\in U\cap W$}. 
\end{enumerate}
וסיימנו.

\textcolor{blue}{מטענה זו אפשר להסיק שאפשר לעשות כמה חיתוכים שנרצה
ועדיין נקבל תת-מרחב.}

\subsection{מסקנה}

יהי \L{$V$} מ"ו מעל \L{$\mathbb{F}$}. אז החיתוך של כל אוסף )אפילו
אינסופי( של תת-מרחבים של \L{$V$} הוא תת-מרחב של \L{$V$}.

\section{טענה {]}קב' הפתרונות של ממ\textquotedblright ל הומוגנית היא תמ\textquotedblright ו{[}}

יהי \L{$\mathbb{F}$} שדה, יהיו \L{$m,n\in\mathbb{N}$}, ותהי \L{$A\in M_{m\times n}(\mathbb{F})$}
)מטריצה עם \L{$m$} שורות ו-\L{$n$} עמודות עם מקדמים משדה \L{$\mathbb{F}$}(.

נגדיר \L{%
\[
U=\left\{ \begin{pmatrix}x_{1}\\
\vdots\\
x_{n}
\end{pmatrix}\in\mathbb{F}^{n}\middle|\:A\begin{pmatrix}x_{1}\\
\vdots\\
x_{n}
\end{pmatrix}=\begin{pmatrix}0\\
\vdots\\
0
\end{pmatrix}\right\} =\text{Null}(A)
\]
}%

זו קבוצת הפתרונות של מערכת משוואות הומוגנית עם מטריצה \L{$A$}.

כלומר, קבוצת כל הפתרונות של ממ\textquotedblright ל הומוגנית שהמטריצה
הלא מורחבת שלה היא \L{$A$}.

מסמנים את הקבוצה הזאת ב- \L{$\text{Null}(A)$}. \\

אז \L{$\text{Null}(A)$} תמ\textquotedblright ו של \L{$\mathbb{F}^{n}$}.

\subsection{\textcolor{teal}{תזכורת}}

\textcolor{teal}{לכתוב ש- \L{$\begin{pmatrix}x_{1}\\
\vdots\\
x_{n}
\end{pmatrix}$} זה פתרון של \L{$(A|b)$} זה אותו דבר כמו לכתוב \L{$A\begin{pmatrix}x_{1}\\
\vdots\\
x_{n}
\end{pmatrix}=b$}}

\subsection{הוכחה}
\begin{enumerate}
\item \L{$\begin{pmatrix}0_{\mathbb{F}}\\
\vdots\\
0_{\mathbb{F}}
\end{pmatrix}\in U$} \quad{}כי\quad{} \L{$A\begin{pmatrix}0_{\mathbb{F}}\\
\vdots\\
0_{\mathbb{F}}
\end{pmatrix}=\begin{pmatrix}0_{\mathbb{F}}\\
\vdots\\
0_{\mathbb{F}}
\end{pmatrix}$} .
\item אם \L{$d,h\in U$} אז \L{%
\[
Ad=Ah=\begin{pmatrix}0_{\mathbb{F}}\\
\vdots\\
0_{\mathbb{F}}
\end{pmatrix}
\]
}%
מכאן \L{%
\[
A(d+h)=Ad+Ah=\begin{pmatrix}0_{\mathbb{F}}\\
\vdots\\
0_{\mathbb{F}}
\end{pmatrix}+\begin{pmatrix}0_{\mathbb{F}}\\
\vdots\\
0_{\mathbb{F}}
\end{pmatrix}=\begin{pmatrix}0_{\mathbb{F}}\\
\vdots\\
0_{\mathbb{F}}
\end{pmatrix}
\]
}%
לכן \L{$d+h\in U$}.
\item אם \L{$d\in U,c\in\mathbb{F}$} אז \L{%
\[
Ad=\begin{pmatrix}0_{\mathbb{F}}\\
\vdots\\
0_{\mathbb{F}}
\end{pmatrix}
\]
}%
נשים לב )לפי תכונות כפל מטריצה בסקלר( \L{%
\[
A(cd)=c(Ad)=c\cdot\begin{pmatrix}0_{\mathbb{F}}\\
\vdots\\
0_{\mathbb{F}}
\end{pmatrix}=\begin{pmatrix}0_{\mathbb{F}}\\
\vdots\\
0_{\mathbb{F}}
\end{pmatrix}
\]
}%
לכן \L{$cd\in U$} 
\end{enumerate}
\begin{flushright}
\L{$\square$}
\par\end{flushright}

\newpage{}

\noindent\fbox{\begin{minipage}[t]{1\columnwidth - 2\fboxsep - 2\fboxrule}%
\begin{center}
אלגברה ליניארית )1( - \L{$80134$} \textbar{} האונ' העברית \textbar{}
תשפ\textquotedblright ו א' \textbar{} מרצה: ד\textquotedblright ר
אלכס גורביץ' \textbar{} מתרגל: מר יואב כהן \textbar{} ליקוט ועריכה:
נעם לב
\par\end{center}%
\end{minipage}}

\part{הרצאה {\beginL 8\endL} - {\beginL 11.11.2025\endL}}

\section{דוגמאות לטענה האחרונה שראינו {]}שקב' הפתרונות של ממ\textquotedblright ל
הומוגנית היא תמ\textquotedblright ו{[}}

\uline{ניזכר בטענה:} יהי \L{$\mathbb{F}$} שדה, יהיו \L{$m,n\in\mathbb{N}$},
ותהי \L{$A\in M_{m\times n}(\mathbb{F})$}.

אז \L{%
\[
U=\left\{ \begin{pmatrix}x_{1}\\
\vdots\\
x_{n}
\end{pmatrix}\in\mathbb{F}^{n}\middle|\:A\begin{pmatrix}x_{1}\\
\vdots\\
x_{n}
\end{pmatrix}=\begin{pmatrix}0\\
\vdots\\
0
\end{pmatrix}\right\} =\text{Null}(A)
\]
}%

)קבוצת הפתרונות של ממ\textquotedblright ל הומוגנית עם מטריצה לא מורחבת
\L{$A$}( תמ\textquotedblright ו של \L{$\mathbb{F}^{n}$}.\\

\textbf{דוגמה }\textbf{{\beginL 1}\endL}\textbf{: \quad{}\L{$A=(0\dots0\:1)\in M_{1\times n}(\mathbb{F}),\quad m=1$}.}
אז \L{%
\[
U=\left\{ \begin{pmatrix}x_{1}\\
\vdots\\
x_{n}
\end{pmatrix}\in\mathbb{F}^{n}\middle|\:x_{n}=0\right\} 
\]
}%

תמ\textquotedblright ו.\\

\textbf{דוגמה }\textbf{{\beginL 2}\endL}\textbf{:\quad{} \L{$m=n\in\mathbb{N}$},\quad{}
}\L{$A=I_{n}=\begin{pmatrix}1 & 0 & \dots & 0\\
0 & 1 & \dots & 0\\
\vdots & \vdots & \ddots & \vdots\\
0 & 0 & \dots & 1
\end{pmatrix}$}. אז

\L{%
\begin{align*}
U & =\left\{ \begin{pmatrix}x_{1}\\
\vdots\\
x_{n}
\end{pmatrix}\in\mathbb{F}^{n}\;\middle|\;I_{n}\begin{pmatrix}x_{1}\\
\vdots\\
x_{n}
\end{pmatrix}=\begin{pmatrix}0\\
\vdots\\
0
\end{pmatrix}\right\} \\
 & =\left\{ \begin{pmatrix}x_{1}\\
\vdots\\
x_{n}
\end{pmatrix}\in\mathbb{F}^{n}\;\middle|\;\begin{pmatrix}x_{1}\\
\vdots\\
x_{n}
\end{pmatrix}=\begin{pmatrix}0\\
\vdots\\
0
\end{pmatrix}\right\} \\
 & =\left\{ \begin{pmatrix}x_{1}\\
\vdots\\
x_{n}
\end{pmatrix}\in\mathbb{F}^{n}\;\middle|\;\begin{cases}
x_{1}=0\\
\vdots\\
x_{n}=0
\end{cases}\right\} =\left\{ \begin{pmatrix}0\\
\vdots\\
0
\end{pmatrix}\right\} 
\end{align*}
}%

תמ\textquotedblright ו.\\

\textcolor{blue}{לגבי הקבוצה שקיבלנו - מסתדר עם מה שאנחנו יודעים על
ממ\textquotedblright ל הומוגניות...} \textcolor{teal}{)אפיון פתרון
יחיד )טריוויאלי(...(}\\

\textbf{דוגמה }\textbf{{\beginL 3}\endL}\textbf{:} \textbf{\quad{}}\L{$A=(0\dots0)\in M_{1\times n}(\mathbb{F}),\quad m=1$}.
אז \L{%
\[
U=\left\{ \begin{pmatrix}x_{1}\\
\vdots\\
x_{n}
\end{pmatrix}\in\mathbb{F}^{n}\;\middle|\;(0\dots0)\begin{pmatrix}x_{1}\\
\vdots\\
x_{n}
\end{pmatrix}=\underset{\begin{array}{c}
\textcolor{green}{\downarrow}\\
\textcolor{green}{(*)}
\end{array}}{0}\right\} 
\]
}%

\textcolor{green}{\L{$(*)$} \L{$m$}-יה}\\

תמ\textquotedblright ו.\\

\textcolor{blue}{\uline{הערה:}}\textcolor{blue}{{} התופעה שבדוגמאות
2 ו-3 מתקיימת בכל מרחב וקטורי, כלומר לכל מ\textquotedblright ו \L{$V$}
מתקיים ש-} \L{$\underset{\begin{array}{c}
\textcolor{green}{\downarrow}\\
\textcolor{green}{(***)}
\end{array}}{V},\underset{\begin{array}{c}
\textcolor{green}{\downarrow}\\
\textcolor{green}{(**)}
\end{array}}{\{0_{V}\}}$} \textcolor{blue}{הם תמ\textquotedblright ו של \L{$V$}.}

\textcolor{green}{\L{$(**)$} התמ\textquotedblright ו \textquotedblleft הכי
קטן\textquotedblright{} של \L{$V$}}

\textcolor{green}{\L{$(***)$} התמ\textquotedblright ו \textquotedblleft הכי
גדול\textquotedblright{} של \L{$V$}}\\

\textcolor{blue}{\uline{הערה:}}\textcolor{blue}{{} למעשה כל תת-מרחב
של \L{$\mathbb{F}^{n}$} ניתן לתאר כמו בטענה, כלומר לכל תת-מרחב \L{$U$}
של \L{$\mathbb{F}^{n}$} קיימים}\textbf{\textcolor{blue}{\quad{}}}\textcolor{blue}{{}
\L{$A\in M_{m\times n}(\mathbb{F}),\quad m\in\mathbb{N}$} כך ש- \L{$U=\text{Null}(A)$}.
)לא נוכיח את זה בשלב זה(. דהיינו, תמיד אפשר למצוא ממ\textquotedblright ל
הומוגנית שאוסף הפתרונות שלה הוא התמ\textquotedblright ו.}

\medskip{}

\rule[0.5ex]{1\columnwidth}{1pt}

\medskip{}


\section{פרוש - \L{$\spn$}}

\subsection{\textcolor{blue}{תיאור במילים}}

\textcolor{blue}{נראה איך מקבוצה שהיא לא תמ\textquotedblright ו אנחנו
יכולים \textquotedblleft ליצור\textquotedblright{} תמ\textquotedblright ו.}

\textcolor{blue}{כלומר למצוא תמ\textquotedblright ו שמכיל את הקבוצה
ושיהיה \textquotedblleft קטן ככל הניתן\textquotedblright .}

\textcolor{blue}{אז גם קודם כל וקטור אפס צריך להיות שם.}

\textcolor{blue}{וכל וקטור של הקבוצה צריך להיות בתמ\textquotedblright ו,
וגם כל מכפלה בסקלר של כל וקטור כזה.}

\textcolor{blue}{חוץ מזה יש דרישה של סגירות ביחס לחיבור. אז גם הסכומים
של כפולות בסקלר של איברים מהקבוצה.}

\textcolor{blue}{כלומר: כל הצירופים הלינאריים האפשריים של איברי הקבוצה.}

\subsection{הגדרה - פרוש}

יהי \L{$V$} מ\textquotedblright ו כלשהו מעל \L{$\mathbb{F}$} ותהא
\L{$\emptyset\neq S\subseteq V$} קבוצה כלשהי. הפרוש של \L{$S$} מוגדר
ע\textquotedblright י \L{%
\[
\spn S=\left\{ c_{1}v_{1}+\dots+c_{n}v_{n}\;\middle|\;n\in\mathbb{N},c_{1},\dots,c_{n}\in\mathbb{F},v_{1},\dots,v_{n}\in S\right\} 
\]
}%
בנוסף מגדירים\L{%
\[
\spn\emptyset=\left\{ 0_{V}\right\} 
\]
}%


\section{טענה {]}פרוש הוא תמ\textquotedblright ו{[}}

יהי \L{$V$} מ\textquotedblright ו מעל \L{$\mathbb{F}$} ותהא \L{$S\subseteq V$}.
אז \L{$\spn S$} הוא תמ\textquotedblright ו של \L{$V$}.

\subsection{הוכחה}

1( {]}וקטור אפס{[}

אם \L{$S=\emptyset$} אז \L{$0_{V}\in\{0_{V}\}=\spn\emptyset=\spn S$}.

אם \L{$S\neq\emptyset$} אז קיים \L{$v\in S$}, לכן \L{$0_{V}=0_{\mathbb{F}}\cdot v\in\spn S$}\\

2( {]}סגירות ביחס לחיבור{[}

אם \L{$S=\emptyset$} אז \L{$\spn\emptyset=\left\{ 0_{V}\right\} $}
סגור ביחס לחיבור.

אם \L{$S\neq\emptyset$}, נבחר \L{$u,v\in\spn S$}.

אז קיימים \L{$v_{1},\dots,v_{n}\in S$} \L{$c_{1},\dots,c_{n}\in\mathbb{F}$}
כך ש- \L{$v=c_{1}v_{1}+\dots+c_{n}v_{n}$} וגם קיימים \L{$u_{1},\dots,u_{m}\in S$},
\L{$d_{1},\dots,d_{m}\in\mathbb{F}$} כך ש- \L{$u=d_{1}u_{1}+\dots+d_{m}u_{m}$}
אז \L{$u+v=d_{1}u_{1}+\dots+d_{m}u_{m}+c_{1}v_{1}+\dots+c_{n}v_{n}\in\spn S$}\\

3( {]}סגירות ביחס לכפל בסקלר{[}

אם \L{$S=\emptyset$} אז \L{$\spn S=\left\{ 0_{V}\right\} $} סגור
ביחס לכפל בסקלר.

אם \L{$S\neq\emptyset$}, נבחר \L{$v\in\spn S$}, \L{$c\in\mathbb{F}$}.

אז קיימים \L{$v_{1},\dots,v_{n}\in S$}, \L{$c_{1},\dots,c_{n}\in\mathbb{F}$}
כך ש- \L{$v=c_{1}v_{1}+\dots+c_{n}v_{n}$} 

אז \L{%
\[
cv=c\left(c_{1}v_{1}+\dots+c_{n}v_{n}\right)=\left(cc_{1}\right)v_{1}+\dots+\left(cc_{n}\right)v_{n}\in\spn S
\]
}%

\begin{flushright}
\L{$\square$}
\par\end{flushright}

\noindent\fbox{\begin{minipage}[t]{1\columnwidth - 2\fboxsep - 2\fboxrule}%
\textbf{\textcolor{teal}{שחזור הוכחה )שלי(}}\textcolor{teal}{}\\

\textcolor{teal}{יהי \L{$V$} מ"ו מעל \L{$\mathbb{F}$} ותהא \L{$S\subseteq V$}.
נוכיח כי \L{$\spn S$} תמ\textquotedblright ו.}

\textcolor{teal}{אם \L{$S=\emptyset$} אז \L{$\spn S=\{0_{V}\}$}
וכפי שלמדנו זה תמ\textquotedblright ו )הטריוויאלי(.}

\textcolor{teal}{נניח כי \L{$S\neq\emptyset$} ונבדוק את קיום שלושת
האקסיומות.}\\

\textcolor{teal}{1( \L{$0_{\mathbb{F}}\in\mathbb{F}$} לפי אקסיומות
שדה וגם \L{$S\neq\emptyset$} אז קיים \L{$v\in S$} אחד לפחות ומתקיים
\L{$0\cdot v=0_{V}$} לכן \L{$0_{V}\in\spn S$}.}\\

\textcolor{teal}{2( יהיו \L{$s_{1},s_{2}\in\spn S$}. אז \L{%
\[
s_{1}=c_{1}v_{1}+\dots+c_{n}v_{n}\quad s.t\quad\forall_{1\le i\le n},\;c_{i}\in\mathbb{F},\;v_{i}\in S
\]
}%
 \L{%
\[
s_{2}=d_{1}k_{1}+\dots+d_{m}k_{m}\quad s.t\quad\forall_{1\le i\le m},\;d_{i}\in\mathbb{F},\;k_{i}\in S
\]
}%
}

\textcolor{teal}{ובסך הכל \L{%
\[
s_{1}+s_{2}=c_{1}v_{1}+\dots+c_{n}v_{n}+d_{1}k_{1}+\dots+d_{m}k_{m}
\]
}%
 זה צ"ל של איברים מ-\L{$S$} ולכן שייך ל-\L{$\spn S$}.}\\

\textcolor{teal}{3( יהיו \L{$d\in\mathbb{F},s_{1}\in\spn S$}. אז
\L{%
\[
s_{1}=c_{1}v_{1}+\dots+c_{n}v_{n}\quad s.t\quad\forall_{1\le i\le n},\;c_{i}\in\mathbb{F},\;v_{i}\in S
\]
}%
 ומתקיים }

\textcolor{teal}{\L{%
\[
d\cdot s_{1}=d\left(c_{1}v_{1}+\dots+c_{n}v_{n}\right)=d\cdot c_{1}v_{1}+\dots+d\cdot c_{n}v_{n}
\]
}%
}

\textcolor{teal}{\L{$\mathbb{F}$} שדה ולכן סגור לכפל ונקבל}

\textcolor{teal}{\L{%
\[
=\left(dc_{1}\right)v_{1}+\dots+\left(dc_{n}\right)v_{n}\quad s.t\quad dc_{1},\dots,dc_{n}\in\mathbb{F},\quad v_{1},\dots,v_{n}\in S
\]
}%
}

\textcolor{teal}{לכן \L{$d\cdot s_{1}\in\spn S$} וסיימנו.}%
\end{minipage}}

\section{תכונה {\beginL 1\endL} של \L{$\spn$} {]}קבוצה מוכלת בפרוש שלה{[}}

יהי \L{$V$} מ\textquotedblright ו ותהא \L{$S\subseteq V$}. אז \L{$S\subseteq\spn S$}.

\subsection{הוכחה}

יהי \L{$w\in S$}. אז \L{$w=1_{\mathbb{F}}\cdot w\in\spn S$}.\\


\section{טענה {]}תכונת המינימליות{[} {]}קבוצה מוכלת בתמ\textquotedblright ו
גורר ספאן של הקבוצה מוכל בתמ\textquotedblright ו{[}}

יהי \L{$V$} מ\textquotedblright ו מעל \L{$\mathbb{F}$} ותהא \L{$S\subseteq V$}.
יהי \L{$U$} תמ\textquotedblright ו של \L{$V$} כך ש- \L{$S\subseteq U$}.

אז \L{$\spn S\subseteq U$}.

\subsection{הוכחה}

{\small\emph{)שחזור שלי, מאוד דומה לרשומות מאלכס( }}{\small\par}

יהי \L{$w\in\spn S$}. אז \L{$w=c_{1}v_{1}+\cdots+c_{n}v_{n}\quad s.t\quad\forall_{1\leq i\leq n}\:c_{i}\in\mathbb{F}\:v_{i}\in S$}.
אבל \L{$S\subseteq U$} ומכך נובע \L{$\forall_{1\leq i\leq n}\:c_{i}\in\mathbb{F}\:v_{i}\in U$}.
\L{$U$} תמ\textquotedblright ו ולכן סגור לחיבור ולכפל בסקלר אז נקבל
\L{$w\in U$} וסיימנו.
\begin{flushright}
\L{$\square$}
\par\end{flushright}

\noindent\fbox{\begin{minipage}[t]{1\columnwidth - 2\fboxsep - 2\fboxrule}%
\textcolor{teal}{תכונה 1 והטענה הזאת, בנוסף לטענה שראינו שפרוש הוא
תמ\textquotedblright ו, מהוות הגדרה מקבילה של ספאן שהיא: \L{$\spn S$}
הוא התמ\textquotedblright ו המינימלי שמכיל את \L{$S$}.}

\textcolor{teal}{מה זה אומר?}

\textcolor{teal}{שהספאן מקיים את התכונה )להיות תמ\textquotedblright ו
ולהכיל את הקבוצה( ומוכל בכל אובייקט אחר שמקיים את התכונה )להיות תמ\textquotedblright ו
ולהכיל את הקבוצה(. )מינימליות זה להיות מוכל בכל אובייקט אחר שמקיים
את התכונה(.}%
\end{minipage}}

\section{תכונה {\beginL 2\endL} של \L{$\spn$} {]}תמ\textquotedblright ו שווה
לפרוש שלו{[}}

יהיו \L{$V$} מ\textquotedblright ו, \L{$U$} תמ\textquotedblright ו
של \L{$V$}. אז \L{$\spn U=U$}.

\subsection{הוכחה}

\L{$U\subseteq\spn U$} לפי תכונה {\beginL 1\endL} {]}קבוצה מוכלת
בפרוש שלה{[}.

בנוסף, מאחר ש- \L{$U$} תמ\textquotedblright ו וגם \L{$U\subseteq U$},

לפי הטענה הקודמת {]}תכונת המינימליות{[} אנו מקבלים ש- \L{$\spn U\subseteq U$}.

מההכלה הדו-כיוונית \L{$\spn U=U$}.
\begin{flushright}
\L{$\square$}
\par\end{flushright}

\section{תכונה {\beginL 3\endL} של \L{$\spn$} {]}ספאן משמר יחס הכלה{[}}

יהי \L{$V$} מ\textquotedblright ו ויהיו \L{$S_{1}\subseteq S_{2}\subseteq V$}.
אז \L{$\spn S_{1}\subseteq\spn S_{2}$}.

\subsection{הוכחה}

לפי תכונה {\beginL 1\endL} {]}קבוצה מוכלת בפרוש שלה{[} מתקיים \L{$S_{2}\subseteq\spn S_{2}$}.

אז מהנתון ומטרנזיטיביות ההכלה \L{$S_{1}\subseteq\spn S_{2}$}.

מאחר ש- \L{$\spn S_{2}$} הוא תמ\textquotedblright ו, לפי הטענה הקודמת
{]}תכונת המינימליות{[} מתקיים \L{$\spn S_{1}\subseteq\spn S_{2}$}.
\begin{flushright}
\L{$\square$}
\par\end{flushright}

\medskip{}

\rule[0.5ex]{1\columnwidth}{1pt}

\medskip{}

\textcolor{blue}{כאשר \L{$S$} היא קבוצה סופית יש דרך נוספת להגדיר
את הפרוש שלה.}

\section{טענה {]}הגדרה מקבילה לספאן כאשר הקבוצה סופית{[}}

יהיו \L{$V$} מ\textquotedblright ו מעל \L{$\mathbb{F}$}, \L{$v_{1},\dots,v_{k}\in V$},
\L{$S=\{v_{1},\dots,v_{k}\}$}. אז \L{%
\[
\spn S=\left\{ c_{1}v_{1}+\dots+c_{k}v_{k}\:\middle|\:c_{1},\dots,c_{k}\in\mathbb{F}\right\} 
\]
}%


\subsection{\textcolor{blue}{הסבר}}

\textcolor{blue}{ההכלה \L{$\supseteq$} מתקיימת כי \L{$c_{1}v_{1}+\dots+c_{k}v_{k}$}
הוא צ\textquotedblright ל של איברים מ- \L{$S$}.}

\textcolor{blue}{ההכלה \L{$\subseteq$} מתקיימת כי כל צ\textquotedblright ל
של איברים מ- \L{$S$} ניתן להפוך ע\textquotedblright י כינוס מחוברים
והוספת מחוברים עם מקדם \L{$0$} לצ\textquotedblright ל מהצורה \L{$c_{1}v_{1}+\dots+c_{k}v_{k}$}
בלי לשנות את ערכו.}

\newpage{}

\noindent\fbox{\begin{minipage}[t]{1\columnwidth - 2\fboxsep - 2\fboxrule}%
\begin{center}
אלגברה ליניארית )1( - \L{$80134$} \textbar{} האונ' העברית \textbar{}
תשפ\textquotedblright ו א' \textbar{} מרצה: ד\textquotedblright ר
אלכס גורביץ' \textbar{} מתרגל: מר יואב כהן \textbar{} ליקוט ועריכה:
נעם לב
\par\end{center}%
\end{minipage}}

\part{הרצאה {\beginL 9\endL} - {\beginL 17.11.2025\endL}}

\textcolor{teal}{בתזכורת דיברנו על הגדרת הפרוש הראשונה שראינו ועל
הטענה {]}הגדרה מקבילה לספאן כאשר הקבוצה סופית{[}.}

\section{הגדרה - ספאן של סדרה של וקטורים}

יהי \L{$V$} מ\textquotedblright ו מעל \L{$\mathbb{F}$}, \L{$(v_{1},\dots,v_{k})$}
סדרה של וקטורים מ- \L{$V$}. נגדיר את הפרוש של הסדרה ע\textquotedblright י
\L{%
\[
\spn(v_{1},\dots,v_{k})=\{c_{1}v_{1}+\dots+c_{k}v_{k}\mid c_{1},\dots,c_{k}\in\mathbb{F}\}
\]
}%


\subsection{דוגמה}

יהי \L{$V$} מ\textquotedblright ו מעל \L{$\mathbb{F}$}, \L{$v\in V$}.
אז \L{$\spn(v)=\{cv\mid c\in\mathbb{F}\}$}.

\section{טענה {]}ספאן עמודות של מטריצה{[}}

תהי \L{$A\in M_{m\times n}(\mathbb{F})$}. לכל \L{$1\leq j\leq n$}
נסמן ב- \L{$v_{j}\in\mathbb{F}^{m}$} את העמודה ה- \L{$j$} של \L{$A$}.
אז \L{%
\[
\spn(v_{1},\dots,v_{n})=\{Ax\mid x\in\mathbb{F}^{n}\}=\{b\in\mathbb{F}^{m}\mid\text{The Linear System }Ax=b\text{ has a solution}\}
\]
}%
\noindent\fbox{\begin{minipage}[t]{1\columnwidth - 2\fboxsep - 2\fboxrule}%
\textcolor{violet}{זוהי אחת הזהויות הבסיסיות והחשובות ביותר באלגברה
לינארית.}

\textcolor{violet}{הסיבה לכך שהשוויון נכון נובעת מהגדרת כפל מטריצה-וקטור.}

\textcolor{violet}{כאשר כופלים מטריצה \L{$A$} בוקטור \L{$x$}, התוצאה
היא בעצם צירוף לינארי של עמודות המטריצה \L{$A$}, כאשר הרכיבים של
\L{$x$} הם המקדמים )הסקלארים( של הצירוף.}

\textcolor{violet}{לכן:}

\textcolor{violet}{הביטוי \L{$\{Ax\mid x\in\mathbb{F}^{n}\}$} הוא
פשוט אוסף כל הצירופים הלינאריים האפשריים של עמודות \L{$A$}, שזו בדיוק
ההגדרה של \L{$\spn$}.}

\textcolor{violet}{מערכת המשוואות \L{$Ax=b$} פתירה אם ורק אם קיים
\L{$x$} כלשהו כך ש-\L{$Ax$} יתן את \L{$b$}. כלומר, אם ורק אם \L{$b$}
נמצא בתוך הספאן של עמודות המטריצה.}\\
\textcolor{violet}{}\\
\textcolor{teal}{ספאן עמודות של מטריצה \L{$A$} הוא קבוצת וקטורי התוצאה
)\L{$b$}( של כל ממ\textquotedblright ל עם פתרון ש-\L{$A$} היא המטריצה
הלא-מורחבת שלה.}%
\end{minipage}}

\subsection{הוכחה}

\L{$\subseteq$}:

אם \L{$v\in\spn(v_{1},\dots,v_{n})$} אז קיימים \L{$c_{1},\dots,c_{n}\in\mathbb{F}$}
כך ש- \L{$v=c_{1}v_{1}+\dots+c_{n}v_{n}$} לכן \L{$v=A\begin{pmatrix}c_{1}\\
\vdots\\
c_{n}
\end{pmatrix}$} ומתקיים \L{$v\in\{Ax\mid x\in\mathbb{F}^{n}\}$}.

\L{$\supseteq$}:

אם \L{$v\in\{Ax\mid x\in\mathbb{F}^{n}\}$} אז קיים \L{$c=\begin{pmatrix}c_{1}\\
\vdots\\
c_{n}
\end{pmatrix}\in\mathbb{F}^{n}$} כך ש- \L{$v=A\begin{pmatrix}c_{1}\\
\vdots\\
c_{n}
\end{pmatrix}$}. מכאן \L{$v=c_{1}v_{1}+\dots+c_{n}v_{n}\in\spn(v_{1},\dots,v_{n})$}
וסיימנו.
\begin{flushright}
\L{$\square$}
\par\end{flushright}

\subsubsection{\textcolor{teal}{שחזור הוכחה שלי}}

נסמן: \L{%
\[
\alpha:=\spn(v_{1},\dots,v_{n}),\quad\beta:=\{Ax\mid x\in\mathbb{F}^{n}\}
\]
}%

\textbf{כיוון ראשון:} נוכיח כי \L{$\alpha\subseteq\beta$}. יהי \L{$v\in\alpha$}.
אז קיימים \L{$c_{1},\dots,c_{n}\in\mathbb{F}$} כך ש: \L{%
\[
v=c_{1}v_{1}+\dots+c_{n}v_{n}
\]
}%
נסתכל על הוקטור \L{$\vec{c}=\begin{pmatrix}c_{1}\\
\vdots\\
c_{n}
\end{pmatrix}\in\mathbb{F}^{n}$}. מהגדרת כפל מטריצה בוקטור מתקיים: \L{%
\[
A\vec{c}=c_{1}v_{1}+\dots+c_{n}v_{n}=v
\]
}%
לכן \L{$v\in\beta$}.

\textbf{כיוון שני:} נוכיח כי \L{$\beta\subseteq\alpha$}. יהי \L{$v\in\beta$}.
דהיינו, קיים \L{$x\in\mathbb{F}^{n}$} כך ש- \L{$v=Ax$}. אז קיימים
\L{$x_{1},\dots,x_{n}\in\mathbb{F}$} כך ש: \L{%
\[
v=Ax=A\begin{pmatrix}x_{1}\\
\vdots\\
x_{n}
\end{pmatrix}=x_{1}C_{1}^{A}+\dots+x_{n}C_{n}^{A}
\]
}%
אבל הגדרנו \L{$\forall_{1\le j\le n},C_{j}^{A}=v_{j}$} לכן: \L{%
\[
v=x_{1}v_{1}+\dots+x_{n}v_{n}
\]
}%
כלומר זהו צ\textquotedblright ל של \L{$(v_{1},\dots,v_{n})$}, ולכן
\L{$v\in\alpha$} וסיימנו.

\hfill{}\L{$\square$}

\subsection{מסקנה {]}ספאן סדרת וקטורי היחידה הסטנדרטיים הוא \L{$\mathbb{F}^{n}$}{[}}

\L{%
\[
\spn(e_{1},\dots,e_{n})=\mathbb{F}^{n}
\]
}%
כאשר \L{$e_{i}=\begin{pmatrix}0\\
\vdots\\
1\\
\vdots\\
0
\end{pmatrix}\in\mathbb{F}^{n}$} )ה-1 במקום ה-\L{$i$}(.

\subsubsection{הוכחה}

\L{%
\[
\spn(e_{1},\dots,e_{n})\underset{(*)}{=}\left\{ I_{n}x\mid x\in\mathbb{F}^{n}\right\} =\left\{ x\mid x\in\mathbb{F}^{n}\right\} =\mathbb{F}^{n}
\]
}%

\L{$(*)$} כפי שראינו, \L{$e_{1},\dots,e_{n}$} הם עמודות המטריצה
\L{$I_{n}$} ואז משתמשים בטענה שהרגע ראינו.

\section{הגדרה - תלות לינארית )תכונה של סדרה(}

יהי \L{$V$} מ\textquotedblright ו מעל \L{$\mathbb{F}$}, \L{$(v_{1},\dots,v_{k})$}
סדרה של וקטורים מ- \L{$V$}.

אומרים כי הסדרה תלויה לינארית )ת\textquotedblright ל( כאשר \L{%
\[
\exists_{1\leq i\leq k}\quad s.t\quad v_{i}\in\spn(v_{1},\dots,v_{i-1},v_{i+1},v_{i+2},\dots,v_{k})
\]
}%

\textcolor{blue}{כלומר, סדרת הוקטורים ללא \L{$v_{i}$}.}

\textcolor{blue}{דהיינו, כשקיים וקטור ששייך לספאן של שאר הוקטורים,
משמע כשקיים וקטור בסדרה שהוא צ\textquotedblright ל של שאר הוקטורים
בסדרה.}

\subsection{תכונות של תלות לינארית}

{\beginL 1\endL}( אם \L{$(v_{1},\dots,v_{k})$} סדרה כך ש-

\L{%
\[
\exists_{1\leq i\leq k,\:1\leq j\leq k}\quad s.t\quad i\neq j\quad\wedge\quad v_{i}=v_{j}
\]
}%

אז היא ת\textquotedblright ל. )\textcolor{blue}{כי הרי \L{$v_{i}\in\spn(v_{1},\dots,v_{i-1},v_{i+1},\dots,v_{j},\dots,v_{k})$}}
(.\\

{\beginL 2\endL}( אם \L{$(v_{1},\dots,v_{k})$} סדרה כך ש-\L{%
\[
\exists_{1\leq i\leq k}\quad s.t\quad v_{i}=0_{V}
\]
}%
אז היא ת"ל. )\textcolor{blue}{כי הרי פרוש הוא תמ\textquotedblright ו
ו-וקטור האפס תמיד שייך לתמ\textquotedblright ו. במילים אחרות \L{$v_{i}=0_{V}\in\spn(v_{1},\dots,v_{i-1},v_{i+1},\dots,v_{k})$}}(.\\

{\beginL 3\endL}( \textbf{\textcolor{teal}{{]}תכונה חשובה{[}}} אם
\L{$(v_{1},\dots,v_{k})$} סדרה ת"ל, אז כל סדרה אחרת שמתקבלת ממנה
ע\textquotedblright י שינוי סדר הוקטורים היא גם ת\textquotedblright ל.\\

{\beginL 4\endL}( אם \L{$(v_{1},\dots,v_{k})$} סדרה ת"ל, ו- \L{$v_{k+1}\in V$},
אז גם הסדרה \L{$(v_{1},\dots,v_{k},v_{k+1})$} ת"ל. )כי הרי אם \L{$v_{i}\in\spn(v_{1},\dots,v_{i-1},v_{i+1},\dots,v_{k})$},
אז \L{$\left(v_{i}\in\spn(v_{1},\dots,v_{i-1},v_{i+1},\dots,v_{k},v_{k+1})\right)$}.\\

{\beginL 5\endL}( אם לסדרה \L{$(v_{1},\dots,v_{k})$} קיימת תת-סדרה
ת"ל, אז היא גם ת"ל.

\section{טענה {]}אפיון ת\textquotedblright ל בעזרת צ\textquotedblright ל מתאפס
לא טריוויאלי{[}}

יהי \L{$V$} מ"ו מעל \L{$\mathbb{F}$}, \L{$(v_{1},\dots,v_{k})$}
סדרה של וקטורים מ-\L{$V$}.

אז \L{$(v_{1},\dots,v_{k})$} ת"ל אם"ם קיימים \L{$c_{1},\dots,c_{k}\in\mathbb{F}$}
כך ש- \L{%
\[
c_{1}v_{1}+\dots+c_{k}v_{k}=0_{V}
\]
}%
ולא כולם שווים ל-\L{$0$}.

\medskip{}

\textcolor{teal}{בכתיב כמתים:} \L{$(v_{1},\dots,v_{k})$} ת\textquotedblright ל
אם\textquotedblright ם \L{%
\[
\exists_{c_{1},\dots,c_{k}\in\mathbb{F}},\exists_{1\leq i\leq k},c_{i}\neq0\quad\wedge\quad c_{1}v_{1}+\dots+c_{k}v_{k}=0_{V}
\]
}%


\subsection{הוכחה}

\L{$\Leftarrow$}: מאחר ש- \L{$(v_{1},\dots,v_{k})$} ת"ל קיים \L{$1\le i\le k$}
כך ש- \L{$v_{i}\in\spn(v_{1},\dots,v_{i-1},v_{i+1},\dots,v_{k})$}.
אז קיימים \L{$c_{1},\dots,c_{i-1},c_{i+1},\dots,c_{k}\in\mathbb{F}$}
כך ש- \L{%
\[
v_{i}=c_{1}v_{1}+\dots+c_{i-1}v_{i-1}+c_{i+1}v_{i+1}+\dots+c_{k}v_{k}
\]
}%
מכאן \L{%
\[
-c_{1}v_{1}-\dots-c_{i-1}v_{i-1}+\underset{\begin{array}{c}
\downarrow\\
{\color{green}c_{i}\ne0}
\end{array}}{1\cdot v_{i}}-c_{i+1}v_{i+1}-\dots-c_{k}v_{k}=0_{V}
\]
}%
\L{$\Rightarrow$}: אם קיימים \L{$c_{1},\dots,c_{k}\in\mathbb{F}$}
שלא כולם שווים ל-\L{$0$} כך ש- \L{%
\[
c_{1}v_{1}+\dots+c_{k}v_{k}=0
\]
}%
אז קיים \L{$1\le i\le k$} כך ש- \L{$c_{i}\ne0$}. מכאן \L{%
\[
c_{i}v_{i}=-c_{1}v_{1}-\dots-c_{i-1}v_{i-1}-c_{i+1}v_{i+1}-\dots-c_{k}v_{k}
\]
}%
\L{$c_{i}\ne0$} לכן קיים לו הופכי ומתקיים \L{%
\[
v_{i}=-\frac{c_{1}}{c_{i}}v_{1}-\dots-\frac{c_{i-1}}{c_{i}}v_{i-1}-\frac{c_{i+1}}{c_{i}}v_{i+1}-\dots-\frac{c_{k}}{c_{i}}v_{k}
\]
}%
\L{%
\[
\in\spn(v_{1},\dots,v_{i-1},v_{i+1},\dots,v_{k})
\]
}%
ובפרט \L{$(v_{1},\dots,v_{k})$} ת"ל.

\section{טענה {]}אפיון ת\textquotedblright ל בעזרת ספאן סדרת הוקטורים הקודמים{[}}

יהי \L{$V$} מ"ו מעל \L{$\mathbb{F}$} ותהי \L{$(v_{1},\dots,v_{k})$}
סדרה של וקטורים מ-\L{$V$}.

אז \L{$(v_{1},\dots,v_{k})$} ת"ל אם"ם \L{%
\[
\exists_{1\leq i\leq k}\quad s.t\quad v_{i}\in\spn(v_{1},\dots,v_{i-1})
\]
}%


\subsection{הוכחה}

\L{$\Rightarrow$}: אם \L{$v_{i}\in\spn(v_{1},\dots,v_{i-1})$} אז
\L{$v_{i}\in\spn(v_{1},\dots,v_{i-1},v_{i+1},\dots,v_{k})$} וסיימנו.\\

\L{$:\Leftarrow$} מאחר ש- \L{$(v_{1},\dots,v_{k})$} ת"ל, לפי הטענה
הקודמת \textcolor{teal}{{]}אפיון ת\textquotedblright ל בעזרת צ\textquotedblright ל
מתאפס לא טריוויאלי{[}} קיימים \L{$c_{1},\dots,c_{k}\in\mathbb{F}$}
כך שלא כולם שווים ל-0 ומתקיים \L{%
\[
c_{1}v_{1}+\dots+c_{k}v_{k}=0_{V}
\]
}%
נגדיר \L{$i=\max\{1\le\ell\le k\mid c_{\ell}\neq0\}$}.

אז \L{$c_{i}\neq0$}, אבל \L{$c_{i+1}=c_{i+2}=\dots=c_{k}=0$}. מכאן
\L{$c_{1}v_{1}+\dots+c_{i}v_{i}=0_{V}$}. לכן \L{$c_{i}v_{i}=-c_{1}v_{1}-\dots-c_{i-1}v_{i-1}$}.

\L{$c_{i}\neq0$} לכן יש לו הופכי ומתקיים \L{%
\[
v_{i}=-\frac{c_{1}}{c_{i}}v_{1}-\dots-\frac{c_{i-1}}{c_{i}}v_{i-1}\in\spn(v_{1},\dots,v_{i-1})
\]
}%
וסיימנו.
\begin{flushright}
\L{$\square$}
\par\end{flushright}

\newpage{}

\noindent\fbox{\begin{minipage}[t]{1\columnwidth - 2\fboxsep - 2\fboxrule}%
\begin{center}
אלגברה ליניארית )1( - \L{$80134$} \textbar{} האונ' העברית \textbar{}
תשפ\textquotedblright ו א' \textbar{} מרצה: ד\textquotedblright ר
אלכס גורביץ' \textbar{} מתרגל: מר יואב כהן \textbar{} ליקוט ועריכה:
נעם לב
\par\end{center}%
\end{minipage}}

\part{הרצאה {\beginL 10\endL} - {\beginL 18.11.2025\endL}}

\textcolor{teal}{בתזכורת דיברנו על ההגדרה של ת\textquotedblright ל
ועל הטענה {]}אפיון ת\textquotedblright ל בעזרת צ\textquotedblright ל
מתאפס לא טריוויאלי{[}. עכשיו נמשיך עם תכונות של ת\textquotedblright ל.}

\section{תכונה {]}סדרת וקטור יחיד אינה ת\textquotedblright ל{[}}

יהיו \L{$V$} מ"ו מעל \L{$\mathbb{F}$}, \L{$0_{V}\neq v\in V$}.
אז \L{$(v)$} אינה ת"ל.

\textcolor{blue}{כי הרי \L{$v\notin\{0_{V}\}=\spn(\ )$}.}

\textcolor{teal}{משתמשים בטענה {]}אפיון ת\textquotedblright ל בעזרת
ספאן סדרת הוקטורים הקודמים{[}, סדרת הוקטורים הקודמים היא הסדרה הריקה
והספאן שלה הוא הקבוצה עם \L{$0_{V}$}.}

\section{תכונה {]}שני וקטורים הם ת\textquotedblright ל אם\textquotedblright ם
אחד מהם הוא כפולה בסקלר של השני{[}}

יהיו \L{$V$} מ"ו מעל \L{$\mathbb{F}$}, \L{$v,w\in V$} כך ש- \L{$w\neq0_{V}\quad\wedge\quad v\neq0_{V}$}.

אז \L{$(v,w)$} ת"ל אם"ם קיים \L{$0\neq c\in\mathbb{F}$} כך ש-
\L{$w=cv$}.

\subsection{הוכחה}

\L{$\Rightarrow$}: אם \L{$w=cv$} אז \L{$w\in\spn(v)$}, לכן \L{$(v,w)$}
ת"ל.\\

\L{$\Leftarrow$}: אם \L{$(v,w)$} ת"ל, אז \L{$w\in\spn(v)$} או
\L{$v\in\spn(w)$}.

אם \L{$w\in\spn(v)$}, אז קיים \L{$c\in\mathbb{F}$} כך ש- \L{$w=cv$}.

כמו כן \L{$c\neq0$} כי אחרת היינו מקבלים \L{$w=0_{V}$}.

אם \L{$v\in\spn(w)$}, אזי קיים \L{$d\in\mathbb{F}$} כך ש- \L{$v=dw$}.

כמו כן \L{$d\neq0$} כי אחרת היינו מקבלים \L{$v=0_{V}$}.

מכאן \L{$\frac{1}{d}v=w$} וברור כי \L{$\frac{1}{d}\neq0$}.

\section{טענה \textcolor{teal}{{]}חשובה{[}} {]}עמודות מטריצה מהוות סדרה ת\textquotedblright ל
אם\textquotedblright ם לממ\textquotedblright ל הומוגנית עם המטריצה
יש פתרון לא טריוויאלי{[}}

תהי \L{$A\in M_{m\times n}(\mathbb{F})$}. לכל \L{$1\leq j\leq n$}
נסמן ב- \L{$v_{j}\in\mathbb{F}^{m}$} את העמודה ה- \L{$j$} של \L{$A$}.
אז \L{$(v_{1},\dots,v_{n})$} ת\textquotedblright ל אם\textquotedblright ם
לממ\textquotedblright ל הומוגנית \L{$Ax=\begin{pmatrix}0\\
\vdots\\
0
\end{pmatrix}$} קיים פתרון לא טריוויאלי.\\

\noindent\fbox{\begin{minipage}[t]{1\columnwidth - 2\fboxsep - 2\fboxrule}%
\textcolor{teal}{כזכור כשכופלים מטריצה בוקטור )\L{$Ax$}( מקבלים צ\textquotedblright ל
של עמודות המטריצה עם ערכי הוקטור כמקדמים. }

\textcolor{teal}{אז כשאומרים ש-לממ\textquotedblright ל הומוגנית \L{$Ax=\begin{pmatrix}0\\
\vdots\\
0
\end{pmatrix}$} יש פתרון, אומרים שיש צ\textquotedblright ל שנותן את וקטור האפס. ואם
יש מקדם שהוא לא שווה לאפס קיבלנו את אפיון ת\textquotedblright ל בעזרת
צ\textquotedblright ל מתאפס לא טריוויאלי.}%
\end{minipage}}

\subsection{הוכחה}

\L{$\Leftarrow$}: אם \L{$(v_{1},\dots,v_{n})$} ת"ל, אז לפי טענה
קודמת {]}אפיון ת\textquotedblright ל בעזרת צ\textquotedblright ל מתאפס
לא טריוויאלי{[} קיימים \L{$c_{1},\dots,c_{n}\in\mathbb{F}$} כך שלא
כולם שווים ל-0 ומתקיים

\L{%
\[
c_{1}v_{1}+\dots+c_{n}v_{n}=0_{V}=\begin{pmatrix}0\\
\vdots\\
0
\end{pmatrix}
\]
}%

מכאן \L{$A\begin{pmatrix}c_{1}\\
\vdots\\
c_{n}
\end{pmatrix}=\begin{pmatrix}0\\
\vdots\\
0
\end{pmatrix}$} לכן \L{$\begin{pmatrix}c_{1}\\
\vdots\\
c_{n}
\end{pmatrix}$} הוא פתרון של \L{$Ax=\begin{pmatrix}0\\
\vdots\\
0
\end{pmatrix}$}.

בנוסף \L{$\begin{pmatrix}c_{1}\\
\vdots\\
c_{n}
\end{pmatrix}\neq\begin{pmatrix}0\\
\vdots\\
0
\end{pmatrix}$} כי \L{$c_{1},\dots,c_{n}$} לא כולם שווים ל-0 אז \L{$\begin{pmatrix}c_{1}\\
\vdots\\
c_{n}
\end{pmatrix}$} הוא לא הפתרון הטריוויאלי.\\

\L{$\Rightarrow$}: אם למערכת \L{$Ax=\begin{pmatrix}0\\
\vdots\\
0
\end{pmatrix}$} קיים פתרון לא טריוויאלי, נסמן אותו ב- \L{$\begin{pmatrix}c_{1}\\
\vdots\\
c_{n}
\end{pmatrix}$}. אז \L{$c_{1},\dots,c_{n}$} לא כולם שווים ל- 0 וגם \L{$A\begin{pmatrix}c_{1}\\
\vdots\\
c_{n}
\end{pmatrix}=\begin{pmatrix}0\\
\vdots\\
0
\end{pmatrix}$}.

אבל \L{$A\begin{pmatrix}c_{1}\\
\vdots\\
c_{n}
\end{pmatrix}=c_{1}v_{1}+\dots+c_{n}v_{n}$} אז נקבל \L{$c_{1}v_{1}+\dots+c_{n}v_{n}=\begin{pmatrix}0\\
\vdots\\
0
\end{pmatrix}=0_{V}$}. כאמור לא כולם שווים ל- {\beginL 0\endL} לכן לפי טענה קודמת {]}אפיון
ת\textquotedblright ל בעזרת צ\textquotedblright ל מתאפס לא טריוויאלי{[},
\L{$(v_{1},\dots,v_{n})$} ת"ל.

\subsection{מסקנה {]}סדרת וקטורי היחידה הסטנדרטיים אינה ת\textquotedblright ל{[}}

סדרת הוקטורים ב-\L{$\mathbb{F}^{n}$} הבאה: \L{$(e_{1},\dots,e_{n})$}
אינה ת"ל, כאשר \L{$e_{i}=\begin{pmatrix}0\\
\vdots\\
1\\
\vdots\\
0
\end{pmatrix}\in\mathbb{F}^{n}$} )ה-1 במקום ה-\L{$i$}(. כי הרי המטריצה שהעמודה ה-\L{$i$} שלה היא
\L{$e_{i}$} היא \L{$I_{n}=\begin{pmatrix}1 & \cdots & 0\\
\vdots & \ddots & \vdots\\
0 & \cdots & 1
\end{pmatrix}$} )מטריצת היחידה( ולממ"ל \L{$I_{n}x=\begin{pmatrix}0\\
\vdots\\
0
\end{pmatrix}$} קיים פתרון יחיד )הפתרון הטריוויאלי( כיוון שמדובר בממ"ל \L{$\begin{cases}
x_{1}=0\\
\vdots\\
x_{n}=0
\end{cases}$} .\\

\textcolor{blue}{סיימנו את התת-פרק על תלות לינארית.}

\section{הגדרות - פורשת, בת\textquotedblright ל ובסיס}

יהי \L{$V$} מ"ו מעל \L{$\mathbb{F}$}.

תהי \L{$(v_{1},\dots,v_{n})$} סדרה של וקטורים מ-\L{$V$}.

\L{$(v_{1},\dots,v_{n})$} נקראת \textbf{פורשת} כאשר \L{%
\[
\spn(v_{1},\dots,v_{n})=V
\]
}%
\L{$(v_{1},\dots,v_{n})$} נקראת \textbf{בלתי תלויה לינארית )בת"ל(}
כאשר \L{$(v_{1},\dots,v_{n})$} אינה ת"ל.

\L{$(v_{1},\dots,v_{n})$} נקראת \textbf{בסיס} כאשר היא פורשת ובת"ל.

\subsection{דוגמה {]}ולמעשה הגדרת הבסיס הסטנדרטי{[}}

\L{$V=\mathbb{F}^{n}$}. אז הסדרה \L{$(e_{1},\dots,e_{n})$} היא פורשת
ובת"ל ולכן היא בסיס.\\

\textcolor{teal}{\emph{פורשת}}\textcolor{teal}{{} כפי שראינו ב- מסקנה
{]}ספאן סדרת וקטורי היחידה הסטנדרטיים הוא \L{$\mathbb{F}^{n}$}{[},}

\textcolor{teal}{\emph{בת\textquotedblright ל}}\textcolor{teal}{{} כפי
שראינו ב- מסקנה {]}סדרת וקטורי היחידה הסטנדרטיים אינה ת\textquotedblright ל{[}.}

\textcolor{teal}{\L{$E=(e_{1},\dots,e_{n})$} זה הבסיס הסטנדרטי. }

\section{טענה {]}אם סדרת עמודות מטריצה היא בסיס של המרחב אז לכל ממ\textquotedblright ל
עם המטריצה יש פתרון יחיד{[}}

תהי \L{$A\in M_{m\times n}(\mathbb{F})$}.

לכל \L{$1\le j\le n$} נסמן ב-\L{$v_{j}\in\mathbb{F}^{m}$} את העמודה
ה-\L{$j$} של \L{$A$}.

אם \L{$(v_{1},\dots,v_{n})$} בסיס של \L{$\mathbb{F}^{m}$}, אז לכל
\L{$b\in\mathbb{F}^{m}$}, לממ"ל \L{$Ax=b$} קיים פתרון יחיד.

\textcolor{teal}{אלכס אמר שאפשר כאן גם אם\textquotedblright ם.}

\subsection{\textcolor{teal}{שחזור הוכחה שלי}}

נניח כי \L{$(v_{1},\dots,v_{n})$} בסיס של \L{$\mathbb{F}^{m}$}.
יהי \L{$b\in\mathbb{F}^{m}$}. נרצה להוכיח קיום של פתרון לממ\textquotedblright ל
\L{$Ax=b$} ולאחר מכן יחידות. מהגדרת בסיס הסדרה היא פורשת ובת"ל.
מכך שהיא פורשת מתקיים \L{$\spn(v_{1},\dots,v_{n})=\mathbb{F}^{m}$}.
אז בפרט \L{$b\in\spn(v_{1},\dots,v_{n})$} )כי \L{$b\in\mathbb{F}^{m}$}(.
כלומר, קיימים \L{$c_{1},\dots,c_{n}\in\mathbb{F}$} כך ש- \L{$c_{1}v_{1}+\dots+c_{n}v_{n}=b$}.
נסתכל על \L{$c=\begin{pmatrix}c_{1}\\
\vdots\\
c_{n}
\end{pmatrix}$}. נשים לב שמתקיים \L{%
\[
Ac=c_{1}v_{1}+\dots+c_{n}v_{n}=b
\]
}%
ובכך הוכחנו קיום. נניח בשלילה שקיים לממ\textquotedblright ל \L{$Ax=b$}
פתרון נוסף. אזי קיים \L{$d=\begin{pmatrix}d_{1}\\
\vdots\\
d_{n}
\end{pmatrix}$} כך ש- \L{$Ad=b$} וגם \L{$\exists_{1\le i\le n}\quad d_{i}\neq c_{i}$}.

נקבל \L{%
\[
Ac-Ad=b-b=0_{V}
\]
}%
כלומר \L{%
\[
c_{1}v_{1}+\dots+c_{n}v_{n}-(d_{1}v_{1}+\dots+d_{n}v_{n})=0_{V}
\]
}%
 כלומר \L{%
\[
(c_{1}-d_{1})v_{1}+\dots+(c_{n}-d_{n})v_{n}=0_{V}
\]
}%
 אבל \L{$d_{i}\neq c_{i}$} לכן \L{$c_{i}-d_{i}\neq0$} ולפי הטענה
האומרת שעמודות מטריצה מהוות סדרה ת"ל אם"ם לממ\textquotedblright ל
ההומוגנית עם המטריצה יש פתרון לא-טריוויאלי, נקבל שהסדרה \L{$(v_{1},\dots,v_{n})$}
ת"ל וזאת סתירה, בכך הוכחנו יחידות.
\begin{flushright}
\L{$\square$}
\par\end{flushright}

\subsection{הוכחה}

מאחר ש- \L{$(v_{1},\dots,v_{n})$} פורשת מתקיים \L{%
\[
\spn(v_{1},\dots,v_{n})=\mathbb{F}^{m}
\]
}%
לכן לפי הטענה מאתמול \textcolor{teal}{)הניחוש שלי הוא שאלכס מתייחס
כאן לטענה {]}ספאן עמודות של מטריצה{[}(} לממ"ל \L{$Ax=b$} קיים פתרון.
נניח כי \L{$d,h\in\mathbb{F}^{n}$} הם פתרונות של הממ"ל \L{$Ax=b$}.
אז \L{$Ad=b$} וגם \L{$Ah=b$}. מכאן \L{%
\[
A(d-h)=Ad-Ah=b-b=\begin{pmatrix}0\\
\vdots\\
0
\end{pmatrix}
\]
}%
כלומר \L{$d-h$} הוא פתרון של הממ"ל \L{$Ax=\begin{pmatrix}0\\
\vdots\\
0
\end{pmatrix}$}.

מאחר ש- \L{$(v_{1},\dots,v_{n})$} בת"ל לממ"ל \L{$Ax=\begin{pmatrix}0\\
\vdots\\
0
\end{pmatrix}$} אין פתרון לא טריוויאלי, לכן \L{$(d-h)$} הוא הפתרון הטריוויאלי, כלומר
\L{$(d-h)=\begin{pmatrix}0\\
\vdots\\
0
\end{pmatrix}$} מכאן \L{$d=h$}.
\begin{flushright}
\L{$\square$}
\par\end{flushright}

\section{טענה {]}יחידות ההצגה לפי בסיס{[}}

יהי \L{$V$} מ"ו מעל \L{$\mathbb{F}$} ויהי \L{$(v_{1},\dots,v_{n})$}
בסיס של \L{$V$}. יהי \L{$v\in V$}.

אז קיימים \L{$c_{1},\dots,c_{n}\in\mathbb{F}$} יחידים כך ש- \L{%
\[
v=c_{1}v_{1}+\dots+c_{n}v_{n}
\]
}%
דהיינו, כל וקטור במ"ו אפשר להציג כצירוף לינארי של וקטורי הבסיס בדרך
יחידה.

\subsection{הוכחה}

\textcolor{teal}{אותו עיקרון של ההוכחה הקודמת.}\\

מאחר ש- \L{$(v_{1},\dots,v_{n})$} פורשת מתקיים \L{$\spn(v_{1},\dots,v_{n})=V$}.

לכן \L{$v\in\spn(v_{1},\dots,v_{n})$} ומכאן קיימים \L{$c_{1},\dots,c_{n}\in\mathbb{F}$}
כך ש- \L{%
\[
v=c_{1}v_{1}+\dots+c_{n}v_{n}
\]
}%

\textcolor{teal}{)הוכחנו קיום(}

נניח כי \L{$d_{1},\dots,d_{n}\in\mathbb{F}$} כך ש- \L{%
\[
v=d_{1}v_{1}+\dots+d_{n}v_{n}
\]
}%
מכאן \L{%
\[
c_{1}v_{1}+\dots+c_{n}v_{n}=d_{1}v_{1}+\dots+d_{n}v_{n}
\]
}%
לכן \L{%
\[
c_{1}v_{1}+\dots+c_{n}v_{n}-d_{1}v_{1}-\dots-d_{n}v_{n}=0_{V}
\]
}%
נקבל \L{%
\[
(c_{1}-d_{1})v_{1}+(c_{2}-d_{2})v_{2}+\dots+(c_{n}-d_{n})v_{n}=0_{V}
\]
}%
מאחר ש- \L{$(v_{1},\dots,v_{n})$} בת"ל, אנו מקבלים \textcolor{teal}{לפי
טענה {]}אפיון ת\textquotedblright ל בעזרת צ\textquotedblright ל מתאפס
לא טריוויאלי{[}}

\L{%
\[
\begin{array}{c}
c_{1}-d_{1}=0\\
\vdots\\
c_{n}-d_{n}=0
\end{array}
\]
}%

מכאן\L{%
\[
\begin{array}{c}
c_{1}=d_{1}\\
\vdots\\
c_{n}=d_{n}
\end{array}
\]
}%

\textcolor{teal}{)הוכחנו יחידות(} וסיימנו.
\begin{flushright}
\L{$\square$}
\par\end{flushright}

\newpage{}

\noindent\fbox{\begin{minipage}[t]{1\columnwidth - 2\fboxsep - 2\fboxrule}%
\begin{center}
אלגברה ליניארית )1( - \L{$80134$} \textbar{} האונ' העברית \textbar{}
תשפ\textquotedblright ו א' \textbar{} מרצה: ד\textquotedblright ר
אלכס גורביץ' \textbar{} מתרגל: מר יואב כהן \textbar{} ליקוט ועריכה:
נעם לב
\par\end{center}%
\end{minipage}}

\part{הרצאה {\beginL 11\endL} - {\beginL 24.11.2025\endL}}

\textcolor{teal}{בתזכורת דיברנו על טענה {]}יחידות ההצגה לפי בסיס{[}.}\textcolor{blue}{{} }

\section{הגדרה - וקטור הקואורדינטות}

יהיו \L{$V$} מ"ו מעל \L{$\mathbb{F}$}, \L{$\mathcal{B}=(v_{1},\dots,v_{n})$}
בסיס של \L{$V$}, \L{$v\in V$}.

יהיו \L{$c_{1},\dots,c_{n}\in\mathbb{F}$} כמו בטענה \textcolor{teal}{)היחידים
כך ש- \L{$c_{1}v_{1}+\dots+c_{n}v_{n}=v$}(}.

ה-\L{$n$}-יה \L{$\begin{pmatrix}c_{1}\\
\vdots\\
c_{n}
\end{pmatrix}\in\mathbb{F}^{n}$} נקראת \L{$n$}-ית הקואורדינטות של \L{$v$} ביחס לבסיס \L{$\mathcal{B}$}
ומסומנת ב- \L{$[v]_{\mathcal{B}}$}.

\section{טענה {]}וקטור הקואורדינטות משמר חיבור וכפל בסקלר{[}}

יהיו \L{$V$} מ"ו מעל \L{$\mathbb{F}$}, \L{$\mathcal{B}=(v_{1},\dots,v_{n})$}
בסיס של \L{$V$}, \L{$c\in\mathbb{F},\quad v,w\in V$}. אז\\

{\beginL 1\endL}( \hspace*{\fill}\L{%
\[
[v+w]_{\mathcal{B}}=[v]_{\mathcal{B}}+[w]_{\mathcal{B}}
\]
}%
\\

{\beginL 2\endL}( \hspace*{\fill}\L{%
\[
[c\cdot v]_{\mathcal{B}}=c\cdot[v]_{\mathcal{B}}
\]
}%


\subsection{הוכחה}

1( נסמן \L{%
\[
[v]_{\mathcal{B}}=\begin{pmatrix}c_{1}\\
\vdots\\
c_{n}
\end{pmatrix},\quad[w]_{\mathcal{B}}=\begin{pmatrix}d_{1}\\
\vdots\\
d_{n}
\end{pmatrix}
\]
}%
אז 

\L{%
\[
v=c_{1}v_{1}+\dots+c_{n}v_{n}
\]
}%
\L{%
\[
w=d_{1}v_{1}+\dots+d_{n}v_{n}
\]
}%
מכאן \L{%
\[
v+w=c_{1}v_{1}+\dots+c_{n}v_{n}+d_{1}v_{1}+\dots+d_{n}v_{n}=
\]
}%
\L{%
\[
(c_{1}+d_{1})v_{1}+\dots+(c_{n}+d_{n})v_{n}
\]
}%
לפיכך \L{%
\[
[v+w]_{\mathcal{B}}=\begin{pmatrix}c_{1}+d_{1}\\
\vdots\\
c_{n}+d_{n}
\end{pmatrix}=\begin{pmatrix}c_{1}\\
\vdots\\
c_{n}
\end{pmatrix}+\begin{pmatrix}d_{1}\\
\vdots\\
d_{n}
\end{pmatrix}
\]
}%
\L{%
\[
=[v]_{\mathcal{B}}+[w]_{\mathcal{B}}
\]
}%

2( נסמן \L{$[v]_{\mathcal{B}}=\begin{pmatrix}c_{1}\\
\vdots\\
c_{n}
\end{pmatrix}$}. אז

\L{%
\[
v=c_{1}v_{1}+\dots+c_{n}v_{n}
\]
}%

מכאן \L{%
\[
c\cdot v=c\cdot(c_{1}v_{1}+\dots+c_{n}v_{n})=
\]
}%
\L{%
\[
(c\cdot c_{1})v_{1}+\dots+(c\cdot c_{n})v_{n}
\]
}%
לפיכך \L{%
\[
[c\cdot v]_{\mathcal{B}}=\begin{pmatrix}c\cdot c_{1}\\
\vdots\\
c\cdot c_{n}
\end{pmatrix}=c\cdot\begin{pmatrix}c_{1}\\
\vdots\\
c_{n}
\end{pmatrix}=c\cdot[v]_{\mathcal{B}}
\]
}%

\begin{flushright}
\L{$\square$}
\par\end{flushright}

\section{הבחנה {]}איך מחשבים וקטור קואורדינטות לפי בסיס{[}}

אם \L{$\mathcal{B}=(v_{1},\dots,v_{n})$} בסיס של \L{$\mathbb{F}^{m}$},

\L{$A$} היא מטריצה שעמודותיה הן \L{$v_{1},\dots,v_{n}$} )וקטורי
הבסיס(, ו- \L{$b\in\mathbb{F}^{m}$}.

אז \L{$[b]_{\mathcal{B}}$} הוא הפתרון היחיד של הממ"ל \L{$Ax=b$}.\\

\textcolor{violet}{במילים פשוטות: כדי למצוא את וקטור הקואורדינטות
של \L{$b$} לפי בסיס \L{$\mathcal{B}$}, אנחנו פשוט צריכים לפתור ממ\textquotedblright ל
שבה עמודות המטריצה הן וקטורי הבסיס.}

\textcolor{teal}{כזכור, \L{$Ax$} הוא צ\textquotedblright ל של עמודות
המטריצה, והוכחנו שיש רק צ\textquotedblright ל אחד כזה לפי הבסיס עבור
כל וקטור )בטענה {]}יחידות ההצגה לפי בסיס{[}(.}

\section{תכונה {]}וקטור הקואורדינטות לפי הבסיס הסטנדרטי שווה לעצמו{[}}

יהי \L{$E=(e_{1},\dots,e_{n})$} הבסיס הסטנדרטי של \L{$\mathbb{F}^{n}$},
ויהי \L{$v\in\mathbb{F}^{n}$}. אז \L{%
\[
[v]_{E}=v
\]
}%


\subsection{דוגמה}

\L{%
\[
V=\mathbb{R}^{2},\quad B=\left(\begin{pmatrix}1\\
0
\end{pmatrix},\begin{pmatrix}0\\
1
\end{pmatrix}\right),\quad v=\begin{pmatrix}1\\
2
\end{pmatrix}
\]
}%

\L{%
\[
\begin{pmatrix}1\\
2
\end{pmatrix}=1\begin{pmatrix}1\\
0
\end{pmatrix}+2\begin{pmatrix}0\\
1
\end{pmatrix}\quad\therefore\quad[v]_{\mathcal{B}}=\begin{pmatrix}1\\
2
\end{pmatrix}
\]
}%

{\small\textcolor{teal}{\emph{\L{$\therefore$} זה הסימון של \textquotedblleft ולכן\textquotedblright .}}}{\small\par}

\section{תכונה {]}שינוי סדר הקואורדינטות מתאים לשינוי סדר וקטורים בבסיס{[}}

אם בסיס \L{$\mathcal{C}$} של מ"ו \L{$V$} מתקבל מבסיס \L{$\mathcal{B}$}
ע"י שינוי סדר הוקטורים, אז לכל \L{$v\in V$}, \L{$[v]_{\mathcal{C}}$}
מתקבל מ- \L{$[v]_{\mathcal{B}}$} ע"י אותו שינוי הסדר של הקואורדינטות.

\subsection{דוגמה}

{\small\textcolor{teal}{\emph{הדוגמה מתקשרת לדוגמה הקודמת.}}}{\small\par}

\L{%
\[
V=\mathbb{R}^{2},\quad B=\left(\begin{pmatrix}0\\
1
\end{pmatrix},\begin{pmatrix}1\\
0
\end{pmatrix}\right),\quad v=\begin{pmatrix}1\\
2
\end{pmatrix}
\]
}%

\L{%
\[
\begin{pmatrix}1\\
2
\end{pmatrix}=2\begin{pmatrix}0\\
1
\end{pmatrix}+1\begin{pmatrix}1\\
0
\end{pmatrix}\quad\therefore\quad[v]_{\mathcal{B}}=\begin{pmatrix}2\\
1
\end{pmatrix}
\]
}%


\section{תכונה {]}וקטורי הקואורדינטות של וקטורי הבסיס{[}}

אם \L{$\mathcal{B}=(v_{1},\dots,v_{n})$} בסיס של מ"ו \L{$V$}, אז
\L{%
\[
\forall_{1\leq i\leq n}\quad[v_{i}]_{\mathcal{B}}=e_{i}=\begin{pmatrix}0\\
\vdots\\
1\\
\vdots\\
0
\end{pmatrix}\leftarrow i
\]
}%


\subsection{דוגמה}

\L{%
\[
V=\mathbb{R}^{2},\quad B=\left(\begin{pmatrix}1\\
1
\end{pmatrix},\begin{pmatrix}1\\
2
\end{pmatrix}\right),\quad v=\begin{pmatrix}1\\
2
\end{pmatrix}
\]
}%

\textcolor{blue}{ניתן לבדוק כי \L{$\mathcal{B}$} אכן בסיס של \L{$\mathbb{R}^{2}$}.}

\L{%
\[
\begin{pmatrix}1\\
2
\end{pmatrix}=0\cdot\begin{pmatrix}1\\
1
\end{pmatrix}+1\cdot\begin{pmatrix}1\\
2
\end{pmatrix}\quad\therefore\quad[v]_{\mathcal{B}}=\begin{pmatrix}0\\
1
\end{pmatrix}
\]
}%

\textcolor{violet}{למה זו דוגמה לתכונה? כי הווקטור \L{$v$} בדוגמה
הוא האיבר השני בבסיס \L{$\mathcal{B}$}, והחישוב מראה שווקטור הקואורדינטות
שלו הוא אכן \L{$e_{2}$} )הווקטור \L{$\begin{pmatrix}0\\
1
\end{pmatrix}$}(, בדיוק כפי שהתכונה מתארת עבור \L{$i=2$}.}

\section{המשפט המרכזי}
יהי \L{$V$} מ"ו מעל \L{$\mathbb{F}$}, ויהיו \L{$v_{1},\dots,v_{m},w_{1},\dots,w_{n}\in V$}
כך ש- \L{%
\[
w_{1},\dots,w_{n}\in\spn(v_{1},\dots,v_{m})
\]
}%
אם \L{$m<n$}, אז \L{$(w_{1},\dots,w_{n})$} ת"ל.

\subsection{הוכחה}

לכל \L{$1\le j\le n$} קיימים \L{$a_{1j},\dots,a_{mj}$} כך ש- \L{%
\[
w_{j}=a_{1j}v_{1}+\dots+a_{mj}v_{m}
\]
}%
נסמן \L{%
\[
A=\begin{pmatrix}a_{11} & \dots & a_{1n}\\
\vdots &  & \vdots\\
a_{m1} & \dots & a_{mn}
\end{pmatrix}
\]
}%
אז \L{$A\in M_{m\times n}(\mathbb{F})$}.

מאחר ש- \L{$m<n$} לממ"ל ההומוגנית \L{%
\[
A\begin{pmatrix}x_{1}\\
\vdots\\
x_{n}
\end{pmatrix}=\begin{pmatrix}0\\
\vdots\\
0
\end{pmatrix}
\]
}%
קיים פתרון לא טריוויאלי, נסמן אותו ב- \L{$\begin{pmatrix}c_{1}\\
\vdots\\
c_{n}
\end{pmatrix}$}.

בפרט לכל \L{$1\le i\le m$} מתקיים \L{%
\[
a_{i1}c_{1}+\dots+a_{in}c_{n}=0
\]
}%
נתבונן בצירוף לינארי לא טריוויאלי \L{$c_{1}w_{1}+\dots+c_{n}w_{n}$}
ונראה שהוא שווה ל- \L{$0_{V}$}. אכן \L{%
\[
c_{1}w_{1}+\dots+c_{n}w_{n}=
\]
}%
\L{%
\[
c_{1}a_{11}v_{1}+c_{1}a_{21}v_{2}+\dots+c_{1}a_{m1}v_{m}+
\]
}%
\L{%
\[
c_{2}a_{12}v_{1}+c_{2}a_{22}v_{2}+\dots+c_{2}a_{m2}v_{m}+
\]
}%
\L{%
\[
\dots+c_{n}a_{1n}v_{1}+c_{n}a_{2n}v_{2}+\dots+c_{n}a_{mn}v_{m}=
\]
}%
\L{%
\[
(a_{11}c_{1}+a_{12}c_{2}+\dots+a_{1n}c_{n})v_{1}+
\]
}%
\L{%
\[
(a_{21}c_{1}+a_{22}c_{2}+\dots+a_{2n}c_{n})v_{2}+
\]
}%
\L{%
\[
\vdots
\]
}%
\L{%
\[
+(a_{m1}c_{1}+a_{m2}c_{2}+\dots+a_{mn}c_{n})v_{m}=
\]
}%
\L{%
\[
0v_{1}+0v_{2}+\dots+0v_{m}=0_{V}
\]
}%
לפיכך \L{$(w_{1},\dots,w_{n})$} ת"ל.
\begin{flushright}
\L{$\square$}
\par\end{flushright}

\subsection{\textcolor{teal}{שחזור הוכחה שלי}}

יהי \L{$V$} מ"ו מעל \L{$\mathbb{F}$} ויהיו \L{%
\[
v_{1},\dots,v_{m},w_{1},\dots,w_{n}\in V
\]
}%
כך ש- \L{$w_{1},\dots,w_{n}\in\spn(v_{1},\dots,v_{m})$} ונניח ש-
\L{$m<n$}. אז קיימים \L{$a_{1},\dots,a_{n}\in\mathbb{F}^{m}$} כך
ש- \L{%
\[
w_{1}=a_{11}v_{1}+a_{12}v_{2}+\dots+a_{1m}v_{m}
\]
}%
\L{%
\[
w_{2}=a_{21}v_{1}+a_{22}v_{2}+\dots+a_{2m}v_{m}
\]
}%
\L{%
\[
\vdots
\]
}%
\L{%
\[
w_{m}=a_{m1}v_{1}+a_{m2}v_{2}+\dots+a_{mm}v_{m}
\]
}%
\L{%
\[
\vdots
\]
}%
\L{%
\[
w_{n}=a_{n1}v_{1}+a_{n2}v_{2}+\dots+a_{nm}v_{m}
\]
}%
לפי טענה שלמדנו, \L{$w_{1},\dots,w_{n}$} ת"ל, אם"ם קיים פתרון לא
טריוויאלי למשוואה \L{$x_{1}w_{1}+\dots+x_{n}w_{n}=0_{V}$}. נציב את
הממ"ל הקודם ונקבל שזה מתקיים אם\textquotedblright ם קיים פתרון לא
טריוויאלי למשוואה:\L{%
\[
x_{1}a_{11}v_{1}+x_{1}a_{12}v_{2}+\dots+x_{1}a_{1m}v_{m}+
\]
}%
\L{%
\[
x_{2}a_{21}v_{1}+x_{2}a_{22}v_{2}+\dots+x_{2}a_{2m}v_{m}+
\]
}%
\L{%
\[
\vdots
\]
}%
\L{%
\[
x_{m}a_{m1}v_{1}+x_{m}a_{m2}v_{2}+\dots+x_{m}a_{mm}v_{m}+
\]
}%
\L{%
\[
\vdots
\]
}%
\L{%
\[
x_{n}a_{n1}v_{1}+x_{n}a_{n2}v_{2}+\dots+x_{n}a_{nm}v_{m}=0_{V}
\]
}%
נקבץ איברים ונקבל \L{%
\[
(x_{1}a_{11}+x_{2}a_{21}+\dots+x_{m}a_{m1}+\dots+x_{n}a_{n1})v_{1}+
\]
}%
\L{%
\[
(x_{1}a_{12}+x_{2}a_{22}+\dots+x_{m}a_{m2}+\dots+x_{n}a_{n2})v_{2}+
\]
}%
\L{%
\[
\vdots
\]
}%
\L{%
\[
(x_{1}a_{1m}+x_{2}a_{2m}+\dots+x_{m}a_{mm}+\dots+x_{n}a_{nm})v_{m}=0_{V}
\]
}%
וזה בפרט מתקיים אם קיים פתרון לא טריוויאלי לממ"ל \L{%
\[
x_{1}a_{11}+x_{2}a_{21}+\dots+x_{m}a_{m1}+\dots+x_{n}a_{n1}=0
\]
}%
\L{%
\[
x_{1}a_{12}+x_{2}a_{22}+\dots+x_{m}a_{m2}+\dots+x_{n}a_{n2}=0
\]
}%
\L{%
\[
\vdots
\]
}%
\L{%
\[
x_{1}a_{1m}+x_{2}a_{2m}+\dots+x_{m}a_{mm}+\dots+x_{n}a_{nm}=0
\]
}%
שמיוצגת ע"י המטריצה המורחבת \L{%
\[
\left(\begin{array}{cccccc|c}
a_{11} & a_{21} & \dots & a_{m1} & \dots & a_{n1} & 0\\
a_{12} & a_{22} & \dots & a_{m2} & \dots & a_{n2} & 0\\
\vdots &  &  &  &  &  & \vdots\\
a_{1m} & a_{2m} & \dots & a_{mm} & \dots & a_{nm} & 0
\end{array}\right)
\]
}%
ונתון כי \L{$n>m$} כלומר בממ"ל ההומוגנית הזאת יש יותר נעלמים מאשר
משוואות ולפי מסקנה שלמדנו יש לה פתרון לא-טריוויאלי ולכן כפי שרצינו
\L{$w_{1},\dots,w_{n}$} ת"ל וסיימנו. 
\begin{flushright}
\L{$\square$}
\par\end{flushright}

\section{טענה {]}לבסיסים של אותו מ\textquotedblright ו יש מספר זהה של וקטורים{[}}

יהי \L{$V$} מ"ו מעל שדה \L{$\mathbb{F}$}, יהיו \L{$(v_{1},\dots,v_{m})$}
, \L{$(w_{1},\dots,w_{n})$} בסיסים של \L{$V$}. אז \L{$m=n$}.

\subsection{הוכחה}

נניח בשלילה כי \L{$m\neq n$}. בלי הגבלת הכלליות ניתן להניח כי \L{$m<n$}.

מאחר ש- \L{$(v_{1},\dots,v_{m})$} בסיס מתקיים \L{$\spn(v_{1},\dots,v_{m})=V$}.

מכאן \L{$w_{1},\dots,w_{n}\in\spn(v_{1},\dots,v_{m})$}. מאחר ש- \L{$m<n$},
מהמשפט המרכזי נובע כי \L{$(w_{1},\dots,w_{n})$} ת"ל בסתירה לכך שהיא
בסיס.

\section{הגדרות - מ\textquotedblright ו נוצר סופית, מימד}

מ"ו \L{$V$} מעל \L{$\mathbb{F}$} נקרא \textbf{נוצר סופית} )נ\textquotedblright ס(
כאשר יש לו בסיס סופי.\\

יהי \L{$V$} מ"ו נוצר סופית. \textbf{המימד} של \L{$V$} זהו מספר
הוקטורים בבסיס של \L{$V$}.

בזכות הטענה שהרגע הוכחנו, המימד מוגדר היטב )כלומר לא תלוי בבחירת הבסיס
הספציפי, לכן הוא יחיד עבור כל מ\textquotedblright ו(.

המימד של \L{$V$} מסומן ב- \L{$\dim V$}.\\

\textcolor{blue}{נשים לב - המושג מימד מוגדר רק על מ\textquotedblright ו
שהוא נ\textquotedblright ס. }

\subsection{דוגמה {]}ולמעשה עובדה חשובה - \L{$\dim\mathbb{F}^{n}=n$}{[}}

מאחר ש- \L{$E=(e_{1},\dots,e_{n})$} הוא תמיד בסיס של \L{$\mathbb{F}^{n}$}
\textcolor{teal}{)ראו - דוגמה {]}ולמעשה הגדרת הבסיס הסטנדרטי{[}(}
מתקיים \L{$\dim\mathbb{F}^{n}=n$}.

\newpage{}

\noindent\fbox{\begin{minipage}[t]{1\columnwidth - 2\fboxsep - 2\fboxrule}%
\begin{center}
אלגברה ליניארית )1( - \L{$80134$} \textbar{} האונ' העברית \textbar{}
תשפ\textquotedblright ו א' \textbar{} מרצה: ד\textquotedblright ר
אלכס גורביץ' \textbar{} מתרגל: מר יואב כהן \textbar{} ליקוט ועריכה:
נעם לב
\par\end{center}%
\end{minipage}}

\part{הרצאה {\beginL 12\endL} - {\beginL 25.11.2025\endL}}

\textcolor{teal}{בתזכורת עברנו על המשפט המרכזי, טענה {]}לבסיסים של
אותו מ\textquotedblright ו יש מספר זהה של וקטורים{[} והגדרות - מ\textquotedblright ו
נוצר סופית, מימד. }

\section{דוגמאות נוספות למימד}

{\beginL 1\endL}( \L{$V=M_{m\times n}(\mathbb{F})$}. לכל \L{$1\leq i\leq m$}
ו- \L{$1\leq j\leq n$} נגדיר מטריצה \L{$E_{ij}\in M_{m\times n}(\mathbb{F})$}
באופן הבא:

בחיתוך של השורה ה-\L{$i$} והעמודה ה-\L{$j$} של \L{$E_{ij}$} מופיע
{\beginL 1\endL}, בשאר המקומות מופיע {\beginL 0\endL}.

אז הסדרה \L{$(E_{11},E_{12},\dots,E_{1n},E_{21},E_{22},\dots,E_{2n},\dots,E_{m1},E_{m2},\dots,E_{mn})$}
מהווה בסיס של \L{$M_{m\times n}(\mathbb{F})$} לכן \L{%
\[
\dim M_{m\times n}(\mathbb{F})=mn
\]
}%

\textcolor{violet}{זה הבסיס הסטנדרטי של מרחב המטריצות.}\\

{\beginL 2\endL}( \L{%
\[
V=\mathbb{F}[x]_{\le n}=\{P\in\mathbb{F}[x]\mid\deg P\le n\}
\]
}%

הסדרה \L{$(1,x,x^{2},\dots,x^{n})$} מהווה בסיס של \L{$\mathbb{F}[x]_{\le n}$}
\textcolor{violet}{למעשה} \textcolor{violet}{הבסיס הסטנדרטי}.

אכן, אם \L{$P\in\mathbb{F}[x]_{\le n}$} אז קיים \L{$k\le n$} כך
ש- \L{$P=a_{k}x^{k}+\dots+a_{1}x+a_{0}$}. אז \L{%
\[
P=a_{0}\cdot1+a_{1}\cdot x+\dots+a_{k}\cdot x^{k}\in\spn(1,x,\dots,x^{n})
\]
}%
לכן \L{$(1,x,x^{2},\dots,x^{n})$} פורשת את \L{$\mathbb{F}[x]_{\le n}$}.

כמו כן, \textbf{אם} \L{$c_{0},c_{1},\dots,c_{n}\in\mathbb{F}$} הם
סקלרים כך ש- \L{%
\[
c_{0}\cdot1+c_{1}\cdot x+\dots+c_{n}x^{n}=0_{V}
\]
}%
\textbf{אז} \L{$c_{0}=c_{1}=\dots=c_{n}=0$}. 

\textcolor{violet}{זה נובע מהגדרת שוויון בין פולינומים: הביטוי שווה
לפולינום האפס )\L{$0_{V}$}(, ופולינום שווה זהותית לאפס אם ורק אם
כל המקדמים )הסקלרים \L{$c_{i}$}( שלו הם אפס.}

לפיכך \L{$(1,x,x^{2},\dots,x^{n})$} היא בת\textquotedblright ל ולכן
\L{$\dim\mathbb{F}[x]_{\le n}=n+1$}.

\section{טענה {]}אם יש בסדרה יותר וקטורים מלבסיס היא ת\textquotedblright ל
ואם פחות היא לא פורשת{[}}

יהי \L{$V$} מ"ו נ"ס מעל \L{$\mathbb{F}$} כך ש- \L{$\dim V=m$}.

תהי \L{$(v_{1},\dots,v_{n})$} סדרה כלשהי של וקטורים מ-\L{$V$}. אזי: 
\begin{enumerate}
\item אם \L{$m<n$} אז \L{$(v_{1},\dots,v_{n})$} ת"ל. 
\item אם \L{$m>n$} אז \L{$(v_{1},\dots,v_{n})$} אינה פורשת. 
\end{enumerate}

\subsection{הוכחה}

קיים בסיס \L{$(w_{1},\dots,w_{m})$} של \L{$V$}. 
\begin{enumerate}
\item אם \L{$m<n$} מאחר ש- \L{$v_{1},\dots,v_{n}\in\spn(w_{1},\dots,w_{m})$}
נובע מהמשפט המרכזי כי \L{$(v_{1},\dots,v_{n})$} ת"ל. 
\item נניח בשלילה כי \L{$(v_{1},\dots,v_{n})$} פורשת, כלומר \L{%
\[
\spn(v_{1},\dots,v_{n})=V
\]
}%
אם \L{$n<m$}, מאחר ש- \L{$w_{1},\dots,w_{m}\in\spn(v_{1},\dots,v_{n})$}
נובע מהמשפט המרכזי כי \L{$(w_{1},\dots,w_{m})$} ת"ל, בסתירה לכך
ש- \L{$(w_{1},\dots,w_{m})$} בסיס. לפיכך \L{$(v_{1},\dots,v_{n})$}
אינה פורשת. 
\end{enumerate}

\subsection{הערה {]}מה אם \L{$m=n$}?{[}}

אם בהנחות הטענה מתקיים \L{$m=n$}, הסדרה \L{$(v_{1},\dots,v_{n})$}
לא חייבת להיות בסיס של \L{$V$}.

לדוגמה:\L{%
\[
V=\mathbb{R}^{2},\quad\dim\mathbb{R}^{2}=2,\quad S=\left(\begin{pmatrix}1\\
1
\end{pmatrix},\begin{pmatrix}2\\
2
\end{pmatrix}\right)
\]
}%

אז יש ב- \L{$S$} {\beginL 2\endL} וקטורים אבל \L{$\spn S\neq\mathbb{R}^{2}$}
וגם \L{$S$} ת"ל.

\section{למה 1 {]}הסרת וקטור תלוי מסדרה פורשת{[}}

יהי \L{$V$} מ"ו מעל \L{$\mathbb{F}$}, \L{$(v_{1},\dots,v_{n})$}
סדרה פורשת של וקטורים מ-\L{$V$} כך ש- \L{$v_{n}\in\spn(v_{1},\dots,v_{n-1})$}.
אזי גם \L{$(v_{1},\dots,v_{n-1})$} פורשת.

\subsection{הוכחה}

לכל \L{$1\le i\le n-1$} מתקיים \L{$v_{i}\in\spn(v_{1},\dots,v_{n-1})$}.

בנוסף \L{$v_{n}\in\spn(v_{1},\dots,v_{n-1})$}, כלומר

\L{%
\[
v_{1},\dots,v_{n-1},v_{n}\in\spn(v_{1},\dots,v_{n-1})
\]
}%

מאחר ש- \L{$\spn(v_{1},\dots,v_{n-1})$} הוא תמ\textquotedblright ו
מתקיים \L{%
\[
V=\spn(v_{1},\dots,v_{n})\subseteq\spn(v_{1},\dots,v_{n-1})
\]
}%
לפיכך \L{$(v_{1},\dots,v_{n-1})$} פורשת.

\section{למה 2 {]}הוספת וקטור לא תלוי לסדרה בת\textquotedblright ל{[}}

יהי \L{$V$} מ"ו מעל \L{$\mathbb{F}$}. תהי \L{$(v_{1},\dots,v_{n})$}
סדרה בת"ל של וקטורים מ-\L{$V$} ויהי \L{$v_{n+1}\in V$} כך ש- \L{$v_{n+1}\notin\spn(v_{1},\dots,v_{n})$}.
אז גם \L{$(v_{1},\dots,v_{n},v_{n+1})$} בת"ל.

\subsection{הוכחה}

נניח בשלילה כי \L{$(v_{1},\dots,v_{n},v_{n+1})$} ת"ל. אז לפי טענה
{]}אפיון ת\textquotedblright ל בעזרת ספאן סדרת הוקטורים הקודמים{[},
קיים \L{$1\le i\le n+1$} כך ש- \L{$v_{i}\in\spn(v_{1},\dots,v_{i-1})$}.

אם \L{$i\le n$}, אז \L{$(v_{1},\dots,v_{n})$} ת\textquotedblright ל
וזו סתירה )לנתון(. 

אם \L{$i=n+1$} אז \L{$v_{n+1}\in\spn(v_{1},\dots,v_{n})$} וזו סתירה
)לנתון(.

לכן \L{$(v_{1},\dots,v_{n+1})$} בת"ל.

\section{טענה {]}סדרה עם מספר וקטורים כלבסיס היא פורשת \L{$\text{\ensuremath{\protect\iff}}$}
בת\textquotedblright ל{[}}

יהי \L{$V$} מ"ו נ"ס מעל \L{$\mathbb{F}$} כך ש- \L{$\dim V=n$}.

תהי \L{$(v_{1},\dots,v_{n})$} סדרה כלשהי של וקטורים מ-\L{$V$}.

אז \L{$(v_{1},\dots,v_{n})$} פורשת אם"ם \L{$(v_{1},\dots,v_{n})$}
בת"ל.

\subsection{הוכחה}

\L{$\Leftarrow$}: נניח כי \L{$(v_{1},\dots,v_{n})$} פורשת.

נניח בשלילה כי \L{$(v_{1},\dots,v_{n})$} ת"ל.

אז קיים \L{$1\le i\le n$} כך ש-\L{%
\[
v_{i}\in\spn(v_{1},\dots,v_{i-1},v_{i+1},\dots,v_{n})
\]
}%
בלי הגבלת הכלליות \L{$i=n$}, כלומר\L{%
\[
v_{n}\in\spn(v_{1},\dots,v_{n-1})
\]
}%
 אז לפי למה 1 גם \L{$(v_{1},\dots,v_{n-1})$} פורשת.

מאחר ש- \L{$\dim V=n>n-1$} לפי הטענה הקודמת \L{$(v_{1},\dots,v_{n-1})$}
אינה פורשת והגענו לסתירה.

\L{$\Rightarrow$}: נניח כי \L{$(v_{1},\dots,v_{n})$} בת"ל.

נניח בשלילה כי \L{$(v_{1},\dots,v_{n})$} אינה פורשת.

כלומר קיים \L{$v\in V$} כך ש- \L{$v\notin\spn(v_{1},\dots,v_{n})$}. 

אז לפי למה 2 גם \L{$(v_{1},\dots,v_{n},v)$} בת"ל.

מאחר ש- \L{$\dim V=n<n+1$}, לפי הטענה הקודמת \L{$(v_{1},\dots,v_{n},v)$}
ת"ל והגענו לסתירה.
\begin{flushright}
\L{$\square$}
\par\end{flushright}

\section{טענה {]}כל סדרה פורשת מכילה תת-סדרה שהיא בסיס{[}}

יהי \L{$V$} מ"ו נ"ס מעל \L{$\mathbb{F}$} כך ש \L{$\dim V=n$}.

תהי \L{$(v_{1},\dots,v_{m})$} סדרה פורשת של וקטורים מ-\L{$V$}.

אז \L{$n\le m$} וקיימת תת-סדרה \L{$(v_{i_{1}},\dots,v_{i_{n}})$}
של סדרה זו אשר מהווה בסיס של \L{$V$}.

\subsection{הוכחה}

אם \L{$m=n$}, הטענה נובעת מהטענה הקודמת.

אם \L{$m>n$} אז \L{$(v_{1},\dots,v_{m})$} ת"ל, כלומר קיים \L{$1\le i\le m$}
כך ש- \L{%
\[
v_{i}\in\spn(v_{1},\dots,v_{i-1},v_{i+1},\dots,v_{m})
\]
}%
בלי הגבלת הכלליות \L{$i=m$}, כלומר \L{$v_{m}\in\spn(v_{1},\dots,v_{m-1})$}.

אז לפי למה 1 הסדרה \L{$(v_{1},\dots,v_{m-1})$} פורשת.

אם \L{$m-1=n$}, אז \L{$(v_{1},\dots,v_{m-1})$} בסיס של \L{$V$}.

אם \L{$m-1>n$} נמשיך בתהליך הורדת הוקטורים מהסדרה, עד שנגיע לסדרה
פורשת באורך \L{$n$}. אז סדרה זו תהיה בסיס של \L{$V$}.

\newpage{}

\noindent\fbox{\begin{minipage}[t]{1\columnwidth - 2\fboxsep - 2\fboxrule}%
\begin{center}
אלגברה ליניארית )1( - \L{$80134$} \textbar{} האונ' העברית \textbar{}
תשפ\textquotedblright ו א' \textbar{} מרצה: ד\textquotedblright ר
אלכס גורביץ' \textbar{} מתרגל: מר יואב כהן \textbar{} ליקוט ועריכה:
נעם לב
\par\end{center}%
\end{minipage}}

\part{הרצאה {\beginL 13\endL} - {\beginL 01.12.2025\endL}}

\textcolor{teal}{בתזכורת עברנו על למה }\textcolor{teal}{{\beginL 1}\endL}\textcolor{teal}{,
למה }\textcolor{teal}{{\beginL 2}\endL}\textcolor{teal}{, והטענה האחרונה
מההרצאה הקודמת.}

\section{מסקנה {]}סדרה פורשת סופית \L{$\Leftarrow$} נוצר סופית{[}}

יהי \L{$V$} מ"ו מעל \L{$\mathbb{F}$}. אם ב-\L{$V$} קיימת סדרה
פורשת סופית, אז \L{$V$} נוצר סופית.

\section{טענה {]}השלמת סדרה בת\textquotedblright ל לבסיס של תמ\textquotedblright ו{[}}

יהיו \L{$V$} מ"ו נ"ס מעל \L{$\mathbb{F}$}, \L{$U$} תת-מרחב של
\L{$V$}, \L{$(v_{1},\dots,v_{m})$} סדרה בת"ל של וקטורים מ-\L{$U$}.
אז קיימים \L{$v_{m+1},\dots,v_{k}\in U$} כך ש- \L{$(v_{1},\dots,v_{m},v_{m+1},\dots,v_{k})$}
בסיס של \L{$U$}.

\subsection{הוכחה}

אם \L{$\spn(v_{1},\dots,v_{m})=U$}, אז \L{$(v_{1},\dots,v_{m})$}
בסיס של \L{$U$}. 

אחרת קיים \L{$v_{m+1}\in U$} כך ש- \L{$v_{m+1}\notin\spn(v_{1},\dots,v_{m})$}.

אז לפי למה 2, \L{$(v_{1},\dots,v_{m},v_{m+1})$} בת"ל.

נחזור על התהליך עבור הסדרה \L{$(v_{1},\dots,v_{m+1})$}.

מאחר שאורך סדרה בת"ל ב-\L{$V$} לא יכול להיות גדול מ- \L{$\dim V$},
התהליך יסתיים בשלב כלשהו וכך נקבל בסיס של \L{$U$}, כנדרש.

\section{מסקנה 1 {]}משפט ההשלמה לבסיס{[}}

יהיו \L{$V$} מ"ו נ"ס מעל \L{$\mathbb{F}$}, \L{$(v_{1},\dots,v_{m})$}
סדרה בת"ל של וקטורים מ-\L{$V$}.

אז קיימים וקטורים \L{$v_{m+1},\dots,v_{n}\in V$} כך ש- \L{$(v_{1},\dots,v_{m},v_{m+1},\dots,v_{n})$}
בסיס של \L{$V$}.

\section{מסקנה 2 {]}תת-מרחב של מ\textquotedblright ו נ\textquotedblright ס
הוא נ\textquotedblright ס{[}}

יהי \L{$V$} מ"ו נ"ס מעל \L{$\mathbb{F}$}, \L{$U$} תת-מרחב של
\L{$V$}. אז \L{$U$} גם נ"ס.

\subsection{הוכחה}

הסדרה הריקה \L{$()$} היא סדרה בת"ל של וקטורים מ-\L{$U$}. אז לפי
הטענה {]}השלמת סדרה בת\textquotedblright ל לבסיס של תמ\textquotedblright ו{[}
קיימים \L{$u_{1},\dots,u_{k}\in U$} כך ש- \L{$(u_{1},\dots,u_{k})$}
בסיס של \L{$U$}, לכן \L{$U$} נ"ס.

\section{טענה {]}מימד של תת-מרחב, ושוויון מימדים גורר שוויון{[}}

יהי \L{$V$} מ"ו נ"ס מעל \L{$\mathbb{F}$}, \L{$U$} תת-מרחב של
\L{$V$}. אז: 
\begin{enumerate}
\item \L{$\dim U\le\dim V$} 
\item אם \L{$\dim U=\dim V$}, אז \L{$U=V$}. \textcolor{blue}{)אפשר לכתוב
גם בתור אם"ם, הכיוון השני טריוויאלי( }
\end{enumerate}

\subsection{הוכחה}
\begin{enumerate}
\item נסמן \L{$\dim U=m,\dim V=n$}. אז קיים בסיס \L{$(u_{1},\dots,u_{m})$}
של \L{$U$}. נניח בשלילה כי \L{$m>n$}. אז \L{$(u_{1},\dots,u_{m})$}
ת"ל, בסתירה לזה שהיא בסיס. 
\item נסמן \L{$\dim U=\dim V=n$}. אז קיים בסיס \L{$(u_{1},\dots,u_{n})$}
של \L{$U$}. נניח בשלילה כי \L{$U\neq V$}. \L{$U$} תת מרחב של \L{$V$}
אז \L{$U\subseteq V$} וגם \L{$V\not\subseteq U$}. אז קיים \L{$v\in V$}
כך ש- \L{$v\notin U$}. מכאן \L{$v\notin\spn(u_{1},\dots,u_{n})$}.
כמו כן \L{$(u_{1},\dots,u_{n})$} בת"ל. אז לפי למה 2 הסדרה \L{$(u_{1},\dots,u_{n},v)$}
גם בת"ל. מצד שני מאחר שאורך הסדרה \L{$(u_{1},\dots,u_{n},v)$} שווה
ל- \L{$n+1$} ו- \L{$\dim V=n<n+1$}, אנו מסיקים כי היא ת"ל וזו סתירה.
\end{enumerate}
\begin{flushright}
\L{$\square$}
\par\end{flushright}

\subsection{הוכחה נוספת ל-2}

נסמן \L{$\dim U=\dim V=n$}. אז קיים בסיס \L{$(u_{1},\dots,u_{n})$}
של \L{$U$}. אז \L{$(u_{1},\dots,u_{n})$} סדרה בת"ל באורך \L{$n$}
של וקטורים מ-\L{$V$}. לכן היא פורשת את \L{$V$}, כלומר \L{$U=\spn(u_{1},\dots,u_{n})=V$}. 
\begin{flushright}
\L{$\square$}\\
\par\end{flushright}

\textcolor{blue}{בזה סיימנו את הפרק על בסיסים ומימדים.}

\section{הגדרה - סכום של קבוצות )במ\textquotedblright ו(}

יהי \L{$V$} מ"ו מעל \L{$\mathbb{F}$}. יהיו \L{$S_{1},S_{2}$} תת-קבוצות
לא ריקות של \L{$V$}. נגדיר \L{%
\[
S_{1}+S_{2}=\{u+v\mid u\in S_{1},v\in S_{2}\}
\]
}%
\textcolor{blue}{\L{$S_{1},S_{2}$} הן תת-קבוצות, לא תת-מרחבים.}

\subsection{דוגמה}

\L{%
\[
\mathbb{F}=\mathbb{R},V=\mathbb{R}
\]
}%
\L{%
\[
S_{1}=\{1,2\},\quad S_{2}=\{3,4\}
\]
}%
\L{%
\[
S_{1}+S_{2}=\{4,5,6\}
\]
}%


\section{הגדרה - ישריה}

\subsection{הערה}

אם \L{$S_{1}=\{v\}$} אז \L{%
\[
S_{1}+S_{2}=\{v\}+S_{2}=\{v+w\mid w\in S_{2}\}
\]
}%


\subsection{ההגדרה}

יהיו \L{$V$} מ"ו מעל \L{$\mathbb{F}$}, \L{$v\in V$}, \L{$U$}
תת-מרחב של \L{$V$}.

קבוצה \L{$\{v\}+U$} נקראת ישריה )\textcolor{blue}{מהמילה ישר}(.

\subsection{דוגמה}

\L{%
\[
\mathbb{F}=\mathbb{R},\quad V=\mathbb{R}^{2},\quad v=\begin{pmatrix}0\\
1
\end{pmatrix}\quad U=\spn\left\{ \begin{pmatrix}1\\
1
\end{pmatrix}\right\} 
\]
}%

\par\noindent\centerline{%
  \L{\includegraphics[width=0.2\textheight]{\detokenize{./gen_sum_vol_2_figs/fig4.png}}}%
}\par

\medskip{}

\textcolor{violet}{)בציור רואים את הישר \L{$U$} עובר בראשית, ואת
הישר המוזז \L{$v+U$} עובר דרך \L{$v$}(.}

\section{טענה {]}הפתרונות של \L{$Ax=b$} הם הזזה של הפתרונות של \L{$Ax=0$}{[}}

יהיו \L{$A\in M_{m\times n}(\mathbb{F})$}, \L{$b\in\mathbb{F}^{m}$},
יהי \L{$v\in\mathbb{F}^{n}$} פתרון של ממ\textquotedblright ל \L{$Ax=b$}.

נסמן ב- \L{$U$} את קבוצת הפתרונות של הממ"ל ההומוגנית \L{$Ax=\begin{pmatrix}0\\
\vdots\\
0
\end{pmatrix}$}.

אזי קבוצת הפתרונות של הממ"ל \L{$Ax=b$} היא \L{$\{v\}+U$}.

\subsection{הוכחה}

נוכיח כי \L{$\{v\}+U\subseteq\{x\in\mathbb{F}^{n}\mid Ax=b\}$}.

יהי \L{$w\in\{v\}+U$}.

אז קיים \L{$u\in U$} כך ש-\L{$w=v+u$}. מתקיים \L{$Au=\begin{pmatrix}0\\
\vdots\\
0
\end{pmatrix}$}. אז \L{%
\[
Aw=A(v+u)=Av+Au=b+\begin{pmatrix}0\\
\vdots\\
0
\end{pmatrix}=b
\]
}%
לכן \L{$w$} הוא פתרון של הממ"ל \L{$Ax=b$}.

נוכיח כי \L{$\{x\in\mathbb{F}^{n}\mid Ax=b\}\subseteq\{v\}+U$}.
\begin{flushright}
יהי \L{$w$} פתרון כלשהו של הממ"ל \L{$Ax=b$}. מתקיים \L{%
\[
A(w-v)=Aw-Av=b-b=\begin{pmatrix}0\\
\vdots\\
0
\end{pmatrix}
\]
}%
לכן \L{$w-v$} הוא פתרון של הממ"ל ההומוגנית \L{$Ax=\begin{pmatrix}0\\
\vdots\\
0
\end{pmatrix}$} ולפיכך \L{$w-v\in U$}. לסיום \L{%
\[
w=v+(w-v)\in\{v\}+U
\]
}%
\L{$\square$}
\par\end{flushright}

\section{מסקנה {]}קבוצת הפתרונות של ממ\textquotedblright ל פתירה היא ישריה{[}}

אם לממ"ל ב-\L{$n$} נעלמים קיים פתרון כלשהו, אז קבוצת הפתרונות שלה
היא ישריה ב-\L{$\mathbb{F}^{n}$}.

\section{למה {]}אם \L{$0_{V}\in S_{1}$} אז \L{$S_{2}\subseteq S_{1}+S_{2}$}{[}}

יהי \L{$V$} מ"ו מעל \L{$\mathbb{F}$} ויהיו \L{$S_{2}\subseteq V$},
\L{$0_{V}\in S_{1}\subseteq V$}. אז \L{$S_{2}\subseteq S_{1}+S_{2}$}.

\subsection{הוכחה}

יהי \L{$v\in S_{2}$}. אז מתקיים \L{%
\[
v=0_{V}+v\in S_{1}+S_{2}
\]
}%


\section{טענה {]}חיבור ספאנים שווה לספאן של האיחוד{[}}

יהי \L{$V$} מ"ו מעל \L{$\mathbb{F}$} ויהיו \L{$S_{1},S_{2}\subseteq V$}
. אז\L{%
\[
\spn S_{1}+\spn S_{2}=\spn(S_{1}\cup S_{2})
\]
}%

\newpage{}

\noindent\fbox{\begin{minipage}[t]{1\columnwidth - 2\fboxsep - 2\fboxrule}%
\begin{center}
אלגברה ליניארית )1( - \L{$80134$} \textbar{} האונ' העברית \textbar{}
תשפ\textquotedblright ו א' \textbar{} מרצה: ד\textquotedblright ר
אלכס גורביץ' \textbar{} מתרגל: מר יואב כהן \textbar{} ליקוט ועריכה:
נעם לב
\par\end{center}%
\end{minipage}}

\part{הרצאה {\beginL 14\endL} - {\beginL 02.12.2025\endL}}

\section{\uline{תזכורת}}

\section{למה מאתמול}

יהי \L{$V$} מ"ו מעל \L{$\mathbb{F}$}, \L{$S_{1},S_{2}\subseteq V$}.
אם \L{$0_{V}\in S_{1}$} אז \L{$S_{2}\subseteq S_{1}+S_{2}$}.

\section{טענה {]}חיבור ספאנים שווה לספאן של האיחוד{[}}

יהי \L{$V$} מ"ו מעל \L{$\mathbb{F}$}, \L{$S_{1},S_{2}\subseteq V$}.
אז \L{%
\[
\spn S_{1}+\spn S_{2}=\spn(S_{1}\cup S_{2})
\]
}%


\subsection{הוכחה}

נוכיח כי \L{$\spn(S_{1}\cup S_{2})\subseteq\spn S_{1}+\spn S_{2}$}.
יהי \L{$w\in\spn(S_{1}\cup S_{2})$}.

אז קיימים \L{$c_{1},\dots,c_{k}\in\mathbb{F}$}, \L{$w_{1},\dots,w_{k}\in S_{1}\cup S_{2}$}
כך ש- \L{$w=c_{1}w_{1}+\dots+c_{k}w_{k}$}. אחרי שינוי סדר של \L{$w_{1},\dots,w_{k}$}
ניתן להגיע למצב שבו \L{$w_{1},\dots,w_{l}\in S_{1},\quad w_{l+1},\dots,w_{k}\in S_{2}$}.
אז \L{%
\[
w=(c_{1}w_{1}+\dots+c_{l}w_{l})+(c_{l+1}w_{l+1}+\dots+c_{k}w_{k})\in\spn S_{1}+\spn S_{2}
\]
}%
נוכיח כי \L{$\spn S_{1}+\spn S_{2}\subseteq\spn(S_{1}\cup S_{2})$}.
יהי \L{$w\in\spn S_{1}+\spn S_{2}$}.

אז קיימים \L{$u\in\spn S_{1},\quad v\in\spn S_{2}$} כך ש- \L{$w=u+v$}.

לכן קיימים \L{$c_{1},\dots,c_{m}\in\mathbb{F}$}, \L{$u_{1},\dots,u_{m}\in S_{1}$}
כך ש- \L{$u=c_{1}u_{1}+\dots+c_{m}u_{m}$}.

בנוסף קיימים \L{$d_{1},\dots,d_{n}\in\mathbb{F}$}, \L{$v_{1},\dots,v_{n}\in S_{2}$}
כך ש- \L{$v=d_{1}v_{1}+\dots+d_{n}v_{n}$}. לכן \L{%
\[
w=c_{1}u_{1}+\dots+c_{m}u_{m}+d_{1}v_{1}+\dots+d_{n}v_{n}\in\spn(S_{1}\cup S_{2})
\]
}%

\begin{flushright}
\L{$\square$}
\par\end{flushright}

\section{מסקנה {]}סכום תתי-מרחבים הוא תמ\textquotedblright ו{[}}

יהי \L{$V$} מ"ו מעל \L{$\mathbb{F}$} ויהיו \L{$U_{1},U_{2}$} תתי-מרחבים
של \L{$V$}.

אז \L{$U_{1}+U_{2}=\spn(U_{1}\cup U_{2})$}. ובפרט סכום של תתי-מרחבים
הוא תת-מרחב. 

\subsection{הוכחה}

לפי הטענה הקודמת מתקיים \L{%
\[
U_{1}+U_{2}=\spn U_{1}+\spn U_{2}=\spn(U_{1}\cup U_{2})
\]
}%


\section{משפט המימדים הראשון \textcolor{green}{)חשוב(}}

יהי \L{$V$} מ"ו נ\textquotedblright ס מעל שדה \L{$\mathbb{F}$}
ויהיו \L{$U_{1},U_{2}$} תתי-מרחבים של \L{$V$}. אז \L{%
\[
\dim(U_{1}+U_{2})=\dim(U_{1})+\dim(U_{2})-\dim(U_{1}\cap U_{2})
\]
}%


\subsection{הוכחה}

\L{$U_{1}\cap U_{2}$} הוא תת-מרחב של \L{$V$} ולכן הוא נ\textquotedblright ס
וקיים בסיס \L{$(v_{1},\dots,v_{k})$} של \L{$U_{1}\cap U_{2}$}.

בנוסף קיימים \L{$u_{1},\dots,u_{m}\in U_{1}$} כך ש- \L{$(v_{1},\dots,v_{k},u_{1},\dots,u_{m})$}
היא בסיס של \L{$U_{1}$}. כמו כן קיימים \L{$w_{1},\dots,w_{n}\in U_{2}$}
כך ש- \L{$(v_{1},\dots,v_{k},w_{1},\dots,w_{n})$} היא בסיס של \L{$U_{2}$}.
נוכיח כי \L{%
\[
(v_{1},\dots,v_{k},u_{1},\dots,u_{m},w_{1},\dots,w_{n})=B
\]
}%
 היא בסיס של \L{$U_{1}+U_{2}$}.

נסמן \L{%
\[
S_{1}=\{v_{1},\dots,v_{k},u_{1},\dots,u_{m}\}
\]
}%
\L{%
\[
S_{2}=\{v_{1},\dots,v_{k},w_{1},\dots,w_{n}\}
\]
}%
אז \L{%
\[
\spn S_{1}=U_{1},\quad\spn S_{2}=U_{2}
\]
}%
. לפי הטענה הקודמת \L{%
\[
U_{1}+U_{2}=\spn S_{1}+\spn S_{2}=\spn(S_{1}\cup S_{2})=\spn(v_{1},\dots,v_{k},u_{1},\dots,u_{m},w_{1},\dots,w_{n})
\]
}%
לכן \L{$B$} פורשת את \L{$U_{1}+U_{2}$}.

יהיו \L{$a_{1},\dots,a_{k},b_{1},\dots,b_{m},c_{1},\dots,c_{n}\in\mathbb{F}$}
כך ש- \L{%
\[
a_{1}v_{1}+\dots+a_{k}v_{k}+b_{1}u_{1}+\dots+b_{m}u_{m}+c_{1}w_{1}+\dots+c_{n}w_{n}=0_{V}
\]
}%
מכאן \L{%
\[
\underbrace{a_{1}v_{1}+\dots+a_{k}v_{k}+b_{1}u_{1}+\dots+b_{m}u_{m}}_{\in U_{1}}=\underbrace{-(c_{1}w_{1}+\dots+c_{n}w_{n})}_{\in U_{2}}
\]
}%
מכאן \L{%
\[
-c_{1}w_{1}-\dots-c_{n}w_{n}\in U_{1}\cap U_{2}
\]
}%
לכן קיימים \L{$d_{1},\dots,d_{k}\in\mathbb{F}$} כך ש- \L{%
\[
-c_{1}w_{1}-\dots-c_{n}w_{n}=d_{1}v_{1}+\dots+d_{k}v_{k}
\]
}%
מכאן \L{%
\[
d_{1}v_{1}+\dots+d_{k}v_{k}+c_{1}w_{1}+\dots+c_{n}w_{n}=0_{V}
\]
}%
ומאחר ש- \L{$(v_{1},\dots,v_{k},w_{1},\dots,w_{n})$} בת"ל אנו מסיקים
כי \L{%
\[
d_{1}=\dots=d_{k}=c_{1}=\dots=c_{n}=0
\]
}%
 מכאן נובע כי \L{%
\[
a_{1}v_{1}+\dots+a_{k}v_{k}+b_{1}u_{1}+\dots+b_{m}u_{m}=0_{V}
\]
}%
מאחר ש- \L{$(v_{1},\dots,v_{k},u_{1},\dots,u_{m})$} בת"ל אנו מסיקים
כי \L{%
\[
a_{1}=\dots=a_{k}=b_{1}=\dots=b_{m}=0
\]
}%
 לכן \L{$B$} בת"ל. לסיום \L{%
\[
\dim U_{1}+\dim U_{2}-\dim(U_{1}\cap U_{2})
\]
}%
\L{%
\[
=(k+m)+(k+n)-k=k+m+n=\dim(U_{1}+U_{2})
\]
}%

\begin{flushright}
\L{$\square$}
\par\end{flushright}

\section{הגדרה - נמצאים בסכום ישר}

יהי \L{$V$} מ"ו מעל \L{$\mathbb{F}$}, יהיו \L{$U_{1},U_{2}$} תת-מרחבים.

אומרים כי \L{$U_{1},U_{2}$} \textbf{נמצאים בסכום ישר} כאשר \L{$U_{1}\cap U_{2}=\{0_{V}\}$}.

\section{מסקנה מהמשפט}

יהי \L{$V$} מ"ו נ"ס מעל \L{$\mathbb{F}$}, יהיו \L{$U_{1},U_{2}$}
תת-מרחבים של \L{$V$}.

אז \L{$U_{1},U_{2}$} נמצאים בסכום ישר אם"ם \L{%
\[
\dim(U_{1}+U_{2})=\dim U_{1}+\dim U_{2}
\]
}%


\subsection{הוכחה}

\L{$U_{1},U_{2}$} נמצאים בסכום ישר אם\textquotedblright ם \L{$U_{1}\cap U_{2}=\{0_{V}\}$}.

זה מתקיים אם\textquotedblright ם \L{$\dim(U_{1}\cap U_{2})=0$}. ולפי
משפט המימדים הראשון זה מתקיים אם\textquotedblright ם\L{%
\[
\dim(U_{1}+U_{2})-\dim U_{1}-\dim U_{2}=0
\]
}%

וזה מתקיים אם\textquotedblright ם \L{%
\[
\dim(U_{1}+U_{2})=\dim U_{1}+\dim U_{2}
\]
}%
וסיימנו.
\begin{flushright}
\L{$\square$}
\par\end{flushright}

\section{טענה {]}יחידות ההצגה בסכום ישר{[}}

יהי \L{$V$} מ"ו מעל \L{$\mathbb{F}$}, יהיו \L{$U_{1},U_{2}$} תת-מרחבים
של \L{$V$}.

אז \L{$U_{1},U_{2}$} נמצאים בסכום ישר אם"ם לכל \L{$v\in U_{1}+U_{2}$}
קיימים \L{$u\in U_{1},w\in U_{2}$} יחידים כך ש- \L{$v=u+w$}.

\subsection{הוכחה}

\L{$\Leftarrow$}: יהי \L{$v\in U_{1}+U_{2}$}.

לפי ההגדרה של הסכום קיימים \L{$u\in U_{1},w\in U_{2}$} כך ש- \L{$v=u+w$}.

נניח כי קיימים \L{$u'\in U_{1},w'\in U_{2}$} כך ש- \L{$v=u'+w'$}.

אז \L{$u+w=u'+w'$} מכאן \L{%
\[
\underbrace{u-u'}_{\in U_{1}}=\underbrace{w'-w}_{\in U_{2}}
\]
}%
לפיכך \L{$w'-w\in U_{1}\cap U_{2}$}. 

אבל הם נמצאים בסכום ישר ולכן \L{$U_{1}\cap U_{2}=\{0_{V}\}$} לכן
\L{$w'-w=0_{V}$} ומכאן \L{$w=w'$}. בנוסף \L{$u-u'=0_{V}$} ומכאן
\L{$u=u'$}.\\

\L{$\Rightarrow$}: נניח בשלילה כי \L{$U_{1},U_{2}$} אינם נמצאים
בסכום ישר.

כלומר קיים \L{$0_{V}\neq w\in U_{1}\cap U_{2}$}. אז ל-\L{$w$} יש
את שתי ההצגות הבאות: \L{%
\[
w=\underset{\in U_{1}}{0_{V}}+\underset{\in U_{2}}{w}
\]
}%
\L{%
\[
w=\underset{\in U_{1}}{w}+\underset{\in U_{2}}{0_{V}}
\]
}%
וזו סתירה להנחה.
\begin{flushright}
\L{$\square$}
\par\end{flushright}

\section{טענה {]}לכל תמ\textquotedblright ו קיים תמ\textquotedblright ו משלים,
המשלים בסכום ישר{[}}

יהיו \L{$V$} מ"ו נ"ס מעל \L{$\mathbb{F}$}, \L{$U$} תת-מרחב של
\L{$V$}.

אז קיים תת-מרחב \L{$W$} של \L{$V$} כך ש- \L{$U,W$} נמצאים בסכום
ישר וגם \L{$V=U+W$}.

\subsection{הוכחה}

\L{$U$} נ"ס לכן קיים בסיס \L{$(v_{1},\dots,v_{k})$} של \L{$U$}.

בנוסף קיימים \L{$v_{k+1},\dots,v_{n}\in V$} )\textcolor{violet}{משפט
ההשלמה לבסיס}( כך ש- \L{$(v_{1},\dots,v_{k},v_{k+1},\dots,v_{n})$}
בסיס של \L{$V$}.

נגדיר \L{$W=\spn(v_{k+1},\dots,v_{n})$}.

הסדרה \L{$(v_{k+1},\dots,v_{n})$} בת"ל ולכן היא בסיס של \L{$W$}.
\L{%
\[
U+W=\spn(v_{1},\dots,v_{k})+\spn(v_{k+1},\dots,v_{n})
\]
}%
\L{%
\[
=\spn(v_{1},\dots,v_{n})=V
\]
}%
\L{%
\[
\dim(U+W)=\dim(V)=n=k+(n-k)
\]
}%
\L{%
\[
=\dim U+\dim W
\]
}%
לכן לפי המסקנה \L{$U,W$} נמצאים בסכום ישר.

\section{סימון של סכום ישר}

כאשר תת-מרחבים \L{$U_{1},U_{2}$} נמצאים בסכום ישר רושמים \L{%
\[
U_{1}+U_{2}=U_{1}\oplus U_{2}
\]
}%

\newpage{}

\noindent\fbox{\begin{minipage}[t]{1\columnwidth - 2\fboxsep - 2\fboxrule}%
\begin{center}
אלגברה ליניארית )1( - \L{$80134$} \textbar{} האונ' העברית \textbar{}
תשפ\textquotedblright ו א' \textbar{} מרצה: ד\textquotedblright ר
אלכס גורביץ' \textbar{} מתרגל: מר יואב כהן \textbar{} ליקוט ועריכה:
נעם לב
\par\end{center}%
\end{minipage}}

\part{הרצאה {\beginL 15\endL} - {\beginL 08.12.2025\endL}}

\textcolor{blue}{היום נתחיל פרק חדש בקורס: העתקות לינאריות. אלה פונקציות
שמקיימות אקסיומות מסוימות.}

\textcolor{blue}{כמו שראינו בתיכון אפשר לפתור משוואות בעזרת פונקציות: }

\medskip{}

\par\noindent\centerline{%
  \L{\includegraphics[width=0.3\textheight]{\detokenize{./gen_sum_vol_2_figs/fig5.png}}}%
}\par

\medskip{}

כלומר, בשביל לפתור את המשווה \L{$2x=6$}, משתמשים בפונקציה \L{$f:\mathbb{R}\to\mathbb{R},f(x)=2x$},
ומגלים ש- \L{$f^{-1}(\{6\})=\{3\}$}. 

באופן דומה, אפשר לעשות את זה בממ"ל \L{$Ax=b$}. אם \L{$A\in M_{m\times n}(\mathbb{F}),\ b\in\mathbb{F}^{m}$}
מסתכלים על הפונקציה \L{$f:\mathbb{F}^{n}\rightarrow\mathbb{F}^{m}$}
המוגדרת ע\textquotedblright י \L{$f(x)=Ax$}, ומקבלים שקבוצת הפתרונות
היא \L{$f^{-1}(\{b\})$}.\\

\textcolor{blue}{אנחנו יודעים שהפונקציה הזאת מקיימת תכונות חשובות.
כמו שראינו:} \L{%
\[
A(x+y)=Ax+Ay
\]
}%
\L{%
\[
A(cx)=c(Ax)
\]
}%

\textcolor{blue}{שתי התכונות האלה אנחנו נבחר בתור האקסיומות של ה"ל.}

\section{העתקות לינאריות}

יהי \L{$\mathbb{F}$} שדה. יהיו \L{$V,W$} מ"ו מעל \L{$\mathbb{F}$}.
העתקה לינארית מ-\L{$V$} ל-\L{$W$} היא פונקציה \L{$T:V\to W$} כך
שמתקיים: 

\L{%
\[
1)\qquad\forall v_{1},v_{2}\in V:\quad T(v_{1}+v_{2})=T(v_{1})+T(v_{2})
\]
}%

\L{%
\[
2)\qquad\forall c\in\mathbb{F},v\in V:\quad T(cv)=cT(v)
\]
}%


\subsection{דוגמאות}

{\beginL 1\endL}( \L{$V,W$} מ"ו מעל \L{$\mathbb{F}$}. נגדיר \L{$T:V\to W$}
ע"י: \L{%
\[
\forall v\in V:\quad T(v)=0_{W}
\]
}%
\textcolor{violet}{העתקה זו נקראת העתקת האפס.}\\

{\beginL 2\endL}( \L{$V$} מ"ו מעל \L{$\mathbb{F}$} )\L{$V=W$}(.
נגדיר \L{$T:V\to V$} ע"י: \L{%
\[
\forall v\in V:\quad T(v)=v
\]
}%
העתקה זו נקראת העתקת הזהות, \L{$Id_{V}$}.\\

{\beginL 3\endL}( \L{$V=\mathbb{F}^{n},W=\mathbb{F}^{m}$}, \L{$A\in M_{m\times n}(\mathbb{F})$}.
נגדיר \L{$T_{A}:\mathbb{F}^{n}\to\mathbb{F}^{m}$} ע"י: \L{%
\[
\forall v\in\mathbb{F}^{n}:\quad T_{A}(v)=Av
\]
}%

\L{$T_{A}$} היא העתקה לינארית כי לכל \L{$v_{1},v_{2}\in\mathbb{F}^{n}$}
מתקיים: \L{%
\[
T_{A}(v_{1}+v_{2})=A(v_{1}+v_{2})=Av_{1}+Av_{2}=T_{A}(v_{1})+T_{A}(v_{2})
\]
}%
ולכל \L{$c\in\mathbb{F},v\in\mathbb{F}^{n}$} מתקיים: \L{%
\[
T_{A}(cv)=A(cv)=c(Av)=cT_{A}(v)
\]
}%

\textcolor{blue}{בהמשך הקורס נראה אפילו שכל ה\textquotedblright ל
ניתן לייצג באמצעות מטריצה.}\\
\\

{\beginL 4\endL}( \L{$\mathbb{F}$} שדה, \L{$S$} קבוצה. \L{$V=\mathbb{F}^{S}$}
זה מרחב הפונקציות מ-\L{$S$} ל-\L{$\mathbb{F}$}. \L{$W=\mathbb{F}$},
\L{$a\in S$}. נגדיר \L{$T_{a}:\mathbb{F}^{S}\to\mathbb{F}$} ע"י:
\L{%
\[
\forall f:S\to\mathbb{F}:\quad T_{a}(f)=f(a)
\]
}%

מתקיים ש-\L{$T_{a}$} היא העתקה לינארית כי לכל \L{$f_{1},f_{2}:S\to\mathbb{F}$}
מתקיים: \L{%
\[
T_{a}(f_{1}+f_{2})=(f_{1}+f_{2})(a)=f_{1}(a)+f_{2}(a)=T_{a}(f_{1})+T_{a}(f_{2})
\]
}%
וגם לכל \L{$c\in\mathbb{F},f:S\to\mathbb{F}$} מתקיים: \L{%
\[
T_{a}(cf)=(cf)(a)=c(f(a))=cT_{a}(f)
\]
}%
\textcolor{blue}{נזכיר: כשדיברנו על דוגמאות של מרחבים וקטוריים אחת
הדוגמאות הייתה מרחב הפונקציות, שם הגדרנו חיבור וכפל בסקלר של פונקציות.}

\section{תכונות של העתקות לינאריות )ה\textquotedblright ל(}

יהיו \L{$V,W$} מ"ו מעל שדה \L{$\mathbb{F}$}, \L{$T:V\to W$} ה"ל.
אז: 

\subsection{{\beginL 1\endL}( \L{$T(0_{V})=0_{W}$}}

\uline{הוכחה}

\L{%
\[
T(0_{V})=T(0\cdot0_{V})\underset{\text{(2)}}{=}0\cdot T(0_{V})=0_{W}
\]
}%


\subsection{{\beginL 2\endL}( \L{$T(-v)=-T(v)$}}

\uline{הוכחה} \L{%
\[
T(-v)=T((-1)v)\underset{\text{ (2)}}{=}(-1)T(v)=-T(v)
\]
}%


\subsection{{\beginL 3\endL}( \L{$\forall v_{1},\dots,v_{n}\in V,\ c_{1},\dots,c_{n}\in\mathbb{F}\quad:\quad T(c_{1}v_{1}+\dots+c_{n}v_{n})=c_{1}T(v_{1})+\dots+c_{n}T(v_{n})$}}

\uline{הוכחה}

\L{%
\[
T(c_{1}v_{1}+\dots+c_{n}v_{n})\underset{\text{(1)}}{=}T(c_{1}v_{1})+\dots+T(c_{n}v_{n})\underset{\text{(2)}}{=}c_{1}T(v_{1})+\dots+c_{n}T(v_{n})
\]
}%

\textcolor{blue}{בינתיים ראינו הגדרה של ה\textquotedblright ל, דוגמאות
וכמה תכונות פשוטות וחשובות.}

\textcolor{blue}{עכשיו נראה שבהינתן ה\textquotedblright ל ממרחב וקטורי
אחד לשני אפשר להגדיר שתי תת קבוצות חשובות, אחת בתחום ההגדרה ואחת בטווח,
המוכרות בשמות הגרעין והתמונה של העתקה לינארית.}

\section{גרעין ותמונה}

יהי \L{$\mathbb{F}$} שדה, \L{$V,W$} מ"ו מעל \L{$\mathbb{F}$},
\L{$T:V\to W$} ה"ל. 
\begin{itemize}
\item \textbf{הגרעין} של \L{$T$} מוגדר ע"י: \L{%
\[
\ker T=T^{-1}(\{0_{W}\})=\{v\in V\mid T(v)=0_{W}\}
\]
}%
\item \textbf{התמונה} של \L{$T$} מוגדרת ע"י: \L{%
\[
\operatorname{Im}(T)=T(V)=\{T(v)\mid v\in V\}
\]
}%
\end{itemize}
\textcolor{blue}{חשוב להבין שהתמונה היא תת קבוצה של \L{$W$}, בניגוד
לגרעין שהוא תת קבוצה של \L{$V$}.}\\

ל\textcolor{blue}{מה שתי הקבוצות האלה כ"כ חשובות? קודם כל כי כל אחת
מהן מהווה תמ"ו )אבל של מרחבים שונים(. עכשיו נציג את זה בתור טענה.}\\


\section{טענה {]}תמונה וגרעין הם תמ\textquotedblright ו{[}}

יהיו \L{$V,W$} מ\textquotedblright ו מעל \L{$\mathbb{F}$}, \L{$T:V\rightarrow W$}
ה\textquotedblright ל.

אז {\beginL 1\endL}( \L{$\ker T$} הוא תת-מרחב של \L{$V$}. 

\ \ \ \ {\beginL 2\endL}( \L{$\operatorname{Im}T$} היא תת-מרחב
של \L{$W$}.

\subsection{הוכחה}

{\beginL 1\endL}( \L{$0_{V}\in\ker T$} כי \L{$T(0_{V})=0_{W}$}.
יהיו \L{$v_{1},v_{2}\in\ker T$}.

אז \L{$T(v_{1})=T(v_{2})=0_{W}$} אז \L{$T(v_{1}+v_{2})\underset{(\text{1})}{=}T(v_{1})+T(v_{2})=0_{W}+0_{W}=0_{W}$}

לכן \L{$v_{1}+v_{2}\in\ker T$}. \\

\textcolor{blue}{בזה בדקנו שהגרעין סגור לחיבור, בדומה אפשר לבדוק שסגור
לכפל בסקלר.}\\
 

יהי \L{$v\in\ker T$}, \L{$c\in\mathbb{F}$}. אז \L{$T(v)=0_{W}$}.
מכאן \L{$T(cv)\underset{(2)}{=}cT(v)=c\cdot0_{W}=0_{W}$}, לכן \L{$cv\in\ker T$}.
\\

{\beginL 2\endL}( \L{$0_{W}\in\operatorname{Im}T$} כי \L{$T(0_{V})=0_{W}$}.

יהיו \L{$w_{1},w_{2}\in\operatorname{Im}T$}. אז קיימים \L{$v_{1},v_{2}\in V$}
כך ש- \L{$T(v_{1})=w_{1},T(v_{2})=w_{2}$}.

מתקיים \L{$T(v_{1}+v_{2})\underset{\text{(1)}}{=}T(v_{1})+T(v_{2})=w_{1}+w_{2}$}

לכן \L{$w_{1}+w_{2}\in\operatorname{Im}T$}.\\

יהי \L{$c\in\mathbb{F}\quad,w\in\operatorname{Im}T$}. אזי קיים \L{$v\in V$}
כך ש- \L{$T(v)=w$}. מתקיים \L{%
\[
T(cv)\underset{\text{(2)}}{=}cT(v)=cw
\]
}%
לכן \L{$cw\in\operatorname{Im}T$}.\\

\textcolor{blue}{בזה הוכחנו שכל האקסיומות של תמ"ו מתקיימות עבור התמונה
והגרעין.}\\

\textcolor{blue}{עכשיו נחזור לכמה דוגמאות של ה"ל ונראה מה הגרעין
והתמונה בכל אחת.}

\section{דוגמאות {]}לגרעין ותמונה בה\textquotedblright ל נפוצות{[}}

{\beginL 1\endL}( \L{$\forall v\in V:T(v)=0_{W}\qquad,T:V\to W$}.
\L{%
\[
\ker T=V\qquad\operatorname{Im}T=\{0_{W}\}
\]
}%

{\beginL 2\endL}( \L{$\forall v\in V:Id_{V}(v)=v\qquad,\text{Id}_{V}:V\to V$}.
\L{%
\[
\ker Id_{V}=\{0_{V}\}\qquad\operatorname{Im}Id_{V}=V
\]
}%

\textcolor{blue}{בדוגמה 2 הגרעין הוא הכי קטן שיכול להיות ובדוגמה 1
הכי גדול שיכול להיות. אבל למעשה יש רק ה"ל אחת בין שני מ"ו שמקיימת
את זה )שהגרעין הוא התחום והתמונה היא וקטור האפס( כמו בדוגמה 1, לעומת
זאת יש הרבה מאוד העתקות שונות מזהות שמקיימות את זה כמו בדוגמה 2.}\\

{\beginL 3\endL}( \L{$A\in M_{m\times n}(\mathbb{F}),V=\mathbb{F}^{n},W=\mathbb{F}^{m}$}.

\L{$T_{A}:\mathbb{F}^{n}\to\mathbb{F}^{m}$},\qquad{} \L{$\forall v\in\mathbb{F}^{n}:T_{A}(v)=Av$}.

\L{%
\[
\ker T_{A}=\{v\in\mathbb{F}^{n}\mid Av=\begin{pmatrix}0\\
\vdots\\
0
\end{pmatrix}\}
\]
}%
זוהי קבוצת הפתרונות של ממ\textquotedblright ל \L{$Ax=\begin{pmatrix}0\\
\vdots\\
0
\end{pmatrix}$}.

\L{%
\[
\operatorname{Im}T_{A}=\{Av\mid v\in\mathbb{F}^{n}\}=\{b\in\mathbb{F}^{m}\mid\exists v\in\mathbb{F}^{n}\text{ s.t }Av=b\}
\]
}%
זהו אוסף ה-\L{$m$}-יות כך שלממ\textquotedblright ל \L{$Ax=b$} קיים
פתרון. 

\newpage{}

\noindent\fbox{\begin{minipage}[t]{1\columnwidth - 2\fboxsep - 2\fboxrule}%
\begin{center}
אלגברה ליניארית )1( - \L{$80134$} \textbar{} האונ' העברית \textbar{}
תשפ\textquotedblright ו א' \textbar{} מרצה: ד\textquotedblright ר
אלכס גורביץ' \textbar{} מתרגל: מר יואב כהן \textbar{} ליקוט ועריכה:
נעם לב
\par\end{center}%
\end{minipage}}

\part{הרצאה {\beginL 16\endL} - {\beginL 09.12.2025\endL}}


\section{תזכורת}

\textcolor{blue}{התחלנו לדבר על ה"ל, גרעין, תמונה. ראינו שכל אחד
מהם תמ"ו )של תחום וטווח בהתאמה(. היום נמשיך לדבר על המושגים האלה
ונראה תכונות נוספות שלהם, ביניהן את משפט הממדים השני. קודם כל נרצה
לראות איך התמ"וים האלה, גרעין ותמונה, קשורים לתכונות של ה"ל כפונקציה,
כלומר פונקציה על או פונקציה חח"ע. נתחיל עכשיו עם תזכורת על המושגים
האלה על פונקציות.}

\section{הגדרות - פונ' חח\textquotedblright ע ועל}

תהי \L{$f:A\to B$} פונקציה. 
\begin{itemize}
\item \L{$f$} נקראת \textbf{על} כאשר לכל \L{$y\in B$} קיים \L{$x\in A$}
כך ש- \L{$f(x)=y$}. 
\item \L{$f$} נקראת \textbf{חח"ע} כאשר לכל \L{$x_{1},x_{2}\in A$} כך
ש- \L{$x_{1}\neq x_{2}$} מתקיים \L{$f(x_{1})\neq f(x_{2})$}. 
\end{itemize}
\noindent\fbox{\begin{minipage}[t]{1\columnwidth - 2\fboxsep - 2\fboxrule}%

\section{\textcolor{violet}{הגדרות שקולות נוספות: חח"ע ועל}}

\textcolor{violet}{תהי \L{$f:A\to B$} פונקציה.}

\subsection{\textcolor{violet}{פונקציה חד-חד-ערכית )\L{$\text{Injective}$}(}}

\textcolor{violet}{התנאים הבאים שקולים לכך ש-\L{$f$} היא חח"ע: }
\begin{enumerate}
\item \textcolor{violet}{לפי הגדרה )בשלילה(: לכל \L{$x_{1},x_{2}\in A$},
אם \L{$x_{1}\neq x_{2}$} אזי \L{$f(x_{1})\neq f(x_{2})$}. }
\item \textcolor{violet}{לפי הגדרה )הקונטרה-פוזיטיב(: לכל \L{$x_{1},x_{2}\in A$},
אם \L{$f(x_{1})=f(x_{2})$} אזי \L{$x_{1}=x_{2}$}. }
\item \textcolor{violet}{קיום הופכית משמאל: קיימת פונקציה \L{$g:B\to A$}
כך ש- \L{$g\circ f=id_{A}$} )בתנאי ש- \L{$A\neq\emptyset$}(. }
\item \textcolor{violet}{מקורות של איברים: לכל \L{$y\in B$}, למשוואה \L{$f(x)=y$}
יש }\textbf{\textcolor{violet}{לכל היותר}}\textcolor{violet}{{} פתרון
אחד ב-\L{$A$}. }
\end{enumerate}

\subsection{\textcolor{violet}{פונקציה על )\L{$\text{Surjective}$}(}}

\textcolor{violet}{התנאים הבאים שקולים לכך ש-\L{$f$} היא על: }
\begin{enumerate}
\item \textcolor{violet}{לפי הגדרה: לכל \L{$y\in B$} קיים \L{$x\in A$}
כך ש- \L{$f(x)=y$}. }
\item \textcolor{violet}{תמונת התחום: \L{$\operatorname{Im}(f)=B$} )או
בכתיב אחר \L{$f(A)=B$}(. }
\item \textcolor{violet}{קיום הופכית מימין: קיימת פונקציה \L{$g:B\to A$}
כך ש- \L{$f\circ g=id_{B}$}. }
\item \textcolor{violet}{מקורות של איברים: לכל \L{$y\in B$}, למשוואה \L{$f(x)=y$}
יש }\textbf{\textcolor{violet}{לפחות}}\textcolor{violet}{{} פתרון אחד
ב-\L{$A$}. }
\end{enumerate}

\subsection{\textcolor{violet}{פונקציה הפיכה )\L{$\text{Bijective}$}(}}

\textcolor{violet}{פונקציה היא הפיכה אם ורק אם מתקיים אחד מהתנאים
הבאים: }
\begin{enumerate}
\item הפונקציה היא חח\textquotedblright ע ועל.
\item \textcolor{violet}{לכל \L{$y\in B$} קיים \L{$x\in A$} יחיד כך ש-
\L{$f(x)=y$}. }
\item \textcolor{violet}{קיימת פונקציה \L{$g:B\to A$} שהיא גם הופכית מימין
וגם הופכית משמאל )כלומר \L{$f$} הפיכה(. }
\end{enumerate}
%
\end{minipage}}

\section{טענה {]}אפיון חח\textquotedblright ע ועל בעזרת גרעין ותמונה{[}}

יהיו \L{$V,W$} מ"ו מעל \L{$\mathbb{F}$}, \L{$T:V\to W$} ה"ל.
אז: 
\begin{enumerate}
\item \L{$T$} היא \textbf{על} אם\textquotedblright ם \L{$\operatorname{Im}(T)=W$}. 
\item \L{$T$} היא \textbf{חח"ע} אם\textquotedblright ם \L{$\ker(T)=\{0_{V}\}$}. 
\end{enumerate}

\subsection{הוכחה}

1( מאחר ש- \L{$\operatorname{Im}(T)=T(V)$}, \L{$T$} היא על אם\textquotedblright ם
\L{$\operatorname{Im}(T)=W$}. \\

{\beginL 2\endL}( 

\L{$\Leftarrow$} נניח \L{$T$} חח"ע.

נניח בשלילה כי \L{$\ker(T)\not=\{0_{V}\}$}, כלומר קיים \L{$0_{V}\neq v\in\ker(T)$}.

אז \L{$T(v)=0_{W}=T(0_{V})$} לכן \L{$T$} אינה חח\textquotedblright ע
- סתירה.

\L{$\Rightarrow$} נניח כי \L{$\ker(T)=\{0_{V}\}$}.

נניח בשלילה ש-\L{$T$} אינה חח"ע. כלומר קיימים \L{$v_{1},v_{2}\in V$}
כך ש- \L{$v_{1}\neq v_{2}$} )לכן \L{$v_{1}-v_{2}\neq0$}( וגם \L{$T(v_{1})=T(v_{2})$}.
מכאן: \L{%
\[
T(v_{1}-v_{2})=T(v_{1})-T(v_{2})=0_{W}
\]
}%
לכן \L{$v_{1}-v_{2}\in\ker(T)$}. כמו כן \L{$v_{1}-v_{2}\neq0_{V}$},
כלומר \L{$\ker(T)\not=\{0_{V}\}$} - סתירה. 

\section{תזכורת {]}תמונה של קבוצה{[}}

תהי \L{$f:A\to B$} פונקציה, \L{$C\subseteq A$}. אז התמונה של \L{$C$}
ביחס ל-\L{$f$} היא: \L{%
\[
f(C)=\{f(x)\mid x\in C\}\subseteq B
\]
}%


\section{טענה {]}תמונת ה\textquotedblright ל היא הספאן של תמונת קבוצה פורשת{[}}

יהיו \L{$V,W$} מ"ו מעל \L{$\mathbb{F}$}, \L{$T:V\to W$} ה"ל,
\L{$S$} תת-קבוצה פורשת של \L{$V$}. אז \L{$\operatorname{Im}(T)=\spn(T(S))$}.

\subsection{הוכחה}

כיוון ראשון \L{$\supseteq$}:

מתקיים \L{$T(S)\subseteq T(V)=\operatorname{Im}(T)$}. בנוסף \L{$\operatorname{Im}(T)$}
הוא תת-מרחב, לכן \L{$\spn(T(S))\subseteq\operatorname{Im}(T)$}.\\

כיוון שני \L{$\subseteq$}:

יהי \L{$w\in\operatorname{Im}(T)$}. אז קיים \L{$v\in V$} כך ש- \L{$T(v)=w$}.

מאחר ש- \L{$S$} פורשת את \L{$V$}, קיימים \L{$v_{1},\dots,v_{n}\in S$}
ו- \L{$c_{1},\dots,c_{n}\in\mathbb{F}$} כך ש: \L{%
\[
v=c_{1}v_{1}+\dots+c_{n}v_{n}
\]
}%
מכאן \L{%
\[
w=T(v)=T(c_{1}v_{1}+\dots+c_{n}v_{n})=c_{1}T(v_{1})+\dots+c_{n}T(v_{n})\in\spn(T(S))
\]
}%


\section{משפט הממדים השני}

יהיו \L{$V,W$} מ"ו מעל \L{$\mathbb{F}$} כך ש- \L{$V$} נ\textquotedblright ס,
\L{$T:V\to W$} ה"ל.

אז \L{$\ker T,\ \operatorname{Im}T$} הם גם נ\textquotedblright ס
ומתקיים \L{%
\[
\dim\ker T+\dim\operatorname{Im}T=\dim V
\]
}%


\subsection{הוכחה}

מאחר ש- \L{$\ker T$} הוא תת-מרחב של \L{$V$}, הוא נ\textquotedblright ס.
נבחר בסיס \L{$(v_{1},\dots,v_{k})$} של \L{$\ker T$}. מאחר שהוא בסיס
הוא בת\textquotedblright ל ולכן קיימים \L{$v_{k+1},\dots,v_{n}\in V$}
כך ש- \L{$(v_{1},\dots,v_{k},v_{k+1},\dots,v_{n})$} בסיס של \L{$V$}.

נוכיח כי \L{$(T(v_{k+1}),\dots,T(v_{n}))$} מהווה בסיס של \L{$\operatorname{Im}T$}.

מאחר ש- \L{$(v_{1},\dots,v_{n})$} פורשת את \L{$V$}, לפי הטענה הקודמת
מתקיים

\L{%
\[
\spn(T(v_{k+1}),\dots,T(v_{n}))=\spn(\underset{0_{W}}{\underbrace{T(v_{1})}},\dots,\underset{0_{W}}{\underbrace{T(v_{k})}},T(v_{k+1}),\dots,T(v_{n}))=\operatorname{Im}T
\]
}%
לכן \L{$(T(v_{k+1}),\dots,T(v_{n}))$} פורשת את \L{$\operatorname{Im}T$}.

יהיו \L{$c_{k+1},\dots,c_{n}\in\mathbb{F}$} כך ש: \L{%
\[
c_{k+1}T(v_{k+1})+\dots+c_{n}T(v_{n})=0_{W}
\]
}%
מכאן \L{%
\[
T(c_{k+1}v_{k+1}+\dots+c_{n}v_{n})=0_{W}
\]
}%
ולכן \L{$c_{k+1}v_{k+1}+\dots+c_{n}v_{n}\in\ker T$}.

מאחר ש-\L{$(v_{1},\dots,v_{k})$} פורשת את \L{$\ker T$}, קיימים \L{$c_{1},\dots,c_{k}\in\mathbb{F}$}
כך ש: \L{%
\[
c_{k+1}v_{k+1}+\dots+c_{n}v_{n}=c_{1}v_{1}+\dots+c_{k}v_{k}
\]
}%
מכאן \L{%
\[
-c_{1}v_{1}-\dots-c_{k}v_{k}+c_{k+1}v_{k+1}+\dots+c_{n}v_{n}=0_{V}
\]
}%
מאחר ש- \L{$(v_{1},\dots,v_{n})$} בת\textquotedblright ל, מתקיים:
\L{%
\[
-c_{1}=\dots=-c_{k}=c_{k+1}=\dots=c_{n}=0
\]
}%
לפיכך \L{$(T(v_{k+1}),\dots,T(v_{n}))$} בת"ל.

לסיום \L{%
\[
\dim\ker T+\dim\operatorname{Im}T=k+(n-k)=n=\dim V
\]
}%


\section{טענה {]}א\textquotedblright ש של מימדים לפי חח\textquotedblright ע
ועל{[}}

יהיו \L{$V,W$} מ"ו נ\textquotedblright ס מעל \L{$\mathbb{F}$},
\L{$T:V\to W$} ה"ל. אז
\begin{enumerate}
\item אם \L{$T$} חח"ע, אז \L{$\dim V\le\dim W$}. 
\item אם \L{$T$} על, אז \L{$\dim V\ge\dim W$}. 
\item אם \L{$T$} חח\textquotedblright ע ועל, אז \L{$\dim V=\dim W$}. 
\end{enumerate}

\subsection{הוכחה}
\begin{enumerate}
\item מאחר ש- \L{$T$} חח"ע, מתקיים \L{$\ker T=\{0_{V}\}$} ולכן \L{$\dim\ker T=0$}.
אז לפי המשפט מתקיים\L{%
\[
\dim V=\dim\ker T+\dim\operatorname{Im}T=\dim\operatorname{Im}T\le\dim W
\]
}%
\textcolor{violet}{)האי-שוויון נובע מכך ש-\L{$\operatorname{Im}T$}
תת-מרחב של \L{$W$}(.}
\item מאחר ש- \L{$T$} על, מתקיים \L{$\operatorname{Im}T=W$} ולכן \L{$\dim\operatorname{Im}T=\dim W$}.
אז לפי המשפט מתקיים\L{%
\[
\dim W=\dim\operatorname{Im}T\le\dim\ker T+\dim\operatorname{Im}T=\dim V
\]
}%
\item נובע מ-)1( ו-)2(. 
\end{enumerate}

\section{טענה {]}ה\textquotedblright ל בין מרחבים נ\textquotedblright ס מאותו
מימד היא חח\textquotedblright ע אם\textquotedblright ם היא על{[}}

יהיו \L{$V,W$} מ"ו נ\textquotedblright ס מעל \L{$\mathbb{F}$} כך
ש-\L{$\dim V=\dim W$}.

תהי \L{$T:V\to W$} ה"ל. אז \L{$T$} חח"ע אם\textquotedblright ם
\L{$T$} על.

\subsection{הוכחה}

\L{$T$} על אם\textquotedblright ם \L{$\operatorname{Im}T=W$} אם\textquotedblright ם
\L{$\dim\operatorname{Im}T=\dim W$} אם\textquotedblright ם )נתון(
\L{$\dim\operatorname{Im}T=\dim V$} אם\textquotedblright ם )לפי משפט
הממדים השני( \L{$\dim\operatorname{Im}T=\dim\ker T+\dim\operatorname{Im}T$}
אם\textquotedblright ם \L{$\dim\ker T=0$} אם\textquotedblright ם
\L{$\ker T=\{0_{V}\}$} אם\textquotedblright ם \L{$T$} חח\textquotedblright ע.

\section{הגדרה - הרכבת פונקציות}

תהיינה \L{$f:A\to B$}, \L{$g:B\to C$} פונ'.

ההרכבה של \L{$g$} על \L{$f$} היא פונקציה \L{$g\circ f:A\to C$}
המוגדרת ע"י: \L{%
\[
\forall x\in A,\quad(g\circ f)(x)=g(f(x))
\]
}%


\section{טענה {]}הרכבת ה\textquotedblright ל היא ה\textquotedblright ל{[}}

יהיו \L{$U,V,W$} מ"ו מעל \L{$\mathbb{F}$}, \L{$T_{1}:U\to V$},
\L{$T_{2}:V\to W$} ה"ל.

אז \L{$T_{2}\circ T_{1}:U\to W$} היא ה"ל.

\subsection{הוכחה}

יהיו \L{$u_{1},u_{2}\in U$}. אז \L{%
\[
(T_{2}\circ T_{1})(u_{1}+u_{2})=T_{2}(T_{1}(u_{1}+u_{2}))
\]
}%
 \L{$T_{1}$} ה\textquotedblright ל לכן \L{%
\[
=T_{2}(T_{1}(u_{1})+T_{1}(u_{2}))
\]
}%
 \L{$T_{2}$} ה\textquotedblright ל לכן: \L{%
\[
=T_{2}(T_{1}(u_{1}))+T_{2}(T_{1}(u_{2}))=(T_{2}\circ T_{1})(u_{1})+(T_{2}\circ T_{1})(u_{2})
\]
}%

יהי \L{$c\in\mathbb{F}$}, \L{$u\in U$}. אז \L{%
\[
(T_{2}\circ T_{1})(cu)=T_{2}(T_{1}(cu))
\]
}%
 \L{$T_{1}$} ה\textquotedblright ל לכן \L{%
\[
=T_{2}(cT_{1}(u))
\]
}%
 \L{$T_{2}$} ה\textquotedblright ל לכן \L{%
\[
=cT_{2}(T_{1}(u))=c(T_{2}\circ T_{1})(u)
\]
}%

וסיימנו.

\newpage{}

\noindent\fbox{\begin{minipage}[t]{1\columnwidth - 2\fboxsep - 2\fboxrule}%
\begin{center}
אלגברה ליניארית )1( - \L{$80134$} \textbar{} האונ' העברית \textbar{}
תשפ\textquotedblright ו א' \textbar{} מרצה: ד\textquotedblright ר
אלכס גורביץ' \textbar{} מתרגל: מר יואב כהן \textbar{} ליקוט ועריכה:
נעם לב
\par\end{center}%
\end{minipage}}

\part{הרצאה {\beginL 17\endL} - {\beginL 11.12.2025\endL}}

\textcolor{teal}{זה היה שיעור השלמה ביום חמישי.}

\section{הגדרה - פונ' הופכיות}

יהיו \L{$f:A\to B$},\quad{} \L{$g:B\to A$} פונקציות.

\L{$f$} ו- \L{$g$} נקראות \textbf{הופכיות} זו לזו כאשר

\L{%
\[
f\circ g=Id_{B}\quad\wedge\quad g\circ f=Id_{A}
\]
}%

כלומר \L{%
\[
\forall a\in A:\ g(f(a))=a\quad\wedge\quad\forall b\in B:\ f(g(b))=b
\]
}%

אז \L{$g,f$} הופכיות זו לזו ועל.

\subsection{עובדה {]}הופכיות \L{$\protect\iff$} חח\textquotedblright ע ועל{[}}

{\beginL 1\endL}( אם \L{$f,g$} הופכיות זו לזו, אז הן חח\textquotedblright ע
ועל.

{\beginL 2\endL}( אם \L{$f:A\rightarrow B$} חח\textquotedblright ע
ועל, אז קיימת פונ' \L{$g:B\rightarrow A$} הופכית ל- \L{$f$}.

\section{טענה {]}הופכית של ה\textquotedblright ל היא ה\textquotedblright ל{[}}

יהי \L{$V,W$} מ"ו מעל \L{$\mathbb{F}$}.

\L{$T:V\to W$} ה"ל חח\textquotedblright ע ועל. נסמן ב- \L{$T':W\to V$}
פונקציה הופכית ל- \L{$T$}. אז \L{$T'$} היא ה\textquotedblright ל.

\subsection{הוכחה}

יהיו \L{$w_{1},w_{2}\in W$}. נסמן \L{$v_{1}=T'(w_{1})$}, \L{$v_{2}=T'(w_{2})$}. 

אז מתקיים \L{$T(v_{1})=T(T'(w_{1}))=w_{1}$} , \L{$T(v_{2})=T(T'(w_{2}))=w_{2}$}.
לכן\L{%
\[
T'(w_{1}+w_{2})=T'(T(v_{1})+T(v_{2}))\underset{\text{\text{T is a Linear Map}}}{=\!=\!=}T'(T(v_{1}+v_{2}))=
\]
}%
\L{%
\[
=v_{1}+v_{2}=T'(w_{1})+T'(w_{2})
\]
}%
יהי \L{$w\in W$}, \L{$c\in\mathbb{F}$}. נסמן \L{$v=T'(w)$}. אז
\L{$T(v)=T(T'(w))=w$}. לכן \L{%
\[
T'(cw)=T'(cT(v))\underset{\text{T is L.M}}{\!=\!}T'(T(cv))=cv=cT'(w)
\]
}%


\section{הגדרה - איזומורפיזם}

ה\textquotedblright ל \L{$T:V\rightarrow W$} כך שקיימת עבורה ה\textquotedblright ל
הופכית \L{$T':W\rightarrow V$} נקראת \textbf{איזומורפיזם.}

יהיו \L{$V,W$} מ"ו מעל השדה \L{$\mathbb{F}$}. אומרים כי \L{$V$}
\textbf{איזומורפי} ל-\L{$W$} כאשר קיים איזומורפיזם \L{$T:V\to W$}.

\subsection{מסקנה {]}ה\textquotedblright ל חח\textquotedblright ע ועל היא איזומורפיזם{[}}

כל ה"ל חח"ע ועל היא איזומורפיזם.

\section{תכונות איזומורפיזם}

יהיו \L{$U,V,W$} מ"ו מעל \L{$\mathbb{F}$}. 
\begin{enumerate}
\item \L{$V$} איזומורפי ל-\L{$V$} )רפלקסיביות( )כי \L{$Id_{V}:V\to V$}
איזומורפיזם(. 
\item אם \L{$V$} איזומורפי ל-\L{$W$} אז \L{$W$} איזומורפי ל-\L{$V$}
)סימטריות(: )כי אם \L{$T:V\to W$} איזומורפיזם, אז \L{$T':W\to V$}
גם איזומורפיזם(. 
\item אם \L{$U$} איזומורפי ל-\L{$V$} ו-\L{$V$} איזומורפי ל-\L{$W$} אז
\L{$U$} איזומורפי ל-\L{$W$} )טרנזיטיביות(: )כי אם \L{$T_{1}:U\rightarrow V,\ T_{2}:V\rightarrow W$}
איזומורפיזם, אז גם \L{$T_{2}\circ T_{1}:U\rightarrow W$} איזומ'(.
\end{enumerate}

\section{טענה {]}פונ' וקטור הקואורדינטות היא איזומורפיזם{[}}

יהי \L{$V$} מ"ו נ\textquotedblright ס מעל \L{$\mathbb{F}$}, \L{$\dim V=n$},
\L{$\mathcal{B}=(v_{1},\dots,v_{n})$} בסיס של \L{$V$}. נגדיר \L{$\omega_{\mathcal{B}}:V\to\mathbb{F}^{n}$}
באופן הבא: \L{%
\[
\forall v\in V\quad\omega_{\mathcal{B}}(v)=[v]_{\mathcal{B}}
\]
}%
אז \L{$\omega_{\mathcal{B}}$} היא איזומורפיזם.

\subsection{הוכחה}

נראה כי \L{$\omega_{\mathcal{B}}$} היא ה"ל.

לכל \L{$v_{1},v_{2}\in V$} מתקיים \L{%
\[
\omega_{\mathcal{B}}(v_{1}+v_{2})=[v_{1}+v_{2}]_{\mathcal{B}}=[v_{1}]_{\mathcal{B}}+[v_{2}]_{\mathcal{B}}=\omega_{\mathcal{B}}(v_{1})+\omega_{\mathcal{B}}(v_{2})
\]
}%
לכל \L{$c\in\mathbb{F}$}, \L{$v\in V$} מתקיים: \L{%
\[
\omega_{\mathcal{B}}(cv)=[cv]_{\mathcal{B}}=c[v]_{\mathcal{B}}=c\omega_{\mathcal{B}}(v)
\]
}%
מאחר ש- \L{$(v_{1},\dots,v_{n})$} פורשת את \L{$V$} מתקיים\L{%
\[
\operatorname{Im}omega_{\mathcal{B}}=\spn(\omega_{\mathcal{B}}((v_{1},\dots,v_{n})))=\spn((\omega_{\mathcal{B}}(v_{1}),\dots,\omega_{\mathcal{B}}(v_{n})))=
\]
}%
\L{%
\[
=\spn((e_{1},\dots,e_{n}))=\mathbb{F}^{n}
\]
}%
לכן \L{$\omega_{\mathcal{B}}$} היא על. מאחר ש- \L{$\dim V=n=\dim\mathbb{F}^{n}$}
היא גם חח"ע, לכן \L{$\omega_{\mathcal{B}}$} איזומורפיזם.

\section{{]}ה-ה\textquotedblright ל ההופכית לפונ' וקטור הקואורדינטות{[}}

תהי \L{$\omega_{\mathcal{B}}:V\to\mathbb{F}^{n}$} כמו בטענה.

נסמן את ההעתקה הליניארית ההופכית ל- \L{$\omega_{\mathcal{B}}$} ב-\L{$\tau_{\mathcal{B}}:\mathbb{F}^{n}\to V$}.
אז\L{%
\[
\forall\begin{pmatrix}c_{1}\\
\vdots\\
c_{n}
\end{pmatrix}\in\mathbb{F}^{n}:\tau_{\mathcal{B}}\begin{pmatrix}c_{1}\\
\vdots\\
c_{n}
\end{pmatrix}=c_{1}v_{1}+\dots+c_{n}v_{n}
\]
}%


\subsection{מסקנה}

יהי \L{$V$} מ"ו נ\textquotedblright ס מעל \L{$\mathbb{F}$} כך ש-
\L{$\dim V=n$}.

אז \L{$V$} איזומורפי ל- \L{$\mathbb{F}^{n}$}.

\section{משפט {]}מ\textquotedblright ו איזומורפים אם\textquotedblright ם הם
מאותו מימד{[}}

יהיו \L{$V,W$} מ"ו נ\textquotedblright ס מעל \L{$\mathbb{F}$}.

אז \L{$V$} איזומורפי ל- \L{$W$} אם"ם \L{$\dim V=\dim W$}. 

\newpage{}

\noindent\fbox{\begin{minipage}[t]{1\columnwidth - 2\fboxsep - 2\fboxrule}%
\begin{center}
אלגברה ליניארית )1( - \L{$80134$} \textbar{} האונ' העברית \textbar{}
תשפ\textquotedblright ו א' \textbar{} מרצה: ד\textquotedblright ר
אלכס גורביץ' \textbar{} מתרגל: מר יואב כהן \textbar{} ליקוט ועריכה:
נעם לב
\par\end{center}%
\end{minipage}}

\part{הרצאה {\beginL 18\endL} - {\beginL 15.12.2025\endL}}

\section{טענה {]}קיום ויחידות ה\textquotedblright ל על בסיס{[}}

יהיו \L{$V,W$} מ"ו מעל \L{$\mathbb{F}$} כך ש- \L{$V$} נ\textquotedblright ס,
\L{$\mathcal{B}=(v_{1},\dots,v_{n})$} בסיס של \L{$V$}, \L{$w_{1},\dots,w_{n}\in W$}.
אזי קיימת ה"ל \L{$T:V\to W$} יחידה כך ש- \L{%
\[
\forall1\le i\le n\quad T(v_{i})=w_{i}
\]
}%


\subsection{הוכחה}

יהי \L{$v\in V$}. נסמן \L{$[v]_{\mathcal{B}}=\begin{pmatrix}c_{1}\\
\vdots\\
c_{n}
\end{pmatrix}$}. אז \L{$v=c_{1}v_{1}+\dots+c_{n}v_{n}$}. 

נגדיר \L{$T(v)=c_{1}w_{1}+\dots+c_{n}w_{n}\in W$}.

באופן זה הגדרנו פונקציה \L{$T:V\to W$}.

נוכיח כי \L{$T$} היא ה"ל.

יהיו \L{$u,v\in V$}. נסמן \L{$[u]_{\mathcal{B}}=\begin{pmatrix}b_{1}\\
\vdots\\
b_{n}
\end{pmatrix},[v]_{\mathcal{B}}=\begin{pmatrix}c_{1}\\
\vdots\\
c_{n}
\end{pmatrix}$}. 

אז \L{$[u+v]_{\mathcal{B}}=\begin{pmatrix}b_{1}+c_{1}\\
\vdots\\
b_{n}+c_{n}
\end{pmatrix}$}. מכאן \L{%
\[
\begin{aligned}T(u+v) & =(b_{1}+c_{1})w_{1}+\dots+(b_{n}+c_{n})w_{n}=\\
 & =(b_{1}w_{1}+\dots+b_{n}w_{n})+(c_{1}w_{1}+\dots+c_{n}w_{n})=T(u)+T(v)
\end{aligned}
\]
}%
יהיו \L{$c\in\mathbb{F},v\in V$}. נסמן \L{$[v]_{\mathcal{B}}=\begin{pmatrix}c_{1}\\
\vdots\\
c_{n}
\end{pmatrix}$}. אז \L{$[cv]_{\mathcal{B}}=\begin{pmatrix}cc_{1}\\
\vdots\\
cc_{n}
\end{pmatrix}$}. מכאן \L{%
\[
T(cv)=(cc_{1})w_{1}+\dots+(cc_{n})w_{n}=c(c_{1}w_{1}+\dots+c_{n}w_{n})=cT(v)
\]
}%
לכל \L{$1\le i\le n$} מתקיים \L{$[v_{i}]_{\mathcal{B}}=e_{i}=\begin{pmatrix}0\\
\vdots\\
1\\
\vdots\\
0
\end{pmatrix}\leftarrow i$}. מכאן \L{%
\[
T(v_{i})=0\cdot w_{1}+\dots+0\cdot w_{i-1}+1\cdot w_{i}+0\cdot w_{i+1}+\dots+0w_{n}=w_{i}
\]
}%
באופן זה בנינו ה\textquotedblright ל המקיימת את התנאי הנדרש.

כדי להוכיח את היחידות, נניח כי קיימות ה"ל \L{$T,T':V\to W$} כך שלכל
\L{$1\le i\le n$} מתקיים \L{$T(v_{i})=w_{i}$} וגם \L{$T'(v_{i})=w_{i}$}.
יהי \L{$v\in V$}. נסמן \L{$[v]_{\mathcal{B}}=\begin{pmatrix}c_{1}\\
\vdots\\
c_{n}
\end{pmatrix}$}. אז \L{$v=c_{1}v_{1}+\dots+c_{n}v_{n}$}. מכאן \L{%
\[
T(v)=T(c_{1}v_{1}+\dots+c_{n}v_{n})=c_{1}T(v_{1})+\dots+c_{n}T(v_{n})=c_{1}w_{1}+\dots+c_{n}w_{n}
\]
}%
באופן דומה מתקיים \L{$T'(v)=c_{1}w_{1}+\dots+c_{n}w_{n}$}.

לכן \L{$T(v)=T'(v)$}. לפיכך ה"ל שמקיימת את התנאי הנדרש יחידה.

\section{טענה {]}כל ה\textquotedblright ל ניתנת לייצוג באמצעות מטריצה{[}}

תהי \L{$T:\mathbb{F}^{n}\to\mathbb{F}^{m}$} ה"ל כלשהי.

אז קיימת \L{$A\in M_{m\times n}(\mathbb{F})$} כך ש- \L{$T=T_{A}$}
כאשר

\L{%
\[
T_{A}:\mathbb{F}^{n}\rightarrow\mathbb{F}^{m},\quad\forall v\in\mathbb{F}^{n}\quad T_{A}(v)=Av
\]
}%


\subsection{הוכחה}

לכל \L{$1\le i\le n$} נסמן \L{$w_{i}=T(e_{i})\in\mathbb{F}^{m}$}.
נגדיר את \L{$A$} להיות המטריצה שהעמודה ה-\L{$i$} שלה היא \L{$w_{i}$},
כלומר \L{$A=(w_{1}\dots w_{n})$}. 

אז \L{$A\in M_{m\times n}(\mathbb{F})$} ולכל \L{$1\le i\le n$} מתקיים
\L{%
\[
T(e_{i})=w_{i}=Ae_{i}=T_{A}(e_{i})
\]
}%
לכן לפי הטענה הקודמת מתקיים \L{$T=T_{A}$}.

\section{מטריצה של העתקה ליניארית {]}\textcolor{teal}{הגדרה חשובה}{[}}

יהיו \L{$V,W$} מ"ו נ\textquotedblright ס מעל \L{$\mathbb{F}$},

\L{$\mathcal{B}=(v_{1},\dots,v_{n})$} בסיס של \L{$V$},

\L{$\mathcal{C}=(w_{1},\dots,w_{m})$} בסיס של \L{$W$}.

תהי \L{$T:V\to W$} ה"ל.

המטריצה של \L{$T$} ביחס לבסיסים \L{$\mathcal{B},\mathcal{C}$} היא
המטריצה \L{$[T]_{\mathcal{C}}^{\mathcal{B}}$} שהעמודה ה-\L{$i$}
שלה היא \L{$[T(v_{i})]_{\mathcal{C}}$} כלומר \L{%
\[
[T]_{\mathcal{C}}^{\mathcal{B}}=\left(\begin{array}{c}
|\\{}
[T(v_{1})]_{\mathcal{C}}\\
|
\end{array}\dots\begin{array}{c}
|\\{}
[T(v_{n})]_{\mathcal{C}}\\
|
\end{array}\right)
\]
}%

לפיכך \L{$[T]_{\mathcal{C}}^{\mathcal{B}}\in M_{m\times n}(\mathbb{F})$}.

\section{טענה {]}ה\textquotedblright ל נקבעת ע\textquotedblright י המטריצה
המייצגת שלה{[}}

יהיו \L{$V,W$} מ"ו נ\textquotedblright ס מעל \L{$\mathbb{F}$},
\L{$\mathcal{B},\mathcal{C}$} בסיסים של \L{$V,W$} בהתאמה, \L{$T,T':V\to W$}
ה"ל כך ש- \L{$[T]_{\mathcal{C}}^{\mathcal{B}}=[T']_{\mathcal{C}}^{\mathcal{B}}$}.
אז \L{$T=T'$}.

\subsection{הוכחה}

מאחר ש- \L{$[T]_{\mathcal{C}}^{\mathcal{B}}=[T']_{\mathcal{C}}^{\mathcal{B}}$},
לכל \L{$1\le i\le n$} מתקיים \L{$[T(v_{i})]_{\mathcal{C}}=[T'(v_{i})]_{\mathcal{C}}$}
ומכאן \L{$T(v_{i})=T'(v_{i})$}. אז לפי היחידות שבטענה הקודמת אנו
מסיקים כי \L{$T=T'$}.

\section{טענה {]}כל מטריצה מייצגת ה\textquotedblright ל{[}}

יהיו \L{$V,W$} מ"ו נ\textquotedblright ס מעל \L{$\mathbb{F}$},
\L{$\dim V=n,\dim W=m$}, \L{$\mathcal{B},\mathcal{C}$} בסיסים של
\L{$V,W$} בהתאמה, \L{$A\in M_{m\times n}(\mathbb{F})$}.

אז קיימת \L{$T:V\to W$} כך ש-\L{$[T]_{\mathcal{C}}^{\mathcal{B}}=A$}.

\subsection{הוכחה}

נסמן \L{$A=\begin{pmatrix}a_{11} & \dots & a_{1n}\\
\vdots &  & \vdots\\
a_{m1} & \dots & a_{mn}
\end{pmatrix}$} ונסמן \L{$\mathcal{B}=(v_{1},\dots,v_{n}),\mathcal{C}=(w_{1},\dots,w_{m})$}.

לכל \L{$1\le i\le n$} נגדיר \L{$u_{i}=a_{1i}w_{1}+\dots+a_{mi}w_{m}\in W$}. 

אז \L{$\begin{pmatrix}a_{1i}\\
\vdots\\
a_{mi}
\end{pmatrix}=[u_{i}]_{\mathcal{C}}$} היא העמודה ה- \L{$i$} של \L{$A$}.

לפי הקיום בטענה {]}קיום ויחידות ה\textquotedblright ל על בסיס{[},
קיימת ה\textquotedblright ל \L{$T:V\to W$} כך ש- \L{%
\[
\forall1\leq i\leq n:T(v_{i})=u_{i}
\]
}%
 אז \L{%
\[
[T]_{\mathcal{C}}^{\mathcal{B}}=\left(\begin{array}{c}
|\\{}
[T(v_{1})]_{\mathcal{C}}\\
|
\end{array}\dots\begin{array}{c}
|\\{}
[T(v_{n})]_{\mathcal{C}}\\
|
\end{array}\right)=\left(\begin{array}{c}
|\\{}
[u_{1}]_{\mathcal{C}}\\
|
\end{array}\dots\begin{array}{c}
|\\{}
[u_{n}]_{\mathcal{C}}\\
|
\end{array}\right)=\begin{pmatrix}a_{11} & \dots & a_{1n}\\
\vdots &  & \vdots\\
a_{m1} & \dots & a_{mn}
\end{pmatrix}=A
\]
}%

\textcolor{teal}{זה ש-\L{$T$} היא ה\textquotedblright ל נובע מהטענה
שבה השתמשנו: היא מבטיחה קיום של העתקה לינארית )ולא סתם פונ'( המקיימת
\L{$T(v_{i})=u_{i}$} לכל \L{$i$}. לכן \L{$T$} היא ה\textquotedblright ל
ואפשר לדבר על המטריצה המייצגת שלה. }

\section{משפט {]}\L{$[T(v)]_{\mathcal{C}}=[T]_{\mathcal{C}}^{\mathcal{B}}[v]_{\mathcal{B}}$}{[}}

יהיו \L{$V,W$} מ"ו נ\textquotedblright ס מעל \L{$\mathbb{F}$}.
\L{$\mathcal{B},\mathcal{C}$} בסיסים של \L{$V,W$} בהתאמה. \L{$T:V\to W$}
ה"ל, \L{$v\in V$}. אז \L{$[T(v)]_{\mathcal{C}}=[T]_{\mathcal{C}}^{\mathcal{B}}[v]_{\mathcal{B}}$}.

\subsection{הוכחה}

נסמן \L{$\mathcal{B}=(v_{1},\dots,v_{n}),[v]_{\mathcal{B}}=\begin{pmatrix}c_{1}\\
\vdots\\
c_{n}
\end{pmatrix}$}. אז \L{$v=c_{1}v_{1}+\dots+c_{n}v_{n}$}. 

\L{%
\[
[T(v)]_{\mathcal{C}}=\omega_{\mathcal{C}}(T(v))=(\omega_{\mathcal{C}}\circ T)(c_{1}v_{1}+\dots+c_{n}v_{n})=
\]
}%

\L{%
\[
=c_{1}(\omega_{\mathcal{C}}\circ T)(v_{1})+\dots+c_{n}(\omega_{\mathcal{C}}\circ T)(v_{n})=
\]
}%

\L{%
\[
=c_{1}[T(v_{1})]_{\mathcal{C}}+\dots+c_{n}[T(v_{n})]_{\mathcal{C}}=
\]
}%

\L{%
\[
=\left([T(v_{1})]_{\mathcal{C}}\dots[T(v_{n})]_{\mathcal{C}}\right)\begin{pmatrix}c_{1}\\
\vdots\\
c_{n}
\end{pmatrix}=[T]_{\mathcal{C}}^{\mathcal{B}}[v]_{\mathcal{B}}
\]
}%

\newpage{}

\noindent\fbox{\begin{minipage}[t]{1\columnwidth - 2\fboxsep - 2\fboxrule}%
\begin{center}
אלגברה ליניארית )1( - \L{$80134$} \textbar{} האונ' העברית \textbar{}
תשפ\textquotedblright ו א' \textbar{} מרצה: ד\textquotedblright ר
אלכס גורביץ' \textbar{} מתרגל: מר יואב כהן \textbar{} ליקוט ועריכה:
נעם לב
\par\end{center}%
\end{minipage}}

\part{הרצאה {\beginL 19\endL} - {\beginL 16.12.2025\endL}}

\section{תזכורת}

יהיו \L{$V,W$} מ\textquotedblright ו נ\textquotedblright ס מעל \L{$\mathbb{F}$},

\L{$\mathcal{B},\mathcal{C}$} הם בסיסים של \L{$V,W$} בהתאמה,

\L{$T:V\rightarrow W$} ה\textquotedblright ל. המטריצה של \L{$T$}
ביחס ל-\L{$\mathcal{B}$} ו-\L{$C$} היא

המטריצה \L{$[T]_{\mathcal{C}}^{\mathcal{B}}$} כך שהעמודה ה-\L{$i$}
שלה היא \L{$[T(v_{i})]_{\mathcal{C}}$} כלומר

\L{%
\[
[T]_{\mathcal{C}}^{\mathcal{B}}=\left([T(v_{1})]_{\mathcal{C}},\dots,[T(v_{n})]_{\mathcal{C}}\right)
\]
}%


\section{משפט}

יהיו \L{$V,W$} מ\textquotedblright ו נ\textquotedblright ס מעל \L{$\mathbb{F}$},

\L{$\mathcal{B},\mathcal{C}$} בסיסים של \L{$V,W$} בהתאמה, \L{$T:V\rightarrow W$}
ה\textquotedblright ל, \L{$v\in V$}.

אז \L{$[T(v)]_{\mathcal{C}}=[T]_{\mathcal{C}}^{\mathcal{B}}[v]_{\mathcal{B}}$}.

\rule[0.5ex]{1\columnwidth}{1pt}

\section{מסקנה}

יהיו \L{$V,W$} מ\textquotedblright ו נ\textquotedblright ס מעל \L{$\mathbb{F}$},

\L{$\mathcal{B},\mathcal{C}$} בסיסים של \L{$V,W$} בהתאמה, \L{$T:V\rightarrow W$}
ה\textquotedblright ל. נסמן \L{$A=[T]_{\mathcal{C}}^{\mathcal{B}}$}.

\L{$\omega_{\mathcal{B}}(v)=[v]_{\mathcal{B}},\quad\omega_{\mathcal{B}}:V\rightarrow\mathbb{F}^{n}$}
לכל \L{$v\in V$}.

\L{$\omega_{\mathcal{C}}(w)=[w]_{\mathcal{C}},\quad\omega_{\mathcal{C}}:W\rightarrow\mathbb{F}^{m}$}
לכל \L{$w\in W$}.

אז מתקיים

\L{%
\begin{align*}
1)\quad\omega_{\mathcal{C}}\circ T & =T_{A}\circ\omega_{\mathcal{B}}\\
2)\quad T & =\tau_{\mathcal{C}}\circ T_{A}\circ\omega_{\mathcal{B}}
\end{align*}
}%


\subsection{הוכחה}

{\beginL 1\endL}( לכל \L{$v\in V$} מתקיים

\L{$(\omega_{\mathcal{C}}\circ T)(v)=\omega_{\mathcal{C}}(T(v))=[T(v)]_{\mathcal{C}}\underset{(i)}{=}[T]_{\mathcal{C}}^{\mathcal{B}}[v]_{\mathcal{B}}=A[v]_{\mathcal{B}}=T_{A}([v]_{\mathcal{B}})=T_{A}(\omega_{\mathcal{B}}(v))=(T_{A}\circ\omega_{\mathcal{B}})v$}

\L{$\qquad$}\L{$(i)$} לפי המשפט.

{\beginL 2\endL}(

לפי ){\beginL 1\endL}( מתקיים \L{$T=Id_{W}\circ T=(\tau_{\mathcal{C}}\circ\omega_{\mathcal{C}})\circ T=\tau_{\mathcal{C}}\circ(\omega_{\mathcal{C}}\circ T)=\tau_{\mathcal{C}}\circ(T_{A}\circ\omega_{\mathcal{B}})$}

\section{טענה}

תהי \L{$A\in M_{m\times n}(\mathbb{F})$}. נגדיר \L{$T_{A}:\mathbb{F}^{n}\rightarrow\mathbb{F}^{m}$}
ע\textquotedblright י \L{$T_{A}(v)=Av$}.

נסמן ב-\L{$\varepsilon_{n},\varepsilon_{m}$} את הבסיסים הסטנדרטיים
של \L{$\mathbb{F}^{n},\mathbb{F}^{m}$} בהתאמה.

אז \L{$[T_{A}]_{\varepsilon_{m}}^{\varepsilon_{n}}=A$}.

\subsection{הוכחה}

נסמן את העמודות של \L{$A$} ב-\L{$C_{1},\dots,C_{n}\in\mathbb{F}^{m}$}.

אז לכל \L{$1\leq i\leq n$} העמודה ה-\L{$i$} של \L{$[T_{A}]_{\varepsilon_{m}}^{\varepsilon_{n}}$}
היא 

\L{$[T_{A}(e_{i})]_{\varepsilon_{m}}=[Ae_{i}]_{\varepsilon_{m}}=[C_{i}]_{\varepsilon_{m}}=C_{i}$}
ולכן היא שווה לעמודה

ה-\L{$i$} של \L{$A$} וסיימנו. 

\section{כפל מטריצות}

יהיו \L{$\ell,m,n\in\mathbb{N}$} ויהי \L{$\mathbb{F}$} שדה,

\L{$A\in M_{m\times n}(\mathbb{F})$}, \L{$B\in M_{\ell\times m}(\mathbb{F})$}.

המכפלה של \L{$A$} ו-\L{$B$} היא המטריצה \L{$BA$} כך שלכל \L{$1\leq i\leq n$}

העמודה ה-\L{$i$} של \L{$BA$} היא \L{$BC_{i}\in\mathbb{F}^{\ell}$}
כאשר \L{$C_{i}\in\mathbb{F}^{m}$} היא העמודה ה-\L{$i$} של \L{$A$}.

לפיכך \L{$BA\in M_{\ell\times n}(\mathbb{F})$}.

\section{טענה}

יהיו \L{$U,V,W$} מ\textquotedblright ו נ\textquotedblright ס מעל
\L{$\mathbb{F}$}, \L{$T_{1}:U\rightarrow V,\quad T_{2}:V\rightarrow W$}
ה\textquotedblright ל,

\L{$\mathcal{B},\mathcal{C},\mathcal{D}$} הם בסיסים של \L{$U,V,W$}
בהתאמה.

אז \L{$[T_{2}\circ T_{1}]_{\mathcal{D}}^{\mathcal{B}}=[T_{2}]_{\mathcal{D}}^{\mathcal{C}}\cdot[T_{1}]_{\mathcal{C}}^{\mathcal{B}}$}.

\subsection{הוכחה}

נסמן \L{$\mathcal{B}=(u_{1},\dots,u_{n})$}. אז לכל \L{$1\leq i\leq n$}
העמודה ה-\L{$i$} של \L{$[T_{2}\circ T_{1}]_{\mathcal{D}}^{\mathcal{B}}$}
היא

\L{%
\[
[(T_{2}\circ T_{1})(u_{i})]_{\mathcal{D}}=[T_{2}(T_{1}(u_{i}))]_{\mathcal{D}}\underset{(i)}{=}[T_{2}]_{\mathcal{D}}^{\mathcal{C}}[T_{1}(u_{i})]_{\mathcal{C}}\underset{(i)}{=}[T_{2}]_{\mathcal{D}}^{\mathcal{C}}([T_{1}]_{\mathcal{C}}^{\mathcal{B}}[u_{i}]_{\mathcal{B}})=[T_{2}]_{\mathcal{D}}^{\mathcal{C}}([T_{1}]_{\mathcal{C}}^{\mathcal{B}}e_{i})
\]
}%

\L{$(i)\quad$} לפי משפט\\
\\
מאחר ש-\L{$[T_{1}]_{\mathcal{C}}^{\mathcal{B}}e_{i}$} היא העמודה
ה-\L{$i$} של \L{$[T_{1}]_{\mathcal{C}}^{\mathcal{B}}$}, קיבלנו בצד
ימין את העמודה ה-\L{$i$} של \L{$[T_{2}]_{\mathcal{D}}^{\mathcal{C}}[T_{1}]_{\mathcal{C}}^{\mathcal{B}}$}
וסיימנו. 

\section{טענה}

יהיו \L{$k,\ell,m,n\in\mathbb{N}$}, יהי \L{$\mathbb{F}$} שדה,

ויהיו \L{$A\in M_{m\times n}(\mathbb{F})$}, \L{$B\in M_{\ell\times m}(\mathbb{F})$},
\L{$C\in M_{k\times\ell}(\mathbb{F})$}.

אז \L{$(CB)A=C(BA)$}.

\subsection{הוכחה}

נתבונן בה\textquotedblright ל הבאות: \L{$T_{A}:\mathbb{F}^{n}\rightarrow\mathbb{F}^{m},T_{B}:\mathbb{F}^{m}\rightarrow\mathbb{F}^{\ell},T_{C}:\mathbb{F}^{\ell}\rightarrow\mathbb{F}^{k}$}.

נסמן ב-\L{$\varepsilon_{n},\varepsilon_{m},\varepsilon_{\ell},\varepsilon_{k}$}
את הבסיסים הסטנדרטיים של \L{$\mathbb{F}^{n},\mathbb{F}^{m},\mathbb{F}^{\ell},\mathbb{F}^{k}$}
בהתאמה.

אז \L{$C(BA)=[T_{C}]_{\varepsilon_{k}}^{\varepsilon_{\ell}}([T_{B}]_{\varepsilon_{\ell}}^{\varepsilon_{m}}[T_{A}]_{\varepsilon_{m}}^{\varepsilon_{n}})=[T_{C}]_{\varepsilon_{k}}^{\varepsilon_{\ell}}[T_{B}\circ T_{A}]_{\varepsilon_{\ell}}^{\varepsilon_{n}}=[T_{C}\circ(T_{B}\circ T_{A})]_{\varepsilon_{k}}^{\varepsilon_{n}}=$}

\L{$=[(T_{C}\circ T_{B})\circ T_{A}]_{\varepsilon_{k}}^{\varepsilon_{n}}=[T_{C}\circ T_{B}]_{\varepsilon_{k}}^{\varepsilon_{m}}[T_{A}]_{\varepsilon_{m}}^{\varepsilon_{n}}=([T_{C}]_{\varepsilon_{k}}^{\varepsilon_{\ell}}[T_{B}]_{\varepsilon_{\ell}}^{\varepsilon_{m}})[T_{A}]_{\varepsilon_{m}}^{\varepsilon_{n}}=(CB)A$}

וסיימנו.

\section{תכונה של כפל מטריצות}

תהיינה \L{$A\in M_{m\times n}(\mathbb{F})$}, \L{$B\in M_{\ell\times m}(\mathbb{F})$}.

אם העמודה ה-\L{$i$} של \L{$A$} היא עמודת האפסים אז העמודה ה-\L{$i$}
של \L{$BA$} היא גם עמודת האפסים. 

\section{תכונה }

תהיינה \L{$A=\left(\begin{array}{ccc}
a_{11} & \dots & a_{1n}\\
\vdots &  & \vdots\\
a_{m1} & \dots & a_{mn}
\end{array}\right)\in M_{m\times n}(\mathbb{F})$}, \L{$\mathcal{B}=\left(\begin{array}{ccc}
b_{11} & \dots & b_{1m}\\
\vdots &  & \vdots\\
b_{\ell1} & \dots & b_{\ell m}
\end{array}\right)\in M_{\ell\times m}(\mathbb{F})$}.

יהיו \L{$1\leq i\leq\ell,\quad1\leq j\leq n$} . אז בחיתוך של השורה
ה-\L{$i$} והעמודה ה-\L{$j$} של \L{$BA$} 

מופיע \L{$b_{i1}a_{1j}+b_{i2}a_{2j}+\dots+b_{im}a_{mj}$}.

\section{מסקנה}

תהיינה \L{$A\in M_{m\times n}(\mathbb{F})$}, \L{$B\in M_{\ell\times m}(\mathbb{F})$}.
אם השורה ה-\L{$i$} של \L{$B$} היא שורת האפסים,

אז השורה ה-\L{$i$} של \L{$BA$} היא גם שורת האפסים. 

\newpage{}

\noindent\fbox{\begin{minipage}[t]{1\columnwidth - 2\fboxsep - 2\fboxrule}%
\begin{center}
אלגברה ליניארית )1( - \L{$80134$} \textbar{} האונ' העברית \textbar{}
תשפ\textquotedblright ו א' \textbar{} מרצה: ד\textquotedblright ר
אלכס גורביץ' \textbar{} מתרגל: מר יואב כהן \textbar{} ליקוט ועריכה:
נעם לב
\par\end{center}%
\end{minipage}}

\part{הרצאה {\beginL 20\endL} - {\beginL 29.12.2025\endL}}

\section{תזכורת}

\subsection{כפל מטריצות}

תהיינה \L{$A\in M_{m\times n}(\mathbb{F})$}, \L{$B\in M_{\ell\times m}(\mathbb{F})$}.

המכפלה \L{$BA$} היא מטריצה בעלת \L{$n$} עמודות כאשר לכל \L{$1\leq i\leq n$}

העמודה ה-\L{$i$} של \L{$BA$} היא \L{$BC_{i}$} כאשר \L{$C_{i}\in\mathbb{F}^{m}$}
היא העמודה ה-\L{$i$} של \L{$A$}.

אז \L{$BA\in M_{\ell\times n}(\mathbb{F})$}.

\subsection{טענה}

תהיינה \L{$A\in M_{m\times n}(\mathbb{F})$}, \L{$B\in M_{\ell\times m}(\mathbb{F})$},
\L{$C\in M_{k\times\ell}(\mathbb{F})$}.

אז \L{$(CB)A=C(BA)$}.\\

\rule[0.5ex]{1\columnwidth}{1pt}

\section{תכונות של כפל מטריצות}

\subsection{{\beginL 1\endL}.}

אם \L{$A\in M_{m\times n}(\mathbb{F})$}, אז \L{$I_{m}A=A$}.

\subsubsection{הוכחה}

לכל \L{$1\leq i\leq n$} נסמן ב-\L{$C_{i}$} את העמודה ה-\L{$i$}
של \L{$A$}.

אז העמודה ה-\L{$i$} של \L{$I_{m}A$} היא \L{$C_{i}=I_{m}C_{i}$}
לכן \L{$I_{m}A=A$}.

\subsection{{\beginL 2\endL}.}

אם \L{$B\in M_{\ell\times m}(\mathbb{F})$}, אז \L{$BI_{m}=B$}.

\subsubsection{הוכחה}

לכל \L{$1\leq i\leq m$} העמודה ה-\L{$i$} של \L{$I_{m}$} היא \L{$e_{i}$}.

אז העמודה ה-\L{$i$} של \L{$BI_{m}$} היא \L{$Be_{i}$}, וזאת העמודה
ה-\L{$i$} של \L{$B$}.

לכן \L{$BI_{m}=B$}.

\section{מטריצות ריבועיות}

\section{הגדרה - מטריצות הופכיות}

תהיינה \L{$A,B\in M_{m\times m}(\mathbb{F})$}.

אומרים כי \L{$A,B$} הופכיות זו לזו כאשר \L{$AB=BA=I_{m}$}.

\section{טענה}

תהיינה \L{$A,B,C\in M_{m\times m}(\mathbb{F})$} כך ש-\L{$AC=BA=I_{m}$}.

אז \L{$B=C$}.

\subsection{הוכחה}

\L{$C=I_{m}C=(BA)C=B(AC)=BI_{m}=B$}.

\section{מסקנה}

תהיינה \L{$A,B,C\in M_{m\times m}(\mathbb{F})$}.

אם \L{$B$} הופכית ל-\L{$A$} וגם \L{$C$} הופכית ל-\L{$A$}, אז \L{$B=C$}.

\section{הגדרה - מטריצה הפיכה}

מטריצה \L{$A\in M_{m\times m}(\mathbb{F})$} נקראת הפיכה כאשר קיימת
\L{$B\in M_{m\times m}(\mathbb{F})$} שהופכית ל-\L{$A$}.

\section{טענה}

תהיינה \L{$A,C\in M_{m\times m}(\mathbb{F})$} הפיכות.

אז \L{$AC$} הפיכה.

בנוסף, אם \L{$B,D\in M_{m\times m}(\mathbb{F})$} הופכיות ל-\L{$A,C$}
בהתאמה,

אז \L{$DB$} היא הופכית ל-\L{$AC$}.

\subsection{הוכחה}

\L{%
\[
(AC)(DB)=((AC)D)B=(A(CD))B=(AI_{m})B=AB=I_{m}
\]
}%

\L{%
\[
(DB)(AC)=((DB)A)C=(D(BA))C=(DI_{m})C=DC=I_{m}
\]
}%


\section{טענה}

תהי \L{$A\in M_{m\times m}(\mathbb{F})$} הפיכה, ויהי \L{$b\in\mathbb{F}^{m}$}.

אז לממ\textquotedblright ל \L{$Ax=b$} קיים פתרון יחיד.

בנוסף, אם \L{$B\in M_{m\times m}(\mathbb{F})$} היא הופכית ל-\L{$A$},

אז הפתרון היחיד של הממ\textquotedblright ל הוא \L{$Bb$}.

\subsection{הוכחה}

{]}ראשית נוכיח כי \L{$Bb$} הוא פתרון של הממ\textquotedblright ל ואז
נוכיח יחידות.{[}\\

מתקיים \L{$A(Bb)=(AB)b=I_{m}b=b$}.

לכן \L{$Bb$} הוא פתרון של הממ\textquotedblright ל.\\
\\
נניח כי \L{$d\in\mathbb{F}^{m}$} הוא פתרון של הממ\textquotedblright ל.
אז \L{$Ad=b$}.

מכאן \L{$d=I_{m}d=(BA)d=B(Ad)=Bb$}.\\
\\
לכן לממ\textquotedblright ל הזאת קיים פתרון יחיד והוא \L{$Bb$}.

\section{תכונה}

יהי \L{$V$} מ\textquotedblright ו נ\textquotedblright ס, \L{$\mathcal{B}=(v_{1},\dots,v_{n})$}
בסיס של \L{$V$}.

אז \L{$[Id_{V}]_{\mathcal{B}}^{\mathcal{B}}=I_{n}$}.

\subsection{הוכחה}

לכל \L{$1\leq i\leq n$} העמודה ה-\L{$i$} של \L{$[Id_{V}]_{\mathcal{B}}^{\mathcal{B}}$}
היא \L{$e_{i}=[v_{i}]_{\mathcal{B}}=[Id_{V}(v_{i})]_{\mathcal{B}}$}.

לכן \L{$[Id_{V}]_{\mathcal{B}}^{\mathcal{B}}=I_{n}$}.

\section{טענה}

יהיו \L{$V,W$} מ\textquotedblright ו נ\textquotedblright ס, \L{$\mathcal{B},\mathcal{C}$}
בסיסים \L{$V,W$} בהתאמה,

\L{$T:V\rightarrow W$} איזומורפיזם, \L{$T':W\rightarrow V$} ה\textquotedblright ל
הופכית ל-\L{$T$}.

אז המטריצות \L{$[T]_{\mathcal{C}}^{\mathcal{B}}$}, \L{$[T']_{\mathcal{B}}^{\mathcal{C}}$}
הופכיות זו לזו.

\subsection{הוכחה}

\L{%
\[
[T']_{\mathcal{B}}^{\mathcal{C}}[T]_{\mathcal{C}}^{\mathcal{B}}=[T'\circ T]_{\mathcal{B}}^{\mathcal{B}}=[Id_{V}]_{\mathcal{B}}^{\mathcal{B}}=I_{n}
\]
}%
\L{%
\[
[T]_{\mathcal{C}}^{\mathcal{B}}[T']_{\mathcal{B}}^{\mathcal{C}}=[T\circ T']_{\mathcal{C}}^{\mathcal{C}}=[Id_{W}]_{\mathcal{C}}^{\mathcal{C}}=I_{n}
\]
}%


\section{טענה}

יהיו \L{$V,W$} מ\textquotedblright ו נ\textquotedblright ס, \L{$\dim V=\dim W$},
\L{$\mathcal{B},\mathcal{C}$} בסיסים של \L{$V,W$} בהתאמה,

\L{$T:V\rightarrow W$} ה\textquotedblright ל.

אז \L{$T$} איזומורפיזם אם\textquotedblright ם \L{$[T]_{\mathcal{C}}^{\mathcal{B}}$}
הפיכה.

\subsection{הוכחה}

\textbf{\L{$\Leftarrow$}:}

אם \L{$T$} איזומורפיזם, אז קיימת ה\textquotedblright ל \L{$T':W\rightarrow V$}
הופכית ל-\L{$T$}.

אז לפי הטענה הקודמת \L{$[T']_{\mathcal{B}}^{\mathcal{C}}$} הופכית
ל-\L{$[T]_{\mathcal{C}}^{\mathcal{B}}$}. מכאן \L{$[T]_{\mathcal{C}}^{\mathcal{B}}$}
הפיכה וסיימנו.

\textbf{\L{$\Rightarrow$}:}

אם \L{$[T]_{\mathcal{C}}^{\mathcal{B}}$} הפיכה, אז קיימת מטריצה \L{$D$}
שהופכית ל-\L{$[T]_{\mathcal{C}}^{\mathcal{B}}$}.

לפי טענה שהוכחנו פעם, קיימת ה\textquotedblright ל \L{$T':W\rightarrow V$}
כך ש-\L{$[T']_{\mathcal{B}}^{\mathcal{C}}=D$}.

אז \L{$[T'\circ T]_{\mathcal{B}}^{\mathcal{B}}=[T']_{\mathcal{B}}^{\mathcal{C}}[T]_{\mathcal{C}}^{\mathcal{B}}=D[T]_{\mathcal{C}}^{\mathcal{B}}=I_{n}$}
כאשר \L{$n=\dim V$}.

נסמן \L{$\mathcal{B}=(v_{1},\dots,v_{n})$}. אז לכל \L{$1\leq i\leq n$}
מתקיים \L{$[(T'\circ T)(v_{i})]_{\mathcal{B}}=e_{i}$}.

מכאן \L{$(T'\circ T)(v_{i})=v_{i}$}. לכן \L{$T'\circ T=Id_{V}$},
ובדומה \L{$T\circ T'=Id_{W}$}.

לכן \L{$T$} ו-\L{$T'$} ה\textquotedblright ל הופכיות זו לזו, כלומר
\L{$T$} איזומורפיזם וסיימנו.

\section{טענה}

תהי \L{$A\in M_{m\times m}(\mathbb{F})$} כך ש-לממ\textquotedblright ל
הומוגנית \L{$Ax=\begin{pmatrix}0\\
\vdots\\
0
\end{pmatrix}$} קיים פתרון יחיד,

אז \L{$A$} הפיכה.

\subsection{הוכחה}

נגדיר \L{$T_{A}:\mathbb{F}^{m}\rightarrow\mathbb{F}^{m}$} ע\textquotedblright י
\L{$T_{A}(v)=Av$}.

אז \L{$\ker T_{A}=\left\{ v\in\mathbb{F}^{m}\,|\,T_{A}(v)=\begin{pmatrix}0\\
\vdots\\
0
\end{pmatrix}\right\} =\left\{ \begin{pmatrix}0\\
\vdots\\
0
\end{pmatrix}\right\} $}.

לכן \L{$T_{A}$} חח\textquotedblright ע, מכאן \L{$T_{A}$} גם על,
כלומר \L{$T_{A}$} איזומורפיזם.

לפי הטענה הקודמת \L{$[T_{A}]_{\varepsilon_{m}}^{\varepsilon_{m}}=A$}
הפיכה כאשר \L{$\varepsilon_{m}=(e_{1},\dots,e_{m})$} הבסיס הסטנדרטי
של \L{$\mathbb{F}^{m}$}.

\newpage{}

\noindent\fbox{\begin{minipage}[t]{1\columnwidth - 2\fboxsep - 2\fboxrule}%
\begin{center}
אלגברה ליניארית )1( - \L{$80134$} \textbar{} האונ' העברית \textbar{}
תשפ\textquotedblright ו א' \textbar{} מרצה: ד\textquotedblright ר
אלכס גורביץ' \textbar{} מתרגל: מר יואב כהן \textbar{} ליקוט ועריכה:
נעם לב
\par\end{center}%
\end{minipage}}

\part{הרצאה {\beginL 21\endL} - {\beginL 30.12.2025\endL}}

\section{תזכורת}

תהיינה \L{$A,B\in M_{m\times m}(\mathbb{F})$}. אומרים כי \L{$A,B$}
הופכיות זו לזו כאשר \L{$AB=BA=I_{m}$}.

\section{טענה}

תהיינה \L{$A,B,C\in M_{m\times m}(\mathbb{F})$} כך ש-\L{$AC=BA=I_{m}$}.

אז \L{$B=C$}.\\
\\
מטענה זו הסקנו את היחידות.

\section{הגדרה - מטריצה הפיכה}

מטריצה \L{$A\in M_{m\times m}(\mathbb{F})$} נקראת הפיכה כאשר קיימת
\L{$B\in M_{m\times m}(\mathbb{F})$} שהופכית ל-\L{$A$}.

\section{טענה {\beginL 1\endL}}

תהי \L{$A\in M_{m\times m}(\mathbb{F})$}, ויהי \L{$b\in\mathbb{F}^{m}$}.

אם \L{$A$} הפיכה, אז לממ\textquotedblright ל \L{$Ax=b$} קיים פתרון
יחיד.

בנוסף, אם \L{$B\in M_{m\times m}(\mathbb{F})$} היא הופכית ל-\L{$A$},

אז הפתרון היחיד של הממ\textquotedblright ל הוא \L{$Bb$}.

\section{טענה}

יהיו \L{$V,W$} מ\textquotedblright ו נ\textquotedblright ס כך ש-
\L{$\dim V=\dim W$}, \L{$\mathcal{B},\mathcal{C}$} בסיסים של \L{$V,W$}
בהתאמה,

\L{$T:V\rightarrow W$} ה\textquotedblright ל.

אז \L{$T$} איזומורפיזם אם\textquotedblright ם \L{$[T]_{\mathcal{C}}^{\mathcal{B}}$}
הפיכה.

\section{טענה {\beginL 2\endL}}

תהי \L{$A\in M_{m\times m}(\mathbb{F})$}. אם לממ\textquotedblright ל
הומוגנית \L{$Ax=\begin{pmatrix}0\\
\vdots\\
0
\end{pmatrix}$} קיים פתרון יחיד,

אז \L{$A$} הפיכה.\\

\rule[0.5ex]{1\columnwidth}{1pt}

\section{טענה}

תהיינה \L{$A,B\in M_{m\times m}(\mathbb{F})$}. אם \L{$BA=I_{m}$},
אז \L{$AB=I_{m}$}.\\
\\
)חשוב לזכור שזה תקף רק למטריצות ריבועיות.(

\subsection{הוכחה}

יהי \L{$d\in\mathbb{F}^{m}$} פתרון של \L{$Ax=\begin{pmatrix}0\\
\vdots\\
0
\end{pmatrix}$}. אז \L{$Ad=\begin{pmatrix}0\\
\vdots\\
0
\end{pmatrix}$}.

מכאן \L{$d=I_{m}d=(BA)d=B(Ad)=B\begin{pmatrix}0\\
\vdots\\
0
\end{pmatrix}=\begin{pmatrix}0\\
\vdots\\
0
\end{pmatrix}$}.

לכן לממ\textquotedblright ל \L{$Ax=\begin{pmatrix}0\\
\vdots\\
0
\end{pmatrix}$} קיים פתרון יחיד.

לפי הטענה הקודמת אנו מסיקים כי \L{$A$} הפיכה.

נסמן ב-\L{$C$} את המטריצה ההופכית ל-\L{$A$}.

אז \L{$CA=AC=I_{m}$}. מכאן \L{$C=B$} ולכן \L{$AB=I_{m}$}.

\section{טענה {\beginL 3\endL}}

תהי \L{$A\in M_{m\times m}(\mathbb{F})$}. אם לכל \L{$b\in\mathbb{F}^{m}$}
לממ\textquotedblright ל \L{$Ax=b$} קיים פתרון,

אז \L{$A$} הפיכה.

\subsection{הוכחה}

נגדיר \L{$T_{A}:\mathbb{F}^{m}\rightarrow\mathbb{F}^{m}$} ע\textquotedblright י
\L{$T_{A}(v)=Av$}. אז
\begin{flushright}
\L{$\operatorname{Im}T=\left\{ T_{A}(v)\,|\,v\in\mathbb{F}^{m}\right\} =\left\{ Av\,|\,v\in\mathbb{F}^{m}\right\} =\left\{ b\in\mathbb{F}^{m}\,|\,Ax=b\quad\text{has a solution}\right\} \underset{(*)}{=}\mathbb{F}^{m}$}
\par\end{flushright}

){*}( נתון.

לכן \L{$T_{A}$} היא על, ומכאן \L{$T_{A}$} היא חח\textquotedblright ע,
כלומר \L{$T_{A}$} איזומורפיזם.

אז \L{$[T_{A}]_{\varepsilon_{m}}^{\varepsilon_{m}}=A$} הפיכה, כאשר
\L{$\varepsilon_{m}=(e_{1},\dots,e_{m})$} הבסיס הסטנדרטי של \L{$\mathbb{F}^{m}$}.

\newpage{}

\section{משפט מסכם של טענות {\beginL 1,2,3\endL}}

תהי \L{$A\in M_{m\times m}(\mathbb{F})$}. אז התנאים הבאים שקולים
זה לזה:
\begin{enumerate}
\item \L{$A$} הפיכה.
\item לכל \L{$b\in\mathbb{F}^{m}$} לממ\textquotedblright ל \L{$Ax=b$}
קיים פתרון יחיד.
\item לממ\textquotedblright ל \L{$Ax=\begin{pmatrix}0\\
\vdots\\
0
\end{pmatrix}$} קיים פתרון יחיד.
\item לכל \L{$b\in\mathbb{F}^{m}$} לממ\textquotedblright ל \L{$Ax=b$}
קיים פתרון.
\end{enumerate}

\section{הוכחה}

\textquotedblleft אני כתבתי פה הוכחה, אבל אני לא הולך לכתוב הוכחה,
אני הולך לכתוב ציור :(\textquotedblright{} - אלכס.\\
\L{%
\begin{align*}
 & (1)\underset{\text{Proposition 1}}{\Rightarrow}(2)\\
 & (2)\underset{\text{trivial}}{\Rightarrow}(3)\\
 & (2)\underset{\text{trivial}}{\Rightarrow}(4)\\
 & (3)\underset{\text{Proposition 2}}{\Rightarrow}(1)\\
 & (4)\underset{\text{Proposition 3}}{\Rightarrow}(1)
\end{align*}
}%


\section{טענה}

תהי \L{$A\in M_{m\times m}(\mathbb{F})$}. אז \L{$A$} הפיכה אם\textquotedblright ם
סדרת העמודות של \L{$A$} בת\textquotedblright ל. 

\subsection{הוכחה}

נסמן את העמודות של \L{$A$} ב-\L{$C_{1},\dots,C_{m}$}.

אז \L{$(C_{1},\dots,C_{m})$} בת\textquotedblright ל אם\textquotedblright ם

\L{$a_{1}C_{1}+\dots+a_{m}C_{m}=\begin{pmatrix}0\\
\vdots\\
0
\end{pmatrix}$} רק כאשר \L{$a_{1}=\dots=a_{m}=0$} אם\textquotedblright ם

\L{$A\begin{pmatrix}a_{1}\\
\vdots\\
a_{m}
\end{pmatrix}=\begin{pmatrix}0\\
\vdots\\
0
\end{pmatrix}$} רק כאשר \L{$\begin{pmatrix}a_{1}\\
\vdots\\
a_{m}
\end{pmatrix}=\begin{pmatrix}0\\
\vdots\\
0
\end{pmatrix}$} אם\textquotedblright ם

לממ\textquotedblright ל \L{$Ax=\begin{pmatrix}0\\
\vdots\\
0
\end{pmatrix}$} קיים פתרון יחיד אם\textquotedblright ם \L{$A$} הפיכה וסיימנו.

\section{טענה}

תהי \L{$A\in M_{m\times m}(\mathbb{F})$}. אז \L{$A$} הפיכה אם\textquotedblright ם
סדרת העמודות של \L{$A$} פורשת את \L{$\mathbb{F}^{m}$}.

\subsection{הוכחה}

נסמן את העמודות של \L{$A$} ב-\L{$C_{1},\dots,C_{m}$}.

\L{$(C_{1},\dots,C_{m})$} פורשת את \L{$\mathbb{F}^{m}$} אם\textquotedblright ם

לכל \L{$b\in\mathbb{F}^{m}$} קיימים \L{$a_{1},\dots,a_{m}\in\mathbb{F}$}
כך ש-\L{$a_{1}C_{1}+\dots+a_{m}C_{m}=b$} אם\textquotedblright ם

לכל \L{$b\in\mathbb{F}^{m}$} לממ\textquotedblright ל \L{$Ax=b$}
קיים פתרון אם\textquotedblright ם \L{$A$} הפיכה וסיימנו.

\section{מטריצות אלמנטריות}

\section{הגדרה של \L{$T_{\varepsilon}$}}

תהי \L{$\varepsilon$} פעולת שורה אלמנטרית על מטריצות בעלות \L{$m$}
שורות.

בהינתן מטריצה \L{$A$} בעלת \L{$m$} שורות נסמן ב-\L{$\varepsilon(A)$}
את המטריצה המתקבלת מ-\L{$A$} ע\textquotedblright י ביצוע של \L{$\varepsilon$}.

נגדיר \L{$T_{\varepsilon}:\mathbb{F}^{m}\rightarrow\mathbb{F}^{m}$}
באופן הבא: \L{$\forall v\in\mathbb{F}^{m}\quad T_{\varepsilon}(v)=\varepsilon(v)$}.

קל לבדוק כי \L{$T_{\varepsilon}$} היא ה\textquotedblright ל.

\section{טענה}

\L{$[T_{\varepsilon}]_{\varepsilon_{m}}^{\varepsilon_{m}}=\varepsilon(I_{m})$}
כאשר \L{$\varepsilon_{m}=(e_{1},\dots,e_{m})$} זהו הבסיס הסטנדרטי
של \L{$\mathbb{F}^{m}$}.

\subsection{הוכחה}

העמודה ה-\L{$i$} של \L{$[T_{\varepsilon}]_{\varepsilon_{m}}^{\varepsilon_{m}}$}
היא \L{$\varepsilon(e_{i})=T_{\varepsilon}(e_{i})=[T_{\varepsilon}(e_{i})]_{\varepsilon_{m}}$}.

העמודה ה-\L{$i$} של \L{$\varepsilon(I_{m})$} היא \L{$\varepsilon(e_{i})$}.
מכאן הנדרש.

\subsection{דוגמה\qquad{} \L{$m=2$}}

{\beginL 1\endL}( \L{$\varepsilon=R_{1}\leftrightarrow R_{2}$}. אז
\L{$\varepsilon(I_{2})=\begin{pmatrix}0 & 1\\
1 & 0
\end{pmatrix}$}.

{\beginL 2\endL}( \L{$\varepsilon=R_{1}\rightarrow2R_{1}$}. אז \L{$\varepsilon(I_{2})=\begin{pmatrix}2 & 0\\
0 & 1
\end{pmatrix}$}.

{\beginL 3\endL}( \L{$\varepsilon=R_{2}\rightarrow R_{2}+R_{1}$}.
אז \L{$\varepsilon(I_{2})=\begin{pmatrix}1 & 0\\
1 & 1
\end{pmatrix}$}.

\section{טענה}

תהי \L{$\varepsilon$} פש\textquotedblright א על מטריצות בעלות \L{$m$}
שורות, \L{$A\in M_{m\times m}(\mathbb{F})$}. אז \L{$\varepsilon(A)=\varepsilon(I_{m})A$}.

\subsection{הוכחה}

לכל \L{$1\leq i\leq m$} נסמן את העמודה ה-\L{$i$} של \L{$A$} ב-\L{$C_{i}$}.

אז העמודה ה-\L{$i$} של \L{$\varepsilon(A)$} היא \L{$\varepsilon(C_{i})$}.

העמודה ה -\L{$i$} של \L{$\varepsilon(I_{m})A$} היא 

\L{%
\[
\varepsilon(C_{i})=T_{\varepsilon}(C_{i})=[T_{\varepsilon}(C_{i})]_{\varepsilon_{m}}=[T_{\varepsilon}]_{\varepsilon_{m}}^{\varepsilon_{m}}[C_{i}]_{\varepsilon_{m}}=\varepsilon(I_{m})C_{i}
\]
}%

מכאן הנדרש.

\subsection{דוגמה\qquad{} \L{$m=2$}, \L{$A=\begin{pmatrix}1 & 2\\
3 & 4
\end{pmatrix}$}}

{\beginL 1\endL}( \L{$\varepsilon=R_{1}\leftrightarrow R_{2}$}. אז
\L{$\varepsilon(A)=\begin{pmatrix}3 & 4\\
1 & 2
\end{pmatrix}$}.

{\beginL 2\endL}( \L{$\varepsilon=R_{1}\rightarrow2R_{1}$}. אז \L{$\varepsilon(A)=\begin{pmatrix}2 & 4\\
3 & 4
\end{pmatrix}$}.

{\beginL 3\endL}( \L{$\varepsilon=R_{2}\rightarrow R_{2}+R_{1}$}.
אז \L{$\varepsilon(A)=\begin{pmatrix}1 & 2\\
4 & 6
\end{pmatrix}$}.

\section{הגדרה - מטריצה אלמנטרית}

מטריצה \L{$E\in M_{m\times m}(\mathbb{F})$} נקראת אלמנטרית כאשר קיימת
פש\textquotedblright א \L{$\varepsilon$} כך ש-\L{$E=\varepsilon(I_{m})$}.

\section{טענה}

כל מטריצה אלמנטרית היא הפיכה.

\subsection{הוכחה}

תהי \L{$E\in M_{m\times m}(\mathbb{F})$} מטריצה אלמנטרית. אז קיימת
פש\textquotedblright א \L{$\varepsilon$} כך ש-\L{$E=\varepsilon(I_{m})$}.

נסמן ב-\L{$\delta$} את הפש\textquotedblright א ההופכית ל-\L{$\varepsilon$}.
נסמן ב-\L{$D=\delta(I_{m})$}.

אז מתקיים \L{$ED=\varepsilon(I_{m})D=\varepsilon(D)=\varepsilon(\delta(I_{m}))=I_{m}$}.

לכן \L{$E,D$} הופכיות זו לזו ומכאן \L{$E$} הפיכה.

\textcolor{violet}{ההוכחה מראה רק ש-\L{$ED=I_{m}$}. הכיוון השני,
\L{$DE=I_{m}$}, נובע מהטענה שהוכחנו קודם בהרצאה זו: אם \L{$BA=I_{m}$}
אז \L{$AB=I_{m}$} )למטריצות ריבועיות(.}

\section{משפט}

תהיינה \L{$A,B\in M_{m\times m}(\mathbb{F})$} שקולות שורה, כלומר
\L{$B$} מתקבלת מ-\L{$A$} ע\textquotedblright י סדרה של פש\textquotedblright א.

אז קיימת מטריצה הפיכה \L{$P\in\sqrm$} כך ש-\L{$B=PA$}. 

\subsection{הוכחה}

נסמן ב-\L{$\varepsilon_{k},\dots,\varepsilon_{1}$} את הפש\textquotedblright א
שבעזרתן \L{$B$} מתקבל מ-\L{$A$},

כלומר \L{$B=\varepsilon_{k}(\dots\varepsilon_{2}(\varepsilon_{1}(A))\dots)$}.
לפי הטענה הקודמת מתקיים
\begin{flushright}
\L{$B=\varepsilon_{k}(\dots\varepsilon_{3}(\varepsilon_{2}(\varepsilon_{1}(I_{m})A))\dots)=\varepsilon_{k}(\dots\varepsilon_{3}(\varepsilon_{2}(I_{m})(\varepsilon_{1}(I_{m})A)))=\dots=\varepsilon_{k}(I_{m})(\dots(\varepsilon_{2}(I_{m})(\varepsilon_{1}(I_{m})A))\dots)$}
\par\end{flushright}

נגדיר \L{$P=\varepsilon_{k}(I_{m})(\dots(\varepsilon_{2}(I_{m})\cdot\varepsilon_{1}(I_{m}))\dots)$}.
אז \L{$B=PA$}.

מאחר שלכל \L{$1\leq i\leq k$} המטריצה \L{$\varepsilon_{i}(I_{m})$}
הפיכה, גם \L{$P$} הפיכה.

\newpage{}

\noindent\fbox{\begin{minipage}[t]{1\columnwidth - 2\fboxsep - 2\fboxrule}%
\begin{center}
אלגברה ליניארית )1( - \L{$80134$} \textbar{} האונ' העברית \textbar{}
תשפ\textquotedblright ו א' \textbar{} מרצה: ד\textquotedblright ר
אלכס גורביץ' \textbar{} מתרגל: מר יואב כהן \textbar{} ליקוט ועריכה:
נעם לב
\par\end{center}%
\end{minipage}}

\part{הרצאה {\beginL 22\endL} - {\beginL 01.01.2026\endL}}

\textcolor{teal}{זה היה שיעור השלמה ביום חמישי.}

\section{תזכורת}

\subsection{טענה}

תהיינה \L{$A,B\in M_{m\times n}(\mathbb{F})$} שקולות שורה.

אז קיימת \L{$P\in M_{m\times m}(\mathbb{F})$} הפיכה כך ש- \L{$B=PA$}.

\subsection{הוכחה}

נסמן ב-\L{$\varepsilon_{k},\dots,\varepsilon_{1}$} סדרה של פש\textquotedblright א
שבעזרתן \L{$B$} מתקבלת מ-\L{$A$}.

אז \L{$B=\varepsilon_{k}(\dots\varepsilon_{2}(\varepsilon_{1}(A))\dots)$}.

מכאן \L{%
\[
B=\varepsilon_{k}(I_{m})(\dots(\varepsilon_{2}(I_{m})\varepsilon_{1}(I_{m})A)\dots)
\]
}%
נגדיר \L{$P=\varepsilon_{k}(I_{m})\dots\varepsilon_{2}(I_{m})\varepsilon_{1}(I_{m})$}.
אז \L{$B=PA$} מאחר שלכל \L{$1\le i\le k$}, \L{$\varepsilon_{i}(I_{m})$}
הפיכה, גם \L{$P$} הפיכה.

\subsubsection{הערה}

בהנחות הוכחת הטענה, מתקיים\L{%
\[
P=\varepsilon_{k}(I_{m})(\dots(\varepsilon_{3}(I_{m})(\varepsilon_{2}(I_{m})\varepsilon_{1}(I_{m})))\dots)=\varepsilon_{k}(I_{m})(\dots(\varepsilon_{3}(I_{m})\varepsilon_{2}(\varepsilon_{1}(I_{m})))\dots)=
\]
}%
\L{%
\[
=\varepsilon_{k}(I_{m})(\dots(\varepsilon_{3}(\varepsilon_{2}(\varepsilon_{1}(I_{m})))\dots)=\dots=\varepsilon_{k}(\dots(\varepsilon_{3}(\varepsilon_{2}(\varepsilon_{1}(I_{m}))))\dots)
\]
}%
כלומר \L{$P$} מתקבלת מ-\L{$I_{m}$} ע\textquotedblright י סדרה של
פש\textquotedblright א \L{$\varepsilon_{k},\dots,\varepsilon_{1}$}.

\section{דוגמא}

נתונה \L{$A=\begin{pmatrix}1 & -1\\
-1 & 3\\
0 & 1
\end{pmatrix}$}.

נמצא מטריצה מדורגת מצומצמת \L{$B$} ש-שקולת שורה ל-\L{$A$} ומטריצה
הפיכה \L{$P$} כך ש- \L{$B=PA$}.

\L{%
\[
\left(\begin{array}{cc|ccc}
1 & -1 & 1 & 0 & 0\\
-1 & 3 & 0 & 1 & 0\\
0 & 1 & 0 & 0 & 1
\end{array}\right)\xrightarrow{R_{2}\to R_{2}+R_{1}}\left(\begin{array}{cc|ccc}
1 & -1 & 1 & 0 & 0\\
0 & 2 & 1 & 1 & 0\\
0 & 1 & 0 & 0 & 1
\end{array}\right)\xrightarrow{R_{2}\to\frac{1}{2}R_{2}}\left(\begin{array}{cc|ccc}
1 & -1 & 1 & 0 & 0\\
0 & 1 & \frac{1}{2} & \frac{1}{2} & 0\\
0 & 1 & 0 & 0 & 1
\end{array}\right)
\]
}%
\L{%
\[
\xrightarrow[R_{3}\to R_{3}-R_{2}]{R_{1}\to R_{1}+R_{2}}\left(\begin{array}{cc|ccc}
1 & 0 & 3/2 & 1/2 & 0\\
0 & 1 & 1/2 & 1/2 & 0\\
0 & 0 & -1/2 & -1/2 & 1
\end{array}\right)
\]
}%
ולכן: \L{%
\[
B=\begin{pmatrix}1 & 0\\
0 & 1\\
0 & 0
\end{pmatrix},\quad P=\begin{pmatrix}3/2 & 1/2 & 0\\
1/2 & 1/2 & 0\\
-1/2 & -1/2 & 1
\end{pmatrix}
\]
}%


\section{למה {]}למטריצה הפיכה אין שורות או עמודות אפסים{[}}

תהי \L{$A\in M_{m\times m}(\mathbb{F})$} הפיכה. אז ל- \L{$A$} אין
שורות אפסים ואין עמודות אפסים.

\subsection{הוכחה}

נסמן ב- \L{$B$} את המטריצה ההופכית ל- \L{$A$}. אז \L{$AB=BA=I_{m}$}.

אם השורה ה- \L{$i$} של \L{$A$} היא שורת אפסים, אז השורה ה- \L{$i$}
של \L{$I_{m}=AB$} היא שורת אפסים, וזו סתירה.

אם העמודה ה- \L{$i$} של \L{$A$} היא עמודת אפסים, אז העמודה ה- \L{$i$}
של \L{$I_{m}=BA$} היא עמודת אפסים, וזו סתירה.

\section{טענה}

תהי \L{$R\in M_{m\times m}(\mathbb{F})$} מטריצה מדורגת מצומצמת והפיכה.

אז \L{$R=I_{m}$}.

\subsection{הוכחה}

ל- \L{$R$} יש בדיוק \L{$m$} איברים מובילים כי אחרת ב- \L{$R$} הייתה
שורה ללא איבר מוביל כלומר שורת אפסים וזאת סתירה ללמה.

לכן בכל שורה של \L{$R$} יש איבר מוביל.

האיבר המוביל של השורה הראשונה חייב להיות בעמודה הראשונה כי אחרת לא
יספיק מקום לשאר האיברים המובילים. בדומה לכל \L{$1\leq i\leq m$} האיבר
המוביל של השורה ה- \L{$i$} חייב להיות בעמודה ה- \L{$i$}. מכאן \L{$R=I_{m}$}.\\


\section{משפט {]}חשוב{[}}

תהי \L{$A\in M_{m\times m}(\mathbb{F})$}. נסמן ב-\L{$R$} מטריצה
מדורגת מצומצמת שקולת שורה ל-\L{$A$}. 

אזי \L{$A$} הפיכה אם\textquotedblright ם \L{$R=I_{m}$}.

\subsection{הוכחה}

קיימת \L{$P\in M_{m\times m}(\mathbb{F})$} הפיכה כך ש-\L{$R=PA$}.

\L{$\Leftarrow$} אם \L{$A$} הפיכה, אז \L{$R=PA$} הפיכה אז לפי הטענה
הקודמת \L{$R=I_{m}$}.

\L{$\Rightarrow$} אם \L{$R=I_{m}$}, אז \L{$I_{m}=PA$} מכאן \L{$I_{m}=AP$}
לכן \L{$A,P$} הופכיות זו לזו ובפרט \L{$A$} הפיכה.

\section{הערה}

אם \L{$A$} הפיכה, אז כדי לחשב את המטריצה ההופכית ל-\L{$A$} )כלומר
\L{$P$}( צריך לבצע על \L{$I_{m}$} את {]}ההקלטה נקטעה כאן{[} {]}ניחוש:
אותן הפעולות שמבצעים על \L{$A$} כדי להביאה לצורה המדורגת קנונית \L{$I_{m}$}{[}.

\newpage{}

\noindent\fbox{\begin{minipage}[t]{1\columnwidth - 2\fboxsep - 2\fboxrule}%
\begin{center}
אלגברה ליניארית )1( - \L{$80134$} \textbar{} האונ' העברית \textbar{}
תשפ\textquotedblright ו א' \textbar{} מרצה: ד\textquotedblright ר
אלכס גורביץ' \textbar{} מתרגל: מר יואב כהן \textbar{} ליקוט ועריכה:
נעם לב
\par\end{center}%
\end{minipage}}

\part{הרצאה {\beginL 23\endL} - {\beginL 05.01.2026\endL}}

\section{דרגה}

יהיו \L{$V,W$} מ\textquotedblright ו נ\textquotedblright ס, \L{$T:V\to W$}
ה"ל. \textbf{הדרגה} של \L{$T$} מוגדרת ע"י: \L{%
\[
rk\ T=\dim\ \operatorname{Im}T
\]
}%


\section{למה}

יהיו \L{$U,V,W$} מ"ו, \L{$T_{1}:U\to V$}, \L{$T_{2}:V\to W$} ה"ל
. אז: 
\begin{enumerate}
\item \L{$\operatorname{Im}(T_{2}\circ T_{1})\subseteq\operatorname{Im}T_{2}$} 
\item \L{$\ker(T_{2}\circ T_{1})\supseteq\ker T_{1}$} 
\end{enumerate}

\subsection{הוכחה}
\begin{enumerate}
\item יהי \L{$w\in\operatorname{Im}(T_{2}\circ T_{1})$}. אז קיים \L{$u\in U$}
כך ש- \L{%
\[
T_{2}(T_{1}(u))=(T_{2}\circ T_{1})(u)=w
\]
}%
מכאן \L{$w\in\operatorname{Im}T_{2}$}.
\item יהי \L{$u\in\ker T_{1}$}. אז \L{$T_{1}(u)=0_{V}$}. מכאן \L{%
\[
(T_{2}\circ T_{1})(u)=T_{2}(T_{1}(u))=T_{2}(0_{V})=0_{W}
\]
}%
לכן \L{$u\in\ker(T_{2}\circ T_{1})$}. 
\end{enumerate}

\section{טענה}

יהיו \L{$U,V,W$} מ"ו נ\textquotedblright ס, \L{$T_{1}:U\to V$},
\L{$T_{2}:V\to W$} ה"ל. אז: 
\begin{enumerate}
\item \L{$rk\ (T_{2}\circ T_{1})\le rk\ T_{2}$} 
\item אם \L{$T_{1}$} על, אז \L{$rk(T_{2}\circ T_{1})=rk\ T_{2}$} 
\item \L{$rk(T_{2}\circ T_{1})\le rkT_{1}$} 
\item אם \L{$T_{2}$} חח"ע, אז \L{$rk(T_{2}\circ T_{1})=rk\ T_{1}$} 
\end{enumerate}

\subsection{הוכחה}

{\beginL 1\endL}( מאחר ש- \L{$\operatorname{Im}(T_{2}\circ T_{1})\subseteq\operatorname{Im}T_{2}$}
מתקיים: \L{%
\[
rk(T_{2}\circ T_{1})=\dim\operatorname{Im}(T_{2}\circ T_{1})\le\dim\operatorname{Im}T_{2}=rkT_{2}
\]
}%

{\beginL 2\endL}( מתקיים \L{$T_{1}(U)=\operatorname{Im}T_{1}=V$}.
אז \L{%
\[
rkT_{2}=\dim\operatorname{Im}T_{2}=\dim T_{2}(V)=\dim T_{2}(T_{1}(U))=\dim(T_{2}\circ T_{1})(U)=\dim\operatorname{Im}(T_{2}\circ T_{1})=rk(T_{2}\circ T_{1})
\]
}%

{\beginL 3\endL}( מאחר ש- \L{$\ker(T_{2}\circ T_{1})\supseteq\ker T_{1}$}
מתקיים \L{$\dim\ker(T_{2}\circ T_{1})\ge\dim\ker T_{1}$}. אז \L{%
\[
rk(T_{2}\circ T_{1})=\dim\operatorname{Im}(T_{2}\circ T_{1})=\dim U-\dim\ker(T_{2}\circ T_{1})
\]
}%
\L{%
\[
=\dim\ \operatorname{Im}T_{1}+\left(\dim\ker T_{1}-\dim\ker(T_{2}\circ T_{1})\right)\le\dim\operatorname{Im}T_{1}=rk\ T_{1}
\]
}%

{\beginL 4\endL}( נוכיח כי \L{$\ker(T_{2}\circ T_{1})\subseteq\ker T_{1}$}.

יהי \L{$u\in\ker(T_{2}\circ T_{1})$}. אז \L{$T_{2}(T_{1}(u))=(T_{2}\circ T_{1})(u)=0_{W}$}.

מאחר ש- \L{$T_{2}$} חח"ע מתקיים \L{$T_{1}(u)\in\ker T_{2}=\{0_{V}\}$}.

לכן \L{$T_{1}(u)=0_{V}$}, כלומר \L{$u\in\ker T_{1}$}.

אז לפי הלמה מתקיים \L{$\ker(T_{2}\circ T_{1})=\ker T_{1}$}.

\section{מסקנה}

יהיו \L{$U,V,W,Z$} מ"ו נ\textquotedblright ס, \L{$T_{1}:U\to V$}
,\L{$T:V\to W$} , \L{$T_{2}:W\to Z$} ה"ל כך ש- \L{$T_{1},T_{2}$}
איזומורפיזמים. אז \L{%
\[
rk(T_{2}\circ T\circ T_{1})=rkT
\]
}%


\subsection{הוכחה}

\L{%
\[
rk(T_{2}\circ T\circ T_{1})\underset{T_{1}\text{ surjective}}{=}rk(T_{2}\circ T)\underset{T_{2}\text{injective}}{=}rk\ T
\]
}%


\section{דרגה של מטריצה}

תהי \L{$A\in M_{m\times n}(\mathbb{F})$}. 

נתבונן בה"ל \L{$T_{A}:\mathbb{F}^{n}\to\mathbb{F}^{m}$} המוגדרת
ע"י \L{$T_{A}(v)=Av$}.

הדרגה של \L{$A$} מוגדרת ע"י: \L{%
\[
rkA=rk\ T_{A}
\]
}%


\section{טענה}

יהיו \L{$V,W$} מ"ו נ\textquotedblright ס, \L{$\mathcal{B},\mathcal{C}$}
בסיסים של \L{$V,W$} בהתאמה, \L{$T:V\to W$} ה"ל. אזי \L{%
\[
rk\ T=rk\ [T]_{\mathcal{C}}^{\mathcal{B}}
\]
}%


\subsection{הוכחה}

נסמן \L{$n=\dim V$}, \L{$m=\dim W$}. נסמן \L{$A=[T]_{\mathcal{C}}^{\mathcal{B}}$}.

אז \L{$T=\tau_{\mathcal{C}}\circ T_{A}\circ\omega_{\mathcal{B}}$}.
לפי המסקנה הקודמת מתקיים \L{%
\[
rk\ T=rk(\tau_{\mathcal{C}}\circ T_{A}\circ\omega_{\mathcal{B}})=rk\ T_{A}=rk\ A=rk[T]_{\mathcal{C}}^{\mathcal{B}}
\]
}%

\noindent\fbox{\begin{minipage}[t]{1\columnwidth - 2\fboxsep - 2\fboxrule}%
דיאגרמה: \L{%
\[
\begin{array}{ccc}
V & \xrightarrow{T} & W\\
\omega_{\mathcal{B}}\downarrow\uparrow\tau_{\mathcal{B}} &  & \tau_{\mathcal{C}}\uparrow\downarrow\omega_{\mathcal{C}}\\
\mathbb{F}^{n} & \xrightarrow[T_{A}]{} & \mathbb{F}^{m}
\end{array}
\]
}%
%
\end{minipage}}

\section{טענה}

תהיינה \L{$A\in M_{m\times n}(\mathbb{F})$}, \L{$A_{1}\in M_{n\times n}(\mathbb{F})$},
\L{$A_{2}\in M_{m\times m}(\mathbb{F})$}. אם \L{$A_{1},A_{2}$} הפיכות,
אז: \L{%
\[
rk(A_{2}AA_{1})=rkA
\]
}%


\subsection{הוכחה}

נתבונן בה"ל \L{$T_{A}:\mathbb{F}^{n}\to\mathbb{F}^{m}$}, \L{$T_{A_{1}}:\mathbb{F}^{n}\to\mathbb{F}^{n}$},
\L{$T_{A_{2}}:\mathbb{F}^{m}\to\mathbb{F}^{m}$}. מאחר ש- \L{$A_{1},A_{2}$}
הפיכות, \L{$T_{A_{1}},T_{A_{2}}$} איזומורפיזמים. אז \L{%
\[
rk(A_{2}AA_{1})=rk([T_{A_{2}}]_{\varepsilon_{m}}^{\varepsilon_{m}}[T_{A}]_{\varepsilon_{m}}^{\varepsilon_{n}}[T_{A_{1}}]_{\varepsilon_{n}}^{\varepsilon_{n}})
\]
}%
\L{%
\[
=rk([T_{A_{2}}\circ T_{A}\circ T_{A_{1}}]_{\varepsilon_{m}}^{\varepsilon_{n}})\underset{\text{prev theorem}}{=}rk(T_{A_{2}}\circ T_{A}\circ T_{A_{1}})=rkT_{A}=rkA
\]
}%


\section{טענה {]}למטריצות שקולות שורה יש אותה דרגה{[}}

תהיינה \L{$A,B\in M_{m\times n}(\mathbb{F})$} שקולות שורה. אז \L{%
\[
rkB=rkA
\]
}%


\subsection{הוכחה}

קיימת \L{$P\in M_{m\times m}(\mathbb{F})$} הפיכה כך ש- \L{$B=PA$}.
אז לפי הטענה הקודמת \L{%
\[
rkB=rk(PAI_{n})=rkA
\]
}%


\section{שחלוף מטריצות )\L{$\text{Transpose}$}(}

תהי \L{$A=\begin{pmatrix}a_{11} & \dots & a_{1n}\\
\vdots &  & \vdots\\
a_{m1} & \dots & a_{mn}
\end{pmatrix}\in M_{m\times n}(\mathbb{F})$}.

\textbf{המטריצה המשוחלפת} של \L{$A$} מוגדרת ע"י: \L{%
\[
A^{t}=\begin{pmatrix}a_{11} & \dots & a_{m1}\\
\vdots &  & \vdots\\
a_{1n} & \dots & a_{mn}
\end{pmatrix}\in M_{n\times m}(\mathbb{F})
\]
}%


\section{טענה}

תהיינה \L{$A\in M_{m\times n}(\mathbb{F})$}, \L{$B\in M_{\ell\times m}(\mathbb{F})$}.
אז \L{%
\[
(BA)^{t}=A^{t}B^{t}
\]
}%


\subsection{הוכחה}

נסמן \L{%
\[
B=\begin{pmatrix}b_{11} & \dots & b_{1m}\\
\vdots &  & \vdots\\
b_{\ell1} & \dots & b_{\ell m}
\end{pmatrix},\quad A=\begin{pmatrix}a_{11} & \dots & a_{1n}\\
\vdots &  & \vdots\\
a_{m1} & \dots & a_{mn}
\end{pmatrix}
\]
}%

יהיו \L{$1\le i\le n$}, \L{$1\le j\le\ell$}. אז בחיתוך של השורה
ה- \L{$i$} והעמודה ה- \L{$j$} של \L{$A^{t}B^{t}$} מופיע \L{%
\[
a_{1i}b_{j1}+a_{2i}b_{j2}+\dots+a_{mi}b_{jm}
\]
}%
מצד שני המספר בשורה ה-\L{$i$} והעמודה ה-\L{$j$} של \L{$(BA)^{t}$}
שווה למספר בשורה ה-\L{$j$} והעמודה ה-\L{$i$} של \L{$BA$} אשר שווה
ל- \L{%
\[
b_{j1}a_{1i}+b_{j2}a_{2i}+\dots+b_{jm}a_{mi}
\]
}%
\textcolor{violet}{שזה שווה לביטוי הקודם )חילוף בכפל סקלרים(}. מכאן
הנדרש.

\section{טענה}

תהי \L{$A\in M_{m\times m}(\mathbb{F})$} הפיכה. אז \L{$A^{t}$} גם
הפיכה.

בנוסף, אם \L{$B$} הופכית ל-\L{$A$}, אז \L{$B^{t}$} הופכית ל-\L{$A^{t}$}.

\subsection{הוכחה}

מתקיים

\L{%
\[
A^{t}B^{t}=(BA)^{t}=(I_{m})^{t}=I_{m}
\]
}%
לכן \L{$B^{t}$} הופכית ל-\L{$A^{t}$}.

\textcolor{violet}{ההוכחה מראה רק ש-\L{$A^{t}B^{t}=I_{m}$}. הכיוון
השני נובע מהטענה מהרצאה 21 שלפיה אם \L{$BA=I_{m}$} אז \L{$AB=I_{m}$}
)למטריצות ריבועיות(, או ישירות: \L{$B^{t}A^{t}=(AB)^{t}=I_{m}^{t}=I_{m}$}.}

\section{מטריצה סימטרית}

מטריצה \L{$A\in M_{m\times m}(\mathbb{F})$} נקראת סימטרית כאשר \L{$A^{t}=A$}. 

\subsection{דוגמה}

\L{$\begin{pmatrix}1 & 2\\
2 & 3
\end{pmatrix}$} סימטרית.

\section{טענה}

תהי \L{$A\in M_{m\times m}(\mathbb{F})$} הפיכה וסימטרית.

אז ההופכית של \L{$A$} גם סימטרית.

\subsection{הוכחה}

נסמן ב-\L{$B$} את המטריצה ההופכית ל-\L{$A$}. אז מתקיים: \L{%
\[
AB^{t}=A^{t}B^{t}=(BA)^{t}=I_{m}^{t}=I_{m}
\]
}%
)\textcolor{violet}{השוויון הראשון נובע מכך ש-\L{$A$} סימטרית ולכן
\L{$A=A^{t}$}}(.

לכן \L{$B^{t}$} הופכית ל-\L{$A$}. מכאן \L{$B^{t}=B$}, כלומר \L{$B$}
סימטרית.

\newpage{}

\noindent\fbox{\begin{minipage}[t]{1\columnwidth - 2\fboxsep - 2\fboxrule}%
\begin{center}
אלגברה ליניארית )1( - \L{$80134$} \textbar{} האונ' העברית \textbar{}
תשפ\textquotedblright ו א' \textbar{} מרצה: ד\textquotedblright ר
אלכס גורביץ' \textbar{} מתרגל: מר יואב כהן \textbar{} ליקוט ועריכה:
נעם לב
\par\end{center}%
\end{minipage}}

\part{הרצאה {\beginL 24\endL} - {\beginL 06.01.2026\endL}}

\section{תזכורת}

יהיו \L{$V,W$} מ\textquotedblright ו נ\textquotedblright ס, \L{$T:V\rightarrow W$}
ה\textquotedblright ל.

הדרגה של \L{$T$} מוגדרת ע\textquotedblright י: \L{$rk\:T=dim\:\operatorname{Im}T$}.

תהי \L{$A\in M_{m\times n}(\mathbb{F})$}. הדרגה של \L{$A$} מוגדרת
ע\textquotedblright י \L{$rk\:A=rk\:T_{A}$}.

\section{טענה}

יהיו \L{$V,W$} מ\textquotedblright ו נ\textquotedblright ס, \L{$\mathcal{B},\mathcal{C}$}
בסיסים של \L{$V,W$} בהתאמה, \L{$T:V\rightarrow W$} ה\textquotedblright ל.

אז \L{$rk\:T=rk\:[T]_{\mathcal{C}}^{\mathcal{B}}$}.

\section{טענה}

תהיינה \L{$A_{2}\in M_{n\times n}(\mathbb{F})$}, \L{$A\in M_{m\times n}(\mathbb{F})$},
\L{$A_{1}\in M_{m\times m}(\mathbb{F})$}.

אם \L{$A_{1},A_{2}$} הפיכות, אז \L{$rk\:A_{1}AA_{2}=rk\:A$}.

\section{טענה}

תהיינה \L{$A,B\in M_{m\times n}(\mathbb{F})$} שקולות שורה. אז \L{$rk\:B=rk\:A$}.

\section{הגדרה - מטריצה משוחלפת}

תהי \L{$A=\begin{pmatrix}a_{11} & \dots & a_{1n}\\
\vdots &  & \vdots\\
a_{m1} & \dots & a_{mn}
\end{pmatrix}\in\matrmn$}. המטריצה המשוחלפת של \L{$A$} היא \L{$A^{t}=\begin{pmatrix}a_{11} & \dots & a_{m1}\\
\vdots &  & \vdots\\
a_{1n} & \dots & a_{mn}
\end{pmatrix}$}.

\section{טענה}

תהיינה \L{$A\in\matrmn$}, \L{$B\in M_{\ell\times m}(\mathbb{F})$}.
אז \L{$(BA)^{t}=A^{t}B^{t}$}.

\section{טענה}

תהי \L{$A\in\sqrm$} הפיכה. אז גם \L{$A^{t}$} הפיכה.\\

\rule[0.5ex]{1\columnwidth}{1pt}

\section{טענה}

תהיינה \L{$A,B\in\matrmn$} שקולות שורה, אז \L{$rk\:B^{t}=rk\:A^{t}$}. 

\subsection{הוכחה}

קיימת \L{$P\in\sqrm$} הפיכה כך ש- \L{$B=PA$}.

מכאן \L{$B^{t}=(PA)^{t}=A^{t}P^{t}$}. בנוסף, \L{$P^{t}$} הפיכה.

לכן \L{$rk\:(B^{t})=rk\:(I_{n}A^{t}P^{t})=rk\ A^{t}$}.

\section{טענה}

תהי \L{$A\in\matrmn$}. לכל \L{$1\leq j\le n$} נסמן את העמודה ה-\L{$j$}
של \L{$A$} ב- \L{$C_{j}$}.

אז \L{$rk\:A=\dim\spn(C_{1},\dots,C_{n})$}.

\subsection{הוכחה}

נתבונן ב- \L{$T_{A}:\mathbb{F}^{n}\rightarrow\mathbb{F}^{m}$}, \L{$T_{A}(v)=Av$}.

מאחר ש- \L{$(e_{1},\dots,e_{n})$} פורשת את \L{$\mathbb{F}^{n}$},
מתקיים
\begin{flushright}
\L{$\operatorname{Im}T_{A}=\spn(T_{A}(e_{1}),\dots,T_{A}(e_{n}))=\spn(Ae_{1},\dots,Ae_{n})=\spn(C_{1},\dots,C_{n})$}
\par\end{flushright}

מכאן \L{$rk\:A=rk\:T_{A}=dim\:\operatorname{Im}T_{A}=dim\:\spn(C_{1},\dots,C_{n})$}.

\section{טענה}

תהי \L{$R\in\matrmn$} מדורגת מצומצמת. אז \L{$rk\:R^{t}=rk\:R$}.

\subsection{הוכחה}

נסמן ב- \L{$r$} את מספר האיברים המובילים ב- \L{$R$}.

לכל \L{$1\leq i\leq r$} נסמן ב- \L{$h_{i}$} את מספר העמודה בה נמצא
האיבר המוביל של השורה ה- \L{$i$}.

אז לכל \L{$1\leq i\leq r$} העמודה ה- \L{$h_{i}$} של \L{$R$} היא
\L{$e_{i}=\begin{pmatrix}0\\
\vdots\\
1\\
\vdots\\
0
\end{pmatrix}$} כאשר \L{$1$} הוא במקום ה-\L{$i$}.

בנוסף \L{$m-r$} שורות אחרונות של \L{$R$} הן שורות אפסים ולכן \L{$m-r$}
קואורדינטות אחרונות של כל עמודה ב- \L{$R$} הן אפסים.

זאת אומרת שכל עמודה של \L{$R$} שייכת ל- \L{$\spn(e_{1},\dots,e_{r})$}.

לכן פרוש העמודות של \L{$R$} הוא \L{$\spn(e_{1},\dots,e_{r})$} והמימד
שלו הוא \L{$r$}.

לכן לפי הטענה הקודמת \L{$rk\:R=r$}.

מאחר ש- \L{$m-r$} העמודות האחרונות של \L{$R^{t}$} הן עמודות אפסים,

פרוש העמודות של \L{$R^{t}$} שווה לפרוש של \L{$r$} העמודות הראשונות
של \L{$R^{t}$}.

בנוסף, \L{$r$} העמודות הראשונות של \L{$R^{t}$} יוצרות סדרה בת\textquotedblright ל,
כי אם נתון צ\textquotedblright ל

של העמודות האלה עם מקדמים \L{$c_{1},\dots,c_{r}$} וגם \L{$c_{j}\neq0$}, 

אז הקואורדינטה ה- \L{$h_{j}$} של הצירוף הלינארי הזה תהיה \L{$c_{j}\neq0$},

כלומר הצ\textquotedblright ל שונה מהעמודה \L{$\begin{pmatrix}0\\
\vdots\\
0
\end{pmatrix}$}.

מכאן הסדרה של \L{$r$} העמודות הראשונות של \L{$R^{t}$} היא בסיס של
פרוש העמודות של \L{$R^{t}$},

מכאן \L{$rk\:R^{t}=r$}.

\section{משפט הדרגה )או משפט המימדים השלישי(}

תהי \L{$A\in\matrmn$}. אז \L{$rk\:A^{t}=rk\:A$}.

\subsection{הוכחה}

נסמן ב- \L{$R$} את המטריצה המדורגת מצומצמת המתקבלת מ- \L{$A$} ע\textquotedblright י
תהליך הדירוג.

אז \L{$A,R$} שקולות שורה.

מכאן \L{$rk\:A^{t}=rk\:R^{t}=rk\:R=rk\:A$}.

\section{טענה}

תהי \L{$A\in\sqrm$}. אז \L{$A$} הפיכה אם\textquotedblright ם \L{$rk\:A=m$}.

\subsection{הוכחה}

לכל \L{$1\le i\leq m$} נסמן את העמודה ה- \L{$i$} של \L{$A$} ב-
\L{$C_{i}$}. אז \L{$A$} הפיכה

אם\textquotedblright ם \L{$\spn(C_{1},\dots,C_{m})=\mathbb{F}^{m}$}
אם\textquotedblright ם \L{$dim\:\spn(C_{1},\dots,C_{m})=dim\:\mathbb{F}^{m}$}
אם\textquotedblright ם \L{$rk\:A=m$}.

\section{מטריצת מעבר בסיס}

יהי \L{$V$} מ\textquotedblright ו נ\textquotedblright ס, \L{$\mathcal{C}$},
\L{$\mathcal{B}=(v_{1},\dots,v_{m})$} בסיסים של \L{$V$}.

מטריצת מעבר בסיס מ-\L{$\mathcal{B}$} ל-\L{$\mathcal{C}$} היא

\L{%
\[
\underset{\mathcal{B}}{V}\overset{Id_{V}}{\longrightarrow}\underset{\mathcal{C}}{V}
\]
}%

\L{%
\[
[Id_{V}]_{\mathcal{C}}^{\mathcal{B}}=\left([v_{1}]_{\mathcal{C}},\dots,[v_{m}]_{\mathcal{C}}\right)\in\sqrm
\]
}%


\section{תכונה}

יהי \L{$V$} מ\textquotedblright ו נ\textquotedblright ס, \L{$\mathcal{B},\mathcal{C}$}
בסיסים של \L{$V$}, \L{$v\in V$}.

אז \L{$[v]_{\mathcal{C}}=[Id_{V}]_{\mathcal{C}}^{\mathcal{B}}[v]_{\mathcal{B}}$}.

\subsection{הוכחה}

\L{%
\[
[v]_{\mathcal{C}}=[Id_{V}(v)]_{\mathcal{C}}=[Id_{V}]_{\mathcal{C}}^{\mathcal{B}}[v]_{\mathcal{B}}
\]
}%


\section{תכונה}

יהי \L{$\mathcal{B}=(v_{1},\dots,v_{n})$} בסיס כלשהו של \L{$\mathbb{F}^{n}$}.

נסמן ב- \L{$A\in\sqrmn$} את המטריצה שהעמודה ה- \L{$i$} שלה היא \L{$v_{i}$}
לכל \L{$1\leq i\leq n$}. 

אז \L{$[Id_{\mathbb{F}^{n}}]_{\varepsilon_{n}}^{B}=A$} כאשר \L{$\varepsilon_{n}=(e_{1},\dots,e_{n})$}
הבסיס הסטנדרטי של \L{$\mathbb{F}^{n}$}.

\subsection{הוכחה}

העמודה ה- \L{$i$} של \L{$[Id_{\mathbb{F}^{n}}]_{\varepsilon_{n}}^{B}$}
היא \L{$[Id_{V}(v_{i})]_{\varepsilon_{n}}=[v_{i}]_{\varepsilon_{n}}=v_{i}$}.

מכאן הנדרש.

\section{תכונה}

יהי \L{$V$} מ\textquotedblright ו נ\textquotedblright ס, \L{$\mathcal{B},\mathcal{C},\mathcal{D}$}
בסיסים של \L{$V$}.

\L{%
\[
[Id_{V}]_{\mathcal{D}}^{\mathcal{B}}=[Id_{V}]_{\mathcal{D}}^{\mathcal{C}}\cdot[Id_{V}]_{\mathcal{C}}^{\mathcal{B}}
\]
}%


\subsection{הוכחה}

\L{%
\[
[Id_{V}]_{\mathcal{D}}^{\mathcal{C}}\cdot[Id_{V}]_{\mathcal{C}}^{\mathcal{B}}=[Id_{V}\circ Id_{V}]_{\mathcal{D}}^{\mathcal{B}}=[Id_{V}]_{\mathcal{D}}^{\mathcal{B}}
\]
}%


\section{תזכורת}

יהי \L{$V$} מ\textquotedblright ו נ\textquotedblright ס, \L{$\mathcal{B}$}
בסיס של \L{$V$}. אז \L{$[Id_{V}]_{\mathcal{B}}^{\mathcal{B}}=I_{n}$}.

\section{תכונה}

יהי \L{$V$} מ\textquotedblright ו נ\textquotedblright ס, \L{$\mathcal{B},\mathcal{C}$}
בסיסים של \L{$V$}.

אז \L{$[Id]_{\mathcal{C}}^{\mathcal{B}}$} הפיכה וההופכית שלה היא
\L{$[Id_{V}]_{\mathcal{B}}^{\mathcal{C}}$}.

\subsection{הוכחה}

לפי התכונה בתזכורת והתכונה הקודמת \L{$I_{n}=[Id_{V}]_{\mathcal{B}}^{\mathcal{B}}=[Id_{V}]_{\mathcal{B}}^{\mathcal{C}}[Id_{V}]_{\mathcal{C}}^{\mathcal{B}}$}.

\section{טענה )חשובה(}

יהי \L{$V$} מ\textquotedblright ו נ\textquotedblright ס מעל \L{$\mathbb{F}$},
\L{$\mathcal{B}=\left(v_{1},\dots,v_{n}\right)$} בסיס של \L{$V$}.

\L{$P\in\sqrmn$} הפיכה. אז קיים בסיס \L{$\mathcal{C}$} של \L{$V$}
כך ש- \L{$[Id]_{\mathcal{B}}^{\mathcal{C}}=P$}.

\subsection{הוכחה}

לכל \L{$1\leq i\leq n$} נסמן ב- \L{$C_{i}$} את העמודה ה-\L{$i$}
של \L{$P$}.

נזכיר כי \L{$\omega_{\mathcal{B}}:V\rightarrow\mathbb{F}^{n}$} ,
\L{$\omega_{\mathcal{B}}(v)=[v]_{\mathcal{B}}$} היא איזומורפיזם.

\L{$\tau_{\mathcal{B}}:\mathbb{F}^{n}\rightarrow V$} הוא ההופכי של
\L{$\omega_{\mathcal{B}}$}.

לכל \L{$1\leq i\leq n$} נגדיר \L{$u_{i}=\tau_{\mathcal{B}}(C_{i})$}.

נגדיר \L{$C=(u_{1},\dots,u_{n})$}. אז \L{$[u_{i}]_{\mathcal{B}}=\omega_{\mathcal{B}}(u_{i})=C_{i}$}.

ולכן \L{$[Id_{V}]_{\mathcal{B}}^{\mathcal{C}}=([u_{1}]_{\mathcal{B}},\dots,[u_{n}]_{\mathcal{B}})=(C_{1},\dots,C_{n})=P$}.

חשוב לציין שמאחר ש-\L{$P$} הפיכה, הסדרה \L{$(C_{1},\dots,C_{n})$}
פורשת את \L{$\mathbb{F}^{n}$},

ולכן \L{$(u_{1},\dots,u_{n})$} פורשת את \L{$V$}, ומכאן \L{$(u_{1},\dots,u_{n})$}
בסיס של \L{$V$}.

\newpage{}

\noindent\fbox{\begin{minipage}[t]{1\columnwidth - 2\fboxsep - 2\fboxrule}%
\begin{center}
אלגברה ליניארית )1( - \L{$80134$} \textbar{} האונ' העברית \textbar{}
תשפ\textquotedblright ו א' \textbar{} מרצה: ד\textquotedblright ר
אלכס גורביץ' \textbar{} מתרגל: מר יואב כהן \textbar{} ליקוט ועריכה:
נעם לב
\par\end{center}%
\end{minipage}}

\part{הרצאה {\beginL 25\endL} - {\beginL 12.01.2026\endL}}

\textcolor{blue}{שיעור שעבר התחלנו לדבר על מטריצות מעבר בסיס.}

\section{תזכורת}

\section{{\beginL 1\endL}. מטריצת מעבר בסיס}

יהי \L{$V$} מ\textquotedblright ו נ\textquotedblright ס מעל \L{$\mathbb{F}$},
\L{$\mathcal{B},\mathcal{C}$} בסיסים של \L{$V$}.

אז המטריצה \L{$[Id_{V}]_{\mathcal{C}}^{\mathcal{B}}$} נקראת מטריצת
מעבר בסיס.

\section{{\beginL 2\endL}. תכונה}

\L{$[Id]_{\mathcal{C}}^{\mathcal{B}}$} הפיכה. המטריצה ההופכית שלה
היא \L{$[Id]_{\mathcal{B}}^{\mathcal{C}}$}.

\section{{\beginL 3\endL}. טענה}

יהי \L{$V$} מ\textquotedblright ו נ\textquotedblright ס מעל \L{$\mathbb{F}$},
\L{$\dim V=n$},

\L{$\mathcal{B}$} בסיס של \L{$V$}, \L{$P\in\sqrmn$} הפיכה.

אז קיים בסיס \L{$\mathcal{C}$} של \L{$V$} כך ש- \L{$[Id]_{\mathcal{B}}^{\mathcal{C}}=P$}.

\medskip{}

\rule[0.5ex]{1\columnwidth}{1pt}

\medskip{}
\textcolor{blue}{מטריצה מייצגת תלויה בבחירת הבסיסים.}

\textcolor{blue}{נרצה לשאול - איך תשתנה המטריצה של ה\textquotedblright ל
אם נשנה את הבסיס של התחום, או של הטווח, או של שניהם?}

\textcolor{blue}{להלן טענה חשובה על כך:}

\section{{\beginL 4\endL}. טענה}

יהיו \L{$V,W$} מ\textquotedblright ו נ\textquotedblright ס מעל \L{$\mathbb{F}$},
\L{$T:V\rightarrow W$} ה\textquotedblright ל, \L{$\mathcal{B},\mathcal{B}'$}
בסיסים של \L{$V$}, \L{$\mathcal{C},\mathcal{C}'$} בסיסים של \L{$W$}.

אז \L{%
\[
[T]_{\mathcal{C}'}^{\mathcal{B}'}=[Id_{W}]_{\mathcal{C}'}^{\mathcal{C}}[T]_{\mathcal{C}}^{\mathcal{B}}[Id_{V}]_{\mathcal{B}}^{\mathcal{B}'}
\]
}%
 

\subsection{הוכחה}

\L{%
\[
[T]_{\mathcal{C}'}^{\mathcal{B}'}=[Id_{W}\circ T\circ Id_{V}]_{\mathcal{C}'}^{\mathcal{B}'}=[Id_{W}]_{\mathcal{C}'}^{\mathcal{C}}[T]_{\mathcal{C}}^{\mathcal{B}}[Id_{V}]_{\mathcal{B}}^{\mathcal{B}'}
\]
}%


\section*{\textcolor{violet}{הסבר הטענה: שינוי בסיס להעתקה ליניארית}}

\textcolor{violet}{הטענה מתארת את הקשר בין המטריצות המייצגות של אותה
העתקה ליניארית \L{$T$}, אך תחת זוגות בסיסים שונים.}

\subsection*{\textcolor{violet}{האינטואיציה מאחורי הנוסחה}}

\textcolor{violet}{נתונה העתקה \L{$T:V\to W$}. יש לנו שתי דרכים להסתכל
על ההעתקה מבחינת קואורדינטות: }
\begin{enumerate}
\item \textcolor{violet}{דרך הבסיסים הישנים \L{$\mathcal{B}$} )בתחום \L{$V$}(
ו-\L{$\mathcal{C}$} )בטווח \L{$W$}(. }
\item \textcolor{violet}{דרך הבסיסים החדשים \L{$\mathcal{B}'$} )בתחום \L{$V$}(
ו-\L{$\mathcal{C}'$} )בטווח \L{$W$}(. }
\end{enumerate}
\textcolor{violet}{הנוסחה אומרת שכדי לקבל את המטריצה בבסיסים החדשים
\L{$[T]_{\mathcal{C}'}^{\mathcal{B}'}$}, ניתן להשתמש במטריצה הישנה
\L{$[T]_{\mathcal{C}}^{\mathcal{B}}$} ולהכפיל אותה משני הצדדים במטריצות
מעבר )מטריצות שינוי בסיס(.}

\textcolor{violet}{סדר הכפל )מימין לשמאל( מתאר את המסלול שעובר וקטור
קואורדינטות: \L{%
\[
[T]_{\mathcal{C}'}^{\mathcal{B}'}=\underbrace{[Id_{W}]_{\mathcal{C}'}^{\mathcal{C}}}_{\text{3}}\cdot\underbrace{[T]_{\mathcal{C}}^{\mathcal{B}}}_{\text{2}}\cdot\underbrace{[Id_{V}]_{\mathcal{B}}^{\mathcal{B}'}}_{\text{1}}
\]
}%
}
\begin{enumerate}
\item \textcolor{violet}{צעד 1 )המטריצה הימנית(: המטריצה \L{$[Id_{V}]_{\mathcal{B}}^{\mathcal{B}'}$}
לוקחת וקטור המיוצג לפי הבסיס החדש \L{$\mathcal{B}'$}, וממירה אותו
לשפה של הבסיס הישן \L{$\mathcal{B}$}. }
\item \textcolor{violet}{צעד 2 )המטריצה האמצעית(: כעת כשהקלט הוא בבסיס הישן,
המטריצה \L{$[T]_{\mathcal{C}}^{\mathcal{B}}$} מבצעת את ההעתקה ומחזירה
תוצאה בבסיס הישן של הטווח \L{$\mathcal{C}$}. }
\item \textcolor{violet}{צעד 3 )המטריצה השמאלית(: המטריצה \L{$[Id_{W}]_{\mathcal{C}'}^{\mathcal{C}}$}
לוקחת את הפלט )שהוא בבסיס \L{$\mathcal{C}$}( וממירה אותו לבסיס החדש
של הטווח \L{$\mathcal{C}'$}. }
\end{enumerate}

\section*{\textcolor{violet}{דוגמה מספרית}}

\textcolor{violet}{נניח כי \L{$V=W=\mathbb{R}^{2}$} ונגדיר העתקה
ליניארית \L{$T:\mathbb{R}^{2}\to\mathbb{R}^{2}$} המוגדרת על ידי:
\L{%
\[
T(x,y)=(x,2y)
\]
}%
}

\subsection*{\textcolor{violet}{1. בחירת בסיסים}}

\textcolor{violet}{נבחר את הבסיס הסטנדרטי כבסיס הישן )\L{$\mathcal{B}=\mathcal{C}$}(:
\L{%
\[
\mathcal{B}=((1,0),(0,1))
\]
}%
נבחר בסיס חדש )\L{$\mathcal{B}'=\mathcal{C}'$}(: \L{%
\[
\mathcal{B}'=((1,1),(0,1))
\]
}%
)לצורך הפשטות, בחרנו \L{$V=W$} וגם \L{$\mathcal{B}=\mathcal{C}$}
ו-\L{$\mathcal{B}'=\mathcal{C}'$}(.}

\subsection*{\textcolor{violet}{2. חישוב המטריצה בבסיס הישן )הסטנדרטי(}}

\textcolor{violet}{נחשב את \L{$[T]_{\mathcal{B}}^{\mathcal{B}}$}:
\L{%
\[
T(1,0)=(1,0),\quad T(0,1)=(0,2)
\]
}%
לכן\L{%
\[
A=[T]_{\mathcal{B}}^{\mathcal{B}}=\begin{pmatrix}1 & 0\\
0 & 2
\end{pmatrix}
\]
}%
}

\subsection*{\textcolor{violet}{3. חישוב מטריצות המעבר}}

\textcolor{violet}{אנו צריכים את המטריצה \L{$[Id]_{\mathcal{B}}^{\mathcal{B}'}$}
)מטריצת המעבר מהחדש לישן(. עמודות המטריצה הן פשוט וקטורי הבסיס \L{$\mathcal{B}'$}
כתובים לפי הבסיס הסטנדרטי: \L{%
\[
P=[Id]_{\mathcal{B}}^{\mathcal{B}'}=\begin{pmatrix}1 & 0\\
1 & 1
\end{pmatrix}
\]
}%
כמו כן אנו צריכים את המטריצה ההפוכה, \L{$[Id]_{\mathcal{B}'}^{\mathcal{B}}$}
)מהישן לחדש(, שהיא \L{$P^{-1}$}: \L{%
\[
P^{-1}=[Id]_{\mathcal{B}'}^{\mathcal{B}}=\begin{pmatrix}1 & 0\\
-1 & 1
\end{pmatrix}
\]
}%
}

\subsection*{\textcolor{violet}{4. שימוש בנוסחה}}

\textcolor{violet}{לפי טענה 4: \L{%
\[
[T]_{\mathcal{B}'}^{\mathcal{B}'}=[Id]_{\mathcal{B}'}^{\mathcal{B}}\cdot[T]_{\mathcal{B}}^{\mathcal{B}}\cdot[Id]_{\mathcal{B}}^{\mathcal{B}'}=P^{-1}AP
\]
}%
נציב ונכפול: \L{%
\[
\begin{pmatrix}1 & 0\\
-1 & 1
\end{pmatrix}\begin{pmatrix}1 & 0\\
0 & 2
\end{pmatrix}\begin{pmatrix}1 & 0\\
1 & 1
\end{pmatrix}
\]
}%
תחילה נכפול את שתי המטריצות הימניות: \L{%
\[
\begin{pmatrix}1 & 0\\
0 & 2
\end{pmatrix}\begin{pmatrix}1 & 0\\
1 & 1
\end{pmatrix}=\begin{pmatrix}1 & 0\\
2 & 2
\end{pmatrix}
\]
}%
כעת נכפול בשמאלית: \L{%
\[
\begin{pmatrix}1 & 0\\
-1 & 1
\end{pmatrix}\begin{pmatrix}1 & 0\\
2 & 2
\end{pmatrix}=\begin{pmatrix}1 & 0\\
1 & 2
\end{pmatrix}
\]
}%
קיבלנו את המטריצה המייצגת בבסיס החדש: \L{%
\[
[T]_{\mathcal{B}'}^{\mathcal{B}'}=\begin{pmatrix}1 & 0\\
1 & 2
\end{pmatrix}
\]
}%
}

\subsection*{\textcolor{violet}{5. בדיקה ישירה )לאימות(}}

\textcolor{violet}{נחשב את \L{$[T]_{\mathcal{B}'}^{\mathcal{B}'}$}
ישירות לפי ההגדרה, ללא שימוש בנוסחת המעבר. נפעיל את \L{$T$} על איברי
הבסיס \L{$\mathcal{B}'=((1,1),(0,1))$} ונציג את התוצאה כצירוף ליניארי
של איברי \L{$\mathcal{B}'$}: }
\begin{itemize}
\item \textcolor{violet}{עבור הוקטור הראשון \L{$(1,1)$}: \L{%
\[
T(1,1)=(1,2)
\]
}%
נציג אותו לפי הבסיס החדש: \L{$(1,2)=1\cdot(1,1)+1\cdot(0,1)$}. לכן
העמודה הראשונה היא \L{$\begin{pmatrix}1\\
1
\end{pmatrix}$}. }
\item \textcolor{violet}{עבור הוקטור השני \L{$(0,1)$}: \L{%
\[
T(0,1)=(0,2)
\]
}%
נציג אותו לפי הבסיס החדש: \L{$(0,2)=0\cdot(1,1)+2\cdot(0,1)$}. לכן
העמודה השנייה היא \L{$\begin{pmatrix}0\\
2
\end{pmatrix}$}. }
\end{itemize}
\textcolor{violet}{סה"כ קיבלנו: \L{%
\[
\begin{pmatrix}1 & 0\\
1 & 2
\end{pmatrix}
\]
}%
זוהי בדיוק אותה מטריצה שקיבלנו בנוסחה, מה שמאשש את נכונות הטענה.}

\textcolor{blue}{לטענה }\textcolor{blue}{{\beginL 4}\endL}\textcolor{blue}{{}
יש מקרה פרטי אחד חשוב: בוחרים \L{$\mathcal{B}=\mathcal{C}$} ו- \L{$\mathcal{B}'=\mathcal{C}'$}
אבל עדיין \L{$\mathcal{B}\neq\mathcal{B}'$}. מכך שהן הופכיות זו לזו
)תכונה }\textcolor{blue}{{\beginL 2}\endL}\textcolor{blue}{( נקבל
משהו מעניין. )בחזקה מינוס אחד על מטריצה הכוונה היא למטריצה ההופכית(.
אפשר לראות איך המסקנה הבאה, וגם הגדרה }\textcolor{blue}{{\beginL 6}\endL}\textcolor{blue}{,
באות לידי ביטוי בדוגמה המספרית שבאפור.}

\section{{\beginL 5\endL}. מסקנה}

יהיו \L{$V$} מ\textquotedblright ו נ\textquotedblright ס מעל \L{$\mathbb{F}$},
\L{$T:V\rightarrow V$} ה\textquotedblright ל, \L{$\mathcal{B},\mathcal{B}'$}
בסיסים של \L{$V$}.

אז מהטענה ){\beginL 4\endL}(

\L{%
\[
[T]_{\mathcal{B}'}^{\mathcal{B}'}=[Id_{V}]_{\mathcal{B}'}^{\mathcal{B}}[T]_{\mathcal{B}}^{\mathcal{B}}[Id_{V}]_{\mathcal{B}}^{\mathcal{B}'}
\]
}%

אבל מכך שהן הופכיות זו לזו ){\beginL 2\endL}( נקבל \textbf{)וזו המסקנה
החשובה:(}

\L{%
\[
[T]_{\mathcal{B}'}^{\mathcal{B}'}=\left([Id_{V}]_{\mathcal{B}}^{\mathcal{B}'}\right)^{-1}[T]_{\mathcal{B}}^{\mathcal{B}}[Id_{V}]_{\mathcal{B}}^{\mathcal{B}'}
\]
}%

\medskip{}
\textcolor{blue}{מה שהמסקנה הזו אומרת למעשה, היא שברגע שיש ה\textquotedblright ל
מ-מ\textquotedblright ו לעצמו עם שני בסיסים שונים, הקשר בין שתי מטריצות
מייצגות במרחב הוא קשר שבו לוקחים את המטריצה הראשונה, כופלים אותה בצד
ימין במטריצה הפיכה, ובצד שמאל במטריצה ההופכית של אותה מטריצה הפיכה.}\\
\textcolor{blue}{}\\
\textcolor{blue}{כאשר קיים קשר כזה בין שתי מטריצות, נקרא להן }\textbf{\textcolor{blue}{מטריצות
דומות}}\textcolor{blue}{.}\\
\textcolor{blue}{}\\
\textcolor{blue}{זו מסקנה חשובה כי בלינארית }\textcolor{blue}{{\beginL 2}\endL}\textcolor{blue}{{}
באה לידי ביטוי באופן מרכזי ה\textquotedblright ל מ-מ\textquotedblright ו
לעצמו, לכן המסקנה הזו קריטית לכל מה שיקרה בקורס ההמשך.}\\


\section{דמיון מטריצות}

\section{{\beginL 6\endL}. הגדרה - דמיון מטריצות}

תהיינה \L{$A,B\in\sqrmn$}. אומרים כי \L{$B$} \textbf{דומה} ל-\L{$A$}
כאשר קיימת \L{$P\in\sqrmn$} הפיכה כך ש- \L{$B=P^{-1}AP$}.

במקרה זה נסמן \L{$B\sim A$}.
\begin{itemize}
\item \textcolor{blue}{נזכיר שכפל מטריצות אינו חילופי, לכן לא בהכרח \L{$A=B$},
אחרת זה לא היה מעניין.}
\item \textcolor{blue}{חשוב לציין, ותכף נראה את זה בתור תכונות, שדמיון מטריצות
מגדיר יחס שקילות )יחס שהוא רפלקסיבי, סימטרי וטרנזיטיבי(.}
\item \textcolor{blue}{עכשיו נכתוב מחדש את המסקנה )5( בעזרת זה.}
\end{itemize}

\section{{\beginL 7\endL}. תכונה}

יהיו \L{$V$} מ\textquotedblright ו נ\textquotedblright ס מעל \L{$\mathbb{F}$},
\L{$T:V\rightarrow V$} ה\textquotedblright ל, \L{$\mathcal{B},\mathcal{B}'$}
בסיסים של \L{$V$}.

אז \L{$[T]_{\mathcal{B}'}^{\mathcal{B}'}$} דומה ל- \L{$[T]_{\mathcal{B}}^{\mathcal{B}}$}.

\section{{\beginL 8\endL}. תכונות של דמיון מטריצות}

\subsection{{\beginL 1\endL}( \L{$A\sim A$} )רפלקסיביות(}

\subsubsection{הוכחה:}

נבחר \L{$P=I_{n}$}. אז \L{$A=\left(I_{n}\right)^{-1}AI_{n}$}.

\subsection{{\beginL 2\endL}( אם \L{$B\sim A$}, אז \L{$A\sim B$} )סימטריה(}

\subsubsection{הוכחה:}

אם \L{$B\sim A$} אז מהגדרת דמיון מטריצות קיימת \L{$P\in\sqrmn$}
הפיכה כך ש- \L{$B=P^{-1}AP$}. אז

\L{%
\[
(P^{-1})^{-1}BP^{-1}\underset{\text{def of \textasciicircum-1}}{=}PBP^{-1}\underset{\text{substitute B}}{=}PP^{-1}APP^{-1}=I_{n}AI_{n}=A
\]
}%
לכן \L{$A\sim B$}.

\subsection{{\beginL 3\endL}( אם \L{$B\sim A$} וגם \L{$C\sim B$} אז \L{$C\sim A$}
)טרנזיטיביות(}

\subsubsection{הוכחה:}

מההנחות קיימות \L{$Q,P$} הפיכות כך ש- \L{$B=P^{-1}AP$} וגם \L{$C=Q^{-1}BQ$}.
אז 

\L{%
\[
C\underset{\text{substitute C}}{=}Q^{-1}BQ\underset{\text{substitute B}}{=}Q^{-1}P^{-1}APQ\underset{*}{=}(PQ)^{-1}A(PQ)
\]
}%
לכן \L{$C\sim A$}.\\

\L{$(*)$} זה נובע מטענה שלמדנו על ההופכי של המכפלה לפיה \L{$(AB)^{-1}=B^{-1}A^{-1}$}.
נשים לב שהסדר מתהפך.

\medskip{}

\textcolor{blue}{הוכחנו פעם שאם כופלים מטריצה הפיכה מצד ימין ומצד
שמאל הדרגה שלה לא משתנה. זה בא לידי ביטוי גם כאן באופן מאוד שימושי
ונציג את זה עכשיו כטענה.}

\section{{\beginL 9\endL}. טענה}

תהיינה \L{$A,B\in\sqrmn$}. אם \L{$B\sim A$}, אז \L{$rk\:B=rk\:A$}.

\subsection{הוכחה}

מהגדרת דמיון מטריצות קיימת \L{$P\in\sqrmn$} הפיכה כך ש- \L{$B=P^{-1}AP$}.
מאחר ש- \L{$P^{-1},P$} הפיכות מתקיים \L{$rk\:B=rk\:A$}.

\textcolor{blue}{נוכיח כי בהינתן בסיס וה\textquotedblright ל מ-מ\textquotedblright ו
לעצמו, כל מטריצה אחרת הדומה למטריצה המייצגת של ההעתקה הזאת, גם היא
מטריצה מייצגת של אותה ההעתקה, פשוט ביחס לבסיס אחר.}\\

\textcolor{blue}{בשביל להוכיח את זה נשתמש בטענה משיעור שעבר )3(.}\\
\textcolor{blue}{}\\
\textcolor{blue}{ראשית נכתוב על סמך ההגדרה של דמיון מטריצות, ואז נרצה
ש- \L{$P$} תהיה מטריצת מעבר. את זה בדיוק אפשר להשיג מהטענה )}\textcolor{blue}{{\beginL 3}\endL}\textcolor{blue}{(
מהשיעור הקודם. }

\section{{\beginL 10\endL}. טענה}

יהי \L{$V$} מ\textquotedblright ו נ\textquotedblright ס מעל \L{$\mathbb{F}$}
כך ש- \L{$\dim V=n$} ויהיו \L{$T:V\rightarrow V$} ה\textquotedblright ל,
\L{$\mathcal{B}$} בסיס של \L{$V$}, \L{$A\in\sqrmn$} מטריצה כך ש-
\L{$A\sim[T]_{\mathcal{B}}^{\mathcal{B}}$}. אז קיים בסיס \L{$\mathcal{C}$}
של \L{$V$} כך ש- \L{$A=[T]_{\mathcal{C}}^{\mathcal{C}}$}.

\subsection{הוכחה}

מההגדרה של דמיון מטריצות קיימת \L{$P\in\sqrmn$} הפיכה כך ש- \L{$A=P^{-1}[T]_{\mathcal{B}}^{\mathcal{B}}P$}.
לפי הטענה מהשיעור הקודם קיים בסיס \L{$\mathcal{C}$} של \L{$V$} כך
ש- \L{$P=[Id_{V}]_{\mathcal{B}}^{\mathcal{C}}$}. מכאן

\L{%
\[
A=P^{-1}[T]_{\mathcal{B}}^{\mathcal{B}}P=\left([Id_{V}]_{\mathcal{B}}^{\mathcal{C}}\right)^{-1}[T]_{\mathcal{B}}^{\mathcal{B}}[Id_{V}]_{\mathcal{B}}^{\mathcal{C}}=[Id_{V}]_{\mathcal{C}}^{\mathcal{B}}[T]_{\mathcal{B}}^{\mathcal{B}}[Id_{V}]_{\mathcal{B}}^{\mathcal{C}}=[T]_{\mathcal{C}}^{\mathcal{C}}
\]
}%


\section{דטרמיננטות}

\textcolor{blue}{ועכשיו חברים מתחילים את הפרק האחרון בקורס, הלוא הוא
דטרמיננטות.}

\section{רעיון כללי ומוטיבציה}

\textcolor{blue}{דטרמיננטה היא פונקציה מ- קבוצה של מטריצות ריבועיות
עם מקדמים מהשדה, ל- שדה. כלומר, לכל מטריצה ריבועית של מקדמים מהשדה
היא מתאימה סקלר מהשדה.}

\textcolor{blue}{יש רק אחת כזאת, מדובר על פונקציה ספציפית, לא על אפיון.}\\

\textcolor{blue}{המשמעות הגיאומטרית שלה:}

\textcolor{blue}{אפשר להסתכל על מטריצה ריבועית כאוסף של n שורות )במקום
עמודות כמו עד עכשיו(. כל שורה במטריצה היא n-יה שוכבת כזאת. נקודה במרחב
\L{$F^{n}$}. אז למשל אם \L{$n=2$} המשמעות הגיאומטרית של המטריצה
היא משהו כמו:}\\

\par\noindent\centerline{%
  \L{\includegraphics[width=0.5\textheight]{\detokenize{./gen_sum_vol_2_figs/Figures_Lec_25/fig_1.jpeg}}}%
}\par

\medskip{}
\textcolor{blue}{\newpage}

\textcolor{blue}{ואם \L{$n=3$} אז למשל:}

\medskip{}

\medskip{}

\par\noindent\centerline{%
  \L{\includegraphics[width=0.5\textheight]{\detokenize{./gen_sum_vol_2_figs/Figures_Lec_25/fig_2.jpeg}}}%
}\par

\medskip{}

\textcolor{blue}{ודטרמיננטה מודדת את ה\textquotedbl גודל\textquotedbl{}
של הגופים האלה.}

\textcolor{blue}{החלק הכי קשה יהיה להוכיח שבכלל קיימת פונקציה שמקיימת
את כל הדרישות שאנחנו נסיק.}

\textcolor{blue}{התכונה העיקרית של הפונקציה הזאת היא שאם נקבע את כל
השורות חוץ מאחת שתהיה המשתנה נקבל פונקציה לינארית.}\\

\par\noindent\centerline{%
  \L{\includegraphics[width=0.5\textheight]{\detokenize{./gen_sum_vol_2_figs/Figures_Lec_25/fig_3.jpeg}}}%
}\par

\medskip{}

\newpage{}

\textbf{הטענה: השטח של צהוב פלוס השטח של אדום שווה לשטח של כחול.}

\medskip{}

\par\noindent\centerline{%
  \L{\includegraphics[width=0.5\textheight]{\detokenize{./gen_sum_vol_2_figs/Figures_Lec_25/fig_4.jpeg}}}%
}\par

\medskip{}
\par\noindent\centerline{%
  \L{\includegraphics[width=0.5\textheight]{\detokenize{./gen_sum_vol_2_figs/Figures_Lec_25/fig_5.jpeg}}}%
}\par

\newpage{}

\subsection{\textcolor{violet}{המחשה ומוטיבציה מהסטודנט האגדי אביב זילברמן פרץ}}

\textcolor{violet}{אל תנסה להבין אינטואיציה גיאומטרית באלגברה )בכלל(.}

\textcolor{violet}{כמה שיותר אינטואיציה גיאומטרית בסיסית ותו לא, לך
ולי אין תפיסה גיאומטרית ב-\L{$n$} מימדים, אנחנו עדיין לא רימן וגאוס.}

\textcolor{violet}{במובן האבסטרקטי, \L{$det$} הוא נפח, ומה שמשתמע
מכך.}

\textcolor{violet}{בסוף, התיאוריה הגיאומטרית שאלכס מפתח לך היא לא
חובה להבנה.}

\textcolor{violet}{מה שחשוב לדעת זה ש- \L{$det$} זו פונקציה רקורסיבית
ושתדע לפתח אותה בעצמך בצורה אלגברית )תלמד בהרצאות הבאות(.}

\textcolor{violet}{כשתגיע לזה כדאי לשים לב שאתה יודע איך לחשב \L{$det$}
ע\textquotedblright י חישוב מטריצה קנונית.}\textcolor{gray}{{} }

\section{לינאריות לפי שורה}

\subsection{סימונים}

בהינתן \L{$A\in\sqrmn$}, \L{$1\leq i\leq n$} נסמן את השורה ה-\L{$i$}
של \L{$A$} ב- \L{$\Lambda_{i}^{A}$}.

\subsection{הגדרה}

תהי \L{$D:M_{n\times n}(\mathbb{F})\rightarrow\mathbb{F}$}, \L{$1\leq i\leq n$}.

אומרים כי \L{$D$} \textbf{לינארית לפי השורה ה- \L{$i$}} כאשר:\\
\\
){\beginL 1\endL}( לכל \L{$A,B,C\in M_{n\times n}(\mathbb{F})$} כך
ש:

\qquad{}לכל \L{$1\leq j\leq n$} המקיים \L{$j\neq i$} מתקיים \L{$\Lambda_{j}^{A}=\Lambda_{j}^{B}=\Lambda_{j}^{C}$}
\qquad{}וגם \L{$\Lambda_{i}^{A}+\Lambda_{i}^{B}=\Lambda_{i}^{C}$}

מתקיים \L{$D(A)+D(B)=D(C)$}.\\
\\
){\beginL 2\endL}( לכל \L{$A,B\in\sqrmn$}, \L{$c\in\mathbb{F}$}
כך ש-

\qquad{}לכל \L{$1\leq j\leq n$} המקיים \L{$j\neq i$} מתקיים \L{$\Lambda_{j}^{A}=\Lambda_{j}^{B}$}

\qquad{}וגם \L{$c\Lambda_{i}^{A}=\Lambda_{i}^{B}$}

מתקיים \L{$c\cdot D(A)=D(B)$}.

\subsection{דוגמה}

\L{$n=2,i=1$}, \L{$D:M_{2\times2}(\mathbb{F})\rightarrow\mathbb{F}$}.

\L{$D\begin{pmatrix}a_{11} & a_{12}\\
a_{21} & a_{22}
\end{pmatrix}=a_{11}$}.\\
\\
\textcolor{violet}{נראה כי \L{$D$} אכן לינארית לפי השורה הראשונה
)\L{$i=1$}(.}

\subsection*{\textcolor{violet}{1. חיבוריות}}

\textcolor{violet}{יהיו \L{$A,B,C\in M_{2\times2}(\mathbb{F})$} מטריצות
בעלות שורה שנייה זהה )\L{$\Lambda_{2}^{A}=\Lambda_{2}^{B}=\Lambda_{2}^{C}$}(,
כך שהשורה הראשונה של \L{$C$} היא סכום השורות הראשונות של \L{$A$}
ו-\L{$B$}. נסמן את השורות הראשונות: \L{$\Lambda_{1}^{A}=(a_{11},a_{12})$}
ו-\L{$\Lambda_{1}^{B}=(b_{11},b_{12})$}. לכן השורה הראשונה של \L{$C$}
היא: \L{%
\[
\Lambda_{1}^{C}=(a_{11}+b_{11},\,a_{12}+b_{12})
\]
}%
לפי הגדרת הפונקציה \L{$D$}, הערך תלוי רק באיבר הראשון בשורה הראשונה:
\L{%
\[
D(C)=(a_{11}+b_{11})=a_{11}+b_{11}=D(A)+D(B)
\]
}%
תנאי החיבוריות מתקיים.}

\subsection*{\textcolor{violet}{2. כפל בסקלר )הומוגניות(}}

\textcolor{violet}{יהי \L{$c\in\mathbb{F}$} סקלר, ותהי \L{$B$} מטריצה
המתקבלת מ-\L{$A$} על ידי הכפלת השורה הראשונה ב-\L{$c$} )כאשר השורה
השנייה נותרת ללא שינוי(. אם השורה הראשונה של \L{$A$} היא \L{$(a_{11},a_{12})$},
אזי השורה הראשונה של \L{$B$} היא \L{$(c\cdot a_{11},c\cdot a_{12})$}.
לפי הגדרת הפונקציה: \L{%
\[
D(B)=c\cdot a_{11}=c\cdot(a_{11})=c\cdot D(A)
\]
}%
תנאי הכפל בסקלר מתקיים.}

\paragraph*{\textcolor{violet}{מסקנה:}}

\textcolor{violet}{הפונקציה \L{$D$} מקיימת את שני התנאים ולכן היא
לינארית לפי השורה הראשונה.}\textcolor{gray}{}\\
\textcolor{gray}{}\\
\textcolor{blue}{באופן דומה יכלנו לכתוב \L{$D\begin{pmatrix}a_{11} & a_{12}\\
a_{21} & a_{22}
\end{pmatrix}=a_{12}$} ולקבל פונקציה לינארית לפי השורה הראשונה אחרת. אבל אם היינו כותבים
\L{$D\begin{pmatrix}a_{11} & a_{12}\\
a_{21} & a_{22}
\end{pmatrix}=a_{21}$} לא היינו מקבלים פונקציה לינארית לפי השורה הראשונה, אלא פונקציה לינארית
לפי השורה השניה.}\\

\textcolor{blue}{אבל מה אם נרצה פונקציה שלינארית גם לפי השורה הראשונה
וגם לפי השורה השניה?}

\section{מולטי-לינאריות}

\subsection{הגדרה}

תהי \L{$D:\sqrmn\rightarrow\mathbb{F}$}. אומרים כי \L{$D$} \textbf{מולטי-לינארית
}כאשר \L{$D$} לינארית לפי השורה ה-\L{$i$} לכל \L{$1\leq i\leq n$}. 

\subsection{דוגמה}

\L{$n=2$}, \L{$D:M_{2\times2}(\mathbb{F})\rightarrow\mathbb{F}$}.

\L{$D\begin{pmatrix}a_{11} & a_{12}\\
a_{21} & a_{22}
\end{pmatrix}=a_{11}a_{21}$}.\\
\medskip{}

\textcolor{violet}{כדי להוכיח שהפונקציה מולטי-לינארית, נראה כי היא
לינארית ביחס לכל שורה בנפרד.}

\subsection*{\textcolor{violet}{1. לינאריות לפי השורה הראשונה )\L{$i=1$}(}}

\textcolor{violet}{נניח כי השורה השנייה קבועה )ולכן \L{$a_{21}$}
הוא סקלר קבוע(. נבדוק שינויים בשורה הראשונה. }
\begin{itemize}
\item \textcolor{violet}{חיבוריות: אם נחבר שתי שורות ראשונות )כאשר האיברים
הראשונים הם \L{$a_{11}$} ו-\L{$a'_{11}$}(, נקבל: \L{%
\[
(a_{11}+a'_{11})a_{21}=a_{11}a_{21}+a'_{11}a_{21}=D(A)+D(A')
\]
}%
}
\item \textcolor{violet}{הומוגניות: אם נכפיל את השורה הראשונה בסקלר \L{$c$}:
\L{%
\[
(c\cdot a_{11})a_{21}=c\cdot(a_{11}a_{21})=c\cdot D(A)
\]
}%
}
\end{itemize}
\textcolor{violet}{הפונקציה לינארית לפי השורה הראשונה.}

\subsection*{\textcolor{violet}{2. לינאריות לפי השורה השנייה )\L{$i=2$}(}}

\textcolor{violet}{נניח כי השורה הראשונה קבועה )ולכן \L{$a_{11}$}
הוא סקלר קבוע(. נבדוק שינויים בשורה השנייה. }
\begin{itemize}
\item \textcolor{violet}{חיבוריות: אם נחבר שתי שורות שניות )כאשר האיברים
הם \L{$a_{21}$} ו-\L{$a'_{21}$}(, נקבל: \L{%
\[
a_{11}(a_{21}+a'_{21})=a_{11}a_{21}+a_{11}a'_{21}=D(A)+D(A')
\]
}%
}
\item \textcolor{violet}{הומוגניות: אם נכפיל את השורה השנייה בסקלר \L{$c$}:
\L{%
\[
a_{11}(c\cdot a_{21})=c\cdot(a_{11}a_{21})=c\cdot D(A)
\]
}%
}
\end{itemize}
\textcolor{violet}{הפונקציה לינארית לפי השורה השנייה.}

\paragraph*{\textcolor{violet}{מסקנה:}}

\textcolor{violet}{הפונקציה \L{$D$} לינארית ביחס לכל אחת משורות המטריצה,
ולכן היא פונקציה מולטי-לינארית.}\textcolor{gray}{}\\
\textcolor{gray}{}\\
\textbf{\textcolor{green}{מולטי-לינאריות זו הגדרה חשובה להמשך הנושא
של דטרמיננטות.}}

\newpage{}

\noindent\fbox{\begin{minipage}[t]{1\columnwidth - 2\fboxsep - 2\fboxrule}%
\begin{center}
אלגברה ליניארית )1( - \L{$80134$} \textbar{} האונ' העברית \textbar{}
תשפ\textquotedblright ו א' \textbar{} מרצה: ד\textquotedblright ר
אלכס גורביץ' \textbar{} מתרגל: מר יואב כהן \textbar{} ליקוט ועריכה:
נעם לב
\par\end{center}%
\end{minipage}}

\part{הרצאה {\beginL 26\endL} - {\beginL 13.01.2026\endL}}

\section{תזכורת}

\section{סימון}

תהי \L{$A\in M_{n\times n}(\mathbb{F})$}. לכל \L{$1\le i\le n$},
נסמן את השורה ה-\L{$i$} של \L{$A$} ב-\L{$\Lambda_{i}^{A}$}. אזי,
ניתן להציג את המטריצה \L{$A$} באופן הבא: \L{%
\[
A=\begin{pmatrix}\Lambda_{1}^{A}\\
\vdots\\
\Lambda_{n}^{A}
\end{pmatrix}
\]
}%


\section{{\beginL 1\endL}. הגדרה: פונקציה ליניארית לפי שורה}

יהי \L{$1\le i\le n$}. תהי \L{$D\colon M_{n\times n}(\mathbb{F})\to\mathbb{F}$}
פונקציה. הפונקציה \L{$D$} נקראת \textbf{ליניארית לפי השורה ה-\L{$i$}}
אם מתקיימים התנאים הבאים:
\begin{enumerate}
\item לכל \L{$A,B,C\in M_{n\times n}(\mathbb{F})$} המקיימות: \L{%
\[
1)\quad\forall j\neq i,1\le j\le n:\quad\Lambda_{j}^{A}=\Lambda_{j}^{B}=\Lambda_{j}^{C}
\]
}%
\L{%
\[
2)\quad\Lambda_{i}^{C}=\Lambda_{i}^{A}+\Lambda_{i}^{B}
\]
}%
מתקיים: \L{%
\[
D(C)=D(A)+D(B)
\]
}%
\item לכל \L{$A,B\in M_{n\times n}(\mathbb{F})$} וסקלר \L{$c\in\mathbb{F}$}
המקיימים: \L{%
\[
1)\quad\forall j\neq i,1\le j\le n:\quad\Lambda_{j}^{A}=\Lambda_{j}^{B}
\]
}%
\L{%
\[
2)\quad\Lambda_{i}^{B}=c\cdot\Lambda_{i}^{A}
\]
}%
מתקיים: \L{%
\[
D(B)=c\cdot D(A)
\]
}%
\end{enumerate}

\section{{\beginL 2\endL}. הגדרה: פונקציה מולטי-ליניארית}

פונקציה \L{$D\colon M_{n\times n}(\mathbb{F})\to\mathbb{F}$} נקראת
\textbf{מולטי-ליניארית} כאשר לכל \L{$1\le i\le n$}, \L{$D$} ליניארית
לפי השורה ה-\L{$i$}.

\subsection{דוגמה}

עבור \L{$n=2$}: \L{%
\[
D\colon M_{2\times2}(\mathbb{F})\to\mathbb{F}
\]
}%
המוגדרת באופן הבא: \L{%
\[
D\begin{pmatrix}a_{11} & a_{12}\\
a_{21} & a_{22}
\end{pmatrix}=a_{11}a_{21}
\]
}%
זוהי פונקציה מולטי-ליניארית )ניתן לוודא שהיא ליניארית לפי השורה הראשונה
וגם לפי השורה השנייה(.

\medskip{}

\rule[0.5ex]{1\columnwidth}{1pt}

\medskip{}

\textcolor{blue}{עכשיו נלמד תכונות של פונקציות מולטי-לינאריות.}

\section{{\beginL 3\endL}. תכונה: התאפסות פונקציה לינארית לפי שורה שהיא שורת
אפסים}

תהי \L{$D\colon M_{n\times n}(\mathbb{F})\to\mathbb{F}$} פונקציה
ליניארית לפי השורה ה-\L{$i$}. תהי \L{$A\in M_{n\times n}(\mathbb{F})$}
מטריצה כך שהשורה ה-\L{$i$} שלה היא שורת אפסים, כלומר: \L{%
\[
\Lambda_{i}^{A}=(0,\dots,0)
\]
}%
אזי מתקיים: \L{%
\[
D(A)=0
\]
}%


\section{{\beginL 4\endL}. תכונה: התאפסות פונקציה מולטי-לינארית לפי שורת אפסים}

תהי \L{$D\colon M_{n\times n}(\mathbb{F})\to\mathbb{F}$} פונקציה
מולטי-ליניארית.

תהי \L{$A\in M_{n\times n}(\mathbb{F})$} מטריצה שיש בה שורת אפסים
)כלומר, קיים אינדקס \L{$k$} שעבורו \L{$\Lambda_{k}^{A}=(0,\dots,0)$}(.
אזי מתקיים: \L{%
\[
D(A)=0
\]
}%

\medskip{}

\textcolor{violet}{המטריצה \L{$A$} היא מהצורה: \L{%
\[
A=\begin{pmatrix}\Lambda_{1}^{A}\\
\vdots\\
\Lambda_{k}^{A}\\
\vdots\\
\Lambda_{n}^{A}
\end{pmatrix}
\]
}%
כאשר \L{$\Lambda_{k}^{A}=(0,\dots,0)$}. מאחר ומתקיים \L{$(0,\dots,0)=0\cdot(0,\dots,0)$},
ומאחר והפונקציה \L{$D$} היא לינארית ביחס לשורה ה-\L{$k$}, נוכל להוציא
את הסקלר \L{$0$} מחוץ לפונקציה: \L{%
\[
D(A)=D\begin{pmatrix}\Lambda_{1}^{A}\\
\vdots\\
(0,\dots,0)\\
\vdots\\
\Lambda_{n}^{A}
\end{pmatrix}=D\begin{pmatrix}\Lambda_{1}^{A}\\
\vdots\\
0\cdot(0,\dots,0)\\
\vdots\\
\Lambda_{n}^{A}
\end{pmatrix}=0\cdot D\begin{pmatrix}\Lambda_{1}^{A}\\
\vdots\\
(0,\dots,0)\\
\vdots\\
\Lambda_{n}^{A}
\end{pmatrix}=0
\]
}%
}

\medskip{}


\section{פונקציה מתחלפת}

\section{אינטואיציה}

\textcolor{blue}{בואו נחשוב מה קורה במטריצה שיש לה שתי שורות זהות.}

\textcolor{blue}{למשל במקרה הדו-מימדי, במקום מקבילית מקבלים מקבילית
מנוונת, פשוט קטע.}

\textcolor{blue}{במקום u+v, מקבלים u+u.}

\medskip{}

\par\noindent\centerline{%
  \L{\includegraphics[width=0.5\textheight]{\detokenize{./gen_sum_vol_2_figs/Figures_Lec_26/fig_1.jpeg}}}%
}\par

\medskip{}

\textcolor{blue}{)המחשה גיאומטרית}

\textcolor{blue}{נשים לב שהנפח של זה הוא למעשה אפס(}\\
\textcolor{blue}{\newpage}

\textbf{\textcolor{violet}{אינטואיציה גיאומטרית: התאפסות פונקציית
נפח על שורות זהות}}

\textcolor{violet}{השורות של המטריצה \L{$A$} מייצגות וקטורים הפורשים
מקבילון במרחב \L{$n$}-ממדי. כאשר שתי שורות במטריצה הן זהות, \L{$\Lambda_{i}^{A}=\Lambda_{j}^{A}$},
המשמעות היא ששני וקטורים הפורשים את הצורה מתלכדים )מצביעים לאותו כיוון
ובעלי אותו אורך(. כתוצאה מכך, המקבילון "קורס" לממד נמוך יותר )לדוגמה:
מנפח תלת-ממדי למישור דו-ממדי שטוח(. במצב זה, ה"נפח" ה-\L{$n$}-ממדי
של הצורה חייב להיות \L{$0$}.}\\

\textcolor{blue}{בעזרת האספקט הזה נגדיר את המושג פונקציה מתחלפת, שבעזרתו
נגדיר פונקציית נפח כפונקציה שהיא מולטי-לינארית ומתחלפת.}

\textcolor{blue}{נשים לב - פונקציית נפח זה עוד לא ההגדרה של דטרמיננטה.}

\textcolor{blue}{יש פונקציות נפח שאינן הדטרמיננטה, אבל הדטרמיננטה
תהיה פונקציית נפח.}\\

\textcolor{violet}{אוסף כל הפונקציות המקיימות את שתי התכונות הנ\textquotedbl ל
)מולטי-לינאריות ומתחלפת( הוא מרחב וקטורי ממימד 1. כלומר, יש אינסוף
פונקציות כאלה, וההבדל ביניהן הוא רק כפל בסקלר )קבוע(.}

\textcolor{violet}{אז מה הופך פונקציית נפח לדטרמיננטה?}

\textcolor{violet}{כדי שפונקציית נפח תהיה הדטרמיננטה, צריך להוסיף
תנאי שלישי )תנאי נרמול(: \L{%
\[
f(I)=1
\]
}%
כלומר, הפונקציה צריכה להחזיר ערך של 1 עבור מטריצת היחידה )קוביה בנפח
יחידה(. זה נכון כי הדטרמיננטה היא מקרה פרטי )ויחיד( של משפחת פונקציות
הנפח, אך לא כל פונקציית נפח היא בהכרח הדטרמיננטה.}

\section{{\beginL 5\endL}. הגדרה: פונקציה מתחלפת}

תהי \L{$D\colon M_{n\times n}(\mathbb{F})\to\mathbb{F}$}. אומרים
כי \L{$D$} \textbf{מתחלפת} כאשר לכל \L{$A\in M_{n\times n}(\mathbb{F})$}
כך ש- \L{$\exists1\leq k\leq n-1:\quad\Lambda_{k+1}^{A}=\Lambda_{k}^{A}$}
מתקיים: \L{$D(A)=0$}.\\

\textcolor{blue}{כלומר, כשיש למטריצה שתי שורות זהות סמוכות, פונקציה
מתחלפת נותנת ערך של אפס כשמפעילים אותה עליה.}

\section{פונקציות נפח}

\section{{\beginL 6\endL}. הגדרה: פונקציית נפח}

אומרים כי \L{$D\colon M_{n\times n}(\mathbb{F})\to\mathbb{F}$} היא
\textbf{פונקציית נפח} )פ\textquotedblright נ( כאשר היא מולטי-ליניארית
ומתחלפת.\\

\textcolor{blue}{עכשיו נלמד תכונות של פונקציות נפח.}

\textcolor{blue}{כשאנחנו חוקרים פ\textquotedbl נ, קו המחשבה שלנו
יהיה איך ערך שלה משתנה על מטריצה כאשר נבצע על המטריצה אחת מפעולות
השורה האלמנטריות.}

\textcolor{blue}{למשל, אם כופלים שורה בסקלר, אז לפי הגדרת הלינאריות
לפי אותה שורה, הערך של פונ' הנפח יהיה מוכפל באותו הסקלר.}

\section{{\beginL 7\endL}. תכונה}
\begin{enumerate}
\item אם \L{$D_{1},D_{2}\colon M_{n\times n}(\mathbb{F})\to\mathbb{F}$}
הן פ\textquotedblright נ, אז \L{$D_{1}+D_{2}$} היא גם פ\textquotedblright נ.
\item אם \L{$D\colon M_{n\times n}(\mathbb{F})\to\mathbb{F}$} היא פ\textquotedblright נ
ו-\L{$c\in\mathbb{F}$} סקלר, אז \L{$c\cdot D$} היא גם פ\textquotedblright נ.
\end{enumerate}

\section{{\beginL 8\endL}. תכונה}

תהי \L{$D\colon M_{n\times n}(\mathbb{F})\to\mathbb{F}$} פ\textquotedblright נ,
ותהי \L{$A\in M_{n\times n}(\mathbb{F})$}. נסמן ב-\L{$\varepsilon$}
את פעולת השורה האלמנטרית המוגדרת על ידי \L{$R_{i}\to c\cdot R_{i}$}
)כפל השורה ה-\L{$i$} בסקלר \L{$c\in\mathbb{F}$}(, עבור \L{$1\le i\le n$}.
אזי מתקיים: \L{%
\[
D(\varepsilon(A))=c\cdot D(A)
\]
}%


\paragraph{\textcolor{violet}{הוכחה}}

\textcolor{violet}{המטריצה \L{$A$} והמטריצה לאחר השינוי \L{$\varepsilon(A)$}
נראות כך: \L{%
\[
A=\begin{pmatrix}\vdots\\
R_{i}\\
\vdots
\end{pmatrix},\quad\varepsilon(A)=\begin{pmatrix}\vdots\\
c\cdot R_{i}\\
\vdots
\end{pmatrix}
\]
}%
}

\textcolor{violet}{מכיוון ש-\L{$D$} היא פונקציה מולטי-ליניארית מתקיים}

\textcolor{violet}{\L{%
\[
D(\varepsilon(A))=D\begin{pmatrix}\vdots\\
c\cdot R_{i}\\
\vdots
\end{pmatrix}=c\cdot D\begin{pmatrix}\vdots\\
R_{i}\\
\vdots
\end{pmatrix}=c\cdot D(A)
\]
}%
}

\section{{\beginL 9\endL}. טענה}

תהי \L{$D\colon M_{n\times n}(\mathbb{F})\to\mathbb{F}$} פ\textquotedblright נ,
ותהי \L{$A\in M_{n\times n}(\mathbb{F})$}. נסמן ב-\L{$\varepsilon$}
את פעולת השורה האלמנטרית המוגדרת על ידי \L{$R_{k}\leftrightarrow R_{k+1}$}
)החלפת שורות סמוכות(, עבור \L{$1\le k\le n-1$}. אזי מתקיים: \L{%
\[
D(\varepsilon(A))=-D(A)
\]
}%


\subsection{הוכחה}

לכל \L{$1\le i\le n$} נסמן ב-\L{$\Lambda_{i}$} את השורה ה-\L{$i$}
של \L{$A$}. אזי: \L{%
\[
A=\begin{pmatrix}\vdots\\
\Lambda_{k-1}\\
\Lambda_{k}\\
\Lambda_{k+1}\\
\Lambda_{k+2}\\
\vdots
\end{pmatrix}
\]
}%
נגדיר מטריצות: \L{%
\[
A_{1}=\begin{pmatrix}\vdots\\
\Lambda_{k-1}\\
\Lambda_{k}+\Lambda_{k+1}\\
\Lambda_{k}+\Lambda_{k+1}\\
\Lambda_{k+2}\\
\vdots
\end{pmatrix},\quad A_{2}=\begin{pmatrix}\vdots\\
\Lambda_{k-1}\\
\Lambda_{k}\\
\Lambda_{k}+\Lambda_{k+1}\\
\Lambda_{k+2}\\
\vdots
\end{pmatrix},\quad A_{3}=\begin{pmatrix}\vdots\\
\Lambda_{k-1}\\
\Lambda_{k+1}\\
\Lambda_{k}+\Lambda_{k+1}\\
\Lambda_{k+2}\\
\vdots
\end{pmatrix}
\]
}%
מאחר ש-\L{$D$} פ\textquotedblright נ היא מתחלפת ולכן המטריצה \L{$A_{1}$}
)שבה יש שתי שורות זהות סמוכות( מקיימת: \L{%
\[
D(A_{1})=0
\]
}%
מאחר ש-\L{$D$} פ\textquotedblright נ היא מולטי-לינארית ובפרט ליניארית
לפי השורה ה-\L{$k$} לכן\L{%
\[
D(A_{1})=D(A_{2})+D(A_{3})
\]
}%

\textcolor{teal}{למעשה נשים לב שהשורה ה-\L{$k$} של \L{$A_{1}$} היא
סכום השורות ה-\L{$k$} של \L{$A_{2}$} ו-\L{$A_{3}$}, ושאר השורות
זהות בשלוש המטריצות.}

\textcolor{teal}{לכן, מאחר ש-\L{$D$} ליניארית לפי השורה ה-\L{$k$},
נקבל את זה.}\\

נגדיר מטריצות נוספות: \L{%
\[
A_{4}=\begin{pmatrix}\vdots\\
\Lambda_{k-1}\\
\Lambda_{k}\\
\Lambda_{k}\\
\Lambda_{k+2}\\
\vdots
\end{pmatrix},\quad A_{5}=\begin{pmatrix}\vdots\\
\Lambda_{k-1}\\
\Lambda_{k}\\
\Lambda_{k+1}\\
\Lambda_{k+2}\\
\vdots
\end{pmatrix}=A
\]
}%
\L{%
\[
A_{6}=\begin{pmatrix}\vdots\\
\Lambda_{k-1}\\
\Lambda_{k+1}\\
\Lambda_{k}\\
\Lambda_{k+2}\\
\vdots
\end{pmatrix}=\varepsilon(A),\quad A_{7}=\begin{pmatrix}\vdots\\
\Lambda_{k-1}\\
\Lambda_{k+1}\\
\Lambda_{k+1}\\
\Lambda_{k+2}\\
\vdots
\end{pmatrix}
\]
}%
מאחר ש-\L{$D$} מתחלפת )ולמטריצות \L{$A_{4},A_{7}$} יש שורות סמוכות
זהות(, מתקיים: \L{%
\[
D(A_{4})=D(A_{7})=0
\]
}%
מאחר ש-\L{$D$} ליניארית לפי השורה ה-\L{$(k+1)$}, מתקיים: \L{%
\[
D(A_{2})=D(A_{4})+D(A_{5})
\]
}%
\L{%
\[
D(A_{3})=D(A_{6})+D(A_{7})
\]
}%
לסיכום: \L{%
\[
0=D(A_{1})=\underbrace{D(A_{2})}_{D(A_{4})+D(A_{5})}+\underbrace{D(A_{3})}_{D(A_{6})+D(A_{7})}=\underbrace{D(A_{4})}_{0}+\underbrace{D(A_{5})}_{D(A)}+\underbrace{D(A_{6})}_{D(\varepsilon(A))}+\underbrace{D(A_{7})}_{0}
\]
}%

אז\L{%
\[
0=0+D(A)+D(\varepsilon(A))+0=D(A)+D(\varepsilon(A))
\]
}%

ולכן

\L{%
\[
D(\varepsilon(A))=-D(A)
\]
}%
\medskip{}

\textcolor{blue}{עכשיו נרצה להרחיב את הטענה הזאת )9( למקרה שבו פעולת
השורה האלמנטרית מחליפה שורות שהן לאו דווקא סמוכות.}

\section{{\beginL 10\endL}. טענה}

תהי \L{$D\colon M_{n\times n}(\mathbb{F})\to\mathbb{F}$} פ\textquotedblright נ,
ותהי \L{$A\in M_{n\times n}(\mathbb{F})$}. נסמן ב-\L{$\varepsilon$}
את פעולת השורה האלמנטרית המוגדרת על ידי \L{$R_{i}\leftrightarrow R_{j}$}
)החלפת שורות( עבור \L{$1\le i<j\le n$}. אזי מתקיים: \L{%
\[
D(\varepsilon(A))=-D(A)
\]
}%


\subsection{הוכחה}

לכל \L{$1\le k\le n-1$} נסמן \L{$\varepsilon_{k}=R_{k}\leftrightarrow R_{k+1}$}. 

אז לכל \L{$A\in M_{n\times n}(\mathbb{F})$} מתקיים השוויון הבא )הרכבת
חילופים סמוכים(: \L{%
\[
\varepsilon(A)=\varepsilon_{i}(\varepsilon_{i+1}(\dots\varepsilon_{j-2}(\varepsilon_{j-1}(\varepsilon_{j-2}(\dots\varepsilon_{i+1}(\varepsilon_{i}(A))\dots)))\dots))
\]
}%
מכאן נובע כי ביצוע של פש\textquotedblright א \L{$\varepsilon$} שקול
לביצוע של מספר אי-זוגי של פש\textquotedblright א מהצורה \L{$\varepsilon_{k}$}
. לכן, לפי הטענה הקודמת ){\beginL 9\endL}(, מתקיים: \L{%
\[
D(\varepsilon(A))=(-1)^{\text{odd}}D(A)=-D(A)
\]
}%


\subsection{\textcolor{violet}{הסבר: מדוע החלפת שורות היא הרכבה של מספר פעמים
אי-זוגי של החלפת שורות סמוכות}}

\subsubsection{\textcolor{violet}{האינטואיציה: "טיול הלוך-חזור"}}

\textcolor{violet}{כדי להחליף בין השורה \L{$i$} לשורה \L{$j$} )כאשר
\L{$i<j$}( בעזרת חילופי שכנים בלבד, עלינו לבצע את המהלכים הבאים: }
\begin{enumerate}
\item \textcolor{violet}{הורדת השורה \L{$i$}: אנו מעבירים את השורה ה-\L{$i$}
מטה, "דילוג" אחרי "דילוג", עד שהיא מגיעה להיות השכנה הצמודה )מלמעלה(
של השורה ה-\L{$j$}. לשם כך אנו מפעילים סדרה של חילופים: \L{$\varepsilon_{i},\varepsilon_{i+1},\dots,\varepsilon_{j-2}$}. }
\item \textcolor{violet}{ההחלפה המרכזית: כעת, כשהשורות סמוכות )במקומות \L{$j-1$}
ו-\L{$j$}(, אנו מחליפים ביניהן. זוהי הפעולה \L{$\varepsilon_{j-1}$}.
פעולה זו מבוצעת פעם אחת בלבד. }
\item \textcolor{violet}{העלאת השורה \L{$j$}: כעת השורה שהייתה במקור ב-\L{$j$}
נמצאת במיקום \L{$j-1$}. עלינו להעלות אותה כל הדרך למעלה למקום ה-\L{$i$}.
אנו עושים זאת על ידי ביצוע אותם חילופים שעשינו קודם, אך בסדר הפוך:
\L{$\varepsilon_{j-2},\dots,\varepsilon_{i}$}. }
\end{enumerate}

\subsubsection{\textcolor{violet}{ספירת החילופים}}

\textcolor{violet}{נסמן ב-\L{$K$} את מספר החילופים שנדרשו בשלב הראשון
)הירידה(. בשלב השלישי )העלייה( נדרש בדיוק אותו מספר חילופים, \L{$K$}.
בשלב השני )האמצעי( התבצע חילוף בודד )\L{$1$}(. סה"כ מספר החילופים
הוא: \L{%
\[
\text{Total}=\underbrace{K}_{\text{(step 1)}}+\underbrace{1}_{(\text{step }2)}+\underbrace{K}_{\text{(step 3)}}=2K+1
\]
}%
המספר \L{$2\text{K}+1$} הוא בהכרח מספר אי-זוגי.}\\

\textcolor{blue}{כלומר,רק אפסילון של \L{$\varepsilon_{j-1}$} חוזר
פעם אחת כי הוא מחליף את שתי השורות החשובות בו זמנית.}

\subsection{\textcolor{violet}{דוגמה: החלפת שורות במטריצה \L{$3\times3$}}}

\textcolor{blue}{{]}אלכס המליץ לעשות לעצמנו דוגמה שכזאת{[}}

\textcolor{violet}{נניח כי \L{$n=3$} ואנו רוצים להחליף את השורה הראשונה
)\L{$R_{1}$}( עם השורה השלישית )\L{$R_{3}$}(. כלומר: \L{$i=1,j=3$}.
המטריצה ההתחלתית: \L{%
\[
A=\begin{pmatrix}R_{1}\\
R_{2}\\
R_{3}
\end{pmatrix}
\]
}%
הנוסחה לביצוע החילוף באמצעות שכנים )לפי \L{$i=1,j=3$}( היא: \L{%
\[
\varepsilon=\varepsilon_{1}\circ\varepsilon_{2}\circ\varepsilon_{1}
\]
}%
)שימו לב: \L{$\varepsilon_{j-1}=\varepsilon_{2}$} מופיע פעם אחת,
ו-\L{$\varepsilon_{1}$} מופיע פעמיים(.}

\subsubsection{\textcolor{violet}{ביצוע שלב אחר שלב:}}
\begin{enumerate}
\item \textcolor{violet}{שלב 1 )הפעלת \L{$\varepsilon_{1}$}(: נחליף את
שורה 1 עם שורה 2. \L{%
\[
\begin{pmatrix}R_{1}\\
R_{2}\\
R_{3}
\end{pmatrix}\xrightarrow{\varepsilon_{1}}\begin{pmatrix}R_{2}\\
R_{1}\\
R_{3}
\end{pmatrix}
\]
}%
)כעת \L{$R_{1}$} ירדה שלב אחד(.}
\item \textcolor{violet}{שלב 2 )הפעלת \L{$\varepsilon_{2}$}(: נחליף את
שורה 2 עם שורה 3. \L{%
\[
\begin{pmatrix}R_{2}\\
R_{1}\\
R_{3}
\end{pmatrix}\xrightarrow{\varepsilon_{2}}\begin{pmatrix}R_{2}\\
R_{3}\\
R_{1}
\end{pmatrix}
\]
}%
)כעת \L{$R_{1}$} הגיעה לתחתית, ו-\L{$R_{3}$} התחילה לעלות(.}
\item \textcolor{violet}{שלב 3 )הפעלת \L{$\varepsilon_{1}$}(: נחליף את
שורה 1 עם שורה 2. \L{%
\[
\begin{pmatrix}R_{2}\\
R_{3}\\
R_{1}
\end{pmatrix}\xrightarrow{\varepsilon_{1}}\begin{pmatrix}R_{3}\\
R_{2}\\
R_{1}
\end{pmatrix}
\]
}%
)כעת \L{$R_{3}$} הגיעה לראש המטריצה(. }
\end{enumerate}
\textcolor{violet}{מסקנה: ביצענו סך הכל 3 חילופים )מספר אי-זוגי(,
והגענו לתוצאה המבוקשת \L{$R_{1}\leftrightarrow R_{3}$}. לכן: \L{%
\[
D(\varepsilon(A))=(-1)^{3}D(A)=-D(A)
\]
}%
}

\textcolor{blue}{מהטענה הזאת אפשר להסיק )זאת תהיה הטענה הבאה, שהיא
חשובה( מה פ\textquotedblright נ מחזירה על מטריצה עם שתי שורות זהות
שאינן סמוכות )\L{$-0=0$}(...}

\textcolor{blue}{)מה שעושים זה שאפשר ליצור ממטריצה כזאת מטריצה עם
שתי שורות זהות וסמוכות ע\textquotedbl י החלפת שורה, ואז לפי מה שהראינו
זה מינוס על הפונקציה על המטריצה הקודמת אבל זה יוצא מינוס על אפס שזה
אפס(}

\section{{\beginL 11\endL}. טענה}

תהי \L{$D\colon M_{n\times n}(\mathbb{F})\to\mathbb{F}$} פ\textquotedblright נ.

תהי \L{$A\in M_{n\times n}(\mathbb{F})$} כך שעבור \L{$1\le i<j\le n$}
מתקיים: \L{%
\[
\Lambda_{i}^{A}=\Lambda_{j}^{A}
\]
}%
אזי מתקיים: \L{%
\[
D(A)=0
\]
}%


\subsection{הוכחה}

אם \L{$j=i+1$} )השורות סמוכות(, אזי \L{$D(A)=0$} כי \L{$D$} מתחלפת.

אם \L{$j\neq i+1$}, נגדיר את פעולת השורה האלמנטרית \L{$\varepsilon=R_{i+1}\leftrightarrow R_{j}$}
ונסמן \L{$B=\varepsilon(A)$}. אזי מתקיים: \L{%
\[
\Lambda_{i}^{B}=\Lambda_{i}^{A}=\Lambda_{j}^{A}=\Lambda_{i+1}^{B}
\]
}%
כלומר, במטריצה \L{$B$} השורות ה-\L{$i$} וה-\L{$(i+1)$} )שהן שורות
סמוכות( הן זהות. מכאן, כי \L{$D$} מתחלפת: \L{%
\[
D(B)=0
\]
}%
לפי הטענה הקודמת ){\beginL 10\endL}(, מתקיים: \L{%
\[
0=D(B)=D(\varepsilon(A))=-D(A)
\]
}%
ולכן \L{$D(A)=0$}.

\medskip{}

\textcolor{blue}{עכשיו נראה הרחבה נוספת - אם שורה אחת היא מכפלה בסקלר
של שורה אחרת, אז פונקציית נפח על המטריצה תחזיר שוב אפס.}

\section{{\beginL 12\endL}. טענה}

תהי \L{$D\colon M_{n\times n}(\mathbb{F})\to\mathbb{F}$} פ\textquotedblright נ.
יהיו \L{$1\le i,j\le n$} עם \L{$i\neq j$} וסקלר \L{$c\in\mathbb{F}$}.
תהי \L{$A\in M_{n\times n}(\mathbb{F})$} כך שמתקיים: \L{%
\[
\Lambda_{j}^{A}=c\Lambda_{i}^{A}
\]
}%
אזי מתקיים: \L{%
\[
D(A)=0
\]
}%


\subsection{הוכחה}

אם \L{$c=0$}, נקבל כי \L{$\Lambda_{j}^{A}$} היא שורת אפסים, ולכן
לפי טענה שראינו ){\beginL 4\endL}( \L{$D(A)=0$}. אם \L{$c\neq0$},
נגדיר את פעולת השורה האלמנטרית \L{$\varepsilon=R_{i}\to cR_{i}$}
ונסמן \L{$B=\varepsilon(A)$}. לכן מתקיים: \L{%
\[
\Lambda_{i}^{B}=c\Lambda_{i}^{A}=\Lambda_{j}^{A}=\Lambda_{j}^{B}
\]
}%
כלומר, למטריצה \L{$B$} יש שתי שורות זהות.

לכן, לפי הטענה הקודמת ){\beginL 11\endL}(, \L{$D(B)=0$}. בנוסף, ידוע
כי: \L{%
\[
0=D(B)=cD(A)
\]
}%
מאחר ש-\L{$c\neq0$}, מתקיים \L{$D(A)=0$}.

\medskip{}

\textcolor{teal}{מכפילים את השורה ה-\L{$i$} בסקלר \L{$c$} )וכך השורה
ה-\L{$i$} הופכת לשווה לשורה ה-\L{$j$}( כדי ליצור מטריצה חדשה \L{$B$}
בעלת שתי שורות זהות - כבר הייתה לנו שורה שהיא כפל בסקלר של השורה המקורית
ועכשיו הכפלנו את השורה המקורית באותו סקלר כדי שתהיה שווה לה. לכן \L{$D(B)=0$},
ומתוך תכונת ההומוגניות )\L{$D(B)=c\cdot D(A)$}( וכך מסיקים שבהכרח
\L{$D(A)=0$}.}

\medskip{}

\textcolor{blue}{עכשיו נראה מה קורה כשעושים את הפש\textquotedbl א
האחרונה שחסרה לנו: הוספה של שורה אחרת כפול סקלר.}

\section{{\beginL 13\endL}. טענה}

תהי \L{$D\colon M_{n\times n}(\mathbb{F})\to\mathbb{F}$} פ\textquotedblright נ.
יהיו \L{$1\le i,j\le n$} עם \L{$i\neq j$} וסקלר \L{$c\in\mathbb{F}$}.
תהי \L{$\varepsilon$} פעולת השורה האלמנטרית המוגדרת על ידי \L{$R_{i}\to R_{i}+cR_{j}$}.
אזי לכל \L{$A\in M_{n\times n}(\mathbb{F})$} מתקיים: \L{%
\[
D(\varepsilon(A))=D(A)
\]
}%


\subsection{הוכחה}

מתקיים: \L{%
\[
A=\begin{pmatrix}\vdots\\
\Lambda_{i-1}\\
\Lambda_{i}\\
\Lambda_{i+1}\\
\vdots\\
\Lambda_{j}\\
\vdots
\end{pmatrix}
\]
}%
נגדיר מטריצת עזר \L{$A'$} כך: \L{%
\[
A'=\begin{pmatrix}\vdots\\
\Lambda_{i-1}\\
c\Lambda_{j}\\
\Lambda_{i+1}\\
\vdots\\
\Lambda_{j}\\
\vdots
\end{pmatrix}
\]
}%
המטריצה המתקבלת לאחר הפעולה היא: \L{%
\[
\varepsilon(A)=\begin{pmatrix}\vdots\\
\Lambda_{i-1}\\
\Lambda_{i}+c\Lambda_{j}\\
\Lambda_{i+1}\\
\vdots\\
\Lambda_{j}\\
\vdots
\end{pmatrix}
\]
}%
לפי הטענה הקודמת ){\beginL 12\endL}( מתקיים \L{$D(A')=0$} )שכן השורה
ה-\L{$i$} היא כפולה בסקלר של השורה ה-\L{$j$}(. בנוסף, לפי הליניאריות
של \L{$D$} לפי השורה ה-\L{$i$} מתקיים: \L{%
\[
D(\varepsilon(A))=D(A')+D(A)=0+D(A)=D(A)
\]
}%

\newpage{}

\noindent\fbox{\begin{minipage}[t]{1\columnwidth - 2\fboxsep - 2\fboxrule}%
\begin{center}
אלגברה ליניארית )1( - \L{$80134$} \textbar{} האונ' העברית \textbar{}
תשפ\textquotedblright ו א' \textbar{} מרצה: ד\textquotedblright ר
אלכס גורביץ' \textbar{} מתרגל: מר יואב כהן \textbar{} ליקוט ועריכה:
נעם לב
\par\end{center}%
\end{minipage}}

\part{הרצאה {\beginL 27\endL} - {\beginL 15.01.2026\endL}}

\textcolor{teal}{הרצאת השלמה בחמישי}

\section{תזכורת}

\section{{\beginL 1\endL}. הגדרה: פונקציית נפח}

פונקציה \L{$D:M_{n\times n}(\mathbb{F})\to\mathbb{F}$} נקראת פ\textquotedblright נ
כאשר היא מולטי-לינארית ומתחלפת.

\section{{\beginL 2\endL}. טענה}

\textcolor{teal}{זו טענה המסכמת את הטענות שהוכחנו בהרצאה הקודמת על
התנהגות פ\textquotedblright נ ביחס לפש\textquotedblright א.}\\

יהיו \L{$D:M_{n\times n}(\mathbb{F})\to\mathbb{F}$} פ\textquotedblright נ,
\L{$A\in\sqrmn$} ו- \L{$\varepsilon$} פש\textquotedblright א. אז:
\begin{enumerate}
\item אם \L{$\varepsilon=R_{i}\to cR_{i}$} עבור \L{$c\neq0$}, אז \L{$D(\varepsilon(A))=cD(A)$}.
\item אם \L{$\varepsilon=R_{i}\to R_{i}+cR_{j}$} עבור \L{$j\neq i$}, אז
\L{$D(\varepsilon(A))=D(A)$}.
\item אם \L{$\varepsilon=R_{i}\leftrightarrow R_{j}$} עבור \L{$j\neq i$},
אז \L{$D(\varepsilon(A))=-D(A)$}.
\end{enumerate}
\medskip{}

\rule[0.5ex]{1\columnwidth}{1pt}

\medskip{}

\textcolor{blue}{עכשיו נגדיר משהו חדש שיפשט את העבודה עם טענות על
פ\textquotedblright נ.}

\section{{\beginL 3\endL}. הגדרה}

\textcolor{teal}{כפי שראינו בטענה }\textcolor{teal}{{\beginL 2}\endL}\textcolor{teal}{{}
בתזכורת, כשמפעילים פש\textquotedblright א על מטריצה מסוימת, }

\textcolor{teal}{אנחנו יודעים להגיד איך הערך שפ\textquotedblright נ
תחזיר עליה ישתנה.}

\textcolor{teal}{וכפי שראינו, הוא ישתנה תמיד באופן מסוים: מקדם כלשהו
כפול הערך שפונקציית הנפח מחזירה כשהיא מופעלת על המטריצה המקורית, כאשר
המקדם הזה נקבע בדיוק לפי פעולת השורה האלמנטרית שעשינו למטריצה המקורית.}

\textcolor{teal}{כאן אנחנו מגדירים את המקדם הזה, לעבודה יותר נוחה
בהמשך.}\\

לכל פש"א \L{$\varepsilon$} נגדיר את \L{$\mu(\varepsilon)\in\mathbb{F}$}
באופן הבא: 
\begin{enumerate}
\item אם \L{$\varepsilon=R_{i}\to cR_{i}$} כאשר \L{$c\neq0$}, אז \L{$\mu(\varepsilon)=c$}. 
\item אם \L{$\varepsilon=R_{i}\to R_{i}+cR_{j}$} כאשר \L{$i\neq j$}, אז
\L{$\mu(\varepsilon)=1$}. 
\item אם \L{$\varepsilon=R_{i}\leftrightarrow R_{j}$} כאשר \L{$i\neq j$},
אז \L{$\mu(\varepsilon)=-1$}. 
\end{enumerate}

\section{{\beginL 4\endL}. טענה}

יהיו \L{$D:M_{n\times n}(\mathbb{F})\to\mathbb{F}$} פ"נ, \L{$A\in M_{n\times n}(\mathbb{F})$},
ו- \L{$\varepsilon$} פש"א. אזי: \L{%
\[
D(\varepsilon(A))=\mu(\varepsilon)\cdot D(A)
\]
}%

\textcolor{blue}{למעשה זה ניסוח מחדש של טענה קודמת )2(, אותה הוכחנו
כטענות מופרדות בהרצאה הקודמת.}

\section{{\beginL 5\endL}. טענה}

תהי \L{$D:M_{n\times n}(\mathbb{F})\to\mathbb{F}$} פ"נ, \L{$A\in M_{n\times n}(\mathbb{F})$}.
נסמן ב-\L{$\varepsilon_{1},\dots,\varepsilon_{k}$} את הפש"א שבעזרתן
מדרגים את \L{$A$}, ונסמן ב- \L{$R$} את המטריצה המדורגת המצומצמת
המתקבלת בתהליך הדירוג. אז: \L{%
\[
D(R)=\mu(\varepsilon_{k})\cdot\cdots\cdot\mu(\varepsilon_{1})\cdot D(A)
\]
}%


\subsection{הוכחה}

מתקיים \L{$R=\varepsilon_{k}(\dots\varepsilon_{2}(\varepsilon_{1}(A))\dots)$}.
אז משימוש חוזר בטענה הקודמת ){\beginL 4\endL}( נובעת הנוסחה הנדרשת.

\medskip{}
\textcolor{violet}{בפירוט:}

\textcolor{violet}{נסמן את סדרת המטריצות בדרך לדירוג באופן הבא: }

\textcolor{violet}{נגדיר \L{$A_{0}=A$}.}

\textcolor{violet}{נגדיר \L{$A_{1}=\varepsilon_{1}(A_{0})$}. אז לפי
טענה }\textcolor{violet}{{\beginL 4}\endL}\textcolor{violet}{: \L{%
\[
D(A_{1})=\mu(\varepsilon_{1})\cdot D(A)
\]
}%
}

\textcolor{violet}{נגדיר \L{$A_{2}=\varepsilon_{2}(A_{1})$}. נציב
את הביטוי של \L{$D(A_{1})$}:\L{%
\[
D(A_{2})=\mu(\varepsilon_{2})D(A_{1})=\mu(\varepsilon_{2})\cdot\mu(\varepsilon_{1})\cdot D(A)
\]
}%
אם נמשיך בתהליך זה \L{$k$} פעמים, עבור המטריצה המדורגת הסופית \L{$R=A_{k}$},
נקבל את מכפלת כל הסקלרים של הפעולות שבוצעו בדרך:\L{%
\[
D(R)=\mu(\varepsilon_{k})\cdot\mu(\varepsilon_{k-1})\cdot\dots\cdot\mu(\varepsilon_{1})D(A)
\]
}%
}\medskip{}

\textcolor{blue}{עכשיו בפעם הראשונה נוכל להגיד בדיוק מה מחזירה פ\textquotedblright נ
על מטריצה מסוימת באופן שאינו יחסי:}

\textcolor{blue}{כפי שלמדנו, במטריצה המדורגת-מצומצמת של מטריצה שאינה
הפיכה יש שורת אפסים.}

\textcolor{blue}{מכך ומהטענה שעכשיו הוכחנו אפשר להסיק שהערך של פ\textquotedblright נ
עליה יהיה אפס. זאת תהיה הטענה הבאה.}

\medskip{}


\section{{\beginL 6\endL}. טענה}

תהי \L{$D:M_{n\times n}(\mathbb{F})\to\mathbb{F}$} פ"נ. תהי \L{$A\in M_{n\times n}(\mathbb{F})$}
מטריצה שאינה הפיכה. אז \L{$D(A)=0$}.

\subsection{הוכחה}

נסמן ב-\L{$R$} את המטריצה המדורגת המצומצמת המתקבלת מ-\L{$A$} ע"י
דירוג, ונסמן ב- \L{$\varepsilon_{1},\dots,\varepsilon_{k}$} את הפש"א
שנעשו במהלך הדירוג. אז )לפי טענה {\beginL 5\endL}(: \L{%
\[
D(R)=\mu(\varepsilon_{k})\cdot\dots\cdot\mu(\varepsilon_{1})\cdot D(A)
\]
}%
מאחר ש-\L{$A$} אינה הפיכה, \L{$R\neq I_{n}$}. לכן ב-\L{$R$} יש
שורת אפסים, ומכאן \textcolor{teal}{)לפי תכונה 4 בסיכום שהכנתי להרצאה
}\textcolor{teal}{{\beginL 26}\endL}\textcolor{teal}{(} \L{$D(R)=0$}.
מכך שלכל \L{$1\le i\le k$} מתקיים \L{$\mu(\varepsilon_{i})\neq0$}
)מההגדרה(, נקבל \L{$D(A)=0$}.\\

\textcolor{blue}{המשפט הבא זה המשפט הכי חשוב על פ\textquotedblright נ.}

\textcolor{blue}{ניזכר שמטריצה הפיכה תמיד שקולת שורה למטריצת היחידה.}

\section{{\beginL 7\endL}. משפט}

תהיינה \L{$D,D':M_{n\times n}(\mathbb{F})\to\mathbb{F}$} פונקציות
נפח, כאשר \L{$D\neq0$} 

)כלומר, קיימת מטריצה \L{$K\in M_{n\times n}(\mathbb{F})$} כך ש-\L{$D(K)\neq0$}(.

אזי:
\begin{enumerate}
\item \L{$D(I_{n})\neq0$}. 
\item \L{$D'=cD$}, כלומר לכל \L{$A\in M_{n\times n}(\mathbb{F})$} מתקיים
\L{$D'(A)=cD(A)$}, כאשר הסקלר \L{$c$} מוגדר ע"י \L{$c=\frac{D'(I_{n})}{D(I_{n})}$}. 
\end{enumerate}

\subsection{הוכחה}

{\beginL 1\endL}. לפי הטענה הקודמת )טענה {\beginL 6\endL}( \L{$K$}
מטריצה הפיכה. לכן ניתן לקבל את \L{$I_{n}$} מ-\L{$K$} ע"י סדרה של
פש"א \L{$\varepsilon_{1},\dots,\varepsilon_{k}$}. אז: \L{%
\[
D(I_{n})=\mu(\varepsilon_{k})\cdot\cdots\cdot\mu(\varepsilon_{1})\cdot D(K)\neq0
\]
}%
)זאת מכיוון ש-\L{$D(K)\neq0$} וגם \L{$\mu(\varepsilon_{i})\neq0$}
לכל \L{$i$}(.

{\beginL 2\endL}. אם \L{$A$} אינה הפיכה, אז לפי הטענה הקודמת )טענה
{\beginL 6\endL}(: \L{%
\[
D'(A)=0=c\cdot0=cD(A)
\]
}%
והשוויון מתקיים.

אם \L{$A$} הפיכה, ניתן לקבל את \L{$I_{n}$} מ-\L{$A$} ע"י סדרה
של פש"א \L{$\varepsilon_{1},\dots,\varepsilon_{k}$}. לכן )לפי טענה
{\beginL 5\endL}( מתקיים: \L{%
\[
D(I_{n})=\mu(\varepsilon_{k})\cdot\dots\cdot\mu(\varepsilon_{1})\cdot D(A)
\]
}%
וגם: \L{%
\[
D'(I_{n})=\mu(\varepsilon_{k})\cdot\dots\cdot\mu(\varepsilon_{1})\cdot D'(A)
\]
}%
מכאן \L{%
\[
D'(A)=\frac{D'(I_{n})}{\mu(\varepsilon_{k})\cdot\dots\cdot\mu(\varepsilon_{1})}=\frac{D'(I_{n})\cdot D(A)}{\mu(\varepsilon_{k})\cdot\dots\cdot\mu(\varepsilon_{1})\cdot D(A)}=\frac{D'(I_{n})\cdot D(A)}{D(I_{n})}=c\cdot D(A)
\]
}%
\medskip{}

\textcolor{teal}{הערכים \L{$\mu(\varepsilon_{k})\cdot\dots\cdot\mu(\varepsilon_{1})$}
זהים בשתי המשוואות מכיוון שבשני המקרים מדובר על הפש\textquotedblright א
שצריך להפעיל על אותה מטריצה \L{$A$} בשביל להעביר אותה לצורה מדורגת-מצומצמת.}

\section{{\beginL 8\endL}. הגדרה: פונקציית נפח מנורמלת}

פונקציית נפח \L{$D:M_{n\times n}(\mathbb{F})\to\mathbb{F}$} נקראת
\textbf{מנורמלת} כאשר: \L{%
\[
D(I_{n})=1
\]
}%

\medskip{}

\textcolor{blue}{בעצם ראינו שפ\textquotedblright נ מאופיינת לפי הערך
שהיא מחזירה כשמפעילים אותה על מטריצת היחידה, וכשהיא מחזירה 1 נקרא
לה מנורמלת.}

\section{{\beginL 9\endL}. מסקנה}

תהיינה \L{$D,D':M_{n\times n}(\mathbb{F})\to\mathbb{F}$} פ"נ מנורמלות.

אז \L{$D'=D$}, כלומר לכל \L{$A\in M_{n\times n}(\mathbb{F})$} מתקיים
\L{$D'(A)=D(A)$}.

\subsection{הוכחה}

מאחר ש-\L{$D$} מנורמלת מתקיים \L{$D\neq0$} )שהרי \L{$D(I_{n})=1$}(.

לפי המשפט הקודם, \L{$D'=cD$} כאשר: \L{%
\[
c=\frac{D'(I_{n})}{D(I_{n})}=\frac{1}{1}=1
\]
}%
מכאן \L{$D'=D$}.

\medskip{}

\textcolor{teal}{כשהוכחנו את המסקנה הזאת הוכחנו למעשה את היחידות של
פ\textquotedblright נ מנורמלת.}\\

\textcolor{teal}{זה מה שרצינו להגדיר כל הזמן הזה: דטרמיננטה היא פ\textquotedblright נ
מנורמלת. }

\textcolor{teal}{\L{%
\[
\backslash(^{\wedge}O^{\wedge})/
\]
}%
}

\textcolor{teal}{המסקנה האחרונה היא בעצם משפט היחידות של הדטרמיננטה.}

\textcolor{teal}{עכשיו רק חסר להוכיח שבכלל קיימת פ\textquotedbl נ
כזאת וזה נעשה בהרצאה הבאה.}

\newpage{}

\noindent\fbox{\begin{minipage}[t]{1\columnwidth - 2\fboxsep - 2\fboxrule}%
\begin{center}
אלגברה ליניארית )1( - \L{$80134$} \textbar{} האונ' העברית \textbar{}
תשפ\textquotedblright ו א' \textbar{} מרצה: ד\textquotedblright ר
אלכס גורביץ' \textbar{} מתרגל: מר יואב כהן \textbar{} ליקוט ועריכה:
נעם לב
\par\end{center}%
\end{minipage}}

\part{הרצאה {\beginL 28\endL} - {\beginL 19.01.2026\endL}}

\textcolor{blue}{מה שעדיין חסר זה הקיום: בשלב זה לא הראינו בשום צורה
שפ\textquotedblright נ מנורמלת בכלל קיימת.}

\textcolor{blue}{למעשה, הדוגמה היחידה לפ\textquotedblright נ שיש לנו
עד כה היא פ\textquotedblright נ הטריוויאלית, ששולחת כל מטריצה לסקלר
\L{$0$}, והיא אמנם מקיימת את כל הדרישות אך היא אינה מנורמלת. המטרה
של הרצאה זאת היא להוכיח שקיימות פ\textquotedblright נ נוספות ובפרט
פ\textquotedblright נ מנורמלת. נוכיח זאת באינדוקציה: נראה איך לבנות
פ\textquotedblright נ מנורמלת על מטריצות בגודל \L{$n\times n$} מתוך
פ\textquotedblright נ מנורמלת על מטריצות בגודל \L{$(n-1)\times(n-1)$},
וכמובן שהבסיס קיים, מה שלא קשה: עבור \L{$n=1$} מטריצות בגודל \L{$1\times1$}
הן פשוט איברים של השדה, דרישת המולטי-לינאריות הופכת לדרישת לינאריות,
דרישת ההתחלפות אינה רלוונטית )אין שתי שורות(, ודרישת הנורמליזציה אומרת
שהמספר \L{$1$} יישלח ל-\L{$1$}. השלב האינדוקטיבי הוא החלק העיקרי
והמסובך.}

\medskip{}

\rule[0.5ex]{1\columnwidth}{1pt}

\medskip{}

\textcolor{blue}{כדי לנסח את הנוסחה שתבצע את שלב האינדוקציה, נזדקק
למושג מינור של מטריצה.}

\section{{\beginL 1\endL}. הגדרה}

תהי \L{$A\in M_{n\times n}(\mathbb{F})$} מטריצה ריבועית, ויהיו \L{$1\le i,j\le n$}. 

נגדיר את \L{$A_{ij}\in M_{(n-1)\times(n-1)}(\mathbb{F})$} להיות המטריצה
המתקבלת מ-\L{$A$} על ידי מחיקת השורה ה-\L{$i$} והעמודה ה-\L{$j$}.
נקרא לה מינור.
\begin{itemize}
\item \textcolor{teal}{לכל זוג \L{$i,j$} מוגדר מינור, ולכן לכל מטריצה בגודל
\L{$n\times n$} מוגדרים \L{$n^{2}$} מינורים, וכל אחד מהם הוא מטריצה
בגודל \L{$(n-1)\times(n-1)$}.}
\item \textcolor{teal}{\L{$a_{ij}$} הוא המספר שנמצא בשורה ה-\L{$i$} ובעמודה
ה-\L{$j$} של \L{$A$}, ואילו \L{$A_{ij}$} )באות גדולה( היא המטריצה
שמתקבלת ממחיקת השורה והעמודה שבהן אותו מספר נמצא.}
\end{itemize}

\section{{\beginL 2\endL}. טענה {]}שלב האינדוקציה בבניית פ\textquotedblright נ
מנורמלת{[}}

תהי \L{$D':M_{(n-1)\times(n-1)}(\mathbb{F})\to\mathbb{F}$} פ\textquotedblright נ
מנורמלת. יהי \L{$1\le j\le n$}. נגדיר \L{$D:M_{n\times n}(\mathbb{F})\to\mathbb{F}$}
ע\textquotedblright י: \L{%
\[
D(A)=(-1)^{j+1}\left(a_{1j}D'(A_{1j})-a_{2j}D'(A_{2j})+\dots+(-1)^{n+1}a_{nj}D'(A_{nj})\right)
\]
}%
לכל \L{$A=\begin{pmatrix}a_{11} & \cdots & a_{1n}\\
\vdots &  & \vdots\\
a_{n1} & \cdots & a_{nn}
\end{pmatrix}\in M_{n\times n}(\mathbb{F})$}.

אז \L{$D$} היא פ\textquotedblright נ מנורמלת.

\medskip{}

\textcolor{teal}{יש \L{$n$} מחוברים, כאשר הסימן שלהם מתחלף, וכל מחובר
הוא מהצורה \L{$a_{ij}\cdot D'(A_{ij})$}. }

\textcolor{teal}{כלומר, עוברים על כל המספרים שבעמודה ה-\L{$j$} של
המטריצה, וכופלים כל אחד מהם ב-\L{$D'$} של המטריצה שמתקבלת ממחיקת
השורה והעמודה שבהן הוא נמצא. כשהאינדקס הראשון אי-זוגי הסימן הוא פלוס,
וכשהוא זוגי הסימן הוא מינוס. לכן המחובר ה-\L{$i$} מופיע עם המקדם
\L{$(-1)^{i+1}$}, ואפשר לרשום את הנוסחה גם כך: \L{%
\[
D(A)=(-1)^{j+1}\sum_{i=1}^{n}(-1)^{i+1}a_{ij}D'(A_{ij})
\]
}%
}

\subsection{הוכחה}

כדי להוכיח שהפונקציה \L{$D$} היא פ\textquotedblright נ מנורמלת, עלינו
להוכיח שלושה דברים: שהיא מולטי-לינארית, שהיא מתחלפת, ושהיא מנורמלת.

\subsubsection{שלב א': \L{$D$} מולטי-לינארית}

לכל \L{$1\le i\le n$} נגדיר \L{$D_{i}:M_{n\times n}(\mathbb{F})\to\mathbb{F}$}
על ידי: \L{%
\[
D_{i}(A)=a_{ij}D'(A_{ij})
\]
}%
לכל \L{$A=\begin{pmatrix}a_{11} & \cdots & a_{1n}\\
\vdots &  & \vdots\\
a_{n1} & \cdots & a_{nn}
\end{pmatrix}\in M_{n\times n}(\mathbb{F})$}.

ראינו בתחילת הפרק שסכום של פונקציות מולטי-לינאריות הוא מולטי-לינארי,
וכן שכפל של פונקציה מולטי-לינארית בסקלר נותן פונקציה מולטי-לינארית.
מאחר ש-\L{$D=(-1)^{j+1}\sum_{i=1}^{n}(-1)^{i+1}D_{i}$}, מספיק להוכיח
כי \L{$D_{i}$} מולטי-לינארית לכל \L{$1\le i\le n$}.

\textcolor{teal}{לא כל מחובר בנפרד הוא פ\textquotedblright נ. אמנם
כל \L{$D_{i}$} בנפרד הוא מולטי-לינארי, אבל הוא לא מתחלף.}

\textcolor{teal}{למשל, עבור \L{$i=j=1$} ומטריצה \L{$A$} ששתי שורותיה
הראשונות זהות: במטריצה \L{$A_{11}$} מחקנו את השורה הראשונה, ולכן
לא נשארו בה שתי שורות זהות ואי אפשר להסיק ש-\L{$D'(A_{11})=0$}. לכן
אנחנו זקוקים לכל המחוברים יחד, שהם אלה שיתנו לנו את ההתחלפות.}

\medskip{}

יהי \L{$1\le k\le n$}. נוכיח כי \L{$D_{i}$} לינארית לפי השורה ה-\L{$k$},
ונפריד לשני מקרים: \L{$k\neq i$} ו-\L{$k=i$}.

\paragraph*{מקרה א': \L{$k\protect\neq i$}}

שמירה על חיבור:

יהי \L{$1\le k\le n$} כך ש-\L{$k\neq i$}. נוכיח כי \L{$D_{i}$}
היא לינארית לפי השורה ה-\L{$k$}.

תהיינה \L{$A,B,C\in M_{n\times n}(\mathbb{F})$} כך שלכל \L{$1\le l\le n$},
\L{$l\neq k$} מתקיים \L{$\Lambda_{l}^{A}=\Lambda_{l}^{B}=\Lambda_{l}^{C}$},
וכן: \L{%
\[
(*)\qquad\Lambda_{k}^{A}+\Lambda_{k}^{B}=\Lambda_{k}^{C}
\]
}%
בפרט, מאחר ש-\L{$i\neq k$}, מתקיים \L{$\Lambda_{i}^{A}=\Lambda_{i}^{B}=\Lambda_{i}^{C}$}
ולכן \L{$a_{ij}=b_{ij}=c_{ij}$}.

מ-\L{$(*)$} נובע שאם \L{$k<i$} אז \L{$\Lambda_{k}^{A_{ij}}+\Lambda_{k}^{B_{ij}}=\Lambda_{k}^{C_{ij}}$},
ואם \L{$k>i$} אז \L{$\Lambda_{k-1}^{A_{ij}}+\Lambda_{k-1}^{B_{ij}}=\Lambda_{k-1}^{C_{ij}}$}.

\textcolor{teal}{כשמוחקים מ-\L{$A,B,C$} את השורה ה-\L{$i$} ואת העמודה
ה-\L{$j$}, מקבלים שלוש מטריצות בגודל \L{$(n-1)\times(n-1)$} שכל
שורותיהן זהות בהתאמה חוץ משורה אחת, שעליה מתקיים אותו יחס כמו במטריצות
המקוריות. אם מחקנו שורה שאחרי השורה ה-\L{$k$} האינדקס נשמר, ואם מחקנו
שורה שלפניה האינדקס יורד ב-\L{$1$}.}

מאחר ש-\L{$D'$} היא מולטי-לינארית, מתקיים \L{$D'(A_{ij})+D'(B_{ij})=D'(C_{ij})$}.
מכאן: \L{%
\begin{align*}
D_{i}(C) & =c_{ij}D'(C_{ij})=c_{ij}D'(A_{ij})+c_{ij}D'(B_{ij})=\\
 & =a_{ij}D'(A_{ij})+b_{ij}D'(B_{ij})=D_{i}(A)+D_{i}(B)
\end{align*}
}%

שמירה על כפל בסקלר:

תהיינה \L{$A,B\in M_{n\times n}(\mathbb{F})$} ו-\L{$c\in\mathbb{F}$}
כך שלכל \L{$1\le l\le n$}, \L{$l\neq k$} מתקיים \L{$\Lambda_{l}^{A}=\Lambda_{l}^{B}$},
וכן \L{$c\Lambda_{k}^{A}=\Lambda_{k}^{B}$}.

אז באופן דומה אפשר להראות כי: \L{%
\[
D_{i}(B)=cD_{i}(A)
\]
}%


\paragraph*{מקרה ב': \L{$k=i$}}

שמירה על חיבור:

נוכיח כי \L{$D_{i}$} לינארית לפי השורה ה-\L{$i$}.

\textcolor{teal}{השורה שבה המטריצות נבדלות זו מזו היא בדיוק השורה
שאותה אנחנו מוחקים, לכן הנימוקים של המקרה הקודם כבר אינם תקפים.}

תהיינה \L{$A,B,C\in M_{n\times n}(\mathbb{F})$} כך שלכל \L{$1\le l\le n$},
\L{$l\neq i$} מתקיים \L{$\Lambda_{l}^{A}=\Lambda_{l}^{B}=\Lambda_{l}^{C}$},
וכן \L{$\Lambda_{i}^{A}+\Lambda_{i}^{B}=\Lambda_{i}^{C}$}.

מכאן \L{$a_{ij}+b_{ij}=c_{ij}$}, ובנוסף \L{$A_{ij}=B_{ij}=C_{ij}$}.

\textcolor{teal}{\L{$A,B,C$} זהות בכל השורות חוץ מהשורה ה-\L{$i$},
ולכן כשמוחקים בשלושתן דווקא את השורה ה-\L{$i$} )ואת אותה עמודה(,
נשארות שלוש מטריצות שכל שורותיהן שוות בהתאמה, כלומר שלוש מטריצות שוות.}

מכאן: \L{%
\begin{align*}
D_{i}(C) & =c_{ij}D'(C_{ij})=a_{ij}D'(C_{ij})+b_{ij}D'(C_{ij})=\\
 & =a_{ij}D'(A_{ij})+b_{ij}D'(B_{ij})=D_{i}(A)+D_{i}(B)
\end{align*}
}%

שמירה על כפל בסקלר:

תהיינה \L{$A,B\in M_{n\times n}(\mathbb{F})$} ו-\L{$c\in\mathbb{F}$}
כך שלכל \L{$1\le l\le n$}, \L{$l\neq i$} מתקיים \L{$\Lambda_{l}^{A}=\Lambda_{l}^{B}$},
וכן \L{$c\Lambda_{i}^{A}=\Lambda_{i}^{B}$}. אז באופן דומה אפשר להראות
כי: \L{%
\[
D_{i}(B)=cD_{i}(A)
\]
}%

לסיכום, \L{$D_{i}$} היא לינארית לפי השורה ה-\L{$k$} לכל \L{$1\le k\le n$},
ולכן \L{$D_{i}$} מולטי-לינארית.

זה נכון לכל \L{$1\le i\le n$}, ולכן גם \L{$D$} מולטי-לינארית.

\subsubsection{שלב ב': \L{$D$} מתחלפת}

\textcolor{teal}{נזכיר שפונקציה נקראת מתחלפת כאשר לכל מטריצה שיש לה
שתי שורות סמוכות וזהות היא מחזירה \L{$0$}.}

תהי \L{$A\in M_{n\times n}(\mathbb{F})$} כך ש-\L{$\Lambda_{k}^{A}=\Lambda_{k+1}^{A}$}.

לכל \L{$1\le i\le n$} כך ש-\L{$i\neq k$} ו-\L{$i\neq k+1$}, למטריצה
\L{$A_{ij}$} יש שתי שורות סמוכות זהות, ולכן \L{$D'(A_{ij})=0$}.

\textcolor{teal}{אם מוחקים מ-\L{$A$} שורה שאיננה השורה ה-\L{$k$}
ואיננה השורה ה-\L{$k+1$}, שתי השורות הזהות נשארות במטריצה והן עדיין
סמוכות )אם \L{$i>k+1$} האינדקסים שלהן נשארים \L{$k,k+1$}, ואם \L{$i<k$}
הם הופכים ל-\L{$k-1,k$}(, ומאחר ש-\L{$D'$} מתחלפת מתקבל \L{$0$}.}

בנוסף, \L{$A_{kj}=A_{k+1\,j}$} וכן \L{$a_{kj}=a_{k+1\,j}$}.

\textcolor{teal}{השורה ה-\L{$k$} והשורה ה-\L{$k+1$} סמוכות וזהות,
ולכן לא משנה איזו מהשתיים נמחק, נקבל את אותה מטריצה, וגם המספרים שבעמודה
ה-\L{$j$} בשתי השורות האלה שווים.}

מכאן, כל המחוברים בנוסחה מתאפסים חוץ משניים, ומתקיים: \L{%
\[
D(A)=(-1)^{j+1}\left((-1)^{k+1}a_{kj}D'(A_{kj})+(-1)^{k+2}a_{k+1\,j}D'(A_{k+1\,j})\right)=0
\]
}%

\textcolor{teal}{שני המחוברים שנשארו שווים זה לזה בערכם המוחלט, אך
הסימנים שלהם, \L{$(-1)^{k+1}$} ו-\L{$(-1)^{k+2}$} , הפוכים, ולכן
הם מבטלים זה את זה והסכום מתאפס.}

\subsubsection{שלב ג': \L{$D$} מנורמלת}

תהי \L{$A=I_{n}$}. אז לכל \L{$1\le i\le n$} כך ש-\L{$i\neq j$}
מתקיים \L{$a_{ij}=0$}, וכן \L{$a_{jj}=1$}. כמו כן \L{$A_{jj}=I_{n-1}$}.

מכאן: \L{%
\[
D(I_{n})=(-1)^{j+1}\left((-1)^{j+1}a_{jj}D'(A_{jj})\right)=(-1)^{2j+2}\cdot1\cdot D'(I_{n-1})=1
\]
}%

\hfill{}\L{$\square$}

\medskip{}

\rule[0.5ex]{1\columnwidth}{1pt}

\medskip{}

\textcolor{blue}{הטענה שהוכחנו מבטיחה לנו את שלב האינדוקציה, ועכשיו
אפשר לנסח את משפט הקיום של הדטרמיננטה.}

\section{{\beginL 3\endL}. משפט {]}קיום פ\textquotedblright נ מנורמלת{[}}

לכל \L{$n\in\mathbb{N}$} קיימת פ\textquotedblright נ מנורמלת \L{$D:M_{n\times n}(\mathbb{F})\to\mathbb{F}$}.

\subsection{הוכחה}

נוכיח באינדוקציה על \L{$n$}.

\textbf{בסיס:} עבור \L{$n=1$} מתקיים \L{$M_{1\times1}(\mathbb{F})=\mathbb{F}$},
ונגדיר \L{$D=\text{Id}_{\mathbb{F}}$}. אז \L{$D$} פ\textquotedblright נ
מנורמלת.

\textcolor{teal}{פונ' הזהות היא לינארית לפי השורה היחידה, היא מתחלפת
באופן ריק, והיא מנורמלת שכן היא מחזירה \L{$1$} עבור מטריצת היחידה.}

\textbf{שלב האינדוקציה:} נניח כי \L{$D':M_{(n-1)\times(n-1)}(\mathbb{F})\to\mathbb{F}$}
היא פ\textquotedblright נ מנורמלת.

נבחר \L{$1\le j\le n$} כלשהו, ונגדיר \L{$D:M_{n\times n}(\mathbb{F})\to\mathbb{F}$}
כמו בטענה הקודמת. אז \L{$D$} היא פ\textquotedblright נ מנורמלת.

\medskip{}


\section{{\beginL 4\endL}. סיכום: הגדרת הדטרמיננטה}

לכל \L{$n\in\mathbb{N}$} קיימת פ\textquotedblright נ מנורמלת \textbf{יחידה}
מ-\L{$M_{n\times n}(\mathbb{F})$} ל-\L{$\mathbb{F}$}. פונקציה זו
נקראת \textbf{דטרמיננטה} ומסומנת ב-\L{$\det$}.

\textcolor{teal}{\L{$\backslash(^{\wedge}O^{\wedge})/$}} )או \L{$\backslash(v\_v)/$}
אם שואלים את אקסלר(

\textcolor{teal}{בסימון \L{$\det$} לא מציינים את גודל המטריצות שעליהן
הפונקציה מוגדרת, למרות שלמעשה לכל \L{$n$} יש פונקציה משלו.}

\textcolor{teal}{בבניית \L{$D$} בחרנו \L{$j$} כלשהו, ולכאורה נראה
שמכל פ\textquotedblright נ מנורמלת על מטריצות בגודל \L{$(n-1)\times(n-1)$}
אנחנו מקבלים \L{$n$} פונקציות שונות על מטריצות בגודל \L{$n\times n$}.
אבל הוכחנו שפ\textquotedblright נ מנורמלת היא יחידה, ולכן לא משנה
איזה \L{$j$} נבחר, הנוסחה תחזיר תמיד את אותה הפונקציה. לאותה פונקציה
יכולות להיות נוסחאות רבות, אבל הערך שהיא מחזירה עבור מטריצה נתונה
הוא תמיד אותו ערך.}

\section{{\beginL 5\endL}. דוגמה: הדטרמיננטה של מטריצה \L{$2\times2$}}

\subsection{חישוב לפי \L{$j=1$}}

ניקח \L{$n=2$}, \L{$D'=\text{Id}_{\mathbb{F}}$} ו-\L{$j=1$}, ותהי
\L{$A=\begin{pmatrix}a_{11} & a_{12}\\
a_{21} & a_{22}
\end{pmatrix}$}. אז: \L{%
\[
D(A)=(-1)^{1+1}\left(a_{11}D'(A_{11})-a_{21}D'(A_{21})\right)=a_{11}a_{22}-a_{21}a_{12}
\]
}%

\textcolor{teal}{\L{$A_{11}=(a_{22})$}, ו- \L{$A_{21}=(a_{12})$}.}

\subsection{חישוב לפי \L{$j=2$}}

ניקח \L{$n=2$}, \L{$D'=\text{Id}_{\mathbb{F}}$} ו-\L{$j=2$}, ותהי
\L{$A=\begin{pmatrix}a_{11} & a_{12}\\
a_{21} & a_{22}
\end{pmatrix}$}. אז: \L{%
\begin{align*}
D(A) & =(-1)^{2+1}\left(a_{12}D'(A_{12})-a_{22}D'(A_{22})\right)=\\
 & =-\left(a_{12}a_{21}-a_{22}a_{11}\right)=a_{11}a_{22}-a_{21}a_{12}
\end{align*}
}%

\textcolor{teal}{קיבלנו את אותה התוצאה בדיוק, כפי שמבטיחה היחידות.}

\textcolor{teal}{גם עבור \L{$n=3$} והבחירות \L{$j=1,2,3$} נקבל שלוש
נוסחאות שנראות שונות אך נותנות אותו ביטוי אחרי סידור.}

\newpage{}

\noindent\fbox{\begin{minipage}[t]{1\columnwidth - 2\fboxsep - 2\fboxrule}%
\begin{center}
אלגברה ליניארית )1( - \L{$80134$} \textbar{} האונ' העברית \textbar{}
תשפ\textquotedblright ו א' \textbar{} מרצה: ד\textquotedblright ר
אלכס גורביץ' \textbar{} מתרגל: מר יואב כהן \textbar{} ליקוט ועריכה:
נעם לב
\par\end{center}%
\end{minipage}}

\part{הרצאה {\beginL 29\endL} ואחרונה - {\beginL 20.01.2026\endL}}

\section{תזכורת}

\section{{\beginL 1\endL}. הגדרה: דטרמיננטה}

פ\textquotedblright נ מנורמלת מ-\L{$M_{n\times n}(\mathbb{F})$} ל-\L{$\mathbb{F}$}
נקראת \textbf{דטרמיננטה} ומסומנת ב-\L{$\det$}.

\section{{\beginL 2\endL}. טענה {]}התנהגות פ\textquotedblright נ תחת פש\textquotedblright א{[}}

תהי \L{$D:M_{n\times n}(\mathbb{F})\to\mathbb{F}$} פ\textquotedblright נ,
\L{$A\in M_{n\times n}(\mathbb{F})$} ו-\L{$\varepsilon$} פש\textquotedblright א.
אז: \L{%
\[
D(\varepsilon(A))=\mu(\varepsilon)\cdot D(A)
\]
}%

\textcolor{teal}{נזכיר את הערכים של \L{$\mu(\varepsilon)$}:}
\begin{enumerate}
\item \textcolor{teal}{אם \L{$\varepsilon=R_{i}\to cR_{i}$} כאשר \L{$c\neq0$},
אז \L{$\mu(\varepsilon)=c$}.}
\item \textcolor{teal}{אם \L{$\varepsilon=R_{i}\to R_{i}+cR_{j}$} כאשר
\L{$i\neq j$}, אז \L{$\mu(\varepsilon)=1$}.}
\item \textcolor{teal}{אם \L{$\varepsilon=R_{i}\leftrightarrow R_{j}$}
כאשר \L{$i\neq j$}, אז \L{$\mu(\varepsilon)=-1$}.}
\end{enumerate}
\medskip{}

\textcolor{blue}{מכיוון שהדטרמיננטה היא פ\textquotedblright נ, כל
הטענות שהוכחנו על פונקציות נפח תקפות לגביה. לכן, מלבד הנוסחה מההרצאה
הקודמת, יש דרך יעילה יותר לחשב דטרמיננטה של מטריצה נתונה: לדרג אותה
בעזרת פש\textquotedblright א, ולעקוב אחרי הסקלר שבו הדטרמיננטה מוכפלת
בכל שלב. בסוף מתקבלת מטריצה מדורגת קנונית, שהיא או מטריצת היחידה )ואז
הדטרמיננטה המקורית מתקבלת מהסקלרים שנצברו(, או מטריצה עם שורת אפסים
)ואז הדטרמיננטה המקורית היא \L{$0$}(.}

\medskip{}

\rule[0.5ex]{1\columnwidth}{1pt}

\medskip{}


\section{{\beginL 3\endL}. למה {]}הקבוע של פש\textquotedblright א הוא הדטרמיננטה
של המטריצה האלמנטרית{[}}

תהי \L{$\varepsilon$} פש\textquotedblright א. אז: \L{%
\[
\mu(\varepsilon)=\det\varepsilon(I_{n})
\]
}%


\subsection{הוכחה}

\L{%
\[
\det(\varepsilon(I_{n}))=\mu(\varepsilon)\det(I_{n})=\mu(\varepsilon)
\]
}%

\textcolor{teal}{הלמה הזאת נכונה אך ורק לדטרמיננטה, ולא לכל פ\textquotedblright נ.
בשוויון הראשון השתמשנו בטענה }\textcolor{teal}{{\beginL 2}\endL}\textcolor{teal}{{}
)שנכונה לכל פ\textquotedblright נ(, אבל בשוויון השני השתמשנו באופן
מהותי בכך שהדטרמיננטה }\textbf{\textcolor{teal}{מנורמלת}}\textcolor{teal}{,
כלומר \L{$\det(I_{n})=1$}.}

\section{{\beginL 4\endL}. טענה}

תהי \L{$A\in M_{n\times n}(\mathbb{F})$} ו-\L{$\varepsilon$} פש\textquotedblright א.
אז: \L{%
\[
\det(\varepsilon(I_{n})A)=\det\varepsilon(I_{n})\cdot\det A
\]
}%


\subsection{הוכחה}

\L{%
\[
\det(\varepsilon(I_{n})\cdot A)=\det(\varepsilon(A))=\mu(\varepsilon)\det A=\det\varepsilon(I_{n})\cdot\det A
\]
}%

\textcolor{teal}{בשוויון הראשון השתמשנו בטענה שראינו כשהתחלנו לדבר
על מטריצות אלמנטריות: כפל של מטריצה במטריצה האלמנטרית המתאימה ל-\L{$\varepsilon$}
משמאל שקול להפעלת \L{$\varepsilon$} על אותה מטריצה. בשוויון השני
השתמשנו בטענה }\textcolor{teal}{{\beginL 2}\endL}\textcolor{teal}{,
ובשלישי בלמה }\textcolor{teal}{{\beginL 3}\endL}\textcolor{teal}{.}

\section{{\beginL 5\endL}. למה}

תהיינה \L{$A,B\in M_{n\times n}(\mathbb{F})$}. אם \L{$A$} אינה הפיכה,
אז \L{$AB$} אינה הפיכה.

\subsection{הוכחה}

נניח בשלילה כי \L{$AB$} הפיכה. כלומר, קיימת \L{$C\in M_{n\times n}(\mathbb{F})$}
כך ש-\L{$C(AB)=(AB)C=I_{n}$}.

מכאן \L{$A(BC)=I_{n}$}, ולכן גם \L{$(BC)A=I_{n}$}, ולכן \L{$A$}
הפיכה -  סתירה.

\section{{\beginL 6\endL}. משפט}

תהיינה \L{$A,B\in M_{n\times n}(\mathbb{F})$}. אז: \L{%
\[
\det(AB)=\det A\cdot\det B
\]
}%


\subsection{הוכחה}

\textbf{מקרה }\textbf{{\beginL 1}\endL}\textbf{:} \L{$A$} אינה הפיכה.

אז לפי למה {\beginL 5\endL} גם \L{$AB$} אינה הפיכה. לכן: \L{%
\[
\det(AB)=0=0\cdot\det B=\det A\cdot\det B
\]
}%

\textbf{מקרה }\textbf{{\beginL 2}\endL}\textbf{:} \L{$A$} הפיכה.

אז קיימת סדרה של פש\textquotedblright א \L{$\varepsilon_{1},\dots,\varepsilon_{k}$}
כך ש: \L{%
\[
I_{n}=\varepsilon_{k}(\dots\varepsilon_{2}(\varepsilon_{1}(A))\dots)
\]
}%

לכל \L{$1\le i\le k$} נסמן ב-\L{$\delta_{i}$} את הפש\textquotedblright א
ההופכית ל-\L{$\varepsilon_{i}$}. אז: \L{%
\[
\delta_{1}(\delta_{2}(\dots(\delta_{k}(I_{n}))\dots))=\delta_{1}(\delta_{2}(\dots(\delta_{k}(\varepsilon_{k}(\dots(\varepsilon_{2}(\varepsilon_{1}(A)))\dots)))\dots))=A
\]
}%

\textcolor{teal}{כל \L{$\delta_{i}$} מבטלת את \L{$\varepsilon_{i}$}
שלידה, וכך כולן מצטמצמות בזוגות ונשארים עם \L{$A$}.}

מכאן, מאחר שהפעלת פש\textquotedblright א על מטריצה שקולה לכפל במטריצה
האלמנטרית המתאימה משמאל: \L{%
\[
\delta_{1}(I_{n})\cdot\delta_{2}(I_{n})\cdot\dots\cdot\delta_{k-1}(I_{n})\cdot\delta_{k}(I_{n})=A
\]
}%

לכן, בעזרת טענה {\beginL 4\endL} )\L{$k-1$} פעמים(: \L{%
\begin{align*}
\det A & =\det\left(\delta_{1}(I_{n})\cdot\delta_{2}(I_{n})\cdot\dots\cdot\delta_{k}(I_{n})\right)=\det\delta_{1}(I_{n})\cdot\det\left(\delta_{2}(I_{n})\cdot\dots\cdot\delta_{k}(I_{n})\right)=\\
 & =\det\delta_{1}(I_{n})\cdot\det\delta_{2}(I_{n})\cdot\dots\cdot\det\delta_{k}(I_{n})
\end{align*}
}%

ובאותו אופן, עבור המכפלה \L{$AB$}: \L{%
\begin{align*}
\det(AB) & =\det\left(\delta_{1}(I_{n})\cdot\dots\cdot\delta_{k}(I_{n})\cdot B\right)=\det\delta_{1}(I_{n})\cdot\det\left(\delta_{2}(I_{n})\cdot\dots\cdot\delta_{k}(I_{n})\cdot B\right)=\dots=\\
 & =\det\delta_{1}(I_{n})\cdot\dots\cdot\det\delta_{k}(I_{n})\cdot\det B=\det A\cdot\det B
\end{align*}
}%


\section{{\beginL 7\endL}. מסקנה {]}דטרמיננטה של מטריצה הפיכה ושל ההופכית
שלה{[}}

אם \L{$A\in M_{n\times n}(\mathbb{F})$} הפיכה, אז \L{$\det A\neq0$}.

בנוסף, אם \L{$B\in M_{n\times n}(\mathbb{F})$} הופכית ל-\L{$A$},
אז: \L{%
\[
\det B=\frac{1}{\det A}
\]
}%


\subsection{הוכחה}

מתקיים \L{$AB=I_{n}$}. מכאן: \L{%
\[
1=\det I_{n}=\det(AB)=\det A\cdot\det B
\]
}%

לכן \L{$\det A\neq0$} )אחרת היינו מקבלים \L{$0=1$}(,ו- \L{$\det B=\frac{1}{\det A}$}.

\textcolor{teal}{ראינו כבר שאם מטריצה אינה הפיכה אז הדטרמיננטה שלה
היא \L{$0$}. יחד עם המסקנה הזאת מתקבל קריטריון נוסף להפיכות: מטריצה
ריבועית \L{$A$} היא הפיכה אם\textquotedblright ם \L{$\det A\neq0$}.}

\section{{\beginL 8\endL}. תזכורת: מטריצות דומות}

מטריצות ריבועיות \L{$A,B\in M_{n\times n}(\mathbb{F})$} נקראות \textbf{דומות}
כאשר קיימת \L{$P\in M_{n\times n}(\mathbb{F})$} הפיכה כך ש: \L{%
\[
B=P^{-1}AP
\]
}%


\section{{\beginL 9\endL}. מסקנה {]}למטריצות דומות יש אותה דטרמיננטה{[}}

תהיינה \L{$A,B\in M_{n\times n}(\mathbb{F})$} דומות. אז \L{$\det A=\det B$}.

\subsection{הוכחה}

קיימת \L{$P\in M_{n\times n}(\mathbb{F})$} הפיכה כך ש-\L{$B=P^{-1}AP$}.
מכאן: \L{%
\[
\det B=\det(P^{-1}AP)=\det(P^{-1})\det(AP)=\frac{1}{\det P}\cdot\det A\cdot\det P=\det A
\]
}%

\textcolor{teal}{המסקנה נותנת דרך חדשה להראות ששתי מטריצות אינן דומות:
אם הדטרמיננטות שלהן שונות.}

\section{{\beginL 10\endL}. למה {]}פש\textquotedblright א המתאימה למטריצה
האלמנטרית המשוחלפת{[}}

תהי \L{$\varepsilon$} פש\textquotedblright א. אז קיימת פש\textquotedblright א
\L{$\varepsilon^{*}$} כך ש: \L{%
\[
\varepsilon^{*}(I_{n})=\varepsilon(I_{n})^{t},\quad\mu(\varepsilon^{*})=\mu(\varepsilon)
\]
}%


\subsection{הוכחה}

נעבור על שלוש האפשרויות לפש\textquotedblright א.
\begin{enumerate}
\item אם \L{$\varepsilon=R_{i}\to cR_{i}$} כאשר \L{$c\neq0$}, אז: \L{%
\[
\varepsilon(I_{n})=\begin{pmatrix}1 &  &  &  & 0\\
 & \ddots\\
 &  & c\\
 &  &  & \ddots\\
0 &  &  &  & 1
\end{pmatrix}
\]
}%
זו מטריצה סימטרית, כלומר \L{$\varepsilon(I_{n})^{t}=\varepsilon(I_{n})$}.
אז נגדיר \L{$\varepsilon^{*}=\varepsilon$}.
\item אם \L{$\varepsilon=R_{i}\leftrightarrow R_{j}$} כאשר \L{$i\neq j$},
אז: \L{%
\[
\varepsilon(I_{n})=\begin{pmatrix}1\\
 & 0 &  & 1\\
 &  & \ddots\\
 & 1 &  & 0\\
 &  &  &  & 1
\end{pmatrix}
\]
}%
גם זו מטריצה סימטרית, כלומר \L{$\varepsilon(I_{n})^{t}=\varepsilon(I_{n})$}.
אז נגדיר \L{$\varepsilon^{*}=\varepsilon$}.
\item אם \L{$\varepsilon=R_{i}\to R_{i}+cR_{j}$} כאשר \L{$i\neq j$}, אז:
\L{%
\[
\varepsilon(I_{n})=\begin{pmatrix}1\\
 & \ddots &  & c\\
 &  & \ddots\\
 &  &  & \ddots\\
 &  &  &  & 1
\end{pmatrix}\qquad\varepsilon(I_{n})^{t}=\begin{pmatrix}1\\
 & \ddots\\
 &  & \ddots\\
 & c &  & \ddots\\
 &  &  &  & 1
\end{pmatrix}
\]
}%
כאן המטריצה אינה סימטרית: ב-\L{$\varepsilon(I_{n})$} הסקלר \L{$c$}
יושב במקום \L{$(i,j)$}, ובמשוחלפת הוא עובר למקום \L{$(j,i)$}. נגדיר
\L{$\varepsilon^{*}=R_{j}\to R_{j}+cR_{i}$}, ואז \L{$\varepsilon^{*}(I_{n})=\varepsilon(I_{n})^{t}$}.
כמו כן: \L{%
\[
\mu(\varepsilon^{*})=1=\mu(\varepsilon)
\]
}%
שהרי שתי הפעולות הן מסוג הוספת כפולה של שורה לשורה אחרת.
\end{enumerate}

\section{{\beginL 11\endL}. משפט {]}דטרמיננטה של מטריצה משוחלפת{[}}

תהי \L{$A\in M_{n\times n}(\mathbb{F})$}. אז: \L{%
\[
\det(A^{t})=\det(A)
\]
}%


\subsection{הוכחה}

\textbf{מקרה }\textbf{{\beginL 1}\endL}\textbf{:} \L{$A$} אינה הפיכה.
אז \L{$A^{t}$} אינה הפיכה )כי אחרת, היינו מקבלים ש- \L{$A=(A^{t})^{t}$}
הפיכה וזו סתירה(. לכן: \L{%
\[
\det A^{t}=0=\det A
\]
}%

\textbf{מקרה }\textbf{{\beginL 2}\endL}\textbf{:} \L{$A$} הפיכה.

אז )כמו בהוכחת משפט {\beginL 6\endL}( קיימת סדרה של פש\textquotedblright א
\L{$\varepsilon_{1},\dots,\varepsilon_{k}$} כך ש- \L{$I_{n}=\varepsilon_{k}(\dots\varepsilon_{2}(\varepsilon_{1}(A))\dots)$}. 

לכל \L{$1\le i\le k$} נסמן ב-\L{$\delta_{i}$} את הפש\textquotedblright א
ההופכית ל-\L{$\varepsilon_{i}$}, ונקבל: \L{%
\[
A=\delta_{1}(I_{n})\cdot\dots\cdot\delta_{k}(I_{n}),\quad\det A=\det\delta_{1}(I_{n})\cdot\dots\cdot\det\delta_{k}(I_{n})
\]
}%

מכאן, מאחר שהשחלוף של מכפלה הוא מכפלת השחלופים בסדר הפוך: \L{%
\[
A^{t}=\left(\delta_{1}(I_{n})\cdot\dots\cdot\delta_{k}(I_{n})\right)^{t}=\delta_{k}(I_{n})^{t}\cdot\dots\cdot\delta_{1}(I_{n})^{t}
\]
}%

לפי למה {\beginL 10\endL}, לכל \L{$1\le i\le k$} קיימת פש\textquotedblright א
\L{$\delta_{i}^{*}$} כך ש- \L{$\delta_{i}^{*}(I_{n})=\delta_{i}(I_{n})^{t}$},
לכן \L{$\mu(\delta_{i}^{*})=\mu(\delta_{i})$}.

מכאן, ממשפט {\beginL 6\endL} נובע כי: \L{%
\begin{align*}
\det A^{t} & =\det\delta_{k}(I_{n})^{t}\cdot\dots\cdot\det\delta_{1}(I_{n})^{t}=\det\delta_{k}^{*}(I_{n})\cdot\dots\cdot\det\delta_{1}^{*}(I_{n})=\\
 & =\mu(\delta_{k}^{*})\cdot\dots\cdot\mu(\delta_{1}^{*})=\mu(\delta_{1})\cdot\dots\cdot\mu(\delta_{k})=\\
 & =\det\delta_{1}(I_{n})\cdot\dots\cdot\det\delta_{k}(I_{n})=\det A
\end{align*}
}%

\textcolor{teal}{במעבר מ-\L{$\mu(\delta_{k}^{*})\cdot\dots\cdot\mu(\delta_{1}^{*})$}
ל-\L{$\mu(\delta_{1})\cdot\dots\cdot\mu(\delta_{k})$} מחקנו את הכוכביות
לפי הלמה וגם שינינו את סדר הגורמים. זה מותר כי כולם סקלרים והכפל בשדה
חילופי. את המעברים מ-\L{$\det$} ל-\L{$\mu$} ובחזרה עשינו לפי למה
}\textcolor{teal}{{\beginL 3}\endL}\textcolor{teal}{.}

\medskip{}

\rule[0.5ex]{1\columnwidth}{1pt}

\medskip{}

\textcolor{blue}{למשפט הזה יש הרבה מאוד תועלת. הגדרנו פ\textquotedblright נ
כפונקציה מולטי-לינארית, כאשר הלינאריות התייחסה לשורות של המטריצה.
אבל השורות של \L{$A^{t}$} הן העמודות של \L{$A$}, ומאחר ש-\L{$\det A^{t}=\det A$},
כל טענה שניסחנו על פעולות שורה תקפה באותה מידה לפעולות עמודה. למשל,
אם נכפול עמודה של מטריצה בסקלר, הדטרמיננטה תוכפל באותו סקלר.}

\section{{\beginL 12\endL}. תזכורת: פיתוח הדטרמיננטה לפי עמודה}

יהי \L{$n\in\mathbb{N}$}, \L{$1\le j\le n$}, ותהי \L{$A=\begin{pmatrix}a_{11} & \cdots & a_{1n}\\
\vdots &  & \vdots\\
a_{n1} & \cdots & a_{nn}
\end{pmatrix}\in M_{n\times n}(\mathbb{F})$}. אז: \L{%
\[
\det A=(-1)^{j+1}\left(a_{1j}\det(A_{1j})-a_{2j}\det(A_{2j})+\dots+(-1)^{n+1}a_{nj}\det(A_{nj})\right)
\]
}%

\textcolor{teal}{הסקלרים \L{$a_{1j},a_{2j},\dots,a_{nj}$} נמצאים
כולם בעמודה ה-\L{$j$} של המטריצה, ולכן היא נקראת נוסחת הפיתוח של
הדטרמיננטה לפי העמודה}\textbf{\textcolor{teal}{{} ה-\L{$j$}}}\textcolor{teal}{.}

\section{{\beginL 13\endL}. מסקנה: פיתוח הדטרמיננטה לפי שורה}

באותן ההנחות מתקיים: \L{%
\[
\det A=\det A^{t}=(-1)^{j+1}\left(a_{j1}\det(A_{j1})-a_{j2}\det(A_{j2})+\dots+(-1)^{n+1}a_{jn}\det(A_{jn})\right)
\]
}%

\textcolor{teal}{במטריצה המשוחלפת כל האינדקסים מתהפכים: מה שעומד בשורה
ה-\L{$i$} ובעמודה ה-\L{$j$} של \L{$A^{t}$} הוא \L{$a_{ji}$} ולא
\L{$a_{ij}$}. }

\textcolor{teal}{לכן מתקבלת אותה נוסחה כאשר כל האינדקסים הפוכים.}

\textcolor{teal}{הסקלרים \L{$a_{j1},a_{j2},\dots,a_{jn}$} נמצאים
כולם בשורה ה-\L{$j$} של המטריצה, ולכן הנוסחה נקראת נוסחת הפיתוח של
הדטרמיננטה לפי השורה ה-\L{$j$}.}

\medskip{}

\rule[0.5ex]{1\columnwidth}{1pt}

\medskip{}


\section{מבנה המבחן}

\textcolor{teal}{בסוף ההרצאה אלכס כתב על הלוח את מבנה המבחן:}
\begin{itemize}
\item \textbf{\textcolor{teal}{חלק א'}}\textcolor{teal}{: }\textcolor{teal}{{\beginL 2}\endL}\textcolor{teal}{{}
שאלות )הוכחות מההרצאה(, צריך לבחור שאלה אחת.}
\item \textbf{\textcolor{teal}{חלק ב'}}\textcolor{teal}{: }\textcolor{teal}{{\beginL 7}\endL}\textcolor{teal}{{}
שאלות )פתוחות, בסגנון תרגילי הבית(, צריך לבחור }\textcolor{teal}{{\beginL 6}\endL}\textcolor{teal}{{}
שאלות.}
\end{itemize}
\textcolor{teal}{משך המבחן הוא שלוש שעות. בהוכחות מההרצאה מותר להסתמך
על טענות קודמות מהקורס בלי להוכיח אותן, אבל חובה לציין במפורש על מה
מסתמכים.}

בהצלחה לכולנו ותודה מיוחדת ליואב המתרגל!
\end{document}
